% ============================================================
% 导学答案册（S5 试产臂汇编草案）—— 源＝导学件 各件 main.tex（只读）；工具/答案抽册器.py 逐件抽册口径，同序同号逐片保持
%% 源件 md5 见 逐件汇总.json；值/详解逐字照源（钉值硬断言）；题面不重印
%% 双档：ansbook-true＝详解印本／ansbook-false＝纯值速查本（\ansbookdetail 件面开关）
% 编译：xelatex <壳> ×2，cwd＝本目录；sty＝qp-m3.sty 副本（md5 7c3930362be8a0a2bdf21bbf8ac16573 锁同）
% ============================================================
\documentclass[fontset=none]{ctexart}
\usepackage{qp-m3}
% —— 册面补钉：pifont（源件值域含 \ding{51} 勾形；qp-m3 不供给，照 ROUTING 承源件同精神挂载，未用件零副作用） ——
\usepackage{pifont}
\xeCJKDeclareCharClass{Default}{"2212}%
\setCJKfamilyfont{nscblk}[Path=C:/提示词/工作区/字替对照-0909/variantF/fonts/,BoldFont=NSC-w600.ttf]{NSC-w500.ttf}%
\xeCJKDeclareSubCJKBlock{nscblk}{"25B1}%
\newcommand{\pxparallelogram}{{\CJKfamily{nscblk}▱}}%
\xeCJKsetup{CJKglue={\hskip 0.10em plus 0.08em minus 0.05em}}%
\ansblockgrayfalse
\newif\ifansbookdetail
\ifdefined\ansbookdetail
  \ifnum\ansbookdetail=0 \ansbookdetailfalse\else\ansbookdetailtrue\fi
\else
  \ansbookdetailtrue
\fi
\renewcommand{\qpzhangming}{第二章\quad 平面解析几何}%
\renewcommand{\qpjianming}{导学答案册}%

\begin{document}
\zhangtitle{导学答案册}
\par\glueguard{1}\addvspace{2pt}\noindent\biaoqian{【纯答案册】}{\fangsong 题面不重印\quad 逐片印面号与正册同号同序}\par\addvspace{2pt}

\begin{multicols}{2}
\emergencystretch=1em

\xjkeshi{2.8 直线与圆锥曲线的位置关系}

\begin{ansblock}[2章-导-衔接-探1]
% ans:2章-导-衔接-探1
\ansitem{1}{C}
\ifansbookdetail
\ansnote{详解}{将\(y=2x+m\)代入\(4x^{2}-y^{2}=1\)，整理得\(4mx=-(m^{2}+1)\)．当\(m=0\)时方程无解，此时直线即为双曲线的渐近线\(y=2x\)，与双曲线无公共点；当\(m\neq 0\)时方程有唯一解\(x=-(m^{2}+1)/(4m)\)，恰有一个交点．故对任意\(m\)，直线与双曲线至多有一个交点，选C．（注意：此处"恰一个交点"源于直线与渐近线平行，并非相切．）}
\fi
\end{ansblock}

\begin{ansblock}[2章-导-衔接-探2]
% ans:2章-导-衔接-探2
\ansitem{2}{3 条}
\ifansbookdetail
\ansnote{详解}{分类讨论：①斜率不存在：直线\(x=0\)与抛物线相切于顶点\((0,0)\)，符合；②斜率\(k=0\)：直线\(y=1\)与抛物线的对称轴平行，恰有一个公共点，符合；③斜率\(k\neq 0\)存在：将\(y=kx+1\)代入\(y^{2}=2x\)得\(k^{2}x^{2}+2(k-1)x+1=0\)，由\(\Delta=4(k-1)^{2}-4k^{2}=4-8k=0\)解得\(k=1/2\)，此时直线\(y=x/2+1\)与抛物线相切，符合．共3条．易错：漏①②两类．}
\fi
\end{ansblock}

\begin{ansblock}[2章-导-衔接-探3]
% ans:2章-导-衔接-探3
\ansitem{3}{2√10}
\ifansbookdetail
\ansnote{详解}{由\(x-2y+1=0\)得\(x=2y-1\)，代入\(y^{2}=2x\)得\(y^{2}-4y+2=0\)，\(\Delta=16-8=8>0\)；由韦达定理\(y_1+y_2=4\)、\(y_1y_2=2\)，故\(|y_1-y_2|=\sqrt{(y_1+y_2)^{2}-4y_1y_2}=2\sqrt{2}\)．直线的斜率\(k=1/2\)，所以\(|AB|=\sqrt{1+1/k^{2}}\,|y_1-y_2|=\sqrt{5}\times 2\sqrt{2}=2\sqrt{10}\)．}
\fi
\end{ansblock}

\begin{ansblock}[2章-导-衔接-探4]
% ans:2章-导-衔接-探4
\ansitem{4}{12}
\ifansbookdetail
\ansnote{详解}{消元\(y\)：将\(y=x+2\)代入\(x^{2}/m+y^{2}/4=1\)，整理得\((4+m)x^{2}+4mx=0\)，一根\(x_A=0\)（直线恰过椭圆的上顶点\((0,2)\)），另一根\(x_B=-4m/(4+m)\)．\(|AB|=\sqrt{1+k^{2}}\,|x_B-x_A|=\sqrt{2}\cdot|4m/(4+m)|=3\sqrt{2}\)，又\(m>0\)，故\(4m/(4+m)=3\)，解得\(m=12\)．回验：\(m=12\)时\(B(-3,-1)\)，\(9/12+1/4=1\)，点\(B\)在椭圆上，符合．}
\fi
\end{ansblock}

\begin{ansblock}[2章-导-衔接-探5]
% ans:2章-导-衔接-探5
\ansitem{5}{4x−y−7=0}
\ifansbookdetail
\ansnote{详解}{点差法：设弦端点\(P_1(x_1,y_1)\)、\(P_2(x_2,y_2)\)，则\(x_1+x_2=4\)、\(y_1+y_2=2\)．两点坐标代入方程相减：\(2(x_1+x_2)(x_1-x_2)-(y_1+y_2)(y_1-y_2)=0\)，故\(k=(y_1-y_2)/(x_1-x_2)=2(x_1+x_2)/(y_1+y_2)=4\)，直线为\(y-1=4(x-2)\)，即\(4x-y-7=0\)．回验\(\Delta\)：联立消\(y\)得\(14x^{2}-56x+51=0\)，\(\Delta=56^{2}-4\times 14\times 51=280>0\)，方为真弦．}
\fi
\end{ansblock}

\begin{ansblock}[2章-导-衔接-探6]
% ans:2章-导-衔接-探6
\ansitem{6}{C}
\ifansbookdetail
\ansnote{详解}{将\(l\colon y=kx+2\)代入\(x^{2}/2+y^{2}=1\)：\((1+2k^{2})x^{2}+8kx+6=0\)，则\(x_1+x_2=-8k/(1+2k^{2})\)、\(x_1x_2=6/(1+2k^{2})\)；\(y_1y_2=k^{2}x_1x_2+2k(x_1+x_2)+4=(4-2k^{2})/(1+2k^{2})\)．\(\angle AOB=90^{\circ}\iff\overrightarrow{OA}\cdot\overrightarrow{OB}=0\iff x_1x_2+y_1y_2=(6+4-2k^{2})/(1+2k^{2})=0\)，解得\(k^{2}=5\)，又\(k>0\)，故\(k=\sqrt{5}\)，选C．}
\fi
\end{ansblock}

\begin{ansblock}[2章-导-衔接-G1]
% ans:2章-导-衔接-G1
\ansitem{7}{A}
\ifansbookdetail
\ansnote{详解}{直线恒过定点\((1,2)\)．由\(1^{2}=1<4\times 2=8\)，点\((1,2)\)在抛物线\(x^{2}=4y\)的内部（\(x^{2}<4y\)一侧）．过抛物线内部点的任意直线必与抛物线相交、不可能相离，故位置关系为相交，选A．}
\fi
\end{ansblock}

\begin{ansblock}[2章-导-衔接-G2]
% ans:2章-导-衔接-G2
\ansitem{8}{B}
\ifansbookdetail
\ansnote{详解}{联立消\(y\)：\(x^{2}/4-((3/2)x+2)^{2}/9=1\)，同乘36：\(9x^{2}-4(9x^{2}/4+6x+4)=36\)，得\(-24x-16=36\)，解得\(x=-13/6\)，唯一解，直线与双曲线恰有一个公共点．而该直线与渐近线\(y=(3/2)x\)平行且不重合，双曲线的切线不可能平行于渐近线，故"恰一个公共点"是平行于渐近线的截断而非相切，判定为\textbf{相交}，选B．（渐近线陷阱：恰一交点\(\neq\)相切．）}
\fi
\end{ansblock}

\begin{ansblock}[2章-导-衔接-G3]
% ans:2章-导-衔接-G3
\ansitem{9}{D}
\ifansbookdetail
\ansnote{详解}{无交点即\(d(O,l)=4/\sqrt{m^{2}+n^{2}}>2\)，得\(m^{2}+n^{2}<4\)，点\(P(m,n)\)在圆\(x^{2}+y^{2}=4\)内．该圆与椭圆\(x^{2}/9+y^{2}/4=1\)内切（\(b=2\)、\(a=3\)），圆内每一点都在椭圆内部，故\(P\)是椭圆内点，过\(P\)的任一直线与椭圆恰有2个公共点，选D．}
\fi
\end{ansblock}

\begin{ansblock}[2章-导-衔接-G4]
% ans:2章-导-衔接-G4
\ansitem{10}{(4,2)}
\ifansbookdetail
\ansnote{详解}{将\(x=y+2\)代入\(y^{2}=4x\)：\(y^{2}-4y-8=0\)，则\(y_1+y_2=4\)，中点纵坐标\(y_0=2\)，横坐标\(x_0=y_0+2=4\)．中点坐标为\((4,2)\)．}
\fi
\end{ansblock}

\begin{ansblock}[2章-导-衔接-G5]
% ans:2章-导-衔接-G5
\ansitem{11}{3}
\ifansbookdetail
\ansnote{详解}{\(a=2\)、\(b=\sqrt{3}\)、\(c=1\)．通径为过焦点且垂直于长轴的弦：令\(x=1\)得\(y^{2}=3(1-1/4)=9/4\)，\(y=\pm 3/2\)，故通径长\(=2b^{2}/a=3\)．}
\fi
\end{ansblock}

\begin{ansblock}[2章-练-衔接-1]
% ans:2章-练-衔接-1
\ansitem{12}{(1) 3；(2) √2；(3) 3+2√2；(4) 1/2}
\ifansbookdetail
\ansnote{详解}{韦达定理：\(x_1+x_2=2\)、\(x_1x_2=1/2\)．(1) \(x_1^{2}+x_2^{2}=(x_1+x_2)^{2}-2x_1x_2=4-1=3\)．\par\noindent (2) \(|x_1-x_2|=\sqrt{(x_1+x_2)^{2}-4x_1x_2}=\sqrt{2}\)．\par\noindent (3) 由求根公式定序：\(x_1=(2+\sqrt{2})/2\)、\(x_2=(2-\sqrt{2})/2\)，故\(x_1/x_2=(2+\sqrt{2})/(2-\sqrt{2})=(2+\sqrt{2})^{2}/2=3+2\sqrt{2}\)．\par\noindent (4) 由\(x_1x_2=1/2\)得\(1/x_1=2x_2\)；又\(x_2\)是原方程的根，\(x_2^{2}=2x_2-1/2\)，故\(1/x_1-x_2^{2}=2x_2-(2x_2-1/2)=1/2\)．}
\fi
\end{ansblock}

\begin{ansblock}[2章-练-衔接-2]
% ans:2章-练-衔接-2
\ansitem{13}{21}
\ifansbookdetail
\ansnote{详解}{由方程得\(x^{2}=x+1\)，即\(x_1^{2}=x_1+1\)、\(x_2^{2}=x_2+1\)，且\(x_1+x_2=1\)．\(2x_1^{5}=2x_1(x_1^{2})^{2}=2x_1(x_1+1)^{2}=2x_1(3x_1+2)=6x_1^{2}+4x_1=10x_1+6\)；\(5x_2^{3}=5x_2(x_2+1)=5x_2^{2}+5x_2=10x_2+5\)．合计\(10(x_1+x_2)+11=10+11=21\)．}
\fi
\end{ansblock}

\begin{ansblock}[2章-练-衔接-3]
% ans:2章-练-衔接-3
\ansitem{14}{另一根 2−√3；c=1}
\ifansbookdetail
\ansnote{详解}{两根之和为4，故另一根\(x_2=4-(2+\sqrt{3})=2-\sqrt{3}\)；两根之积为\(c\)，故\(c=(2+\sqrt{3})(2-\sqrt{3})=4-3=1\)．双路验算：把\(2+\sqrt{3}\)代入方程得\(7+4\sqrt{3}-8-4\sqrt{3}+c=0\)，亦得\(c=1\)．}
\fi
\end{ansblock}

\begin{ansblock}[2章-练-衔接-4]
% ans:2章-练-衔接-4
\ansitem{15}{(1) 1；(2) k=2/(c−1)，k>0 或 k<−2}
\ifansbookdetail
\ansnote{详解}{(1) 两根之积\(=-ck/k=-c\)，故另一根\((-c)/(-c)=1\)．(2) 两根之和\(1-c=-2/k\)，故\(k=2/(c-1)\)．由\(c>0\)且\(c\neq 1\)：\(c>1\)时\(c-1>0\)，\(k>0\)；\(0<c<1\)时\(-1<c-1<0\)，\(k=2/(c-1)<2/(-1)=-2\)，即\(k<-2\)．故\(k>0\)或\(k<-2\)．（另"两不等实根"前提恒相容：\(\Delta=4+4ck^{2}>0\)恒成立．）}
\fi
\end{ansblock}

\begin{ansblock}[2章-练-衔接-5]
% ans:2章-练-衔接-5
\ansitem{16}{(1) 证毕（Δ=4m²+9>0）；(2) m=8}
\ifansbookdetail
\ansnote{详解}{(1) \(\Delta=(2m+1)^{2}-4(m-2)=4m^{2}+4m+1-4m+8=4m^{2}+9\geqslant 9>0\)，故无论\(m\)取何值，方程总有两个不相等的实数根．\par\noindent (2) 韦达定理：\(x_1+x_2=-(2m+1)\)、\(x_1x_2=m-2\)，代入\(x_1+x_2+3x_1x_2=1\)得\(-(2m+1)+3(m-2)=1\)，即\(m-7=1\)，解得\(m=8\)．}
\fi
\end{ansblock}

\begin{ansblock}[2章-练-衔接-6]
% ans:2章-练-衔接-6
\ansitem{17}{(1) 不存在；(2) k=−2，−3，−5}
\ifansbookdetail
\ansnote{详解}{两实根须\(4k\neq 0\)且\(\Delta=(-4k)^{2}-4\cdot 4k\cdot(k+1)=-16k\geqslant 0\)，故\(k<0\)；韦达定理：\(x_1+x_2=1\)、\(x_1x_2=(k+1)/(4k)\)．\par\noindent (1) \((2x_1-x_2)(x_1-2x_2)=2x_1^{2}+2x_2^{2}-5x_1x_2=2[(x_1+x_2)^{2}-2x_1x_2]-5x_1x_2=2-9x_1x_2=2-9(k+1)/(4k)\)．令其等于\(-3/2\)：\(18(k+1)=28k\)，得\(10k=18\)，\(k=9/5>0\)，与实根前提\(k<0\)矛盾，故不存在．\par\noindent (2) \(x_1/x_2+x_2/x_1-2=(x_1+x_2)^{2}/(x_1x_2)-4=4k/(k+1)-4=-4/(k+1)\)．\(k\)为负整数且\(k\neq -1\)，需\(k+1\)整除4，\(k+1\in\{-1,-2,-4\}\)，即\(k\in\{-2,-3,-5\}\)．}
\fi
\end{ansblock}

\begin{ansblock}[2章-练-衔接-7]
% ans:2章-练-衔接-7
\ansitem{18}{30}
\ifansbookdetail
\ansnote{详解}{设两直角边为\(a\)、\(b\)：\(a+b=2m+7\)、\(ab=4m(m-2)\)．勾股定理：\(169=a^{2}+b^{2}=(a+b)^{2}-2ab=(2m+7)^{2}-8m(m-2)\)，展开得\(4m^{2}+28m+49-8m^{2}+16m=169\)，整理得\(m^{2}-11m+30=0\)，解得\(m=5\)或\(m=6\)．\par\noindent 实根检验：\(\Delta=(2m+7)^{2}-16m(m-2)\)，\(m=6\)时\(\Delta=361-384=-23<0\)，舍；\(m=5\)时\(\Delta=289-240=49>0\)，且\(a+b=17>0\)、\(ab=60>0\)均正，符合．故\(S_{\triangle ABC}=ab/2=30\)．}
\fi
\end{ansblock}

\begin{ansblock}[2章-练-衔接-8]
% ans:2章-练-衔接-8
\ansitem{19}{(1) x=0,y=−1 或 x=1,y=2；(2) x=1,y=2 或 x=−1,y=0；(3) x=2,y=1}
\ifansbookdetail
\ansnote{详解}{(1) 由\(3x-y=1\)得\(y=3x-1\)，代入得\(9x^{2}-6x+1=3x+1\)，即\(9x^{2}-9x=0\)，解得\(x=0\)或\(x=1\)，回代得\(y=-1\)或\(y=2\)．\par\noindent (2) 整体设元\(t=x+y\)：\(t^{2}-2t-3=0\)，得\(t=3\)或\(t=-1\)；配\(x-y=-1\)：\(t=3\)时\(x=1\)、\(y=2\)；\(t=-1\)时\(x=-1\)、\(y=0\)．\par\noindent (3) 平方差：\(9x^{2}-4y^{2}=(3x-2y)(3x+2y)=32\)，而\(3x-2y=4\)，故\(3x+2y=8\)；与\(3x-2y=4\)联立得\(x=2\)、\(y=1\)．（注意：消元所得根必须回代求另一未知数，防丢解．）}
\fi
\end{ansblock}

\begin{ansblock}[2章-练-衔接-9]
% ans:2章-练-衔接-9
\ansitem{20}{x=−2,y=2 或 x=2,y=−2}
\ifansbookdetail
\ansnote{详解}{\(xy\neq 0\)，第一式同乘\(xy\)：\(x+y=0\)，即\(x=-y\)．代入\(xy=-4\)：\(-y^{2}=-4\)，\(y=\pm 2\)，故\((x,y)=(-2,2)\)或\((2,-2)\)．}
\fi
\end{ansblock}

\begin{ansblock}[2章-练-衔接-10]
% ans:2章-练-衔接-10
\ansitem{21}{x=(25+√601)/2,\allowbreak y=(−25+√601)/2 或 x=(25−√601)/2,\allowbreak y=(−25−√601)/2}
\ifansbookdetail
\ansnote{详解}{设\(A=\sqrt{x-y}\geqslant 0\)：\(A^{2}-3A-10=0\)，\((A-5)(A+2)=0\)，得\(A=5\)（负根舍），故\(x-y=25\)．配\(xy=-6\)：\(x(x-25)=-6\)，即\(x^{2}-25x+6=0\)，解得\(x=(25\pm\sqrt{625-24})/2=(25\pm\sqrt{601})/2\)；\(y=x-25=(-25\pm\sqrt{601})/2\)（\(\pm\)号对应取同号）．}
\fi
\end{ansblock}

\begin{ansblock}[2章-练-衔接-11]
% ans:2章-练-衔接-11
\ansitem{22}{x=2,y=3 或 x=3,y=2}
\ifansbookdetail
\ansnote{详解}{法一（代入）：\(y=5-x\)代入\(xy=6\)得\(x^{2}-5x+6=0\)，\(x=2\)或\(x=3\)，对应\(y=3\)或\(y=2\)．法二（对称构造）：\(x\)、\(y\)是\(t^{2}-5t+6=0\)的两根，同解．}
\fi
\end{ansblock}

\begin{ansblock}[2章-练-衔接-12]
% ans:2章-练-衔接-12
\ansitem{23}{−√5<m<√5}
\ifansbookdetail
\ansnote{详解}{把\(y=x+m\)代入\(x^{2}/4+y^{2}=1\)：\(x^{2}/4+(x+m)^{2}=1\)，整理得\(5x^{2}+8mx+4m^{2}-4=0\)．方程组有两解\(\iff\)该二次方程有两不等实根\(\iff\Delta=64m^{2}-20(4m^{2}-4)=80-16m^{2}>0\)，即\(m^{2}<5\)，故\(-\sqrt{5}<m<\sqrt{5}\)．}
\fi
\end{ansblock}

\begin{ansblock}[2章-练-衔接-13]
% ans:2章-练-衔接-13
\ansitem{24}{√5（此时 x=−4√5/5、y=√5/5）}
\ifansbookdetail
\ansnote{详解}{设\(m=y-x\)，即\(y=x+m\)，代入\(x^{2}/4+y^{2}=1\)：\(5x^{2}+8mx+4(m^{2}-1)=0\)．实数解存在\(\iff\Delta=64m^{2}-80(m^{2}-1)\geqslant 0\)，即\(m^{2}\leqslant 5\)，故\(m\)的最大值为\(\sqrt{5}\)．取等时关于\(x\)的方程为\(5x^{2}+8\sqrt{5}x+16=0\)，解得\(x=-8\sqrt{5}/10=-4\sqrt{5}/5\)，\(y=x+\sqrt{5}=\sqrt{5}/5\)，点在椭圆上，符合．}
\fi
\end{ansblock}

\begin{ansblock}[2章-练-衔接-14]
% ans:2章-练-衔接-14
\ansitem{25}{B（x²−y²=9）}
\ifansbookdetail
\ansnote{详解}{设双曲线方程为\(x^{2}-y^{2}=\lambda\)（\(\lambda>0\)，焦点在\(x\)轴的等轴双曲线）．与\(y=x/2\)联立：\((3/4)x^{2}=\lambda\)，\(x=\pm 2\sqrt{\lambda/3}\)，\(|\Delta x|=4\sqrt{\lambda/3}\)；\(|AB|=\sqrt{1+k^{2}}\,|\Delta x|=(\sqrt{5}/2)\cdot 4\sqrt{\lambda/3}=2\sqrt{5\lambda/3}=2\sqrt{15}\)，得\(5\lambda/3=15\)，\(\lambda=9\)，方程为\(x^{2}-y^{2}=9\)，选B．}
\fi
\end{ansblock}

\begin{ansblock}[2章-练-衔接-15]
% ans:2章-练-衔接-15
\ansitem{26}{D（2x−y−1=0）}
\ifansbookdetail
\ansnote{详解}{设切线为\(2x-y+m=0\)，即\(y=2x+m\)，代入\(y=x^{2}\)：\(x^{2}-2x-m=0\)．相切\(\iff\Delta=4+4m=0\)，得\(m=-1\)，切线为\(2x-y-1=0\)（切点\((1,1)\)），选D．}
\fi
\end{ansblock}

\begin{ansblock}[2章-练-衔接-16]
% ans:2章-练-衔接-16
\ansitem{27}{选①②：2<t<3；选①③：1<t<2}
\ifansbookdetail
\ansnote{详解}{条件①表示椭圆\(\iff\)两分母同正且不等：\(3-t>0\)、\(t-1>0\)、\(3-t\neq t-1\)，即\(1<t<3\)且\(t\neq 2\)．选①②（长轴在\(y\)轴）：\(t-1>3-t>0\)，解得\(2<t<3\)；选①③（长轴在\(x\)轴）：\(3-t>t-1>0\)，解得\(1<t<2\)．}
\fi
\end{ansblock}

\jietitle{2.1 坐标法}

\begin{ansblock}[2章-练-课时01-E4]
% ans:2章-练-课时01-E4
\ansitem{1}{|AB|＝5；中点坐标为(3/2,−2)．}
\ifansbookdetail
\ansnote{详解}{\(|AB|=\sqrt{(0-3)^{2}+(-4-0)^{2}}=\sqrt{9+16}=5\)；中点坐标为\(((3+0)/2,(0-4)/2)=(3/2,-2)\)．}
\fi
\end{ansblock}

\begin{ansblock}[2章-练-课时01-E1]
% ans:2章-练-课时01-E1
\ansitem{2}{D}
\ifansbookdetail
\ansnote{详解}{四边形\(MNPQ\)的对角线为\(MP\)与\(NQ\)．\(MP\)的中点为\(((1+4)/2,(1+0)/2)=(5/2,1/2)\)，\(NQ\)的中点为\(((3+2)/2,(-1+2)/2)=(5/2,1/2)\)，两对角线互相平分，故\(MNPQ\)是平行四边形．又 \(|MP|^{2}=(4-1)^{2}+(0-1)^{2}=10\)，\(|NQ|^{2}=(2-3)^{2}+(2+1)^{2}=10\)，对角线相等，故\pxparallelogram\(MNPQ\)为矩形．（若担心为正方形：\(|MN|^{2}=(3-1)^{2}+(-1-1)^{2}=8\)，\(|NP|^{2}=(4-3)^{2}+(0+1)^{2}=2\)，邻边不等，非正方形．）故选：D．}
\fi
\end{ansblock}

\begin{ansblock}[2章-练-课时01-E5]
% ans:2章-练-课时01-E5
\ansitem{3}{C的坐标为5．}
\ifansbookdetail
\ansnote{详解}{\(B\)是线段\(AC\)的中点，设\(C(x)\)，则 \((-1+x)/2=2\)，解得\(x=5\)，即\(C\)的坐标为5．}
\fi
\end{ansblock}

\begin{ansblock}[2章-练-课时01-E6]
% ans:2章-练-课时01-E6
\ansitem{4}{P的坐标为3或−5．}
\ifansbookdetail
\ansnote{详解}{设\(P(x)\)，依题意 \(|x+9|=2|x+3|\)．两边均非负，平方得 \((x+9)^{2}=4(x+3)^{2}\)，即 \(x^{2}+18x+81=4x^{2}+24x+36\)，整理得 \(3x^{2}+6x-45=0\)，即 \(x^{2}+2x-15=0\)，解得\(x=3\)或\(x=-5\)．检验：\(x=3\)时\(|3+9|=12=2|3+3|\)；\(x=-5\)时\(|-5+9|=4=2|-5+3|\)．所以点\(P\)的坐标为3或\(-5\)．}
\fi
\end{ansblock}

\begin{ansblock}[2章-练-课时01-E9]
% ans:2章-练-课时01-E9
\ansitem{5}{AB边上的中线长3；BC边上的中线长3；CA边上的中线长3√2．}
\ifansbookdetail
\ansnote{详解}{设\(AB\)、\(BC\)、\(CA\)的中点分别为\(M\)、\(N\)、\(P\)．\(M((2-2)/2,(1+3)/2)=M(0,2)\)，中线\(|CM|=\sqrt{(0-0)^{2}+(2+1)^{2}}=3\)；\(N((-2+0)/2,(3-1)/2)=N(-1,1)\)，中线\(|AN|=\sqrt{(-1-2)^{2}+(1-1)^{2}}=3\)；\(P((0+2)/2,(-1+1)/2)=P(1,0)\)，中线\(|BP|=\sqrt{(1+2)^{2}+(0-3)^{2}}=\sqrt{18}=3\sqrt{2}\)．所以三条中线的长分别为3、3、\(3\sqrt{2}\)．}
\fi
\end{ansblock}

\begin{ansblock}[2章-练-课时01-E8]
% ans:2章-练-课时01-E8
\ansitem{6}{证明见详解．}
\ifansbookdetail
\ansnote{详解}{由两点间距离公式：\par\noindent\(|AB|=\sqrt{(3+1)^{2}+(3-3)^{2}}=\sqrt{16}=4\)；\par\noindent\(|AC|=\sqrt{(1+1)^{2}+(2\sqrt{3}+3-3)^{2}}=\sqrt{4+12}=4\)；\par\noindent\(|BC|=\sqrt{(1-3)^{2}+(2\sqrt{3}+3-3)^{2}}=\sqrt{4+12}=4\)．\par\noindent 所以\(|AB|=|AC|=|BC|\)，故\(\triangle ABC\)是等边三角形．}
\fi
\end{ansblock}

\begin{ansblock}[2章-练-课时01-E7]
% ans:2章-练-课时01-E7
\ansitem{7}{2}
\ifansbookdetail
\ansnote{详解}{曲线方程中\(x=0\)时\(2=0\)不成立，故\(x\neq 0\)，且 \(y=1+2/x\)．设曲线上任意一点\(A(x,1+2/x)\)，则\(|AP|=\sqrt{x^{2}+(1+2/x-1)^{2}}=\sqrt{x^{2}+4/x^{2}}\)．由基本不等式 \(x^{2}+4/x^{2}\geqslant 2\sqrt{x^{2}\cdot 4/x^{2}}=4\)，当且仅当 \(x^{2}=4/x^{2}\)，即\(x=\pm\sqrt{2}\) 时等号成立，此时\(|AP|\)最小值为\(\sqrt{4}=2\)．故所求距离为2．}
\fi
\end{ansblock}

\begin{ansblock}[2章-练-课时01-E2]
% ans:2章-练-课时01-E2
\ansitem{8}{BCD}
\ifansbookdetail
\ansnote{详解}{设\(P_1(x_1,y_1)\)，\(P_2(x_2,y_2)\)．A：\(d_1=d_2=1\)，则 \(ax_1+by_1+c=ax_2+by_2+c=\sqrt{a^{2}+b^{2}}\)，即 \(P_1\)、\(P_2\) 都在直线 \(ax+by+c-\sqrt{a^{2}+b^{2}}=0\) 上，该直线与\(l\)平行，故直线\(P_1P_2\parallel l\)，A正确；B：\(d_1=1\)，\(d_2=-1\) 只说明\(P_1\)、\(P_2\)在\(l\)的两侧且到\(l\)的距离相等，线段\(P_1P_2\)的中点落在\(l\)上，但\(P_1P_2\)与\(l\)不一定垂直，B错误；C：\(d_1+d_2=0\) 包含\(d_1=d_2=0\)的情形，此时\(P_1\)、\(P_2\)都在\(l\)上，直线\(P_1P_2\)与\(l\)重合，谈不上垂直，C错误；D：\(d_1\cdot d_2\leqslant 0\) 表示\(P_1\)、\(P_2\)分居\(l\)两侧或至少一点在\(l\)上，此时直线\(P_1P_2\)与\(l\)相交或重合，D把“重合”漏掉，D错误．故选：BCD．}
\fi
\end{ansblock}

\begin{ansblock}[2章-练-课时01-E3]
% ans:2章-练-课时01-E3
\ansitem{9}{ABD}
\ifansbookdetail
\ansnote{详解}{A：\((|x_1-x_2|+|y_1-y_2|)^{2}=(x_1-x_2)^{2}+2|x_1-x_2||y_1-y_2|+(y_1-y_2)^{2}\geqslant(x_1-x_2)^{2}+(y_1-y_2)^{2}\)，开方即得 \(d(A,B)\geqslant|AB|\)，A正确；B：由绝对值三角不等式，\(|x_1-x_2|+|x_2-x_3|\geqslant|x_1-x_3|\)，\(|y_1-y_2|+|y_2-y_3|\geqslant|y_1-y_3|\)，两式相加即得 \(d(A,B)+d(B,C)\geqslant d(A,C)\)（将\(A\)、\(B\)、\(C\)重排即题设形式），B正确；C：设\(Q(x,2x-1)\)是\(l\)上任一点，\(d(P,Q)=|x-3|+|2x-2|=|x-3|+|x-1|+|x-1|\geqslant|(x-3)-(x-1)|+0=2\)，当且仅当\(x=1\)时等号成立，故 \(d(P,l)=2\neq 4\sqrt{5}/5\)，C错误；D：\(d(P,O)=|x|+|y|=1\) 的图象是以\((1,0)\)、\((0,1)\)、\((-1,0)\)、\((0,-1)\)为顶点的正方形，对角线长均为2，面积\(=(1/2)\times 2\times 2=2\)，D正确．故选：ABD．}
\fi
\end{ansblock}

\begin{ansblock}[2章-练-课时01-E11]
% ans:2章-练-课时01-E11
\ansitem{10}{证明见详解．}
\ifansbookdetail
\ansnote{详解}{以长方形的顶点\(A\)为原点、\(AB\)所在直线为\(x\)轴建立平面直角坐标系．设\(AB=a\)，\(AD=b\)，则\(A(0,0)\)，\(B(a,0)\)，\(C(a,b)\)，\(D(0,b)\)．再设\(M(x,y)\)．\(|AM|^{2}+|CM|^{2}=x^{2}+y^{2}+(x-a)^{2}+(y-b)^{2}\)；\(|BM|^{2}+|DM|^{2}=(x-a)^{2}+y^{2}+x^{2}+(y-b)^{2}\)．两式展开后均为 \(x^{2}+y^{2}+(x-a)^{2}+(y-b)^{2}\)，恒相等．所以 \(|AM|^{2}+|CM|^{2}=|BM|^{2}+|DM|^{2}\)，即对平面上任意一点\(M\)结论都成立．}
\fi
\end{ansblock}

\begin{ansblock}[2章-练-课时01-E14]
% ans:2章-练-课时01-E14
\ansitem{11}{证明见详解．}
\ifansbookdetail
\ansnote{详解}{以平行四边形相邻两边为坐标轴方向建立平面直角坐标系：取顶点\(A\)为原点、\(AB\)所在直线为\(x\)轴．设\(B(a,0)\)，\(D(b,c)\)，则由平行四边形性质知 \(C(a+b,c)\)．两条对角线：\(|AC|^{2}=(a+b)^{2}+c^{2}\)，\(|BD|^{2}=(b-a)^{2}+c^{2}\)，所以 \(|AC|^{2}+|BD|^{2}=\left(a^{2}+2ab+b^{2}+c^{2}\right)+\left(a^{2}-2ab+b^{2}+c^{2}\right)=2a^{2}+2b^{2}+2c^{2}\)．四条边：\(|AB|^{2}=a^{2}\)，\(|CD|^{2}=a^{2}\)，\(|AD|^{2}=b^{2}+c^{2}\)，\(|BC|^{2}=b^{2}+c^{2}\)，所以 \(|AB|^{2}+|BC|^{2}+|CD|^{2}+|DA|^{2}=2a^{2}+2b^{2}+2c^{2}\)．两式相同，故平行四边形两条对角线的平方和等于四条边的平方和．}
\fi
\end{ansblock}

\begin{ansblock}[2章-练-课时01-E13]
% ans:2章-练-课时01-E13
\ansitem{12}{证明见详解．}
\ifansbookdetail
\ansnote{详解}{以\(A\)为原点、\(AB\)所在直线为\(x\)轴建立平面直角坐标系，则\(A(0,0)\)，\(B(4,0)\)，\(C(4,4)\)，\(D(0,4)\)；由\(E\)、\(F\)分别是\(BC\)、\(CD\)的中点得 \(E(4,2)\)，\(F(2,4)\)．于是 \(\overrightarrow{AE}=(4,2)\)，\(\overrightarrow{BF}=(2-4,4-0)=(-2,4)\)．\(\overrightarrow{AE}\cdot\overrightarrow{BF}=4\times(-2)+2\times 4=-8+8=0\)，所以 \(\overrightarrow{AE}\perp\overrightarrow{BF}\)，即 \(BF\perp AE\)．}
\fi
\end{ansblock}

\begin{ansblock}[2章-练-课时01-E10]
% ans:2章-练-课时01-E10
\ansitem{13}{P(0,4)或P(0,−1)．}
\ifansbookdetail
\ansnote{详解}{设\(P(0,y)\)．\(PA\perp PB\) 即\(\triangle PAB\)是以\(AB\)为斜边的直角三角形，\(|PA|^{2}+|PB|^{2}=|AB|^{2}\)．\(|PA|^{2}=9+(y-1)^{2}\)，\(|PB|^{2}=4+(y-2)^{2}\)，\(|AB|^{2}=(-2-3)^{2}+(2-1)^{2}=26\)，代入：\(9+y^{2}-2y+1+4+y^{2}-4y+4=26\)，整理得 \(2y^{2}-6y-8=0\)，即 \(y^{2}-3y-4=0\)，解得\(y=4\)或\(y=-1\)．所以点\(P\)的坐标为\((0,4)\)或\((0,-1)\)．}
\fi
\end{ansblock}

\begin{ansblock}[2章-练-课时01-E12]
% ans:2章-练-课时01-E12
\ansitem{14}{f(x)＝−(x−4)²−1．}
\ifansbookdetail
\ansnote{详解}{在\(y=f(x)\)的图象上任取一点\(P(x,y)\)，它关于点\(M(2,0)\)的对称点为 \(P'(4-x,-y)\)．由对称性，\(P'\)在函数\(y=x^{2}+1\)的图象上，所以 \(-y=(4-x)^{2}+1\)，即 \(y=-(x-4)^{2}-1\)．故 \(f(x)=-(x-4)^{2}-1\)．}
\fi
\end{ansblock}

\begin{ansblock}[2章-练-课时01-E15]
% ans:2章-练-课时01-E15
\ansitem{15}{B}
\ifansbookdetail
\ansnote{详解}{由中点坐标公式，\(M((b+0)/2,(0+a)/2)\)，即 \(M(b/2,a/2)\)。由两点间距离公式：\(|CM|=\sqrt{(b/2)^{2}+(a/2)^{2}}=\sqrt{a^{2}+b^{2}}/2\)；\(|AB|=\sqrt{(b-0)^{2}+(0-a)^{2}}=\sqrt{a^{2}+b^{2}}\)。故 \(|CM|=|AB|/2\)，选 B。这就用坐标法证明了：直角三角形斜边上的中线等于斜边的一半。}
\fi
\end{ansblock}

\begin{ansblock}[2章-练-课时01-E16]
% ans:2章-练-课时01-E16
\ansitem{16}{B}
\ifansbookdetail
\ansnote{详解}{由中点坐标公式，\(M(b/2,c/2)\)，\(N((a+b)/2,c/2)\)。两点纵坐标相等且 \(|MN|=a/2>0\)（\(M\neq N\)），故 \(MN\parallel x\) 轴；而 \(BC\) 在 \(x\) 轴上，所以 \(MN\parallel BC\)。又 \(|MN|=|(a+b)/2-b/2|=a/2\)，\(|BC|=a\)，故 \(|MN|=|BC|/2\)。选 B。这就用坐标法证明了：三角形中位线平行于第三边，且等于第三边的一半。}
\fi
\end{ansblock}

\begin{ansblock}[2章-拓-课时01-T1]
% ans:2章-拓-课时01-T1
\ansitem{17}{(1)见详解；(2) 等号成立$\iff$C的坐标x∈[−3,2]；等号不成立$\iff$x<−3或x>2．}
\ifansbookdetail
\ansnote{详解}{(1) 设\(C\)的坐标为\(x_3\)．因为 \(x_1-x_2=(x_1-x_3)+(x_3-x_2)\)，由绝对值三角不等式\(d(A,B)=|x_1-x_2|=|(x_1-x_3)+(x_3-x_2)|\leqslant|x_1-x_3|+|x_3-x_2|=d(A,C)+d(B,C)\)，得证．(2) \(|x_1-x_2|=|(x_1-x_3)+(x_3-x_2)|\)取等号\(\iff(x_1-x_3)(x_3-x_2)\geqslant 0\)，即\(x_3\)介于\(x_1\)与\(x_2\)之间（含端点）．此处\(x_1=-3\)，\(x_2=2\)，故等号成立\(\iff -3\leqslant x\leqslant 2\)；等号不成立（即\(d(A,B)<d(A,C)+d(B,C)\)）\(\iff x<-3\)或\(x>2\)．}
\fi
\end{ansblock}

\begin{ansblock}[2章-拓-课时01-T2]
% ans:2章-拓-课时01-T2
\ansitem{18}{(1)见详解；(2) 等号成立$\iff$C在以AB为对角线的矩形内（含边界），即$(x_C-x_1)(x_C-x_2)\leqslant 0$且$(y_C-y_1)(y_C-y_2)\leqslant 0$；等号不成立$\iff$C在该矩形外．}
\ifansbookdetail
\ansnote{详解}{(1) 设\(C(x_3,y_3)\)．\(x_1-x_2=(x_1-x_3)+(x_3-x_2)\)，\(y_1-y_2=(y_1-y_3)+(y_3-y_2)\)，两次用绝对值三角不等式：\(d(A,B)=|x_1-x_2|+|y_1-y_2|\leqslant(|x_1-x_3|+|x_3-x_2|)+(|y_1-y_3|+|y_3-y_2|)=d(A,C)+d(B,C)\)，得证．(2) 横、纵两个方向分别取等：第一个不等式取等\(\iff(x_1-x_3)(x_3-x_2)\geqslant 0\)，第二个取等\(\iff(y_1-y_3)(y_3-y_2)\geqslant 0\)，故总等号成立\(\iff x_3\)介于\(x_1\)、\(x_2\)之间且\(y_3\)介于\(y_1\)、\(y_2\)之间，即\(C\)落在以\(AB\)为对角线（边平行于坐标轴）的矩形内部或边界上；只要有一个方向不取等，就有\(d(A,B)<d(A,C)+d(B,C)\)，即\(C\)在该矩形外．}
\fi
\end{ansblock}

\begin{ansblock}[2章-导-课时01-G1]
% ans:2章-导-课时01-G1
\ansitem{19}{B}
\ifansbookdetail
\ansnote{详解}{设\(M(x,y)\)，由中点坐标公式得 \((x-2)/2=1\)，\((y+5)/2=0\)，解得\(x=4\)，\(y=-5\)，所以点\(M(4,-5)\)．则 \(|OM|=\sqrt{4^{2}+(-5)^{2}}=\sqrt{41}\)．故选：B．}
\fi
\end{ansblock}

\begin{ansblock}[2章-导-课时01-G2]
% ans:2章-导-课时01-G2
\ansitem{20}{A}
\ifansbookdetail
\ansnote{详解}{由两点间距离公式：\(|AB|^{2}=(2+1)^{2}+(-1-1)^{2}=9+4=13\)；\(|AC|^{2}=(1+1)^{2}+(4-1)^{2}=4+9=13\)；\(|BC|^{2}=(1-2)^{2}+(4+1)^{2}=1+25=26\)．因为\(|AB|^{2}+|AC|^{2}=13+13=26=|BC|^{2}\)，由勾股定理逆定理知\(\triangle ABC\)是以\(A\)为直角顶点的直角三角形（且\(|AB|=|AC|\)，还是等腰直角三角形）．故选：A．}
\fi
\end{ansblock}

\begin{ansblock}[2章-导-课时01-G3]
% ans:2章-导-课时01-G3
\ansitem{21}{D}
\ifansbookdetail
\ansnote{详解}{设\(C(x_0,0)\)．\(\angle ACB=\pi/2\) 即\(\triangle ACB\)是以\(AB\)为斜边的直角三角形，故 \(|CA|^{2}+|CB|^{2}=|AB|^{2}\)．\(|CA|^{2}=(x_0-4)^{2}+4\)，\(|CB|^{2}=(x_0-1)^{2}+4\)，\(|AB|^{2}=(4-1)^{2}+(2+2)^{2}=9+16=25\)，代入得 \((x_0-4)^{2}+(x_0-1)^{2}+8=25\)，展开整理得 \(x_0^{2}-5x_0=0\)，解得\(x_0=0\)或\(x_0=5\)，经检验均使\(A\)、\(B\)、\(C\)构成三角形．所以点\(C\)的坐标为\((0,0)\)或\((5,0)\)．故选：D．}
\fi
\end{ansblock}

\begin{ansblock}[2章-导-课时01-G4]
% ans:2章-导-课时01-G4
\ansitem{22}{距离为12；中点坐标为4．}
\ifansbookdetail
\ansnote{详解}{数轴上两点距离 \(d(A,B)=|10-(-2)|=12\)；中点坐标为 \((-2+10)/2=4\)．}
\fi
\end{ansblock}

\begin{ansblock}[2章-导-课时01-G5]
% ans:2章-导-课时01-G5
\ansitem{23}{a＝±3√21．}
\ifansbookdetail
\ansnote{详解}{由两点间距离公式 \(|AB|=\sqrt{a^{2}+10^{2}}=17\)，故 \(a^{2}+100=289\)，\(a^{2}=189\)，\(a=\pm\sqrt{189}=\pm 3\sqrt{21}\)．}
\fi
\end{ansblock}

\xjkeshi{2.2.1 直线的倾斜角与斜率（第1课时）}

\begin{ansblock}[2章-练-课时02-E1]
% ans:2章-练-课时02-E1
\ansitem{1}{C}
\ifansbookdetail
\ansnote{详解}{\(O\)、\(O_3\) 都为五角星的中心点，故 \(OO_3\) 平分第三颗小星的一个角；又五角星的内角为 \(36^{\circ}\)，可知\(\angle BAO_3=18^{\circ}\)．过 \(O_3\) 作 \(x\) 轴的平行线 \(O_3E\)，则\(\angle OO_3E=\alpha\approx 16^{\circ}\)，所以直线 \(AB\) 的倾斜角约为 \(18^{\circ}-16^{\circ}=2^{\circ}\)，选C．}
\fi
\end{ansblock}

\begin{ansblock}[2章-练-课时02-E9]
% ans:2章-练-课时02-E9
\ansitem{2}{B}
\ifansbookdetail
\ansnote{详解}{\(\sin 210^{\circ}=-1/2\)，直线方程即\(-x-y-2=0\)，即\(y=-x-2\)，斜率\(k=2\sin 210^{\circ}=-1\)，故倾斜角为 \(135^{\circ}\)，选B．}
\fi
\end{ansblock}

\begin{ansblock}[2章-练-课时02-E15]
% ans:2章-练-课时02-E15
\ansitem{3}{120°}
\ifansbookdetail
\ansnote{详解}{\(l_2\) 与 \(l_1\) 垂直，则 \(l_2\) 的向上方向与 \(x\) 轴正向所成的角为\(30^{\circ}+90^{\circ}=120^{\circ}\)（在 \(l_2\) 向上方向取与 \(l_1\) 垂直且指向左上的一侧，另一直线方向与之共线）．由倾斜角的定义及\(0^{\circ}\leqslant\alpha<180^{\circ}\)，得 \(l_2\) 的倾斜角为\(120^{\circ}\)．（说明：两直线垂直时倾斜角关系的斜率判据在§2.2.3学习，本题只需用平面几何角的运算作答．）}
\fi
\end{ansblock}

\begin{ansblock}[2章-练-课时02-E16]
% ans:2章-练-课时02-E16
\ansitem{4}{(1) 60°；(2) 150°；(3) 90°；(4) 0°．}
\ifansbookdetail
\ansnote{详解}{（1）\(k=\sqrt{3}\)，\(\tan\alpha=\sqrt{3}\)，\(\alpha=60^{\circ}\)；（2）化为\(y=-(\sqrt{3}/3)x-\sqrt{3}/3\)，\(k=-\sqrt{3}/3\)，\(\tan\alpha=-\sqrt{3}/3\)，\(\alpha=150^{\circ}\)；（3）直线垂直于 \(x\) 轴，斜率不存在，倾斜角\(\alpha=90^{\circ}\)；（4）直线与 \(x\) 轴平行，倾斜角\(\alpha=0^{\circ}\)．}
\fi
\end{ansblock}

\begin{ansblock}[2章-练-课时02-E7]
% ans:2章-练-课时02-E7
\ansitem{5}{C}
\ifansbookdetail
\ansnote{详解}{三点共线\(\iff k_{AB}=k_{AC}\)．\(k_{AB}=(2-3)/(3+2)=-1/5\)；\(k_{AC}=(m-3)/(1/2+2)=(m-3)/(5/2)=2(m-3)/5\)．令\(2(m-3)/5=-1/5\)，得\(2(m-3)=-1\)，\(m=5/2\)，选C．}
\fi
\end{ansblock}

\begin{ansblock}[2章-练-课时02-E14]
% ans:2章-练-课时02-E14
\ansitem{6}{a＝1}
\ifansbookdetail
\ansnote{详解}{由斜率公式\(k=[(a+1)-(-1)]/(2-a)=(a+2)/(2-a)=3\)（须\(a\neq 2\)），去分母得\(a+2=3(2-a)=6-3a\)，\(4a=4\)，\(a=1\)，且\(a=1\neq 2\)满足条件．所以\(a=1\)．}
\fi
\end{ansblock}

\begin{ansblock}[2章-练-课时02-E10]
% ans:2章-练-课时02-E10
\ansitem{7}{(−∞,1)∪(1,+∞)}
\ifansbookdetail
\ansnote{详解}{三点能构成三角形\(\iff\)三点不共线．由\(k_{AB}\neq k_{AC}\)，即\((k-1)/(-2-3)\neq(1-1)/(8-3)\)，得\((k-1)/(-5)\neq 0\)，\(k-1\neq 0\)，\(k\neq 1\)．故\(k\)的取值范围为\((-\infty,1)\cup(1,+\infty)\)．}
\fi
\end{ansblock}

\begin{ansblock}[2章-练-课时02-E6]
% ans:2章-练-课时02-E6
\ansitem{8}{BCD}
\ifansbookdetail
\ansnote{详解}{倾斜角为钝角\(\iff\)斜率存在且为负，即\(A\)、\(B\)横坐标不等且\(k_{AB}<0\)．\(k_{AB}=(1+a-a)/[(1-a)-3]=1/(-2-a)<0\)，得\(a+2>0\)，即\(a>-2\)．又\(a=-2\)时\(A(3,-1)\)、\(B(3,-2)\)横坐标相同、斜率不存在，倾斜角\(90^{\circ}\)，不在"钝角"取值内且\(a>-2\)已排除．故\(a>-2\)，选项中\(0\)、\(1\)、\(2\)满足，选BCD．}
\fi
\end{ansblock}

\begin{ansblock}[2章-练-课时02-E2]
% ans:2章-练-课时02-E2
\ansitem{9}{A}
\ifansbookdetail
\ansnote{详解}{直线\(PA\)的斜率为\(k_{PA}=(-2+1)/(1-0)=-1\)，直线\(PB\)的斜率为\(k_{PB}=(1+1)/(2-0)=1\)．数形结合：当过\(P\)的直线\(l\)绕\(P\)从\(PA\)旋转到\(PB\)（保持与线段\(AB\)有公共点）时，其斜率介于\(k_{PA}\)与\(k_{PB}\)之间，故\(k\in[-1,1]\)，选A．}
\fi
\end{ansblock}

\begin{ansblock}[2章-练-课时02-E8]
% ans:2章-练-课时02-E8
\ansitem{10}{B}
\ifansbookdetail
\ansnote{详解}{\(\tan\alpha=k_{OA}=\sqrt{2}\)．\(m'\)的倾斜角为\(2\alpha\)，故\(k=\tan 2\alpha=2\tan\alpha/(1-\tan^{2}\alpha)=2\sqrt{2}/(1-2)=-2\sqrt{2}\)，选B．（合理性：\(\tan\alpha=\sqrt{2}>1\)，\(\alpha\in(\pi/4,\pi/2)\)，\(2\alpha\in(\pi/2,\pi)\)确为钝角，\(\tan 2\alpha<0\)．）}
\fi
\end{ansblock}

\begin{ansblock}[2章-练-课时02-E3]
% ans:2章-练-课时02-E3
\ansitem{11}{C}
\ifansbookdetail
\ansnote{详解}{\(\alpha\in[0^{\circ},180^{\circ})\)且\(\sin\alpha=3/5>0\)，故\(\alpha\)为锐角或钝角．当\(\alpha\)为锐角时\(\cos\alpha=4/5\)，\(k=\tan\alpha=3/4\)；当\(\alpha\)为钝角时\(\cos\alpha=-4/5\)，\(k=\tan\alpha=-3/4\)．故直线\(l\)的斜率是\(3/4\)或\(-3/4\)，选C．}
\fi
\end{ansblock}

\begin{ansblock}[2章-练-课时02-E4]
% ans:2章-练-课时02-E4
\ansitem{12}{B}
\ifansbookdetail
\ansnote{详解}{\(\alpha\in[0,\pi)\)，\(-1\leqslant k\leqslant\sqrt{3}\)即\(-1\leqslant\tan\alpha\leqslant\sqrt{3}\)．当\(0\leqslant\tan\alpha\leqslant\sqrt{3}\)时\(\alpha\in[0,\pi/3]\)；当\(-1\leqslant\tan\alpha<0\)时\(\alpha\in[3\pi/4,\pi)\)．所以\(\alpha\in[0,\pi/3]\cup[3\pi/4,\pi)\)，选B．}
\fi
\end{ansblock}

\begin{ansblock}[2章-练-课时02-E5]
% ans:2章-练-课时02-E5
\ansitem{13}{A}
\ifansbookdetail
\ansnote{详解}{\(k_l=(t^{2}-1)/(3-2)=t^{2}-1\)．由\(-\sqrt{2}\leqslant t\leqslant\sqrt{2}\)得\(t^{2}\in[0,2]\)，故\(k\in[-1,1]\)，即\(\tan\theta\in[-1,1]\)（\(\theta\)为倾斜角），得\(\theta\in[0,\pi/4]\cup[3\pi/4,\pi)\)，选A．（注：本选项组为本卷补写——源题为无选项填空，为入选择槽按原题数据配四选项，答案为\(A\)不变．）}
\fi
\end{ansblock}

\begin{ansblock}[2章-练-课时02-E11]
% ans:2章-练-课时02-E11
\ansitem{14}{[√3/3, √3)}
\ifansbookdetail
\ansnote{详解}{正切函数\(k=\tan\alpha\)在\([\pi/6,\pi/3)\)上单调递增，故\(\alpha\in[\pi/6,\pi/3)\)时\(\tan\alpha\in[\tan(\pi/6),\tan(\pi/3))=[\sqrt{3}/3,\sqrt{3})\)，即斜率\(k\in[\sqrt{3}/3,\sqrt{3})\)．}
\fi
\end{ansblock}

\begin{ansblock}[2章-练-课时02-E12]
% ans:2章-练-课时02-E12
\ansitem{15}{[1,+∞)∪(−∞,−√3]}
\ifansbookdetail
\ansnote{详解}{\(\theta\in[\pi/4,\pi/2)\)时，\(k=\tan\theta\in[\tan(\pi/4),+\infty)\)\allowbreak\(=[1,+\infty)\)\par\noindent \(\theta\in(\pi/2,2\pi/3]\)时，\(k=\tan\theta\in(-\infty,\tan(2\pi/3)]=(-\infty,-\sqrt{3}]\)\par\noindent 故\(k\)的取值范围为\([1,+\infty)\cup(-\infty,-\sqrt{3}]\)．}
\fi
\end{ansblock}

\begin{ansblock}[2章-练-课时02-E13]
% ans:2章-练-课时02-E13
\ansitem{16}{−2−√3}
\ifansbookdetail
\ansnote{详解}{\(l_1\)的斜率为\(-1\)，其倾斜角为\(135^{\circ}\)，故\(l_2\)的倾斜角为\(135^{\circ}-30^{\circ}=105^{\circ}\)．\(k_2=\tan 105^{\circ}=\tan(60^{\circ}+45^{\circ})=(\tan 60^{\circ}+\tan 45^{\circ})/(1-\tan 60^{\circ}\cdot\tan 45^{\circ})=(\sqrt{3}+1)/(1-\sqrt{3})=(\sqrt{3}+1)^{2}/[(1-\sqrt{3})(1+\sqrt{3})]=(4+2\sqrt{3})/(-2)=-2-\sqrt{3}\)．所以直线\(l_2\)的斜率为\(-2-\sqrt{3}\)．}
\fi
\end{ansblock}

\begin{ansblock}[2章-拓-课时02-T1]
% ans:2章-拓-课时02-T1
\ansitem{17}{(1) θ∈[0,π/2)时斜率由0递增并趋向＋∞；θ∈(π/2,π)时斜率由−∞递增趋向0；(2) 不能，理由见详解．}
\ifansbookdetail
\ansnote{详解}{（1）\(k=\tan\theta\)在\([0,\pi/2)\)上单调递增，\(k\)从\(0\)增大到正无穷大；在\((\pi/2,\pi)\)上也单调递增，\(k\)从负无穷大增大到\(0\)．（2）不能．例如取\(\theta_1=45^{\circ}\)，\(\theta_2=120^{\circ}\)，\(\theta_1<\theta_2\)，但\(k_1=1>k_2=-\sqrt{3}\)，即倾斜角大者斜率反而小；斜率关于倾斜角只在每一段\([0,\pi/2)\)与\((\pi/2,\pi)\)内分别递增，整体上不具有单调性．}
\fi
\end{ansblock}

\begin{ansblock}[2章-拓-课时02-T2]
% ans:2章-拓-课时02-T2
\ansitem{18}{(1) \(k_{AB}=0\)（倾斜角\(0\)）；\(k_{BC}=\sqrt{3}\)（倾斜角\(\pi/3\)）；\(k_{AC}=\sqrt{3}/3\)（倾斜角\(\pi/6\)）；(2) \([\pi/6,\ \pi/3]\)．}
\ifansbookdetail
\ansnote{详解}{（1）\(k_{AB}=(1-1)/(1+1)=0\)，\(k_{BC}=(\sqrt{3}+1-1)/(2-1)=\sqrt{3}\)，\(k_{AC}=(\sqrt{3}+1-1)/(2+1)=\sqrt{3}/3\)．因斜率等于倾斜角的正切值且倾斜角\(\alpha\in[0,\pi)\)，故\(AB\)、\(BC\)、\(AC\)的倾斜角依次为\(0\)、\(\pi/3\)、\(\pi/6\)．（2）当直线\(CD\)绕点\(C\)由\(CA\)逆时针旋转到\(CB\)时，\(CD\)与线段\(AB\)恒有交点（即\(D\)在\(AB\)上），此时\(k_{CD}\)由\(k_{AC}\)增大到\(k_{BC}\)，\(k_{CD}\in[\sqrt{3}/3,\sqrt{3}]\)，对应倾斜角\(\alpha\in[\pi/6,\pi/3]\)．}
\fi
\end{ansblock}

\begin{ansblock}[2章-拓-课时02-T3]
% ans:2章-拓-课时02-T3
\ansitem{19}{(1) 0°；(2) 90°；(3) 45°．}
\ifansbookdetail
\ansnote{详解}{（1）纵坐标同为\(c\)，直线与\(x\)轴平行，倾斜角为\(0^{\circ}\)（\(k=0\)）；（2）横坐标同为\(a\)，直线垂直于\(x\)轴，斜率不存在，倾斜角为\(90^{\circ}\)；（3）\(k=[(c+a)-(b+c)]/(a-b)=(a-b)/(a-b)=1\)（\(a\neq b\)），\(\tan\alpha=1\)，\(\alpha\in[0^{\circ},180^{\circ})\)，故倾斜角为\(45^{\circ}\)．}
\fi
\end{ansblock}

\begin{ansblock}[2章-拓-课时02-T4]
% ans:2章-拓-课时02-T4
\ansitem{20}{[−1/6, 5/3]}
\ifansbookdetail
\ansnote{详解}{\((y+1)/(x+1)=[y-(-1)]/[x-(-1)]\)的几何意义是过\(M(x,y)\)与定点\(N(-1,-1)\)两点的直线的斜率，其中点\(M\)在线段\(y=-2x+8\)（\(x\in[2,5]\)）上运动．线段两端点为\((2,4)\)与\((5,-2)\)．记\(k(x)=(y+1)/(x+1)=(-2x+9)/(x+1)=-2+11/(x+1)\)，在\(x\in[2,5]\)上随\(x\)增大而单调递减，故最大值在\(x=2\)处取得：\(k=(4+1)/(2+1)=5/3\)；最小值在\(x=5\)处取得：\(k=(-2+1)/(5+1)=-1/6\)．所以\((y+1)/(x+1)\)的取值范围为\([-1/6,5/3]\)．}
\fi
\end{ansblock}

\begin{ansblock}[2章-导-课时02-G1]
% ans:2章-导-课时02-G1
\ansitem{21}{D}
\ifansbookdetail
\ansnote{详解}{直线\(l\)的向上方向与\(y\)轴正方向成\(30^{\circ}\)角，\(l\)有两种位置（偏左或偏右）：与\(x\)轴正方向所成的角可能为\(90^{\circ}-30^{\circ}=60^{\circ}\)或\(90^{\circ}+30^{\circ}=120^{\circ}\)，故\(l\)的倾斜角为\(60^{\circ}\)或\(120^{\circ}\)，选D．}
\fi
\end{ansblock}

\begin{ansblock}[2章-导-课时02-G2]
% ans:2章-导-课时02-G2
\ansitem{22}{D}
\ifansbookdetail
\ansnote{详解}{A：当\(\alpha=90^{\circ}\)时，直线的斜率不存在，故A不正确；B：虽然直线的斜率为\(\tan\alpha\)，但只有当\(0^{\circ}\leqslant\alpha<180^{\circ}\)时，\(\alpha\)才是此直线的倾斜角，故B不正确；C：当直线与\(x\)轴平行或重合时，\(\alpha=0^{\circ}\)，\(\sin\alpha=0\)，故C不正确；由倾斜角的定义（任意直线都有倾斜角，范围\(0^{\circ}\leqslant\alpha<180^{\circ}\)）与斜率的定义（\(\alpha\neq 90^{\circ}\)时\(k=\tan\alpha\)）知D正确，选D．}
\fi
\end{ansblock}

\begin{ansblock}[2章-导-课时02-G3]
% ans:2章-导-课时02-G3
\ansitem{23}{D}
\ifansbookdetail
\ansnote{详解}{A：取\(\alpha_1=45^{\circ}\)，\(\alpha_2=135^{\circ}\)，则\(k_1=1\)，\(k_2=-1\)，\(k_1>k_2\)，故A错误；B：取\(\alpha_1=\alpha_2=90^{\circ}\)，则\(k_1\)、\(k_2\)都不存在，故B错误；C：取\(k_1=-1\)，\(k_2=1\)，则\(\alpha_1=135^{\circ}\)，\(\alpha_2=45^{\circ}\)，\(\alpha_1>\alpha_2\)，故C错误；D：斜率相等即\(\tan\alpha_1=\tan\alpha_2\)，而\(\alpha_1\)、\(\alpha_2\in[0^{\circ},180^{\circ})\)且均不为\(90^{\circ}\)，正切值相等则角相等，D正确，选D．}
\fi
\end{ansblock}

\begin{ansblock}[2章-导-课时02-G4]
% ans:2章-导-课时02-G4
\ansitem{24}{(1) 90°；(2) 0°；(3) 45°；(4) 135°．}
\ifansbookdetail
\ansnote{详解}{（1）与\(x\)轴垂直，倾斜角为\(90^{\circ}\)；（2）与\(x\)轴平行或重合，规定倾斜角为\(0^{\circ}\)；（3）角平分线\(y=x\)与\(x\)轴正方向成\(45^{\circ}\)角，倾斜角为\(45^{\circ}\)；（4）角平分线\(y=-x\)向上的方向与\(x\)轴正方向成\(135^{\circ}\)角，倾斜角为\(135^{\circ}\)．}
\fi
\end{ansblock}

\begin{ansblock}[2章-导-课时02-G5]
% ans:2章-导-课时02-G5
\ansitem{25}{10°≤α＜100°}
\ifansbookdetail
\ansnote{详解}{由倾斜角的取值范围\(0^{\circ}\leqslant 2\alpha-20^{\circ}<180^{\circ}\)，各边加\(20^{\circ}\)再除以\(2\)，得\(10^{\circ}\leqslant\alpha<100^{\circ}\)．}
\fi
\end{ansblock}

\xjkeshi{2.2.1 直线的方向向量与法向量（第2课时）}

\begin{ansblock}[2章-练-课时03-E3]
% ans:2章-练-课时03-E3
\ansitem{1}{方向向量如(1,−1)（答案不唯一）；k＝−1；θ＝135°}
\ifansbookdetail
\ansnote{详解}{\(\overrightarrow{AB}=(-5+2,\ 3-0)=(-3,3)\)是直线\(l\)的一个方向向量（与之共线的非零向量如\((1,-1)\)也可）．斜率\(k=(3-0)/(-5-(-2))=3/(-3)=-1\)（也可由方向向量\((1,-1)\)得\(k=-1/1=-1\)）．由\(\tan\theta=-1\)且\(\theta\in[0^{\circ},180^{\circ})\)，得\(\theta=135^{\circ}\)．}
\fi
\end{ansblock}

\begin{ansblock}[2章-练-课时03-E1]
% ans:2章-练-课时03-E1
\ansitem{2}{2；4/3}
\ifansbookdetail
\ansnote{详解}{由方向向量\((2,4)\)得直线\(l\)的斜率\(k=4/2=2\)．由两点斜率公式\(k=(m-(-2))/(3-m)=(m+2)/(3-m)\)（\(m\neq 3\)），令\((m+2)/(3-m)=2\)，得\(m+2=6-2m\)，\(3m=4\)，\(m=4/3\)．故斜率为\(2\)，\(m=4/3\)．}
\fi
\end{ansblock}

\begin{ansblock}[2章-练-课时03-E2]
% ans:2章-练-课时03-E2
\ansitem{3}{(1) 1或2；(2) 4/3}
\ifansbookdetail
\ansnote{详解}{(1) \(k_{AC}=[4-(-m-3)]/(-1-m)=(m+7)/(-m-1)\)（\(m\neq -1\)），\(k_{BC}=(4-m+1)/(-1-2)=(5-m)/(-3)\)．由\(k_{AC}=3k_{BC}\)：\((m+7)/(-m-1)=3\cdot(5-m)/(-3)=m-5\)，去分母得\(m+7=(m-5)(-m-1)=-m^{2}+4m+5\)，即\(m^{2}-3m+2=0\)，解得\(m=1\)或\(m=2\)，经检验均使两斜率存在，符合题意．\par\noindent (2) 设\(l_1\)的倾斜角为\(\alpha\)，则\(l_2\)的倾斜角为\(2\alpha\)．由方向向量\(\overrightarrow{n}=(2,1)\)得\(\tan\alpha=1/2\)，故\(l_2\)的斜率\(\tan 2\alpha=2\tan\alpha/(1-\tan^{2}\alpha)=1/(1-1/4)=4/3\)．}
\fi
\end{ansblock}

\begin{ansblock}[2章-练-课时03-E4]
% ans:2章-练-课时03-E4
\ansitem{4}{A、B、C不共线；A、B、D共线．}
\ifansbookdetail
\ansnote{详解}{\(\overrightarrow{AB}=(0+1,\ -1+3)=(1,2)\)，\(\overrightarrow{AC}=(1+1,\ 2+3)=(2,5)\)．因为\(1\times 5-2\times 2=1\neq 0\)，\(\overrightarrow{AB}\)与\(\overrightarrow{AC}\)不共线，且两向量有公共起点\(A\)，所以\(A\)、\(B\)、\(C\)不共线．\(\overrightarrow{AD}=(3+1,\ 5+3)=(4,8)=4\overrightarrow{AB}\)，\(\overrightarrow{AB}\)与\(\overrightarrow{AD}\)共线且有公共起点\(A\)，所以\(A\)、\(B\)、\(D\)共线．}
\fi
\end{ansblock}

\begin{ansblock}[2章-练-课时03-E5]
% ans:2章-练-课时03-E5
\ansitem{5}{真命题．证明见详解．}
\ifansbookdetail
\ansnote{详解}{必要性：若\(A\)、\(B\)、\(C\)共线，则直线上的向量\(\overrightarrow{AB}\)、\(\overrightarrow{BC}\)及与直线平行的向量都是该直线的方向向量，故\(\overrightarrow{AB}\)与\(\overrightarrow{BC}\)共线．\par\noindent 充分性：若\(\overrightarrow{AB}\)与\(\overrightarrow{BC}\)共线，则直线\(AB\)与直线\(BC\)的方向相同或相反；又两直线过公共点\(B\)，由过一点且方向确定的直线唯一，直线\(AB\)与直线\(BC\)是同一条直线，故\(A\)、\(B\)、\(C\)三点共线．综上命题为真．}
\fi
\end{ansblock}

\begin{ansblock}[2章-练-课时03-E12]
% ans:2章-练-课时03-E12
\ansitem{6}{B}
\ifansbookdetail
\ansnote{详解}{\(\overrightarrow{AB}=(3-1,\ 6-2)=(2,4)=2(1,2)\)，故\((1,2)\)为\(l\)的一个方向向量，选B．（辨析：A 中\((2,4)=k(2,1)\)无解，与\(\overrightarrow{AB}\)不成比例；C 中\((-2,1)\cdot(2,4)=-4+4=0\)，是\(l\)的法向量方向而非方向向量；D 中\((1,-2)\)与\(\overrightarrow{AB}\)不成比例，均排除．）}
\fi
\end{ansblock}

\begin{ansblock}[2章-练-课时03-E15]
% ans:2章-练-课时03-E15
\ansitem{7}{m=0}
\ifansbookdetail
\ansnote{详解}{直线\(Ax+By+C=0\)的方向向量可取\((B,\ -A)\)，故\(l\)的方向向量\(t=(2,\ -(m+1))\)．\(t\parallel v\)（横坐标已同为\(2\)）\(\Rightarrow -(m+1)=m-1\)\(\Rightarrow 2m=0\)\(\Rightarrow m=0\)．（验证：\(m=0\)时\(l\)为\(x+2y-6=0\)，方向向量\((2,-1)=v\)；又法向量\((1,2)\)，\(v\cdot n=2-2=0\)，两法一致．）}
\fi
\end{ansblock}

\begin{ansblock}[2章-练-课时03-E8]
% ans:2章-练-课时03-E8
\ansitem{8}{一定．理由见详解．}
\ifansbookdetail
\ansnote{详解}{设\((x_0,y_0)\)是\(l\)的一个方向向量．由方向向量非零知\((x_0,y_0)\neq(0,0)\)，从而\((-y_0,x_0)\neq(0,0)\)．又\((x_0,y_0)\cdot(-y_0,x_0)=-x_0y_0+y_0x_0=0\)，即\((-y_0,x_0)\)与\(l\)的方向向量垂直，于是以\((-y_0,x_0)\)为方向向量的有向线段所在直线与\(l\)垂直，由法向量的定义，\((-y_0,x_0)\)是\(l\)的一个法向量．反之，若\((-y_0,x_0)\)为\(l\)的方向向量，同理\((-y_0,x_0)\neq(0,0)\)即\((x_0,y_0)\neq(0,0)\)，且两向量数量积为\(0\)，故\((x_0,y_0)\)一定是\(l\)的法向量．所以结论成立．}
\fi
\end{ansblock}

\begin{ansblock}[2章-练-课时03-E11]
% ans:2章-练-课时03-E11
\ansitem{9}{证明见详解．}
\ifansbookdetail
\ansnote{详解}{\(\overrightarrow{AB}=(3-1,\ 3-4)=(2,-1)\)，故直线\(AB\)的一个方向向量为\(a=(2,-1)\)；\(\overrightarrow{BC}=(4-3,\ 5-3)=(1,2)\)．取\(v=(1,2)\)，验证\(v\)与\(a\)垂直：\(v\cdot a=1\times 2+2\times(-1)=0\)，且\(v\neq 0\)，所以\(v\)是直线\(AB\)的一个法向量．而\(\overrightarrow{BC}=v\)，即直线\(BC\)的方向向量就是直线\(AB\)的法向量，由法向量的定义，直线\(BC\)与直线\(AB\)垂直，故\(AB\perp BC\)．}
\fi
\end{ansblock}

\begin{ansblock}[2章-练-课时03-E14]
% ans:2章-练-课时03-E14
\ansitem{10}{(1) n₁=(3,−1)（任一非零倍数均可）　(2) a=2/3}
\ifansbookdetail
\ansnote{详解}{(1) 由一般式系数直接读出\(n_1=(3,-1)\)．(2) \(n_2=(a,2)\)．同一平面内每条直线的法向量与其方向向量互相垂直，将两直线绕任一点同旋\(90^{\circ}\)，方向向量恰转到各自法向量方向且夹角不变，故\(l_1\perp l_2\)\(\iff n_1\perp n_2\)\(\iff n_1\cdot n_2=0\)\(\iff 3a+(-1)\times 2=0\)\(\iff a=2/3\)．}
\fi
\end{ansblock}

\begin{ansblock}[2章-练-课时03-E13]
% ans:2章-练-课时03-E13
\ansitem{11}{2x−y−5=0}
\ifansbookdetail
\ansnote{详解}{法向量\((2,-1)\)\(\Rightarrow l\)具形式\(2x-y+C=0\)．代\(P(3,1)\)：\(6-1+C=0\)\(\Rightarrow C=-5\)．故\(2x-y-5=0\)．（验证：\(P\)代入\(6-1-5=0\)；\(l\)的方向向量\((1,2)\)，\(n\cdot(1,2)=2-2=0\)．）}
\fi
\end{ansblock}

\begin{ansblock}[2章-练-课时03-E6]
% ans:2章-练-课时03-E6
\ansitem{12}{(1) y＝−2x−4（即2x＋y＋4＝0）；(2) 当u≠0时为y−y₀＝(v/u)(x−x₀)；当u＝0时为x＝x₀．}
\ifansbookdetail
\ansnote{详解}{(1) 方向向量\(a=(1,-2)\)的横坐标\(1\neq 0\)，故斜率\(k=-2/1=-2\)，由点斜式得\(y-(-2)=-2[x-(-1)]\)，即\(y=-2x-4\)．\par\noindent (2) 设\(Q(x,y)\)是直线\(l\)上任意一点，则\(\overrightarrow{PQ}=(x-x_0,\ y-y_0)\)与方向向量\(a=(u,v)\)共线，得\(v(x-x_0)-u(y-y_0)=0\)．当\(u\neq 0\)时化为\(y-y_0=(v/u)(x-x_0)\)；当\(u=0\)时（此时\(v\neq 0\)）得\(x-x_0=0\)，即\(x=x_0\)（斜率不存在的竖直线）．}
\fi
\end{ansblock}

\begin{ansblock}[2章-练-课时03-E7]
% ans:2章-练-课时03-E7
\ansitem{13}{(1) y＝3x；(2) u(x−x₀)＋v(y−y₀)＝0，即ux＋vy−ux₀−vy₀＝0．}
\ifansbookdetail
\ansnote{详解}{(1) 设\(Q(x,y)\)是直线\(l\)上任意一点，则\(\overrightarrow{PQ}=(x-1,\ y-3)\)与法向量\(a=(-3,1)\)垂直，\(a\cdot\overrightarrow{PQ}=0\)，即\(-3(x-1)+1\cdot(y-3)=0\)，化简得\(y=3x\)．\par\noindent (2) 同法：\(a\cdot\overrightarrow{PQ}=u(x-x_0)+v(y-y_0)=0\)（\(u\)、\(v\)不全为\(0\)，因法向量是非零向量）．这就是直线的点法式方程；当\(v\neq 0\)时可化为斜截形\(y=-(u/v)x+(ux_0+vy_0)/v\)，斜率\(k=-u/v\)（与"方向向量\((v,-u)\)读斜率\(-u/v\)"一致）．}
\fi
\end{ansblock}

\begin{ansblock}[2章-练-课时03-E9]
% ans:2章-练-课时03-E9
\ansitem{14}{(1) y＝2x−11；(2) y＝−(4/3)x＋5（即4x＋3y−15＝0）}
\ifansbookdetail
\ansnote{详解}{(1) \(n=(1,2)\)横坐标不为\(0\)，\(k=2/1=2\)，由点斜式\(y+5=2(x-3)\)，即\(y=2x-11\)．\par\noindent (2) \(n=(3,-4)\)横坐标不为\(0\)，\(k=-4/3\)，由点斜式\(y-5=-(4/3)(x-0)\)，即\(y=-(4/3)x+5\)．}
\fi
\end{ansblock}

\begin{ansblock}[2章-练-课时03-E10]
% ans:2章-练-课时03-E10
\ansitem{15}{(1) 3x−4y＋5＝0；(2) 3x＋4y−5＝0}
\ifansbookdetail
\ansnote{详解}{(1) 设\(Q(x,y)\)在\(l\)上，\(v\cdot\overrightarrow{PQ}=3(x-1)-4(y-2)=0\)，化简得\(3x-4y+5=0\)．\par\noindent (2) 同法：\(3(x+1)+4(y-2)=0\)，化简得\(3x+4y-5=0\)．}
\fi
\end{ansblock}

\begin{ansblock}[2章-练-课时03-E16]
% ans:2章-练-课时03-E16
\ansitem{16}{−√3/3}
\ifansbookdetail
\ansnote{详解}{斜率\(k=\sqrt{3}/1=\sqrt{3}\)，\(l\)：\(y=\sqrt{3}x+1\)．令\(y=0\)：\(x=-1/\sqrt{3}=-\sqrt{3}/3\)．（验证：点\((-\sqrt{3}/3,\ 0)\)代入\(\sqrt{3}\times(-\sqrt{3}/3)+1=-1+1=0\)，斜率与\(v\)一致．）}
\fi
\end{ansblock}

\begin{ansblock}[2章-导-课时03-G1]
% ans:2章-导-课时03-G1
\ansitem{17}{D}
\ifansbookdetail
\ansnote{详解}{直线\(l\)的一个方向向量为\(\overrightarrow{PQ}=(-2-1,\ -2-2)=(-3,-4)\)，\(|\overrightarrow{PQ}|=\sqrt{(-3)^{2}+(-4)^{2}}=5\)，故单位方向向量为\(\pm\overrightarrow{PQ}/|\overrightarrow{PQ}|=\pm(1/5)(-3,-4)=\pm(3/5,4/5)\)（方向相同或相反的两个单位向量都是方向向量）．C 写成"\(\pm\)"作用于每个分量的四种组合，其中\((3/5,-4/5)\)不与\(\overrightarrow{PQ}\)共线，故C错．选D．}
\fi
\end{ansblock}

\begin{ansblock}[2章-导-课时03-G2]
% ans:2章-导-课时03-G2
\ansitem{18}{B}
\ifansbookdetail
\ansnote{详解}{\(\overrightarrow{PQ}=(-m-3,\ 2-2m)\)．两直线平行\(\iff\)方向向量共线（此处即向量共线，注意\(P\)、\(Q\)不重合需\(2m\neq 2\)即\(m\neq 1\)），由共线条件\((-m-3)\times 5-(2-2m)\times(-5)=0\)，得\(-5m-15+10-10m=0\)，\(-15m-5=0\)，\(m=-1/3\)．经检验\(m=-1/3\)符合题意．选B．}
\fi
\end{ansblock}

\begin{ansblock}[2章-导-课时03-G3]
% ans:2章-导-课时03-G3
\ansitem{19}{A}
\ifansbookdetail
\ansnote{详解}{法向量\(v=(2,-4)\)垂直于直线\(l\)，故\(l\)的一个方向向量为与\(v\)垂直的非零向量\(a=(4,2)\)（因\(v\cdot a=2\times 4+(-4)\times 2=0\)），其横坐标不为\(0\)，斜率\(k=2/4=1/2\)．故选A．（另解：法向量为\((A,B)\)的直线可写成\(Ax+By+C=0\)形，\(2x-4y+C=0\)即\(y=x/2+C/4\)，\(k=1/2\)．）}
\fi
\end{ansblock}

\begin{ansblock}[2章-导-课时03-G4]
% ans:2章-导-课时03-G4
\ansitem{20}{(3,4)（答案不唯一）；5}
\ifansbookdetail
\ansnote{详解}{直线\(AB\)的一个方向向量为\(\overrightarrow{AB}=(3-0,\ 0+4)=(3,4)\)（与之平行的任何非零向量均可）；线段\(AB\)的长度为\(|AB|=\sqrt{3^{2}+4^{2}}=5\)．}
\fi
\end{ansblock}

\begin{ansblock}[2章-导-课时03-G5]
% ans:2章-导-课时03-G5
\ansitem{21}{[0,3]}
\ifansbookdetail
\ansnote{详解}{方向向量的横坐标为\(1\neq 0\)，故直线\(AB\)的斜率存在，设为\(k\)．由方向向量\((1,\sin\alpha)\)知\(k=\sin\alpha/1=\sin\alpha\in[-1,1]\)；又由两点坐标\(k=(2m-1-2)/(-2-1)=(3-2m)/3\)．所以\(-1\leqslant(3-2m)/3\leqslant 1\)，解得\(0\leqslant m\leqslant 3\)．故\(m\)的取值范围为\([0,3]\)．}
\fi
\end{ansblock}

\jietitle{2.2.2 直线的点斜式与斜截式}

\begin{ansblock}[2章-练-课时04-E7]
% ans:2章-练-课时04-E7
\ansitem{1}{(1) y−5＝4(x−2)；(2) y−2＝(√3/3)(x＋2)}
\ifansbookdetail
\ansnote{详解}{(1) 由点斜式\(y-y_0=k(x-x_0)\)，得\(y-5=4(x-2)\)．\par\noindent (2) 斜率\(k=\tan30^\circ=\sqrt{3}/3\)，故\(y-2=(\sqrt{3}/3)[x-(-2)]=(\sqrt{3}/3)(x+2)\)．}
\fi
\end{ansblock}

\begin{ansblock}[2章-练-课时04-E8]
% ans:2章-练-课时04-E8
\ansitem{2}{(1) y＝(√3/2)x−2；(2) y＝−x＋3}
\ifansbookdetail
\ansnote{详解}{(1) \(k=\sqrt{3}/2\)，\(b=-2\)，由斜截式\(y=kx+b\)得\(y=(\sqrt{3}/2)x-2\)．\par\noindent (2) \(k=\tan135^\circ=-1\)，\(b=3\)，故\(y=-x+3\)．}
\fi
\end{ansblock}

\begin{ansblock}[2章-练-课时04-E1]
% ans:2章-练-课时04-E1
\ansitem{3}{C}
\ifansbookdetail
\ansnote{详解}{将\(y-b=2(x-a)\)化为斜截式：\(y=2x-2a+b\)，故直线在\(y\)轴上的截距为\(b-2a\)（\(x=0\)时的纵坐标值，可正可负，不加绝对值）．选：C．}
\fi
\end{ansblock}

\begin{ansblock}[2章-练-课时04-E9]
% ans:2章-练-课时04-E9
\ansitem{4}{(1) k＝1，b＝3；(2) k＝√3，b＝−2√3−4；(3) k＝−4，b＝3}
\ifansbookdetail
\ansnote{详解}{(1) \(y-2=x+1\)化为\(y=x+3\)，故\(k=1\)，截距\(b=3\)（定点\((-1,2)\)）；\par\noindent (2) 化为\(y=\sqrt{3}x-2\sqrt{3}-4\)，故\(k=\sqrt{3}\)，截距\(b=-2\sqrt{3}-4\)；\par\noindent (3) 化为\(y=-4x+3\)，故\(k=-4\)，截距\(b=3\)．}
\fi
\end{ansblock}

\begin{ansblock}[2章-练-课时04-E3]
% ans:2章-练-课时04-E3
\ansitem{5}{D}
\ifansbookdetail
\ansnote{详解}{直线\(\sqrt{3}x-y-\sqrt{3}=0\)的斜率为\(\sqrt{3}\)，倾斜角\(\alpha\)满足\(\tan\alpha=\sqrt{3}\)，\(\alpha=60^\circ\)；所求直线倾斜角\(\beta=2\alpha=120^\circ\)，斜率\(k=\tan120^\circ=-\sqrt{3}\)；又在\(y\)轴上的截距为\(-1\)，故斜截式方程为\(y=-\sqrt{3}x-1\)，即\(\sqrt{3}x+y+1=0\)，比对\(ax+by-1=0\)的常数项须为\(-1\)，须将\(\sqrt{3}x+y+1=0\)两边乘\(-1\)得\(-\sqrt{3}x-y-1=0\)，所以\(a=-\sqrt{3}\)，\(b=-1\)．选：D．}
\fi
\end{ansblock}

\begin{ansblock}[2章-练-课时04-E5]
% ans:2章-练-课时04-E5
\ansitem{6}{A在，B不在，C在．}
\ifansbookdetail
\ansnote{详解}{\(x=0\)时\(y=-6\times0+7=7\)，故\(A(0,7)\)在直线上；\par\noindent\(x=2\)时\(y=-12+7=-5\neq5\)，故\(B(2,5)\)不在直线上；\par\noindent\(x=1\)时\(y=-6+7=1\)，故\(C(1,1)\)在直线上．}
\fi
\end{ansblock}

\begin{ansblock}[2章-练-课时04-E6]
% ans:2章-练-课时04-E6
\ansitem{7}{a＝0；b＝4}
\ifansbookdetail
\ansnote{详解}{两点都满足直线方程：\(1=-3a+1\)，得\(a=0\)；\(b=-3\times(-1)+1=4\)．\par\noindent（检验：\(a=0\)时两点为\((0,1)\)、\((-1,4)\)，确为\(y=-3x+1\)上两点，两点横坐标不等，直线唯一．）}
\fi
\end{ansblock}

\begin{ansblock}[2章-练-课时04-E10]
% ans:2章-练-课时04-E10
\ansitem{8}{(1) 真；(2) 假；(3) 真}
\ifansbookdetail
\ansnote{详解}{(1) 斜率不存在时\(l\)为竖直线\(x=m\)，它与\(x\)轴恰有一个交点\((m,0)\)，\(x\)轴截距为\(m\)，唯一，真命题；\par\noindent (2) 斜率存在时\(l\)与\(y\)轴必只有一个交点（\(y\)轴截距唯一），但与\(x\)轴未必相交：水平直线\(y=b\)（\(b\neq0\)，如\(y=1\)）与\(x\)轴无交点，\(x\)轴截距不存在，故“截距唯一”为假命题；\par\noindent (3) 过原点的直线（如\(y=x\)）与两轴交点都为原点\((0,0)\)，两截距都为0，真命题．}
\fi
\end{ansblock}

\begin{ansblock}[2章-练-课时04-E11]
% ans:2章-练-课时04-E11
\ansitem{9}{B}
\ifansbookdetail
\ansnote{详解}{\(M\)、\(N\)纵坐标相等（都为2），直线\(MN\)与\(x\)轴平行，斜率为0，方程为\(y=2\)（同纵坐标的两点连线与\(x\)轴平行）．选：B．}
\fi
\end{ansblock}

\begin{ansblock}[2章-练-课时04-E15]
% ans:2章-练-课时04-E15
\ansitem{10}{(1) y＝−2x−1；(2) y＝−1；(3) x＝1}
\ifansbookdetail
\ansnote{详解}{(1) 两点横、纵坐标均不等，斜率\(k=(1+1)/(-1-0)=-2\)，点斜式（用\(A\)）：\(y+1=-2(x-0)\)，即\(y=-2x-1\)；\par\noindent (2) \(C\)、\(D\)纵坐标同为\(-1\)，\(l_2\parallel x\)轴，方程为\(y=-1\)（纵坐标相同，与\(x\)轴平行）；\par\noindent (3) \(E\)、\(F\)横坐标同为1，\(l_3\perp x\)轴，方程为\(x=1\)（横坐标相同，为竖直线\(x=1\)）．}
\fi
\end{ansblock}

\begin{ansblock}[2章-练-课时04-E13]
% ans:2章-练-课时04-E13
\ansitem{11}{B}
\ifansbookdetail
\ansnote{详解}{令\(y=0\)得\(x=2\)；令\(x=0\)得\(y=-5\)．即\(a=2\)，\(b=-5\)（化截距式\(x/2+y/(-5)=1\)同判）．选：B．}
\fi
\end{ansblock}

\begin{ansblock}[2章-练-课时04-E4]
% ans:2章-练-课时04-E4
\ansitem{12}{x＋2y−2＝0或2x＋3y−6＝0}
\ifansbookdetail
\ansnote{详解}{设直线在\(x\)轴上的截距为\(a\)（\(a\neq0\)且\(a-1\neq0\)），则在\(y\)轴上的截距为\(a-1\)，由截距式：\(x/a+y/(a-1)=1\)．\par\noindent 将\((6,-2)\)代入：\(6/a-2/(a-1)=1\)，去分母：\(6(a-1)-2a=a(a-1)\)，即\(4a-6=a^{2}-a\)，\(a^{2}-5a+6=0\)，解得\(a=2\)或\(a=3\)．\par\noindent \(a=2\)时：\(x/2+y/1=1\)，即\(x+2y-2=0\)；\(a=3\)时：\(x/3+y/2=1\)，即\(2x+3y-6=0\)．\par\noindent （另须检验截距为0的情形：若过原点，设\(y=mx\)，代\((6,-2)\)得\(m=-1/3\)，此时两截距都为0，差为\(0\neq1\)，不合．）所以直线\(l\)的方程为\(x+2y-2=0\)或\(2x+3y-6=0\)．}
\fi
\end{ansblock}

\begin{ansblock}[2章-练-课时04-E2]
% ans:2章-练-课时04-E2
\ansitem{13}{A}
\ifansbookdetail
\ansnote{详解}{设\(l\)与\(x\)轴交于\(M_0(t,0)\)，则\(t\in(-1,3)\)且\(t\neq2\)（\(t=2\)时直线过\(P\)与\(x\)轴交点同横坐标，即\(x=2\)，斜率不存在，截距为\(2\in(-1,3)\)——此情形\(k\)不存在，题设“斜率\(k\)”已默认存在，排除）．\(k=(3-0)/(2-t)=3/(2-t)\)．当\(t\in(-1,2)\)时\(2-t\in(0,3)\)，\(k=3/(2-t)\in(1,+\infty)\)；当\(t\in(2,3)\)时\(2-t\in(-1,0)\)，\(k\in(-\infty,-3)\)．故\(k\in(-\infty,-3)\cup(1,+\infty)\)．选：A．（临界图法同件源：取\(x\)轴上两端点\(M(-1,0)\)、\(N(3,0)\)，\(k_{PM}=3/3=1\)，\(k_{PN}=3/(-1)=-3\)，\(l\)与开线段\(MN\)（挖去\((2,0)\)上方对应竖线）相交即\(k>1\)或\(k<-3\)．）}
\fi
\end{ansblock}

\begin{ansblock}[2章-练-课时04-E12]
% ans:2章-练-课时04-E12
\ansitem{14}{C}
\ifansbookdetail
\ansnote{详解}{分类：①过原点（两截距均为0）：\(y=4x\)，1条；\par\noindent ②不过原点，设\(x/a+y/b=1\)，\(|a|=|b|\neq0\)：代入\((1,4)\)得\(1/a+4/b=1\)．\(b=a\)时\(5/a=1\)，\(a=5\)（\(x+y=5\)）；\(b=-a\)时\(1/a-4/a=1\)，\(a=-3\)（\(-x/3+y/3=1\)，即\(y-x=3\)）．共2条．\par\noindent 合计3条．选：C．}
\fi
\end{ansblock}

\begin{ansblock}[2章-练-课时04-E14]
% ans:2章-练-课时04-E14
\ansitem{15}{x＋2y−6＝0}
\ifansbookdetail
\ansnote{详解}{设\(l:x/a+y/b=1\)（\(a>0\)，\(b>0\)），代入\(P(4,1)\)：\(4/a+1/b=1\)．\par\noindent\(|OA|+|OB|=a+b=(a+b)(4/a+1/b)=5+4b/a+a/b\geqslant5+2\sqrt{(4b/a)\cdot(a/b)}=9\)，当且仅当\(4b/a=a/b\)即\(a=2b\)时取等；联立\(4/a+1/b=1\)代入\(a=2b\)：\(4/(2b)+1/b=1\)\(\Rightarrow3/b=1\)\(\Rightarrow b=3\)，\(a=6\)．\par\noindent 此时\(l\)为\(x/6+y/3=1\)，即\(x+2y-6=0\)．}
\fi
\end{ansblock}

\begin{ansblock}[2章-练-课时04-E16]
% ans:2章-练-课时04-E16
\ansitem{16}{(1) y＝−(4/3)x−2；(2) x/3＋y/2＝1（即2x＋3y−6＝0）}
\ifansbookdetail
\ansnote{详解}{(1) 截距\(b=-2\)，直线过\((0,-2)\)与\((-3,2)\)，斜率\(k=(2+2)/(-3-0)=-4/3\)，斜截式\(y=-(4/3)x-2\)；\par\noindent (2) 直线过\((3,0)\)，故\(x\)轴截距\(a=3\)（非零），\(y\)轴截距\(b=5-3=2\)（非零），由截距式\(x/3+y/2=1\)，即\(2x+3y-6=0\)．}
\fi
\end{ansblock}

\begin{ansblock}[2章-拓-课时04-T1]
% ans:2章-拓-课时04-T1
\ansitem{17}{(1) k≥0；(2) S的最小值为16，此时直线l的方程为x−2y＋8＝0}
\ifansbookdetail
\ansnote{详解}{(1) 将方程整理为\(k(x+4)+(-y+2)=0\)，令\(x+4=0\)且\(y=2\)，得直线恒过定点\((-4,2)\)（在第二象限）．斜率\(k=\tan\theta\)．过\((-4,2)\)的直线不经过第四象限，当且仅当它“向右上或水平张开”：由临界位置（过\((-4,2)\)与原点方向）逆时针转到竖直位置均不进入第四象限，即倾斜角\(\theta\in[0^\circ,90^\circ]\)，故\(k\geqslant0\)（\(k\)不存在时直线\(x=-4\)也不过第四象限，但它不在\(y=kx+b\)的参数范围内——题设方程本身取遍\(k\in\mathbf{R}\)时不含竖直线，而\(x=-4\)无法由该方程表示，故不另计）；\(k<0\)时直线向右下倾斜，自第二象限的点\((-4,2)\)出发必穿过第四象限，不合．所以\(k\geqslant0\)．\par\noindent (2) 由题意\(A\)、\(B\)存在且\(a<0\)、\(b>0\)（竖直线\(x=-4\)不与\(y\)轴相交，此时\(k=0\)给\(y=2\)与\(x\)轴不相交，均已在端点情形排除，直线可设为截距式\(x/a+y/b=1\)，\(a<0\)，\(b>0\)）．\((-4,2)\)在直线上：\(-4/a+2/b=1\)．因\(a<0\)，\(b>0\)，\(-4/a>0\)，\(2/b>0\)，\(1=-4/a+2/b\geqslant2\sqrt{(-4/a)\cdot(2/b)}=2\sqrt{-8/(ab)}\)，两边平方得\(1\geqslant-32/(ab)\)，即\(-ab\geqslant32\)，\(ab\leqslant-32\)．\(S=(1/2)|OA|\cdot|OB|=(1/2)(-a)b=-ab/2\geqslant16\)，当且仅当\(-4/a=2/b=1/2\)，即\(a=-8\)，\(b=4\)时取等号．此时直线为\(x/(-8)+y/4=1\)，化简得\(x-2y+8=0\)．所以\(S\)的最小值为16，此时直线\(l\)的方程为\(x-2y+8=0\)．}
\fi
\end{ansblock}

\begin{ansblock}[2章-导-课时04-G1]
% ans:2章-导-课时04-G1
\ansitem{18}{D}
\ifansbookdetail
\ansnote{详解}{对任意\(k\in\mathbf{R}\)，\(x=1\)时\(y=k(1-1)=0\)，故直线恒过点\((1,0)\)；又方程形式本身要求斜率\(k\)存在，故不含与\(y\)轴平行的直线\(x=1\)；反之，过\((1,0)\)且斜率存在的每条直线都可写成该形．所以它表示过点\((1,0)\)且除直线\(x=1\)外的一切直线．选：D．}
\fi
\end{ansblock}

\begin{ansblock}[2章-导-课时04-G2]
% ans:2章-导-课时04-G2
\ansitem{19}{A}
\ifansbookdetail
\ansnote{详解}{由点斜式\(y-y_0=k(x-x_0)\)的结构知直线过定点\((4,3)\)，斜率\(k=\sqrt{3}\)，即\(\tan\alpha=\sqrt{3}\)，倾斜角\(\alpha=60^\circ\)．选：A．}
\fi
\end{ansblock}

\begin{ansblock}[2章-导-课时04-G3]
% ans:2章-导-课时04-G3
\ansitem{20}{D}
\ifansbookdetail
\ansnote{详解}{直线的斜率为\(\tan60^\circ=\sqrt{3}\)，又\(y\)轴上的截距\(b=-2\)，由斜截式得\(y=\sqrt{3}x-2\)．故选：D．}
\fi
\end{ansblock}

\begin{ansblock}[2章-导-课时04-G4]
% ans:2章-导-课时04-G4
\ansitem{21}{y＝√3x−6或y＝−√3x−6}
\ifansbookdetail
\ansnote{详解}{直线与\(y\)轴相交成\(30^\circ\)角，则其倾斜角为\(90^\circ-30^\circ=60^\circ\)或\(90^\circ+30^\circ=120^\circ\)，斜率为\(\sqrt{3}\)或\(-\sqrt{3}\)；又\(b=-6\)，由斜截式得所求方程为\(y=\sqrt{3}x-6\)或\(y=-\sqrt{3}x-6\)．}
\fi
\end{ansblock}

\begin{ansblock}[2章-导-课时04-G5]
% ans:2章-导-课时04-G5
\ansitem{22}{y＝−3x＋4}
\ifansbookdetail
\ansnote{详解}{所求直线与\(y=3x+4\)关于\(y\)轴对称：斜率互为相反数（\(k=-3\)），与\(y\)轴交点\((0,4)\)不变（截距仍为4），故方程为\(y=-3x+4\)．}
\fi
\end{ansblock}

\jietitle{2.2.2 直线的两点式与一般式}

\begin{ansblock}[2章-练-课时05-E2]
% ans:2章-练-课时05-E2
\ansitem{1}{A}
\ifansbookdetail
\ansnote{详解}{两点式：\((y+1)/(5+1)=(x+1)/(2+1)\)，即\((y+1)/6=(x+1)/3\)，化简得\(y=2x+1\)．\par\noindent 代入\((1009,b)\)：\(b=2\times 1009+1=2019\)．选：A．}
\fi
\end{ansblock}

\begin{ansblock}[2章-练-课时05-E3]
% ans:2章-练-课时05-E3
\ansitem{2}{C}
\ifansbookdetail
\ansnote{详解}{A、B要求\(x_1\neq x_2\)且\(y_1\neq y_2\)（或至少分母不为0），不能表示与坐标轴平行（或重合）的直线；\par\noindent C式即\((y_2-y_1)(x-x_1)=(x_2-x_1)(y-y_1)\)：当\(x_1\neq x_2\)时化为点斜式，当\(x_1=x_2\)时得\(x=x_1\)（竖直线）也成立，\(y_1=y_2\)时得\(y=y_1\)同样成立，对任意两点均恰当表示直线；D是把两坐标“同向”组合，一般不表示过A、B的直线（它表示到A、B横、纵坐标差相等点的轨迹型式子）．选：C．}
\fi
\end{ansblock}

\begin{ansblock}[2章-练-课时05-E11]
% ans:2章-练-课时05-E11
\ansitem{3}{2x−y＋1＝0}
\ifansbookdetail
\ansnote{详解}{AB中点\(D((3-1)/2,(2+4)/2)=(1,3)\)．\(C\)、\(D\)横坐标不等（\(2\neq1\)）、纵坐标不等（\(5\neq3\)），用两点式：\((y-5)/(3-5)=(x-2)/(1-2)\)，即\((y-5)/(-2)=(x-2)/(-1)\)，\par\noindent 去分母：\(y-5=2(x-2)\)，化简得\(2x-y+1=0\)．}
\fi
\end{ansblock}

\begin{ansblock}[2章-练-课时05-E15]
% ans:2章-练-课时05-E15
\ansitem{4}{不是（见详解）}
\ifansbookdetail
\ansnote{详解}{方程\((y-2)/(x-3)=1\)要求\(x\neq3\)，其解集对应的点集是直线\(y-2=x-3\)（即\(y=x-1\)）上去掉点\((3,2)\)后的全体；\par\noindent 而“直线\(y=x-1\)的方程”应使以该直线上的点的坐标为解、且以方程的解为坐标的点都在这条直线上——点\((3,2)\)在直线\(y=x-1\)上却不是该方程的解，两集合不同．\par\noindent 所以\((y-2)/(x-3)=1\)不是（这条）直线的方程，它表示去掉一个点的“直线型”曲线；应改写成\((y-2)=(x-3)\)即\(x-y-1=0\)．}
\fi
\end{ansblock}

\begin{ansblock}[2章-练-课时05-E16]
% ans:2章-练-课时05-E16
\ansitem{5}{(1) C＝0；(2) A≠0且B≠0；(3) A＝0且C≠0；(4) B＝0且C≠0；(5) B＝0（且A≠0）；(6) A＝0（且B≠0）}
\ifansbookdetail
\ansnote{详解}{(1) 原点坐标满足方程：\(C=0\)（\(A\)、\(B\)不同时为0由前提保证）；\par\noindent (2) 与两坐标轴都相交\(\iff\)既不与\(x\)轴平行也不与\(y\)轴平行\(\iff A\neq0\)且\(B\neq0\)；\par\noindent (3) 与\(x\)轴无交点\(\iff\)\(l\)为水平线\(y=m\)（\(A=0\)，\(B\neq0\)）且\(m=-C/B\neq0\)，即\(A=0\)且\(C\neq0\)；\par\noindent (4) 与\(y\)轴无交点\(\iff\)\(l\)为竖直线\(x=n\)（\(B=0\)，\(A\neq0\)）且\(n=-C/A\neq0\)，即\(B=0\)且\(C\neq0\)；\par\noindent (5) 与\(x\)轴垂直\(\iff\)\(l\)为竖直线\(x=n\)（\(n\)可为0）：\(B=0\)且\(A\neq0\)；\par\noindent (6) 与\(y\)轴垂直\(\iff\)\(l\)为水平线\(y=m\)（\(m\)可为0）：\(A=0\)且\(B\neq0\)．\par\noindent （说明：(3)(4)是(5)(6)中去掉与坐标轴重合情形的细分——(3)即(6)中\(C\neq0\)（排除\(y=0\)），(4)即(5)中\(C\neq0\)（排除\(x=0\)）．）}
\fi
\end{ansblock}

\begin{ansblock}[2章-练-课时05-E5]
% ans:2章-练-课时05-E5
\ansitem{6}{C}
\ifansbookdetail
\ansnote{详解}{二元一次方程表示直线\(\iff\)x、y系数不全为0．令\(2m^{2}+m-3=(2m+3)(m-1)=0\)且\(m^{2}-m=m(m-1)=0\)同时成立，公共解只有\(m=1\)；\par\noindent 故方程表示直线\(\iff m\neq1\)（\(m=-3/2\)时仅\(x\)系数为0、\(y\)系数\(15/4\neq0\)，仍表示直线；\(m=0\)时\(x\)系数\(-3\neq0\)）．选：C．}
\fi
\end{ansblock}

\begin{ansblock}[2章-练-课时05-E6]
% ans:2章-练-课时05-E6
\ansitem{7}{C}
\ifansbookdetail
\ansnote{详解}{\(P\)在\(l\)外\(\Rightarrow Ax_0+By_0+C\neq0\)．新方程即\(Ax+By+2C+Ax_0+By_0=0\)，\(x\)、\(y\)系数与\(l\)完全相同而常数项不同，故新直线与\(l\)平行（不重合）；\par\noindent 将\(P\)代入新方程左边得\((Ax_0+By_0+C)+(Ax_0+By_0+C)=2(Ax_0+By_0+C)\neq0\)，故新直线不过\(P\)．选：C．}
\fi
\end{ansblock}

\begin{ansblock}[2章-练-课时05-E7]
% ans:2章-练-课时05-E7
\ansitem{8}{A}
\ifansbookdetail
\ansnote{详解}{把\(A(2,1)\)代入两线方程：\(2a_1+b_1+1=0\)且\(2a_2+b_2+1=0\)，\par\noindent 即\(P_1\)、\(P_2\)两点都满足方程\(2x+y+1=0\)．又两已知直线不同（否则\(a_1:a_2=b_1:b_2=1:1\)，题中\(P_1\)、\(P_2\)为两点确定直线），所以\(2x+y+1=0\)即所求直线．选：A．}
\fi
\end{ansblock}

\begin{ansblock}[2章-练-课时05-E9]
% ans:2章-练-课时05-E9
\ansitem{9}{A}
\ifansbookdetail
\ansnote{详解}{\(M\)在直线上：\(Ax_0+By_0+C=0\)，故\(C=-Ax_0-By_0\)，代入\(Ax+By+C=0\)得\(A(x-x_0)+B(y-y_0)=0\)．选：A．}
\fi
\end{ansblock}

\begin{ansblock}[2章-练-课时05-E10]
% ans:2章-练-课时05-E10
\ansitem{10}{B}
\ifansbookdetail
\ansnote{详解}{\(l\parallel y\)轴\(\iff\)\(y\)的系数为0且\(x\)的系数不为0：\(b+2=0\)且\(a-1\neq0\)，即\(b=-2\)、\(a\neq1\)；\par\noindent \(l\)与\(y\)轴（方程\(x=0\)）不重合\(\iff c\neq0\)（\(b+2=0\)时\(l\)为\((a-1)x+c=0\)即\(x=-c/(a-1)\)，与\(y\)轴重合\(\iff c=0\)）．\par\noindent 所以\(a\neq1\)，\(b=-2\)，\(c\neq0\)．选：B．}
\fi
\end{ansblock}

\begin{ansblock}[2章-练-课时05-E17]
% ans:2章-练-课时05-E17
\ansitem{11}{A（它们都经过定点(−1,1)；且除直线x＝−1外，过(−1,1)的直线都在其中）}
\ifansbookdetail
\ansnote{详解}{对任意实数\(k\)，把\(x=-1\)、\(y=1\)代入方程两边：左\(=1-1=0\)，右\(=k(-1+1)=0\)恒成立，故每条直线都过点\((-1,1)\)；\par\noindent 反之，过\((-1,1)\)且斜率存在（设为\(k\)）的直线方程即\(y-1=k(x+1)\)；斜率不存在的直线\(x=-1\)无法由该方程表示．选：A．（注：本选项组为本卷补写——源题为问答题，为入单选槽按原题数据配四选项，结论为A项不变．）}
\fi
\end{ansblock}

\begin{ansblock}[2章-练-课时05-E18]
% ans:2章-练-课时05-E18
\ansitem{12}{A（定点(9,−4)，证明见详解）}
\ifansbookdetail
\ansnote{详解}{将方程按\(m\)整理：\(m(x+2y-1)+(-x-y+5)=0\)对一切实数\(m\)恒成立，\par\noindent 须\(x+2y-1=0\)且\(-x-y+5=0\)，联立解得\(x=9\)，\(y=-4\)．\par\noindent 即点\((9,-4)\)使方程对任意\(m\)都成立，故直线\(l\)恒过定点\((9,-4)\)．\par\noindent （取\(m=1\)得\(y=-4\)、取\(m=1/2\)得\(x=9\)，交点\((9,-4)\)代回原方程恒成立，亦可验证．）选：A．（注：本选项组为本卷补写——源题为「求证恒过定点并求坐标」解答形，为入单选槽按原题数据配四选项，定点结论为A项不变，证明详解全保留．）}
\fi
\end{ansblock}

\begin{ansblock}[2章-练-课时05-E21]
% ans:2章-练-课时05-E21
\ansitem{13}{3x＋4y−12＝0，3x−4y＋12＝0，3x−4y−12＝0，3x＋4y＋12＝0（即3x±4y±12＝0的四种符号组合）}
\ifansbookdetail
\ansnote{详解}{对角线在坐标轴上且长分别为8、6，则四个顶点为\((\pm4,0)\)、\((0,\pm3)\)．\par\noindent 四边即依次连接\((4,0)\)、\((0,3)\)、\((-4,0)\)、\((0,-3)\)的回字斜方形，各用截距式：\par\noindent \((4,0)\)-\((0,3)\)：\(x/4+y/3=1\)\(\rightarrow\)\(3x+4y-12=0\)；\((0,3)\)-\((-4,0)\)：\(x/(-4)+y/3=1\)\(\rightarrow\)\(3x-4y+12=0\)；\par\noindent \((-4,0)\)-\((0,-3)\)：\(x/(-4)+y/(-3)=1\)\(\rightarrow\)\(3x+4y+12=0\)；\((0,-3)\)-\((4,0)\)：\(x/4+y/(-3)=1\)\(\rightarrow\)\(3x-4y-12=0\)．}
\fi
\end{ansblock}

\begin{ansblock}[2章-练-课时05-E22]
% ans:2章-练-课时05-E22
\ansitem{14}{(1) x−2y−4＝0；(2) x＋2y＝0或x−y−3＝0}
\ifansbookdetail
\ansnote{详解}{(1) 与\(2x+y+3=0\)（法向量\((2,1)\)）垂直的直线，其法向量与该直线方向共线，可设\(l\)：\(x-2y+m=0\)（即以\((1,-2)\)为法向量，\((1,-2)\perp(2,1)\)……由课时03点法式知识合法）；\par\noindent 代入\(P(2,-1)\)：\(2+2+m=0\)，\(m=-4\)，故\(l\)：\(x-2y-4=0\)．\par\noindent (2) 当\(l\)过原点时，设\(y=kx\)，代\(P(2,-1)\)得\(k=-1/2\)，即\(x+2y=0\)，此时两截距都是0，互为相反数；\par\noindent 当\(l\)不过原点时，设两截距为\(a\)、\(-a\)（\(a\neq0\)），\(x/a+y/(-a)=1\)，即\(x-y=a\)，代\(P(2,-1)\)：\(a=3\)，\(l\)：\(x-y-3=0\)．\par\noindent 综上，所求直线为\(x+2y=0\)或\(x-y-3=0\)．}
\fi
\end{ansblock}

\begin{ansblock}[2章-拓-课时05-T6]
% ans:2章-拓-课时05-T6
\ansitem{15}{(1) 2x＋y−6＝0或8x＋y−12＝0；(2) 8，4x＋y−8＝0}
\ifansbookdetail
\ansnote{详解}{设\(A(a,0)\)、\(B(0,b)\)，\(a>0\)，\(b>0\)，\(l\)：\(x/a+y/b=1\)，代入\(P(1,4)\)得\(1/a+4/b=1\)（*）．\par\noindent (1) 又\((1/2)ab=9\)即\(ab=18\)；由(*)去分母\(b+4a=ab=18\)，\(b=18-4a\)，\(a(18-4a)=18\)，\(2a^{2}-9a+9=0\)，\(a=3\)或\(3/2\)；\par\noindent 对应\(b=6\)或12，\(l\)为\(x/3+y/6=1\)或\(x/(3/2)+y/12=1\)，即\(2x+y-6=0\)或\(8x+y-12=0\)．\par\noindent (2) \(S=(1/2)ab=(1/2)ab\cdot(1/a+4/b)^{2}=(1/2)(8+b/a+16a/b)\geqslant(1/2)(8+2\sqrt{(b/a)\cdot(16a/b)})=8\)，\par\noindent 当且仅当\(b/a=16a/b\)即\(b=4a\)，代(*)：\(1/a+1/a=1\)，\(a=2\)、\(b=8\)时取等，\(S\)最小值8，此时\(l\)：\(x/2+y/8=1\)即\(4x+y-8=0\)．}
\fi
\end{ansblock}

\begin{ansblock}[2章-拓-课时05-T7]
% ans:2章-拓-课时05-T7
\ansitem{16}{(1) 最小值8，x/4＋y/2＝1；(2) 最小值4，x/3＋y/3＝1}
\ifansbookdetail
\ansnote{详解}{(1) 设\(l\)：\(x/a+y/b=1\)（\(a>0\)，\(b>0\)），过\(P\)得\(2/a+1/b=1\)．\par\noindent \(1=2/a+1/b\geqslant2\sqrt{2/(ab)}\)，得\(ab\geqslant8\)，当且仅当\(2/a=1/b\)即\(a=4\)、\(b=2\)时取等；\par\noindent \(|OA|\cdot|OB|=ab\)的最小值为8，此时\(l\)：\(x/4+y/2=1\)．\par\noindent (2) 由(1)得\(b=a/(a-2)\)（\(a>2\)）．\(A(a,0)\)、\(B(0,b)\)，\par\noindent \(|PA|=\sqrt{(a-2)^{2}+1}\)，\(|PB|=\sqrt{4+(b-1)^{2}}\)，而\(b-1=a/(a-2)-1=2/(a-2)\)，\par\noindent \(|PA|\cdot|PB|=\sqrt{(a-2)^{2}+1}\cdot\sqrt{4+4/(a-2)^{2}}\)\par\noindent\(=2\sqrt{(a-2)^{2}+1/(a-2)^{2}+2}\geqslant2\sqrt{2+2}=4\)，\par\noindent 当且仅当\((a-2)^{2}=1/(a-2)^{2}\)即\(a-2=1\)、\(a=3\)时取等，此时\(b=3\)；\par\noindent \(|PA|\cdot|PB|\)的最小值为4，此时\(l\)：\(x/3+y/3=1\)．}
\fi
\end{ansblock}

\begin{ansblock}[2章-拓-课时05-T8]
% ans:2章-拓-课时05-T8
\ansitem{17}{(1) y＝(3/4)x±3；(2) m≠1时y＝(1/(m−1))(x−1)，m＝1时x＝1；(3) x＋y−1＝0或x−y−7＝0或y＝−(3/4)x}
\ifansbookdetail
\ansnote{详解}{(1) 设\(l\)：\(y=(3/4)x+b\)，则\(y\)轴截距\(b\)，\(x\)轴截距\(-4b/3\)，面积\((1/2)|b\cdot(-4b/3)|=2b^{2}/3=6\)，\(b^{2}=9\)，\(b=\pm3\)，\(l\)：\(y=(3/4)x\pm3\)；回代验证：\(b=3\) 代回，面积\((1/2)\cdot|3\cdot(-4)|=6\)；\(b=-3\) 代回，面积\((1/2)\cdot|(-3)\cdot4|=6\)，双向验证成立；\par\noindent (2) \(m\neq1\)时两点横、纵坐标均不等，斜率\(k=(1-0)/(m-1)=1/(m-1)\)，点斜式\(y=(1/(m-1))(x-1)\)；\(m=1\)时\(A\)、\(B\)同在直线\(x=1\)上，\(l\)：\(x=1\)；\par\noindent (3) 设截距\(a\)、\(b\)．①过原点（\(a=b=0\)）：\(l\)过\((0,0)\)、\((4,-3)\)，\(y=-(3/4)x\)；\par\noindent ②不过原点，\(x/a+y/b=1\)，\(|a|=|b|\neq0\)且\(4/a-3/b=1\)：\(a=b\)时\(1/a=1\)、\(a=1\)（\(x+y-1=0\)）；\(a=-b\)时\(4/a+3/a=1\)、\(a=7\)（\(x-y-7=0\)）．\par\noindent 共3条．\par\noindent （按源卷核对注：源答案(1)写作“\(y=(4/3)x\pm3\)”，与其题面“斜率\(3/4\)”不符——本卷按题面取修正值 \(y=(3/4)x\pm3\)；面积回代双向验证已并入(1)详解正文，详见亲算账④“源误更正”．）}
\fi
\end{ansblock}

\begin{ansblock}[2章-练-课时05-E19]
% ans:2章-练-课时05-E19
\ansitem{18}{②④}
\ifansbookdetail
\ansnote{详解}{根据点到直线距离公式可得点\((1,0)\)到直线\((x-1)\cos\theta+y\sin\theta=1\)的距离\par\noindent \(d=|0\times\cos\theta+0\times\sin\theta-1|/\sqrt{\cos^{2}\theta+\sin^{2}\theta}=1\)，即点\((1,0)\)到直线的距离为定值1，\par\noindent 因此直线\((x-1)\cos\theta+y\sin\theta=1\)是圆\((x-1)^{2}+y^{2}=1\)的切线（\(\theta\)取遍\([0,2\pi)\)恰给出该圆全体切线）．\par\noindent 对于①，\(M\)中所有直线到点\((1,0)\)的距离为定值1，但并不满足经过一个定点，错误；\par\noindent 对于②，当定点\(P\)在该圆内部（不含圆周）时，\(P\)不在\(M\)中的任意一条直线上，正确；\par\noindent 对于③，\(M\)中的直线所能围成的正三角形，既可以是该圆的内接方向的外切正三角形（每边都与圆相切、圆心即内心），也可以让圆恰为三角形的旁切圆型位置（三边仍为切线但圆在三角形外一侧），二者面积不等，错误；\par\noindent 对于④，当正六边形为该圆的外切正六边形时，其六条边所在直线均在\(M\)中，符合要求，正确．\par\noindent 综上所述，真命题的序号是②④．}
\fi
\end{ansblock}

\begin{ansblock}[2章-练-课时05-E20]
% ans:2章-练-课时05-E20
\ansitem{19}{4√3/3或12√3}
\ifansbookdetail
\ansnote{详解}{直线系\(A\)：\((x-3)\cos\alpha+y\sin\alpha=2\)可变形为\(\sqrt{(x-3)^{2}+y^{2}}\cdot\sin(\alpha+\varphi)=2\)（辅助角公式，\(\tan\varphi=(x-3)/y\)），\par\noindent 故\(\sin(\alpha+\varphi)=2/\sqrt{(x-3)^{2}+y^{2}}\)，而\(|\sin(\alpha+\varphi)|\leqslant1\)，得\((x-3)^{2}+y^{2}\geqslant4\)，\par\noindent 其几何意义为：直线系\(A\)是圆\((x-3)^{2}+y^{2}=4\)的切线的包络（圆上所有点的切线集合）．\par\noindent 把圆心平移至原点，圆\(x^{2}+y^{2}=r^{2}\)上一点\(P(x_0,y_0)\)的切线为\(x_0x+y_0y=r^{2}\)；令\(x_0=r\cos\alpha\)、\(y_0=r\sin\alpha\)，切线即\(x\cos\alpha+y\sin\alpha=r\)，\par\noindent 圆心\((0,0)\)到此直线距离恰为\(r\)，\par\noindent （\(|-r|/\sqrt{\cos^{2}\alpha+\sin^{2}\alpha}=r\)）故\(\alpha\in[0,2\pi)\)时该直线系为圆的全体切线．\par\noindent 直线系中的直线构成正三角形有无数个，但面积的值只有两个．取\(r=2\)：\par\noindent 设一边所在切线为\(y=2\)（上切点水平切线），另两边取斜率\(\pm\sqrt{3}\)的切线\(y=\sqrt{3}x+b\)，圆心到其距离\(|b|/\sqrt{3+1}=2\)，\(b=\pm4\)．\par\noindent ①取\(b\)使三边围成含圆心的正三角形：由\(y=2\)与\(y=\sqrt{3}x-4\)、\(y=-\sqrt{3}x-4\)（即顶点在下方）围成，此时底边\(y=2\)、\(B(-2/\sqrt{3},2)\)，边长\(|BC|=4/\sqrt{3}\)，\(S=(\sqrt{3}/4)\times(4/\sqrt{3})^{2}=4\sqrt{3}/3\)；\par\noindent ②取\(y=-2\)为底、两斜切线为\(y=\sqrt{3}x+4\)、\(y=-\sqrt{3}x+4\)围成另一正三角形：\(B'(-6/\sqrt{3},-2)\)，边长\(|B'C'|=12/\sqrt{3}=4\sqrt{3}\)，\(S=(\sqrt{3}/4)\times(12/\sqrt{3})^{2}=12\sqrt{3}\)．\par\noindent 将圆\(x^{2}+y^{2}=4\)向右平移3个单位即为\((x-3)^{2}+y^{2}=4\)，平移不改变正三角形面积，\par\noindent 故直线系\(A\)中能组成的正三角形面积为\(4\sqrt{3}/3\)或\(12\sqrt{3}\)．}
\fi
\end{ansblock}

\begin{ansblock}[2章-导-课时05-G1]
% ans:2章-导-课时05-G1
\ansitem{20}{D}
\ifansbookdetail
\ansnote{详解}{令\(y=0\)得\(-x/2=-1\)，\(x=2\)，即\(x\)轴上的截距为2；令\(x=0\)得\(y/3=-1\)，\(y=-3\)，即\(y\)轴上的截距为\(-3\)．选：D．\par\noindent （化为标准截距式\(x/2+y/(-3)=1\)同判．）}
\fi
\end{ansblock}

\begin{ansblock}[2章-导-课时05-G2]
% ans:2章-导-课时05-G2
\ansitem{21}{C}
\ifansbookdetail
\ansnote{详解}{直线过二、三、四象限，必与\(y\)轴相交（\(B\neq0\)），化斜截式\(y=-(A/B)x-(C/B)\)：\par\noindent 不过第一象限\(\iff\)斜率\(-A/B<0\)且纵截距\(-C/B<0\)\(\iff AB>0\)且\(BC>0\)，即\(A\)、\(B\)同号且\(B\)、\(C\)同号，亦即\(A\)，\(B\)，\(C\)同号．选：C．\par\noindent （检验\(C<0\)情形？\(A\)、\(B\)、\(C\)同为负与同为正化出的方程相同（两边乘\(-1\)），故“同号”完整覆盖）}
\fi
\end{ansblock}

\begin{ansblock}[2章-导-课时05-G3]
% ans:2章-导-课时05-G3
\ansitem{22}{A}
\ifansbookdetail
\ansnote{详解}{由题意 \(-A/B=5\) 且 \(A-2B+3C=0\)，取\(B\neq0\)，得\(A=-5B\)，\(C=(2B-A)/3=7B/3\)．\par\noindent 直线为 \(-5Bx+By+(7/3)B=0\)，两边除以\((-B/3)\)：\(15x-3y-7=0\)．选：A．}
\fi
\end{ansblock}

\begin{ansblock}[2章-导-课时05-G4]
% ans:2章-导-课时05-G4
\ansitem{23}{x＋2y−1＝0（答案不唯一）}
\ifansbookdetail
\ansnote{详解}{直线与两轴围成面积为\(1/4\)，只需两截距之积的绝对值为\(1/2\)（\(S=(1/2)||a||b||=(1/2)|ab|=1/4\to|ab|=1/2\)）．\par\noindent 取截距\(a=1\)、\(b=1/2\)，得\(x/1+y/(1/2)=1\)，即\(x+2y-1=0\)（满足题意；凡\(|ab|=1/2\)的截距式方程均可）．}
\fi
\end{ansblock}

\begin{ansblock}[2章-导-课时05-G5]
% ans:2章-导-课时05-G5
\ansitem{24}{−10}
\ifansbookdetail
\ansnote{详解}{在\(l_1\)上取\(M(-1,0)\)、\(N(0,-2)\)，它们关于\(P(1,0)\)的对称点分别为\(M_1(2\times1-(-1),2\times0-0)=(3,0)\)、\(N_1(2,2\times0-(-2))=(2,2)\)，两点都在\(l_2\)上：\par\noindent 代入\(l_2\)：\(4\times3+b\cdot0+c=0\)且\(4\times2+2b+c=0\)，得\(c=-12\)，\(8+2b-12=0\)，\(b=2\)．\par\noindent 所以\(b+c=2-12=-10\)．}
\fi
\end{ansblock}

\jietitle{2.2.3 两条直线的位置关系}

\begin{ansblock}[2章-练-课时06-T2]
% ans:2章-练-课时06-T2
\ansitem{1}{B}
\ifansbookdetail
\ansnote{详解}{①\(k_{AB}=(8-2)/(4-1)=2=k_1\)，但只知斜率相等，\(l_1\)与\(l_2\)可能重合，①不入选；\par\noindent ②\(k_{PQ}=0\)，\(l_1\)为水平线；\(l_2\)平行于\(x\)轴且不过\(P\)，两水平线不重合，平行，②入选；\par\noindent ③\(k_1=(-2-0)/(-5+1)=1/2\)，\(k_2=(5-3)/(0+4)=1/2\)；检验\(M(-1,0)\)不在\(l_2\)（\(l_2: y=x/2+5\)，\(x=-1\)时\(y=9/2\)）上，故平行，③入选．选：B．}
\fi
\end{ansblock}

\begin{ansblock}[2章-练-课时06-T3]
% ans:2章-练-课时06-T3
\ansitem{2}{D}
\ifansbookdetail
\ansnote{详解}{\(k_1=(-3-0)/(1-0)=-3\)；\(l_2: ax+y-2=0\)的斜率\(k_2=-a\)．\(l_1\parallel l_2\iff k_1=k_2\)，得\(-a=-3\)，\(a=3\)．重合检验：\(a=3\)时\(l_2: 3x+y-2=0\)，取\(l_1\)上点\((0,0)\)代入得\(-2\neq0\)，\((0,0)\)不在\(l_2\)上，两直线不重合，平行成立．选：D．}
\fi
\end{ansblock}

\begin{ansblock}[2章-练-课时06-T4]
% ans:2章-练-课时06-T4
\ansitem{3}{C}
\ifansbookdetail
\ansnote{详解}{\(k_1=\tan60^\circ=\sqrt{3}\)；\(k_2=(2\sqrt{3}-\sqrt{3})/(-2-1)=-\sqrt{3}/3\)．\(k_1\cdot k_2=\sqrt{3}\times(-\sqrt{3}/3)=-1\)，且\(k_1\neq k_2\)（两线不平行、更不重合），故\(l_1\perp l_2\)．选：C．}
\fi
\end{ansblock}

\begin{ansblock}[2章-练-课时06-T6]
% ans:2章-练-课时06-T6
\ansitem{4}{C}
\ifansbookdetail
\ansnote{详解}{垂直\(\iff k(k-1)+(1-k)(2k+3)=0\)．展开：\(k^2-k+2k+3-2k^2-3k=-k^2-2k+3=0\)，即\(k^2+2k-3=0\)，解得\(k=1\)或\(k=-3\)．选：C．}
\fi
\end{ansblock}

\begin{ansblock}[2章-练-课时06-T11]
% ans:2章-练-课时06-T11
\ansitem{5}{−1}
\ifansbookdetail
\ansnote{详解}{\(l_1\parallel l_2\iff 2\times1-2m=0\)，得\(m=1\)（此时\(l_1: 2x+y+2=0\)与\(l_2\)不重合）；\(l_1\perp l_3\iff 2\times1+n=0\)，得\(n=-2\)．故\(m+n=-1\)．}
\fi
\end{ansblock}

\begin{ansblock}[2章-练-课时06-T14]
% ans:2章-练-课时06-T14
\ansitem{6}{(1) m=1 或 6；(2) m=3 或 −4}
\ifansbookdetail
\ansnote{详解}{\(k_1=(2-m)/(m-4)\)，\(k_2=(m+2-2)/(-2-1)=-m/3\)．\par\noindent (1) \(k_1=k_2\)：\((2-m)/(m-4)=-m/3\)，即\(3(2-m)=-m(m-4)\)，\(m^2-7m+6=0\)，\(m=1\)或\(6\)．检验不重合：\(m=1\)时\(l_1\)过\((3,1)\)、\((0,2)\)，斜率为\(-1/3\)，\(C(1,2)\)不在\(l_1\)上；\(m=6\)时\(l_1: y=-2x+12\)，\(C(1,2)\)不在其上，均不重合，故\(m=1\)或\(6\)；\par\noindent (2) \(m=0\)时\(k_2=0\)（水平），需\(l_1\)铅直，但\(A(3,0)\)、\(B(-1,2)\)横坐标不同，斜率存在，舍；\(k_2\neq0\)时\(k_1\cdot k_2=-1\)：\((2-m)/(m-4)\cdot(-m/3)=-1\)，即\(m(2-m)=-3(m-4)\)，\(m^2-m-12=0\)，\(m=3\)或\(-4\)．验：\(m=3\)时\(k_1=1\)，\(k_2=-1\)；\(m=-4\)时\(k_1=-3/4\)，\(k_2=4/3\)，乘积均为\(-1\)，故\(m=3\)或\(-4\)．}
\fi
\end{ansblock}

\begin{ansblock}[2章-练-课时06-T1]
% ans:2章-练-课时06-T1
\ansitem{7}{B}
\ifansbookdetail
\ansnote{详解}{\(A_1A_2+B_1B_2=2a\times1+1\times[-(a+1)]=a-1=0\)，得\(a=1\)（\(a=-1\)时\(l_2: x+2=0\)铅直而\(l_1\)斜率为2，不垂直，已排除）．此时\(l_1: 2x+y-2=0\)，\(l_2: x-2y+2=0\)．由\(y=2-2x\)代入：\(x-2(2-2x)+2=0\)，\(5x=2\)，\(x=2/5\)，\(y=6/5\)．交点\((2/5, 6/5)\)，选：B．}
\fi
\end{ansblock}

\begin{ansblock}[2章-练-课时06-T5]
% ans:2章-练-课时06-T5
\ansitem{8}{A}
\ifansbookdetail
\ansnote{详解}{\(k_{BC}=2/(-8)=-1/4\)，则\(k_{AH}=4\)；\(k_{BH}=1/(-5)=-1/5\)，则\(k_{AC}=5\)．设\(A(x,y)\)：\((y-2)/(x+3)=4\)且\((y-3)/(x+6)=5\)，即\(y=4x+14\)且\(y=5x+33\)，解得\(x=-19\)，\(y=-62\)．验第三组垂直：\(k_{CH}=(2-3)/(-3+6)=-1/3\)，而\(k_{AB}=(1+62)/(2+19)=3\)，乘积为\(-1\)，成立．选：A．}
\fi
\end{ansblock}

\begin{ansblock}[2章-练-课时06-T7]
% ans:2章-练-课时06-T7
\ansitem{9}{B}
\ifansbookdetail
\ansnote{详解}{垂直\(\iff a\cdot a+(b+2)(b-2)=0\)，即\(a^2+b^2=4\)．由\(a^2+b^2\geqslant2|ab|\)得\(|ab|\leqslant2\)；取\(a=b=\sqrt{2}\)（满足\(a^2+b^2=4\)，此时两直线为\(\sqrt{2}x+(2+\sqrt{2})y+4=0\)与\(\sqrt{2}x+(\sqrt{2}-2)y-3=0\)，\(A_1A_2+B_1B_2=2+(2+\sqrt{2})(\sqrt{2}-2)=2+(2-4)=0\)，垂直成立），\(ab=2\)．故\(ab\)的最大值为2，选：B．}
\fi
\end{ansblock}

\begin{ansblock}[2章-练-课时06-T9]
% ans:2章-练-课时06-T9
\ansitem{10}{B}
\ifansbookdetail
\ansnote{详解}{联立\(2x+y-3=0\)与\(x-3y+2=0\)：由第一式\(y=3-2x\)代入第二式：\(x-3(3-2x)+2=0\)，得\(x=1\)，\(y=1\)，交点\((1,1)\)．\(k_1=-2\)，所求直线斜率\(k=1/2\)．\(y-1=(1/2)(x-1)\)，即\(x-2y+1=0\)，选：B．}
\fi
\end{ansblock}

\begin{ansblock}[2章-练-课时06-T12]
% ans:2章-练-课时06-T12
\ansitem{11}{1}
\ifansbookdetail
\ansnote{详解}{四边形由\(l_1\)、\(l_2\)与\(x\)轴、\(y\)轴围成，坐标轴交点处的内角为\(90^\circ\)，四边形有外接圆\(\iff\)对角互补\(\iff l_1\)与\(l_2\)的交角为\(90^\circ\)，即\(l_1\perp l_2\)．\(k_1=-1/3\)，\(k_2=3k\)，垂直\(\iff(-1/3)\times3k=-1\)，得\(k=1\)．（检验：\(l_1\)不过原点、不平行于坐标轴；\(k=1\)时\(l_2: 3x-y+1=0\)同样，四边形存在且对角互补，有外接圆．）}
\fi
\end{ansblock}

\begin{ansblock}[2章-练-课时06-T15]
% ans:2章-练-课时06-T15
\ansitem{12}{(1) D(−1,6)；(2) 是菱形}
\ifansbookdetail
\ansnote{详解}{(1) 平行四边形对角线互相平分：\(AC\)中点与\(BD\)中点重合．设\(D(x,y)\)：\(((1+3)/2, (2+4)/2)=((5+x)/2, (0+y)/2)\)，得\(x=-1\)，\(y=6\)，即\(D(-1,6)\)；\par\noindent (2) 邻边相等即菱形：\(|AB|^2=16+4=20\)，\(|BC|^2=4+16=20\)，\(|AB|=|BC|\)，故\pxparallelogram{}ABCD 为菱形．（或对角线：\(k_{AC}=1\)，\(k_{BD}=6/(-6)=-1\)，\(AC\perp BD\)且互相平分，也为菱形．）}
\fi
\end{ansblock}

\begin{ansblock}[2章-练-课时06-T13]
% ans:2章-练-课时06-T13
\ansitem{13}{(1,1)；√10}
\ifansbookdetail
\ansnote{详解}{按\(\lambda\)整理：\((x+y-2)+\lambda(3x+2y-5)=0\)，令两因式同时为零：\(x+y-2=0\)且\(3x+2y-5=0\)，解得\(x=y=1\)，定点\(Q(1,1)\)．\(\lambda\)变化时\(l\)绕\(Q\)转动（斜率可取除\(-3/2\)外的一切值，且\(\lambda=-1/2\)时为铅直线\(x=1\)，方向覆盖齐全）．\(d\)最大\(\iff l\perp PQ\)，最大距离\(|PQ|=\sqrt{(-2-1)^2+(0-1)^2}=\sqrt{10}\)．}
\fi
\end{ansblock}

\begin{ansblock}[2章-练-课时06-T8]
% ans:2章-练-课时06-T8
\ansitem{14}{ACD}
\ifansbookdetail
\ansnote{详解}{A：\(l_2\)按\(a\)整理：\(a(x-2y)+(3y-1)=0\)，令\(x-2y=0\)且\(3y-1=0\)，得\((2/3, 1/3)\)，恒过，A对；\par\noindent B：\(a=1\)时\(l_1: x+y-1=0\)，\(l_2: x+y-1=0\)，重合非平行；\(a=-3\)时\(k_1=-1/3\)，\(l_2: -3x+9y-1=0\)的斜率为\(1/3\)，不平行，B错；\par\noindent C：\(A_1A_2+B_1B_2=a+a(3-2a)=2a(2-a)=0\)，得\(a=0\)或\(2\)，C对；\par\noindent D：\(a>0\)时\(l_1: y=-x/a+1\)，斜率为负、纵截距\(1>0\)、横截距\(a>0\)，图象过一、二、四象限，不过第三象限，D对．选：ACD．}
\fi
\end{ansblock}

\begin{ansblock}[2章-练-课时06-T10]
% ans:2章-练-课时06-T10
\ansitem{15}{BCD}
\ifansbookdetail
\ansnote{详解}{\(AB\)中点\((-2,2)\)，\(k_{AB}=1\)，中垂线\(y-2=-(x+2)\)，即\(x+y=0\)．与欧拉线联立：\(x=-1\)，\(y=1\)，外心\(O'(-1,1)\)，\(R^2=9+1=10\)．\par\noindent 设\(C(m,n)\)：重心\(((m-4)/3, (n+4)/3)\)在欧拉线上，得\(m-n-2=0\)；\(|O'C|^2=10\)，即\((m+1)^2+(n-1)^2=10\)，展开\(m^2+n^2+2m-2n=8\)．代入\(n=m-2\)：\(2m^2-4m=0\)，\(m=0\)或\(2\)，\(C(2,0)\)或\((0,-2)\)，C对；\par\noindent 重心为\((-2/3, 4/3)\)或\((-4/3, 2/3)\)，A所给数值错，A错；\par\noindent 垂心：由欧拉关系\(H=3G-2O'\)（\(G\)为重心）：\(C(2,0)\)时\(H=3(-2/3, 4/3)-2(-1,1)=(0,2)\)；\(C(0,-2)\)时\(H=3(-4/3, 2/3)-2(-1,1)=(-2,0)\)．（直接验证\(C(2,0)\)情形：\(AC\)水平，过\(B\)的高为\(x=0\)；过\(C\)垂直\(AB\)的高为\(y=-x+2\)，交点\((0,2)\)．）B对；\par\noindent 面积比：\(AB\)即\(x-y+4=0\)与欧拉线平行，\(C(2,0)\)到欧拉线距离\(4/\sqrt{2}\)，到\(AB\)距离\(6/\sqrt{2}\)，相似比\(2/3\)，小三角形与大三角形面积比\(4/9\)，故两部分之比\(4:5\)，即\(4/5\)，D对．选：BCD．}
\fi
\end{ansblock}

\begin{ansblock}[2章-练-课时06-T16]
% ans:2章-练-课时06-T16
\ansitem{16}{(1) 5x+y−13=0；(2) 7/2}
\ifansbookdetail
\ansnote{详解}{(1) \(k_{BC}=(-2-3)/(3-2)=-5\)，\(y-3=-5(x-2)\)，即\(5x+y-13=0\)；\par\noindent (2) \(|BC|=\sqrt{1+25}=\sqrt{26}\)，\(A\)到\(BC\)的距离\(d=|5\times1+1-13|/\sqrt{26}=7/\sqrt{26}\)，故\(S=(1/2)\times\sqrt{26}\times(7/\sqrt{26})=7/2\)．}
\fi
\end{ansblock}

\begin{ansblock}[2章-导-课时06-G1]
% ans:2章-导-课时06-G1
\ansitem{17}{D}
\ifansbookdetail
\ansnote{详解}{\(k_1=\tan135^\circ=-1\)；\(k_2=(-6-(-1))/(3-(-2))=-5/5=-1\)．两斜率相等，但仅凭斜率相等不能排除重合（题面未给两点异于\(l_1\)的信息），故\(l_1\)与\(l_2\)平行或重合．选：D．}
\fi
\end{ansblock}

\begin{ansblock}[2章-导-课时06-G2]
% ans:2章-导-课时06-G2
\ansitem{18}{B}
\ifansbookdetail
\ansnote{详解}{平行\(\iff a(a-1)-2\times1=0\)，即\(a^2-a-2=0\)，\(a=2\)或\(a=-1\)．检验：\(a=-1\)时两直线分别为\(-x+2y-1=0\)与\(x-2y+1=0\)，系数成比例，重合，舍；\(a=2\)时两直线\(2x+2y-1=0\)与\(x+y+4=0\)平行不重合．故\(a=2\)，选：B．}
\fi
\end{ansblock}

\begin{ansblock}[2章-导-课时06-G3]
% ans:2章-导-课时06-G3
\ansitem{19}{C}
\ifansbookdetail
\ansnote{详解}{联立\(y=2x+10\)与\(y=x+1\)：\(x=-9\)，\(y=-8\)，交点\((-9,-8)\)．代入\(y=ax-2\)：\(-8=-9a-2\)，得\(a=2/3\)，选：C．}
\fi
\end{ansblock}

\begin{ansblock}[2章-导-课时06-G5]
% ans:2章-导-课时06-G5
\ansitem{20}{3x+y+1=0}
\ifansbookdetail
\ansnote{详解}{设\(l\)与\(l_1\)、\(l_2\)分别交于\(A\)、\(B\)，\(P(-1,2)\)为\(AB\)中点．设\(A=P+(u,v)\)，则\(B=P-(u,v)\)．\(A\in l_1\)：\(4(-1+u)+(2+v)+3=0\)，即\(4u+v+1=0\)；\(B\in l_2\)：\(3(-1-u)-5(2-v)-5=0\)，即\(-3u+5v-18=0\)．由\(v=-4u-1\)代入第二式：\(-23u-23=0\)，\(u=-1\)，\(v=3\)．故\(l\)的斜率\(k=v/u=-3\)，过\(P\)：\(y-2=-3(x+1)\)，即\(3x+y+1=0\)．（验：\(A(-2,5)\)代入\(l_1\)：\(-8+5+3=0\)，\(B(0,-1)\)代入\(l_2\)：\(0+5-5=0\)，中点为\((-1,2)=P\)，且\(A\)、\(B\)均在\(l\)上．）}
\fi
\end{ansblock}

\begin{ansblock}[2章-导-课时06-G4]
% ans:2章-导-课时06-G4
\ansitem{21}{−2}
\ifansbookdetail
\ansnote{详解}{\(x=1-2y\)即\(x+2y-1=0\)；\(2x+4y+m=0\)即\(x+2y+m/2=0\)．重合\(\iff\)两方程完全相同，故\(m/2=-1\)，\(m=-2\)．}
\fi
\end{ansblock}

\jietitle{2.2.4 点到直线距离与平行线距离（选学）}

\begin{ansblock}[2章-练-课时06B-T1]
% ans:2章-练-课时06B-T1
\ansitem{1}{C}
\ifansbookdetail
\ansnote{详解}{设点\((0,t)\)，\(d=|4t+5|/\sqrt{3^{2}+4^{2}}=|4t+5|/5=5\)，即\(|4t+5|=25\)，解得\(t=5\)或\(t=-15/2\)。所求点为\((0,5)\)或\((0,-15/2)\)，故选：C．（D 中\((5,0)\)、\((-5,0)\)到直线的距离分别为4与2，不合，系轴别混淆干扰项．）}
\fi
\end{ansblock}

\begin{ansblock}[2章-练-课时06B-T3]
% ans:2章-练-课时06B-T3
\ansitem{2}{D}
\ifansbookdetail
\ansnote{详解}{以\(BC\)为底：\(|BC|=\sqrt{(-1-2)^{2}+(0-1)^{2}}=\sqrt{10}\)；直线\(BC\)过\(B(2,1)\)、\(C(-1,0)\)，斜率为\(1/3\)，方程\(x-3y+1=0\)。\(A(1,3)\)到\(BC\)的高\(h=|1-9+1|/\sqrt{1^{2}+(-3)^{2}}=7/\sqrt{10}\)。故\(S=(1/2)\times\sqrt{10}\times 7/\sqrt{10}=7/2\)，选：D．}
\fi
\end{ansblock}

\begin{ansblock}[2章-练-课时06B-T9]
% ans:2章-练-课时06B-T9
\ansitem{3}{C}
\ifansbookdetail
\ansnote{详解}{\(k=\tan 120^\circ=-\sqrt{3}\)，直线\(l\)：\(y-\sqrt{3}=-\sqrt{3}(x-2)\)，即\(\sqrt{3}x+y-3\sqrt{3}=0\)。\(d=|\sqrt{3}\times 1-\sqrt{3}-3\sqrt{3}|/\sqrt{(\sqrt{3})^{2}+1^{2}}=3\sqrt{3}/2\)，故选：C．}
\fi
\end{ansblock}

\begin{ansblock}[2章-练-课时06B-T12]
% ans:2章-练-课时06B-T12
\ansitem{4}{4/25}
\ifansbookdetail
\ansnote{详解}{\(m^{2}+n^{2}=|OM|^{2}\)，其最小值在\(OM\)垂直于\(l\)时取得，即原点到\(l\)的距离：\(d=|0+0+2|/\sqrt{3^{2}+4^{2}}=2/5\)。故\(m^{2}+n^{2}\)的最小值为\(d^{2}=4/25\)．}
\fi
\end{ansblock}

\begin{ansblock}[2章-练-课时06B-T13]
% ans:2章-练-课时06B-T13
\ansitem{5}{2+√2}
\ifansbookdetail
\ansnote{详解}{\(d=|\cos\theta+\sin\theta-2|/\sqrt{\cos^{2}\theta+\sin^{2}\theta}\)\par\noindent \(=|\sqrt{2}\sin(\theta+\pi/4)-2|\)。\par\noindent 因\(\sqrt{2}\sin(\theta+\pi/4)\leqslant\sqrt{2}<2\)，绝对值内恒为负。\par\noindent 故\(d=2-\sqrt{2}\sin(\theta+\pi/4)\)，当\(\sin(\theta+\pi/4)=-1\)时取最大值\(2+\sqrt{2}\)．}
\fi
\end{ansblock}

\begin{ansblock}[2章-练-课时06B-T5]
% ans:2章-练-课时06B-T5
\ansitem{6}{C}
\ifansbookdetail
\ansnote{详解}{两线过\(P\)、\(Q\)且始终平行。当两线均与\(PQ\)垂直时距离最大，最大值\(=|PQ|=\sqrt{(2+1)^{2}+(-1-3)^{2}}=5\)；当两线向重合方向旋转时距离趋于0，但重合时不能分别过两不同点，0取不到。故\(d\in(0,5]\)，故选：C．}
\fi
\end{ansblock}

\begin{ansblock}[2章-练-课时06B-T6]
% ans:2章-练-课时06B-T6
\ansitem{7}{C}
\ifansbookdetail
\ansnote{详解}{两线平行，\(d=|a-b|/\sqrt{2}\)。由韦达定理\(a+b=-1\)，\(ab=c\)，故\(|a-b|=\sqrt{(a+b)^{2}-4ab}=\sqrt{1-4c}\)。由\(0\leqslant c\leqslant 1/8\)得\(1/2\leqslant 1-4c\leqslant 1\)，故\(\sqrt{2}/2\leqslant|a-b|\leqslant 1\)。于是\(d\in[1/2,\ \sqrt{2}/2]\)：最大值\(\sqrt{2}/2\)，最小值\(1/2\)，故选：C．（实根存在性：\(\Delta=1-4c\geqslant 0\)由\(c\leqslant 1/8\)保证．）}
\fi
\end{ansblock}

\begin{ansblock}[2章-练-课时06B-T8]
% ans:2章-练-课时06B-T8
\ansitem{8}{B}
\ifansbookdetail
\ansnote{详解}{两线平行，\(|PQ|\)等于平行线间的距离，为定值\(=|2-(-1)|/\sqrt{2}=3\sqrt{2}/2\)；且向量\(\overrightarrow{QP}\)恒为\(l_1\)指向\(l_2\)的法向位移\((3/2,\ 3/2)\)。作平移\(A'=A+(3/2,\ 3/2)=(-3/2,\ -3/2)\)，则\(AA'\parallel QP\)且\(|AA'|=|QP|\)，四边形\(AA'QP\)为平行四边形，故\(|AP|=|A'Q|\)。于是\(|AP|+|PQ|+|QB|=|A'Q|+3\sqrt{2}/2+|QB|\geqslant|A'B|+3\sqrt{2}/2\)，当\(A'\)、\(Q\)、\(B\)共线（\(Q\)在线段上）时取等。\(|A'B|=\sqrt{(3/2+3/2)^{2}+(1/2+3/2)^{2}}=\sqrt{13}\)；取等点验证：线段\(A'B\)上参数\(t=4/5\)处的点\((9/10,\ 1/10)\)满足\(9/10+1/10-1=0\)，确在\(l_2\)上，等号可取。故最小值\(=\sqrt{13}+3\sqrt{2}/2\)，故选：B．}
\fi
\end{ansblock}

\begin{ansblock}[2章-练-课时06B-T10]
% ans:2章-练-课时06B-T10
\ansitem{9}{AB}
\ifansbookdetail
\ansnote{详解}{平行\(\iff 1/2=(-2)/n\)（\(n\neq 0\)），得\(n=-4\)。\(l_2\)：\(2x-4y-6=0\)，即\(x-2y-3=0\)。\(d=|m-(-3)|/\sqrt{1^{2}+(-2)^{2}}=|m+3|/\sqrt{5}=2\sqrt{5}\)，得\(|m+3|=10\)，\(m=7\)或\(m=-13\)。故\(m+n=3\)或\(-17\)，故选：AB．}
\fi
\end{ansblock}

\begin{ansblock}[2章-练-课时06B-T11]
% ans:2章-练-课时06B-T11
\ansitem{10}{2x+3y−8=0 或 6x+9y−10=0}
\ifansbookdetail
\ansnote{详解}{\(l_1\)化为\(4x+6y-2=0\)。\(l\)与两线均平行，设\(l\)：\(4x+6y+C=0\)（\(C\neq -2\)且\(C\neq -9\)）。由\(|C+2|/\sqrt{52}=2\times|C+9|/\sqrt{52}\)，即\(|C+2|=2|C+9|\)：\(C+2=2(C+9)\)得\(C=-16\)；\(C+2=-2(C+9)\)得\(C=-20/3\)。故\(l\)：\(4x+6y-16=0\)即\(2x+3y-8=0\)（带域外侧），或\(4x+6y-20/3=0\)即\(6x+9y-10=0\)（带域内侧，此时到\(l_1\)、\(l_2\)的距离分别为\(14/3\)与\(7/3\)，比值确为\(2:1\)）。两解均保留．}
\fi
\end{ansblock}

\begin{ansblock}[2章-练-课时06B-T14]
% ans:2章-练-课时06B-T14
\ansitem{11}{1}
\ifansbookdetail
\ansnote{详解}{设\(A(m,n)\)，\(B(a,b)\)，则\(A\)在直线\(3x+4y-6=0\)上，\(B\)在直线\(3x+4y-1=0\)上，所求即\(|AB|\)。两线平行，当\(AB\)沿公垂方向时最短，最小值为平行线间的距离：\(d=|-6-(-1)|/\sqrt{3^{2}+4^{2}}=5/5=1\)．}
\fi
\end{ansblock}

\begin{ansblock}[2章-练-课时06B-T2]
% ans:2章-练-课时06B-T2
\ansitem{12}{C}
\ifansbookdetail
\ansnote{详解}{\(A\)关于反射面\(x\)轴的对称点为\(A'(-3,-5)\)。由反射角等于入射角，光路\(A\to M\to B\)的总长等于直线段\(|A'B|\)的最短值\(|A'B|=\sqrt{(2+3)^{2}+(10+5)^{2}}=\sqrt{250}=5\sqrt{10}\)，故选：C．}
\fi
\end{ansblock}

\begin{ansblock}[2章-练-课时06B-T4]
% ans:2章-练-课时06B-T4
\ansitem{13}{B}
\ifansbookdetail
\ansnote{详解}{\(A(3,1)\)关于\(y=x\)的对称点\(A_1(1,3)\)（对称式\((x,y)\to(y,x)\)），关于\(x\)轴的对称点\(A_2(3,-1)\)。则\(|AC|=|A_1C|\)，\(|AB|=|A_2B|\)，周长\(=|A_1C|+|CB|+|BA_2|\geqslant|A_1A_2|\)。取等检验：直线\(A_1A_2\)为\(y=-2x+5\)，与\(y=x\)交于\((5/3,\ 5/3)\)、与\(x\)轴交于\((5/2,\ 0)\)，均落在允许动点范围内，等号可取。最小值\(=|A_1A_2|=\sqrt{(3-1)^{2}+(-1-3)^{2}}=\sqrt{20}=2\sqrt{5}\)，故选：B．}
\fi
\end{ansblock}

\begin{ansblock}[2章-练-课时06B-T7]
% ans:2章-练-课时06B-T7
\ansitem{14}{D}
\ifansbookdetail
\ansnote{详解}{对称\(\iff AB\perp l\)且\(AB\)的中点在\(l\)上。\(k_{AB}=(-b-(b+2))/((b-a)-(a+2))=(-2b-2)/(b-2a-2)=3/4\)，整理得\(-8b-8=3b-6a-6\)，即\(6a-11b-2=0\)①。中点坐标\(((b+2)/2,\ 1)\)，代入\(4x+3y-11=0\)：\(2(b+2)+3-11=0\)，得\(b=2\)；代回①：\(6a-22-2=0\)，\(a=4\)。故\(a=4\)，\(b=2\)，故选：D．}
\fi
\end{ansblock}

\begin{ansblock}[2章-练-课时06B-T15]
% ans:2章-练-课时06B-T15
\ansitem{15}{（1）3√5/5；（2）C(−13/5, 31/5)。}
\ifansbookdetail
\ansnote{详解}{（1）\(d=|2\times 1+2-1|/\sqrt{2^{2}+1^{2}}=3/\sqrt{5}=3\sqrt{5}/5\)。\par\noindent （2）内角平分线性质：\(A\)关于\(CD\)的对称点\(A'\)落在边\(BC\)所在直线上。设\(A'(x,y)\)：中点\(((x+1)/2,\ (y+2)/2)\)在\(CD\)上，得\((x+1)+(y+2)/2-1=0\)，即\(2x+y+2=0\)；\(AA'\perp CD\)，\(k_{CD}=-2\)，故\(k_{AA'}=1/2\)，得\((y-2)/(x-1)=1/2\)，即\(x-2y+3=0\)。\par\noindent 联立解得\(x=-7/5\)，\(y=4/5\)，即\(A'(-7/5,\ 4/5)\)。直线\(BC\)过\(B(-1,-1)\)与\(A'\)，斜率\(=(4/5+1)/(-7/5+1)=(9/5)/(-2/5)=-9/2\)，方程\(y+1=-(9/2)(x+1)\)，即\(9x+2y+11=0\)。\par\noindent \(C\)为\(BC\)与\(CD\)的交点：把\(y=1-2x\)代入，\(9x+2(1-2x)+11=0\)，\(5x+13=0\)，\(x=-13/5\)，\(y=1-2\times(-13/5)=31/5\)。故\(C(-13/5,\ 31/5)\)．}
\fi
\end{ansblock}

\begin{ansblock}[2章-练-课时06B-T16]
% ans:2章-练-课时06B-T16
\ansitem{16}{能，P(1/9, 37/18)。}
\ifansbookdetail
\ansnote{详解}{设\(P(x_0,\ y_0)\)，记\(u=2x_0-y_0\)。\par\noindent 条件（3）：\(|u+3|/\sqrt{5}=(\sqrt{2}/\sqrt{5})\cdot|x_0+y_0-1|/\sqrt{2}\)，即\(|u+3|=|x_0+y_0-1|\)，平方去绝对值：\((u+3)^{2}-(x_0+y_0-1)^{2}=(3x_0+2)(x_0-2y_0+4)=0\)。\(x_0=-2/3<0\)与第一象限矛盾，舍；故\(x_0-2y_0+4=0\)①。\par\noindent 条件（2）：\(d_1=|u+3|/\sqrt{5}\)，\(d_2=|-2u+1|/(2\sqrt{5})\)，由\(d_1=(1/2)d_2\)得\(4|u+3|=|1-2u|\)，平方：\((6u+11)(2u+13)=0\)，\(u=-11/6\)或\(u=-13/2\)，即\(2x_0-y_0+11/6=0\)②或\(2x_0-y_0+13/2=0\)③。\par\noindent 联立逐一检验：②与①：\(y_0=2x_0+11/6\)，代入①：\(x_0-4x_0-11/3+4=0\)，\(x_0=1/9>0\)，\(y_0=2/9+11/6=37/18>0\)，在第一象限；③与①：\(y_0=2x_0+13/2\)，代入①：\(x_0-4x_0-13+4=0\)，\(x_0=-3<0\)，舍。\par\noindent 终检\(P(1/9,\ 37/18)\)：\(d_1=(21/18)/\sqrt{5}=7/(6\sqrt{5})\)，\(d_2=(14/3)/(2\sqrt{5})=7/(3\sqrt{5})\)，\(d_1=(1/2)d_2\)成立；\(d_3=(21/18)/\sqrt{2}=7/(6\sqrt{2})\)，\(d_1:d_3=\sqrt{2}:\sqrt{5}\)成立。故存在满足三条件的点\(P(1/9,\ 37/18)\)．}
\fi
\end{ansblock}

\begin{ansblock}[2章-导-课时06B-G1]
% ans:2章-导-课时06B-G1
\ansitem{17}{D}
\ifansbookdetail
\ansnote{详解}{\(d=|m\cdot 1+4-1|/\sqrt{m^{2}+1}=|m+3|/\sqrt{m^{2}+1}=3\)，两边非负可平方：\(m^{2}+6m+9=9m^{2}+9\)，即\(8m^{2}-6m=0\)，\(m=0\)或\(m=3/4\)。检验：\(m=0\)时直线为\(y-1=0\)，\(d=|4-1|=3\)成立；\(m=3/4\)时\(d=(15/4)/(5/4)=3\)成立。两根均保留，故选：D．}
\fi
\end{ansblock}

\begin{ansblock}[2章-导-课时06B-G2]
% ans:2章-导-课时06B-G2
\ansitem{18}{A}
\ifansbookdetail
\ansnote{详解}{\(d=|3\times 0+4\times 0-5|/\sqrt{3^{2}+4^{2}}=5/5=1\)，故选：A．}
\fi
\end{ansblock}

\begin{ansblock}[2章-导-课时06B-G3]
% ans:2章-导-课时06B-G3
\ansitem{19}{B}
\ifansbookdetail
\ansnote{详解}{平行\(\iff 1\times 3-a(a-2)=0\)且系数不成比例（排除重合），由\(a^{2}-2a-3=0\)得\(a=3\)或\(a=-1\)；\(a=3\)时两方程同为\(x+3y+6=0\)（重合），舍，故\(a=-1\)。此时\(l_1\)：\(x-y+6=0\)，\(l_2\)：\(-3x+3y-2=0\)即\(x-y+2/3=0\)，\(d=|6-2/3|/\sqrt{1^{2}+(-1)^{2}}=(16/3)/\sqrt{2}=8\sqrt{2}/3\)，故选：B．}
\fi
\end{ansblock}

\begin{ansblock}[2章-导-课时06B-G4]
% ans:2章-导-课时06B-G4
\ansitem{20}{√13}
\ifansbookdetail
\ansnote{详解}{两线\(x\)、\(y\)系数已相同，直接套平行线间的距离公式：\(d=|5-(-8)|/\sqrt{2^{2}+(-3)^{2}}=13/\sqrt{13}=\sqrt{13}\)．}
\fi
\end{ansblock}

\begin{ansblock}[2章-导-课时06B-G5]
% ans:2章-导-课时06B-G5
\ansitem{21}{m＝±√10}
\ifansbookdetail
\ansnote{详解}{\(y=3x+6\)化为一般式\(3x-y+6=0\)，\(d=|3m-6+6|/\sqrt{3^{2}+(-1)^{2}}=3|m|/\sqrt{10}=3\)，得\(|m|=\sqrt{10}\)，即\(m=\pm\sqrt{10}\)．}
\fi
\end{ansblock}

\jietitle{2.3 圆的方程}

\begin{ansblock}[2章-练-课时07-T1]
% ans:2章-练-课时07-T1
\ansitem{1}{D}
\ifansbookdetail
\ansnote{详解}{\((5a+1-1)^{2}+(12a)^{2}=25a^{2}+144a^{2}=169a^{2}<1\)，即\(|a|<1/13\)，故选：D．}
\fi
\end{ansblock}

\begin{ansblock}[2章-练-课时07-T2]
% ans:2章-练-课时07-T2
\ansitem{2}{D}
\ifansbookdetail
\ansnote{详解}{\((4a-1+1)^{2}+(3a+2-2)^{2}=16a^{2}+9a^{2}=25a^{2}\leqslant 25\)，即\(a^{2}\leqslant 1\)，\(-1\leqslant a\leqslant 1\)，故选：D．}
\fi
\end{ansblock}

\begin{ansblock}[2章-练-课时07-T3]
% ans:2章-练-课时07-T3
\ansitem{3}{C}
\ifansbookdetail
\ansnote{详解}{表圆需\(m>0\)。原点在圆外\(\iff(0-3)^{2}+(0+4)^{2}=25>m\)。故\(0<m<25\)，故选：C．}
\fi
\end{ansblock}

\begin{ansblock}[2章-练-课时07-T9]
% ans:2章-练-课时07-T9
\ansitem{4}{A}
\ifansbookdetail
\ansnote{详解}{圆心\((3,-5)\)。\(d=|12+15-2|/5=25/5=5\)。圆上点到直线的最短距离\(=d-r\)（\(d>r\)时）\(=1\)，得\(r=4\)，故选：A．}
\fi
\end{ansblock}

\begin{ansblock}[2章-练-课时07-T12]
% ans:2章-练-课时07-T12
\ansitem{5}{(x−1)²+y²=9（答案不唯一）}
\ifansbookdetail
\ansnote{详解}{圆心\((a,0)\)，\(r^{2}=(1-a)^{2}+9\)。取\(a=1\)：\(r=3\)，得\((x-1)^{2}+y^{2}=9\)。（一般形式\((x-a)^{2}+y^{2}=(1-a)^{2}+9\)，\(a\in\mathbf{R}\)均合．）}
\fi
\end{ansblock}

\begin{ansblock}[2章-练-课时07-T7]
% ans:2章-练-课时07-T7
\ansitem{6}{ABC}
\ifansbookdetail
\ansnote{详解}{圆心\((2,0)\)。圆关于圆心中心对称，A对；圆心在\(x\)轴上，故关于\(y=0\)对称，B对；\(x+3y-2=0\)过\((2,0)\)（\(2+0-2=0\)成立），过圆心的直线是圆的对称轴，C对；\(x-y+2=0\)：\(2-0+2=4\neq 0\)，不过圆心，D错。故选：ABC．}
\fi
\end{ansblock}

\begin{ansblock}[2章-练-课时07-T10]
% ans:2章-练-课时07-T10
\ansitem{7}{CD}
\ifansbookdetail
\ansnote{详解}{圆心\((1,-2)\)。\(l\)平分圆\(\iff l\)过圆心。截距相等分两类：①截距均为0：\(l\)过原点，\(k=-2/1=-2\)，\(l\)：\(2x+y=0\)（D对）；②截距\(a\neq 0\)：\(x/a+y/a=1\)，代入\((1,-2)\)：\((1-2)/a=1\)，\(a=-1\)，\(l\)：\(x+y+1=0\)（C对）。检验：A：\(1+2+1=4\neq 0\)，不过圆心；B：\(1+4=5\neq 0\)，不过圆心。故选：CD．}
\fi
\end{ansblock}

\begin{ansblock}[2章-练-课时07-T13]
% ans:2章-练-课时07-T13
\ansitem{8}{64；4；[3−2√2, 3+2√2]}
\ifansbookdetail
\ansnote{详解}{\((m+2)^{2}+(n+1)^{2}=|P-(-2,\ -1)|^{2}\)。圆心\(C(2,2)\)，\(r=3\)；\(|C-(-2,\ -1)|=\sqrt{16+9}=5\)。\par\noindent 距离范围\([5-3,\ 5+3]=[2,8]\)，平方范围\([4,64]\)，最大64、最小4。\par\noindent \(\sqrt{m^{2}+n^{2}}=|PO|\)，\(|OC|=2\sqrt{2}<3\)，故范围为\([3-2\sqrt{2},\ 3+2\sqrt{2}]\)．}
\fi
\end{ansblock}

\begin{ansblock}[2章-练-课时07-T4]
% ans:2章-练-课时07-T4
\ansitem{9}{B}
\ifansbookdetail
\ansnote{详解}{\(x^{2}\)与\(y^{2}\)系数均为1。\(D^{2}+E^{2}-4F=4a^{2}+16a^{2}+40a=20a(a+2)>0\)，即\(a>0\)或\(a<-2\)，故选：B．}
\fi
\end{ansblock}

\begin{ansblock}[2章-练-课时07-T5]
% ans:2章-练-课时07-T5
\ansitem{10}{A}
\ifansbookdetail
\ansnote{详解}{表圆：\(m^{2}+(-2)^{2}-4\times 2>0\)，即\(m^{2}>4\)，\(m>2\)或\(m<-2\)。点在圆外：\(1+4+m-4+2=m+3>0\)，即\(m>-3\)。取交：\((-3,-2)\cup(2,+\infty)\)，故选：A．}
\fi
\end{ansblock}

\begin{ansblock}[2章-练-课时07-T6]
% ans:2章-练-课时07-T6
\ansitem{11}{C}
\ifansbookdetail
\ansnote{详解}{过原点\(\iff 2m^{2}-6m+4=0\)，即\(m=1\)或\(m=2\)。表圆条件：\(D^{2}+E^{2}-4F=8(m-1)^{2}-4(2m^{2}-6m+4)=8m-8>0\)，即\(m>1\)。舍\(m=1\)，得\(m=2\)，故选：C．}
\fi
\end{ansblock}

\begin{ansblock}[2章-练-课时07-T8]
% ans:2章-练-课时07-T8
\ansitem{12}{D}
\ifansbookdetail
\ansnote{详解}{配方：\((x+1)^{2}+(y-m)^{2}=m^{2}+4m+5=(m+2)^{2}+1\)。\(r\) 最小为1当\(m=-2\)，此时圆心\((-1,-2)\)。\(|OC|=\sqrt{5}>r\)，圆上点到\(O\)距离最大\(=|OC|+r=\sqrt{5}+1\)，故选：D．}
\fi
\end{ansblock}

\begin{ansblock}[2章-练-课时07-T14]
% ans:2章-练-课时07-T14
\ansitem{13}{(−2,−4)；5}
\ifansbookdetail
\ansnote{详解}{表圆需\(a^{2}=a+2\)，即\(a=-1\)或\(a=2\)。\(a=2\)：\(4x^{2}+4y^{2}+4x+8y+10=0\)，除以4配方得半径平方为负，舍。\par\noindent \(a=-1\)：\(x^{2}+y^{2}+4x+8y-5=0\)，即\((x+2)^{2}+(y+4)^{2}=25\)。圆心\((-2,-4)\)，半径5．}
\fi
\end{ansblock}

\begin{ansblock}[2章-练-课时07-T11]
% ans:2章-练-课时07-T11
\ansitem{14}{x²+y²−4x+6y+8=0}
\ifansbookdetail
\ansnote{详解}{圆心在\(AB\)的中垂线\(y=-3\)上，又在\(2x-y-7=0\)上：\(2x+3-7=0\)，\(x=2\)，圆心\((2,-3)\)。\(r^{2}=4+(-3+4)^{2}=5\)。标准方程\((x-2)^{2}+(y+3)^{2}=5\)，展开：\(x^{2}+y^{2}-4x+6y+8=0\)。（验：\(A(0,-4)\)：\(16-24+8=0\)成立；\(B(0,-2)\)：\(4-12+8=0\)成立．）}
\fi
\end{ansblock}

\begin{ansblock}[2章-练-课时07-T16]
% ans:2章-练-课时07-T16
\ansitem{15}{x²+y²+2x−4y+3=0}
\ifansbookdetail
\ansnote{详解}{圆心\((-D/2,\ -E/2)\)在\(x+y-1=0\)上：\(-D/2-E/2-1=0\)，即\(D+E=-2\)。\par\noindent \(r^{2}=(D^{2}+E^{2}-12)/4=2\)，即\(D^{2}+E^{2}=20\)。由\(E=-2-D\)：\(D^{2}+(D+2)^{2}=20\)，\(2D^{2}+4D-16=0\)，\(D^{2}+2D-8=0\)，\(D=2\)或\(D=-4\)。\par\noindent \(D=2\)：\(E=-4\)，圆心\((-1,2)\)在第二象限成立；\(D=-4\)：\(E=2\)，圆心\((2,-1)\)在第四象限，舍。故\(D=2\)，\(E=-4\)，圆的一般方程为\(x^{2}+y^{2}+2x-4y+3=0\)。（验：\(r^{2}=(4+16-12)/4=2\)成立．）}
\fi
\end{ansblock}

\begin{ansblock}[2章-练-课时07-T15]
% ans:2章-练-课时07-T15
\ansitem{16}{(1) a≠−1；(2) 恒过 (0,0) 与 (16/5, −8/5)；(3) a=1/4}
\ifansbookdetail
\ansnote{详解}{(1) \(a\neq -1\)时\(x^{2}\)、\(y^{2}\)系数同为\(1+a\neq 0\)，两边除以\((1+a)\)化一般式，\(D^{2}+E^{2}-4F=(16+64a^{2})/(1+a)^{2}>0\)恒成立，表圆；\(a=-1\)时二次项消失，方程退化为直线\(-4x-8y=0\)，不合。故\(a\neq -1\)。\par\noindent (2) 方程按\(a\)整理：\((x^{2}+y^{2}-4x)+a(x^{2}+y^{2}+8y)=0\)。两定点须对一切\(a\)成立\(\iff x^{2}+y^{2}-4x=0\)且\(x^{2}+y^{2}+8y=0\)。相减：\(4x+8y=0\)，\(x=-2y\)。\par\noindent 代入\(x^{2}+y^{2}-4x=0\)：\(5y^{2}=0\)，得\((0,0)\)；代入\(x^{2}+y^{2}+8y=0\)：\(5y^{2}+8y=0\)，\(y=0\)（已得）或\(y=-8/5\)，\(x=16/5\)。定点\((0,0)\)与\((16/5,\ -8/5)\)。（验\((16/5,\ -8/5)\)：\(x^{2}+y^{2}=320/25\)，\(4x=64/5=320/25\)，\(8y=-64/5\)，两式皆零。）\par\noindent (3) 表圆时\(r^{2}=(4+16a^{2})/(1+a)^{2}\)。令\(f(a)=(4+16a^{2})/(1+a)^{2}\)（\(a\neq -1\)），求导得\(f'(a)=(32a-8)/(1+a)^{3}\)，驻点\(a=1/4\)。符号按分支定号：①\(a>1/4\)：分子\(32a-8>0\)、分母\((1+a)^{3}>0\)，\(f'>0\)；②\(-1<a<1/4\)：分子\(<0\)、分母\(>0\)，\(f'<0\)；③\(a<-1\)：分子\(<0\)，但分母\((1+a)^{3}<0\)，\(f'>0\)——该支\(f\)单调递增且\(a\to -1^{-}\)时\(f\to+\infty\)、\(f\)恒大于\(16>16/5\)，不产生最小值。\par\noindent 故\(f\)仅在\(a=1/4\)处取得最小值（\(16/5\)）。（几何：圆族中过两定点的圆以定点弦为直径时半径最小，与上互证。）故\(a=1/4\)．}
\fi
\end{ansblock}

\begin{ansblock}[2章-导-课时07-G1]
% ans:2章-导-课时07-G1
\ansitem{17}{A}
\ifansbookdetail
\ansnote{详解}{圆心\((0,b)\)：\(x^{2}+(y-b)^{2}=1\)。代入\((1,2)\)：\(1+(2-b)^{2}=1\)，得\(b=2\)。圆的方程为\(x^{2}+(y-2)^{2}=1\)，故选：A．}
\fi
\end{ansblock}

\begin{ansblock}[2章-导-课时07-G2]
% ans:2章-导-课时07-G2
\ansitem{18}{C}
\ifansbookdetail
\ansnote{详解}{\((3-2)^{2}+(2-3)^{2}=1+1=2<4=r^{2}\)，点在圆内，故选：C．}
\fi
\end{ansblock}

\begin{ansblock}[2章-导-课时07-G3]
% ans:2章-导-课时07-G3
\ansitem{19}{D}
\ifansbookdetail
\ansnote{详解}{配方：\((x+1)^{2}+(y-2)^{2}=6+1+4=11\)。圆心\((-1,2)\)，半径\(\sqrt{11}\)，故选：D．}
\fi
\end{ansblock}

\begin{ansblock}[2章-导-课时07-G4]
% ans:2章-导-课时07-G4
\ansitem{20}{(−1, 1/2)；√7/2}
\ifansbookdetail
\ansnote{详解}{两端除以2：\(x^{2}+y^{2}+2x-y-1/2=0\)，配方\((x+1)^{2}+(y-1/2)^{2}=1/2+1+1/4=7/4\)。圆心\((-1,\ 1/2)\)，半径\(\sqrt{7}/2\)．}
\fi
\end{ansblock}

\begin{ansblock}[2章-导-课时07-G5]
% ans:2章-导-课时07-G5
\ansitem{21}{−2}
\ifansbookdetail
\ansnote{详解}{圆上任意点关于\(l\)的对称点仍在圆上\(\iff\)圆关于\(l\)对称\(\iff l\)过圆心。圆心\((-1,\ -a/2)\)代入\(l\)：\(-1+a/2+2=0\)，得\(a=-2\)。（此时\(r^{2}=1+a^{2}/4+3=5>0\)，表圆成立．）}
\fi
\end{ansblock}

\jietitle{2.3.3 直线与圆的位置关系}

\begin{ansblock}[2章-练-课时08-T5]
% ans:2章-练-课时08-T5
\ansitem{1}{A}
\ifansbookdetail
\ansnote{详解}{圆\(x^{2}+(y-1)^{2}=1\)，圆心\((0,1)\)。直线\(y=k(x-1)+1\)过定点\((1,1)\)（\(x=1\)时\(y=1\)），恰为圆心，故\(d=0<r\)，恒相交，故选：A．}
\fi
\end{ansblock}

\begin{ansblock}[2章-练-课时08-T11]
% ans:2章-练-课时08-T11
\ansitem{2}{相切}
\ifansbookdetail
\ansnote{详解}{圆心\((0,0)\)，半径\(r\)。\(d=|-r|/\sqrt{\cos^{2}\theta+\sin^{2}\theta}=r\)，恒相切．}
\fi
\end{ansblock}

\begin{ansblock}[2章-练-课时08-T12]
% ans:2章-练-课时08-T12
\ansitem{3}{[√2, +∞)}
\ifansbookdetail
\ansnote{详解}{表圆：\(\sqrt{m}\)有意义即\(m\geqslant 0\)，\(r^{2}=m^{2}+m-2>0\)，在\(m\geqslant 0\)下\(m>1\)。\par\noindent 与\(x\)轴有公共点\(\iff|\sqrt{m}|\leqslant r\)，即\(m\leqslant m^{2}+m-2\)，即\(m^{2}\geqslant 2\)，得\(m\geqslant\sqrt{2}\)。故\(m\in[\sqrt{2},\ +\infty)\)．}
\fi
\end{ansblock}

\begin{ansblock}[2章-练-课时08-T13]
% ans:2章-练-课时08-T13
\ansitem{4}{(−∞, −5√2/4)∪(5√2/4, +∞)}
\ifansbookdetail
\ansnote{详解}{视线不被挡\(\iff\)直线\(AB\)与圆相离。\(k_{AB}=a/5\)，直线：\(ax-5y+2a=0\)。\par\noindent \(d=|3a|/\sqrt{a^{2}+25}>1\)，即\(9a^{2}>a^{2}+25\)，\(8a^{2}>25\)，\(|a|>5/(2\sqrt{2})=5\sqrt{2}/4\)．}
\fi
\end{ansblock}

\begin{ansblock}[2章-练-课时08-T8]
% ans:2章-练-课时08-T8
\ansitem{5}{C}
\ifansbookdetail
\ansnote{详解}{\(x\geqslant 0\)：\(y=x-1\)，\(x^{2}+(x-1)^{2}=1\)，即\(2x^{2}-2x=0\)，\(x=0\)（\(y=-1\)）或\(x=1\)（\(y=0\)）；\(x<0\)：\(y=-x-1\)，\(x^{2}+(x+1)^{2}=1\)，即\(2x^{2}+2x=0\)，\(x=-1\)（\(y=0\)）（\(x=0\)不在此段）。\par\noindent 交点\((0,-1)\)、\((1,0)\)、\((-1,0)\)共3个，真子集\(2^{3}-1=7\)个，故选：C．}
\fi
\end{ansblock}

\begin{ansblock}[2章-练-课时08-T1]
% ans:2章-练-课时08-T1
\ansitem{6}{A}
\ifansbookdetail
\ansnote{详解}{\((x-1)^{2}+(y-2)^{2}=9\)，圆心\((1,2)\)，\(r=3\)。\(d=|3+8-1|/5=2\)。弦长\(=2\sqrt{9-4}=2\sqrt{5}\)，故选：A．}
\fi
\end{ansblock}

\begin{ansblock}[2章-练-课时08-T2]
% ans:2章-练-课时08-T2
\ansitem{7}{A}
\ifansbookdetail
\ansnote{详解}{过圆内点的弦以垂直于\(OP\)的弦最短。\(k_{OP}=2\)，故所求斜率\(-1/2\)：\(y-2=-(1/2)(x-1)\)，即\(x+2y-5=0\)，故选：A．}
\fi
\end{ansblock}

\begin{ansblock}[2章-练-课时08-T3]
% ans:2章-练-课时08-T3
\ansitem{8}{C}
\ifansbookdetail
\ansnote{详解}{圆心\((1,-2)\)，\(r^{2}=5-m\)。\(d=|1-2-3|/\sqrt{2}=2\sqrt{2}\)。半弦长的平方\(=r^{2}-d^{2}\)，即\(1=5-m-8\)，得\(m=-4\)，故选：C．}
\fi
\end{ansblock}

\begin{ansblock}[2章-练-课时08-T4]
% ans:2章-练-课时08-T4
\ansitem{9}{B}
\ifansbookdetail
\ansnote{详解}{\(|AB|\geqslant 2\sqrt{3}\)，即\(d=\sqrt{r^{2}-(|AB|/2)^{2}}\leqslant\sqrt{4-3}=1\)。又\(d=|m|/\sqrt{2}\)，故\(|m|\leqslant\sqrt{2}\)，故选：B．}
\fi
\end{ansblock}

\begin{ansblock}[2章-练-课时08-T6]
% ans:2章-练-课时08-T6
\ansitem{10}{D}
\ifansbookdetail
\ansnote{详解}{圆\((x-2)^{2}+y^{2}=1\)，圆心\((2,0)\)，\(r=1\)。①过原点（两截距同为0）：\(y=kx\)，\(d=|2k|/\sqrt{1+k^{2}}=1\)，即\(k^{2}=1/3\)，2条；②不过原点：\(x/a+y/a=1\)即\(x+y=a\)，\(d=|2-a|/\sqrt{2}=1\)，\(a=2\pm\sqrt{2}\)（均非0），2条。\par\noindent 共4条，故选：D．}
\fi
\end{ansblock}

\begin{ansblock}[2章-练-课时08-T7]
% ans:2章-练-课时08-T7
\ansitem{11}{B}
\ifansbookdetail
\ansnote{详解}{圆与两轴都切且过\((2,1)\)（第一象限），则圆心\((a,a)\)，\(r=a\)，\(a>0\)。\((2-a)^{2}+(1-a)^{2}=a^{2}\)，即\(a^{2}-6a+5=0\)，\(a=1\)或\(a=5\)。\par\noindent \(a=1\)：圆心\((1,1)\)，\(d=|2-1-3|/\sqrt{5}=2/\sqrt{5}=2\sqrt{5}/5\)；\(a=5\)：圆心\((5,5)\)，\(d=|10-5-3|/\sqrt{5}=2/\sqrt{5}\)，相同。故\(d=2\sqrt{5}/5\)，故选：B．}
\fi
\end{ansblock}

\begin{ansblock}[2章-练-课时08-T9]
% ans:2章-练-课时08-T9
\ansitem{12}{A}
\ifansbookdetail
\ansnote{详解}{圆\((x-2)^{2}+(y+3)^{2}=1\)，圆心\(C(2,-3)\)，\(r=1\)。切线长\(=\sqrt{|PC|^{2}-r^{2}}\)，最小即\(|PC|\)最小\(=C\)到直线距离\(=|4+3-3|/\sqrt{5}=4/\sqrt{5}\)。\par\noindent 最小切线长\(=\sqrt{16/5-1}=\sqrt{11/5}=\sqrt{55}/5\)，故选：A．}
\fi
\end{ansblock}

\begin{ansblock}[2章-练-课时08-T10]
% ans:2章-练-课时08-T10
\ansitem{13}{C}
\ifansbookdetail
\ansnote{详解}{圆\((x-3)^{2}+(y-1)^{2}=9\)，圆心\((3,1)\)，\(r=3\)。对称轴过圆心：\(3+a-1=0\)，得\(a=-2\)，\(A(-1,-2)\)。\(|AC|=\sqrt{16+9}=5\)，切线长\(|AB|=\sqrt{25-9}=4\)，故选：C．}
\fi
\end{ansblock}

\begin{ansblock}[2章-练-课时08-T14]
% ans:2章-练-课时08-T14
\ansitem{14}{√3}
\ifansbookdetail
\ansnote{详解}{设直线\(y=\sqrt{3}x+b\)，\(A(0,b)\)。相切：\(d=|b-1|/2=1\)，得\(b=3\)或\(b=-1\)，\(|AC|=2\)（\(C(0,1)\)）。\(|AB|=\sqrt{|AC|^{2}-r^{2}}=\sqrt{4-1}=\sqrt{3}\)．}
\fi
\end{ansblock}

\begin{ansblock}[2章-练-课时08-T15]
% ans:2章-练-课时08-T15
\ansitem{15}{(1) x=3 或 3x−4y−5=0；(2) a=0 或 a=4/3；(3) a=−3/4}
\ifansbookdetail
\ansnote{详解}{\(C(1,2)\)，\(r=2\)；\(|MC|=\sqrt{4+1}=\sqrt{5}>2\)，\(M\)在圆外。\par\noindent (1) 斜率不存在时\(x=3\)：\(d=|3-1|=2=r\)，是切线。斜率存在时设\(y-1=k(x-3)\)，即\(kx-y-3k+1=0\)：\(d=|k-2-3k+1|/\sqrt{k^{2}+1}=|2k+1|/\sqrt{k^{2}+1}=2\)，即\(4k^{2}+4k+1=4k^{2}+4\)，得\(k=3/4\)，切线\(3x-4y-5=0\)。\par\noindent (2) \(d=|a-2+4|/\sqrt{a^{2}+1}=|a+2|/\sqrt{a^{2}+1}=2\)，即\(a^{2}+4a+4=4a^{2}+4\)，\(3a^{2}-4a=0\)，得\(a=0\)或\(a=4/3\)。\par\noindent (3) 弦长\(2\sqrt{3}\)，半弦\(\sqrt{3}\)，\(d=\sqrt{4-3}=1\)：\(|a+2|/\sqrt{a^{2}+1}=1\)，即\(a^{2}+4a+4=a^{2}+1\)，得\(a=-3/4\)．}
\fi
\end{ansblock}

\begin{ansblock}[2章-练-课时08-T16]
% ans:2章-练-课时08-T16
\ansitem{16}{(1) x²+y²=36；(2) 一定进入，航行时间 1 小时}
\ifansbookdetail
\ansnote{详解}{(1) \(O\)为原点，\(A(3,0)\)，\(B(12,0)\)。由\(|PA|/|PB|=1/2\)：\(4[(x-3)^{2}+y^{2}]=(x-12)^{2}+y^{2}\)，展开：\(4x^{2}-24x+36+4y^{2}=x^{2}-24x+144+y^{2}\)，即\(3x^{2}+3y^{2}=108\)，得\(x^{2}+y^{2}=36\)。\par\noindent (2) \(C(0,-2\sqrt{11})\)。北偏东\(30^{\circ}\)即航线倾角\(60^{\circ}\)，斜率\(\sqrt{3}\)，航线：\(y+2\sqrt{11}=\sqrt{3}x\)，即\(\sqrt{3}x-y-2\sqrt{11}=0\)。\(d=|2\sqrt{11}|/\sqrt{3+1}=\sqrt{11}<6\)，相交，必进入。\par\noindent 弦长\(=2\sqrt{36-11}=10\)，\(t=10/10=1\)（小时）．}
\fi
\end{ansblock}

\begin{ansblock}[2章-导-课时08-G1]
% ans:2章-导-课时08-G1
\ansitem{17}{D}
\ifansbookdetail
\ansnote{详解}{圆心\(O(0,0)\)，\(r=2\)。\(O\)到直线距离\(d=|2|/\sqrt{3+1}=1\)。弦心距\(OM=1\)，\(|AM|=\sqrt{r^{2}-d^{2}}=\sqrt{3}\)，故\(\sin\angle AOM=\sqrt{3}/2\)，\(\angle AOM=60^{\circ}\)，圆心角\(\angle AOB=120^{\circ}\)，故选：D．}
\fi
\end{ansblock}

\begin{ansblock}[2章-导-课时08-G2]
% ans:2章-导-课时08-G2
\ansitem{18}{B}
\ifansbookdetail
\ansnote{详解}{\(d=|-1|/\sqrt{\sin^{2}\theta+1}\)。由\(\theta\neq k\pi+\pi/2\)知\(\sin\theta\neq\pm 1\)，故\(0\leqslant\sin^{2}\theta<1\)，\(d=1/\sqrt{1+\sin^{2}\theta}\in(1/\sqrt{2},\ 1]\)，恒大于\(r=\sqrt{2}/2=1/\sqrt{2}\)，故相离，故选：B．}
\fi
\end{ansblock}

\begin{ansblock}[2章-导-课时08-G3]
% ans:2章-导-课时08-G3
\ansitem{19}{C}
\ifansbookdetail
\ansnote{详解}{\((\sqrt{3})^{2}+(3-2)^{2}=4\)，点在圆上，切线唯一。圆心\((0,2)\)，\(k_{OC}=(3-2)/(\sqrt{3}-0)=\sqrt{3}/3\)，切线斜率\(k=-\sqrt{3}\)，倾斜角\(120^{\circ}\)，故选：C．}
\fi
\end{ansblock}

\begin{ansblock}[2章-导-课时08-G4]
% ans:2章-导-课时08-G4
\ansitem{20}{[−1−2√2, −1+2√2]}
\ifansbookdetail
\ansnote{详解}{圆心\((a,0)\)，\(r=2\)。有公共点\(\iff d=|a+1|/\sqrt{2}\leqslant 2\)，即\(|a+1|\leqslant 2\sqrt{2}\)，得\(-1-2\sqrt{2}\leqslant a\leqslant -1+2\sqrt{2}\)．}
\fi
\end{ansblock}

\begin{ansblock}[2章-导-课时08-G5]
% ans:2章-导-课时08-G5
\ansitem{21}{x−y−3=0}
\ifansbookdetail
\ansnote{详解}{圆心\(C(1,0)\)。\(k_{CP}=(-1-0)/(2-1)=-1\)。\(CP\perp AB\)，故\(k_{AB}=1\)。\(y+1=1\cdot(x-2)\)，即\(x-y-3=0\)．}
\fi
\end{ansblock}

\jietitle{2.3.4 圆与圆的位置关系}

\begin{ansblock}[2章-练-课时09-T4]
% ans:2章-练-课时09-T4
\ansitem{1}{D}
\ifansbookdetail
\ansnote{详解}{满足条件的直线\(l\)即分别以\(A\)、\(B\)为圆心，1、2为半径的两圆的公切线．\(d=|AB|=3=1+2\)，两圆外切，公切线有3条（2外1内），选D．}
\fi
\end{ansblock}

\begin{ansblock}[2章-练-课时09-T1]
% ans:2章-练-课时09-T1
\ansitem{2}{D}
\ifansbookdetail
\ansnote{详解}{恰有两条公切线\(\iff\)两圆相交．圆心\((0,-m)\)、\((m,0)\)，半径\(\sqrt{2}\)、\(2\sqrt{2}\)，圆心距\(d=\sqrt{2}|m|\)．\par\noindent 相交\(\iff 2\sqrt{2}-\sqrt{2}<\sqrt{2}|m|<2\sqrt{2}+\sqrt{2}\)，即\(1<|m|<3\)，选D．}
\fi
\end{ansblock}

\begin{ansblock}[2章-练-课时09-T2]
% ans:2章-练-课时09-T2
\ansitem{3}{C}
\ifansbookdetail
\ansnote{详解}{\(M\cap N=N\iff N\subseteq M\)，即圆\(N\)内含或内切于圆\(M\)：\(d+r\leqslant 2\)，其中\(d=\sqrt{2}\)．故\(\sqrt{2}+r\leqslant 2\)，得\(0<r\leqslant 2-\sqrt{2}\)，选C．}
\fi
\end{ansblock}

\begin{ansblock}[2章-练-课时09-T7]
% ans:2章-练-课时09-T7
\ansitem{4}{ACD}
\ifansbookdetail
\ansnote{详解}{圆心\(M(2,1)\)、\(N(-2,-1)\)，\(d=2\sqrt{5}>2\)，两圆相离，公切线4条．逐项验圆心到直线的距离是否等于1：\par\noindent A：\(y=0\)到\(M\)距1，到\(N\)距1，是；B：\(3x-4y=0\)到\(M\)距\(|6-4|/5=2/5\neq 1\)，不是；\par\noindent C：\(x-2y+\sqrt{5}=0\)到\(M\)距\(\sqrt{5}/\sqrt{5}=1\)，到\(N\)距1，是；D：同理距离均为1，是．故选ACD．}
\fi
\end{ansblock}

\begin{ansblock}[2章-练-课时09-T10]
% ans:2章-练-课时09-T10
\ansitem{5}{AD}
\ifansbookdetail
\ansnote{详解}{四条公切线\(\iff\)两圆外离．圆心\((0,a)\)、\((a,0)\)，\(d=\sqrt{2}|a|\)，半径3、1．外离\(\iff\sqrt{2}|a|>4\)，即\(|a|>2\sqrt{2}\)．\(-4\)合，\(3>2\sqrt{2}\)合，\(-2\)与\(2\sqrt{2}\)均不合．选AD．}
\fi
\end{ansblock}

\begin{ansblock}[2章-练-课时09-T8]
% ans:2章-练-课时09-T8
\ansitem{6}{ACD}
\ifansbookdetail
\ansnote{详解}{\(C\): \((x-3)^{2}+y^{2}=9\)，圆心\((3,0)\)，\(r=3\)．\par\noindent A：\(r=3\)，对；B：点\((1,2\sqrt{2})\)到圆心距离\(\sqrt{4+8}=2\sqrt{3}>3\)，在圆外，错；\par\noindent C：\(d=|3+0+3|/2=3=r\)，相切，对；D：圆心距\(4\)，半径3、2，\(1<4<5\)，相交，对．故选ACD．}
\fi
\end{ansblock}

\begin{ansblock}[2章-练-课时09-T3]
% ans:2章-练-课时09-T3
\ansitem{7}{A}
\ifansbookdetail
\ansnote{详解}{\(\overrightarrow{PA}\cdot\overrightarrow{PB}=0\iff\angle APB=90^{\circ}\)，即\(P\)在以\(AB\)为直径的圆\(x^{2}+y^{2}=4\)上（除\(A\)、\(B\)）．存在这样的\(P\iff\)两圆相交：\(|r-2|<3<r+2\)，得\(r>1\)且\(r<5\)，选A．（恰过\(A\)、\(B\)的情形对应相切\(r=1\)或\(r=5\)，已由“不同于\(A\)，\(B\)”排除．）}
\fi
\end{ansblock}

\begin{ansblock}[2章-练-课时09-T12]
% ans:2章-练-课时09-T12
\ansitem{8}{5}
\ifansbookdetail
\ansnote{详解}{\(\angle APB=90^{\circ}\iff P\)在以\(AB\)为直径的圆\(x^{2}+y^{2}=m^{2}\)上．存在这样的\(P\iff\)两圆有公共点：\(|m-1|\leqslant|OC|\leqslant m+1\)，而\(|OC|=\sqrt{9+7}=4\)．由\(4\geqslant m-1\)得\(m\leqslant 5\)（另一侧\(4\leqslant m+1\)给\(m\geqslant 3\)），故\(m\)的最大值为5．}
\fi
\end{ansblock}

\begin{ansblock}[2章-练-课时09-T13]
% ans:2章-练-课时09-T13
\ansitem{9}{相交}
\ifansbookdetail
\ansnote{详解}{\(C_2\)的圆心\((2,-m/2)\)在对称轴上：\(2-(\sqrt{3}/2)m+1=0\)，得\(m=2\sqrt{3}\)．\(r_2^{2}=4+(m/2)^{2}-3=4\)，\(r_2=2\)，圆心\(C_2(2,-\sqrt{3})\)；圆\(C_1\)圆心\((4,4)\)，\(r_1=5\)．\par\noindent \(d^{2}=(4-2)^{2}+(4+\sqrt{3})^{2}=23+8\sqrt{3}\approx 36.9\)，而\((r_1-r_2)^{2}=9\)，\((r_1+r_2)^{2}=49\)，\(9<23+8\sqrt{3}<49\)，即\(3<d<7\)，两圆相交．}
\fi
\end{ansblock}

\begin{ansblock}[2章-练-课时09-T9]
% ans:2章-练-课时09-T9
\ansitem{10}{AD}
\ifansbookdetail
\ansnote{详解}{\(|OC|=\sqrt{4+9}=\sqrt{13}>2+1=3\)，两圆外离，公切线4条，A对．\par\noindent B：\(x+y=5\)过\(C(2,3)\)且两截距均为5，合；但\(x-y+1=0\)过\(C\)，两截距为\(-1\)与\(1\)，不相等，B的表述错误，B不对．\par\noindent C：\(C(2,3)\)在圆\(O\)外，过\(C\)与圆\(O\)相切的切线应有两条，选项只给一条且按唯一解表述，不成立（\(9x-16y+30=0\)过\(C\)但只是其中一条），C不对．\par\noindent D：\(|PQ|_{\max}=|OC|+2+1=\sqrt{13}+3\)，\(|PQ|_{\min}=|OC|-2-1=\sqrt{13}-3\)，D对．故选AD．}
\fi
\end{ansblock}

\begin{ansblock}[2章-练-课时09-T5]
% ans:2章-练-课时09-T5
\ansitem{11}{C}
\ifansbookdetail
\ansnote{详解}{两圆方程相减得公共弦\(AB\)所在直线：\(x=-2/3\)．\(|AB|=2\sqrt{r_1^{2}-d_1^{2}}=2\sqrt{4-4/9}=8\sqrt{2}/3\)，\(|O_1O_2|=3\)．连心线垂直平分公共弦，四边形\(AO_1BO_2\)的对角线互相垂直，其面积\(=(1/2)\times(8\sqrt{2}/3)\times 3=4\sqrt{2}\)，选C．}
\fi
\end{ansblock}

\begin{ansblock}[2章-练-课时09-T11]
% ans:2章-练-课时09-T11
\ansitem{12}{1}
\ifansbookdetail
\ansnote{详解}{两圆方程相减得公共弦所在直线：\(2ay-2=0\)，即\(y=1/a\)，圆心到它的距离\(d=1/a\)．由勾股关系\((\sqrt{3})^{2}+d^{2}=r^{2}\)：\(3+1/a^{2}=4\)，得\(a^{2}=1\)，又\(a>0\)，故\(a=1\)．}
\fi
\end{ansblock}

\begin{ansblock}[2章-练-课时09-T6]
% ans:2章-练-课时09-T6
\ansitem{13}{C}
\ifansbookdetail
\ansnote{详解}{设\(P(x_0,y_0)\)．以\(MP\)为直径的圆过切点\(B\)、\(C\)（\(MB\perp PB\)），其方程为\((x-1)(x-x_0)+y(y-y_0)=0\)．与圆\(M\): \(x^{2}+y^{2}-2x-3=0\)展开相减，得公共弦（即切点弦）\(BC\)所在直线：\((x_0-1)x+y_0y-x_0-3=0\)．\par\noindent \(A(2,1)\)在\(BC\)上：\(2(x_0-1)+y_0-x_0-3=0\)，即\(x_0+y_0=5\)．故点\(P\)的轨迹方程为\(x+y-5=0\)，选C．}
\fi
\end{ansblock}

\begin{ansblock}[2章-练-课时09-T14]
% ans:2章-练-课时09-T14
\ansitem{14}{x+2y+1=0}
\ifansbookdetail
\ansnote{详解}{圆\(C\): \((x-1)^{2}+(y-1)^{2}=4\)，圆心\((1,1)\)，\(r=2\)．\(S=2\times(1/2)r\cdot|MA|=2\sqrt{|MC|^{2}-4}\)，最小\(\iff|MC|\)最小，即\(C\)到\(l\)的距离\(|1+2+2|/\sqrt{5}=\sqrt{5}\)．\par\noindent 此时\(M\)为垂足：\(CM\)的方向\((1,2)\)，直线\(CM\): \(y-1=2(x-1)\)，与\(l\)联立\(x+2(2x-1)+2=0\)，得\(x=0\)，\(M(0,-1)\)．以\(CM\)为直径的圆圆心\((1/2,0)\)、半径\(\sqrt{5}/2\)，方程为\(x^{2}+y^{2}-x-1=0\)，与圆\(C\): \(x^{2}+y^{2}-2x-2y-2=0\)相减，得切点弦\(AB\)所在直线：\(x+2y+1=0\)．}
\fi
\end{ansblock}

\begin{ansblock}[2章-练-课时09-T15]
% ans:2章-练-课时09-T15
\ansitem{15}{猜想正确，理由见解析}
\ifansbookdetail
\ansnote{详解}{设交点\(A(x_1,y_1)\)、\(B(x_2,y_2)\)，则两交点坐标同时满足两圆方程：\(x_i^{2}+y_i^{2}+D_1x_i+E_1y_i+F_1=0\)且\(x_i^{2}+y_i^{2}+D_2x_i+E_2y_i+F_2=0\)（\(i=1,2\)）．\par\noindent 两式相减：\((D_1-D_2)x_i+(E_1-E_2)y_i+(F_1-F_2)=0\)，即\(A\)、\(B\)都在直线\((D_1-D_2)x+(E_1-E_2)y+(F_1-F_2)=0\)上．又\(D_1\neq D_2\)或\(E_1\neq E_2\)（否则两圆方程相同），该直线系数不全为零，而过两个不同交点的直线有且仅有一条，故它就是公共弦\(AB\)所在直线的方程，猜想正确．}
\fi
\end{ansblock}

\begin{ansblock}[2章-练-课时09-T16]
% ans:2章-练-课时09-T16
\ansitem{16}{(x−4)²+(y−3)²=25}
\ifansbookdetail
\ansnote{详解}{设所求圆的圆心\(D(a,b)\)．切点\(B\)在两圆连心线上（连心线过切点且与公切线垂直）：\(k_{CB}=k_{DB}\)，即\(b/a=6/8=3/4\)，得\(b=3a/4\)．又\(|DA|=|DB|\)：\((a-1)^{2}+(b+1)^{2}=(a-8)^{2}+(b-6)^{2}\)，展开整理得\(14a+14b=98\)，即\(a+b=7\)．\par\noindent 联立\(b=3a/4\)得\(a=4\)，\(b=3\)，\(r=|DB|=\sqrt{16+9}=5\)．故所求圆的方程为\((x-4)^{2}+(y-3)^{2}=25\)（验：\(|DA|=5\)合；\(|DC|=5=R-r\)（大圆半径\(10\)减所求圆半径\(5\)），所求圆与大圆内切于\(B\)）．}
\fi
\end{ansblock}

\begin{ansblock}[2章-导-课时09-G1]
% ans:2章-导-课时09-G1
\ansitem{17}{B}
\ifansbookdetail
\ansnote{详解}{\(C_1\): \((x-1)^{2}+(y+2)^{2}=9\)，圆心\((1,-2)\)，\(r_1=3\)；\(C_2\): \((x+3)^{2}+(y-4)^{2}=25\)，圆心\((-3,4)\)，\(r_2=5\)．\(d=\sqrt{16+36}=2\sqrt{13}\approx 7.2\)，\(5-3<2\sqrt{13}<5+3\)，两圆相交，公切线2条，故选：B．}
\fi
\end{ansblock}

\begin{ansblock}[2章-导-课时09-G2]
% ans:2章-导-课时09-G2
\ansitem{18}{B}
\ifansbookdetail
\ansnote{详解}{面积被直线平分\(\iff\)直线过圆心：\(1-m+1=0\)，得\(m=2\)．\(C_1\): \((x-1)^{2}+(y+1)^{2}=1\)，圆心\((1,-1)\)，\(r_1=1\)；\(C_2\)圆心\((-2,3)\)，\(r_2=5\)．\(d=\sqrt{9+16}=5\)，\(4<5<6\)，两圆相交，故选：B．}
\fi
\end{ansblock}

\begin{ansblock}[2章-导-课时09-G3]
% ans:2章-导-课时09-G3
\ansitem{19}{A}
\ifansbookdetail
\ansnote{详解}{已知圆：\((x+2)^{2}+(y-3)^{2}=9\)，圆心\((-2,3)\)，\(r=3\)．\(|AC|=\sqrt{16+9}=5=r+R\)，得\(R=2\)．圆\(C\): \((x-2)^{2}+y^{2}=4\)，即\(x^{2}+y^{2}-4x=0\)，故选：A．}
\fi
\end{ansblock}

\begin{ansblock}[2章-导-课时09-G4]
% ans:2章-导-课时09-G4
\ansitem{20}{(0,−1)、(−1,0)}
\ifansbookdetail
\ansnote{详解}{两圆方程相减得公共弦所在直线：\(x+y+1=0\)．代入\(x^{2}+y^{2}=1\)：\(x^{2}+(-x-1)^{2}=1\)，即\(2x^{2}+2x=0\)，解得\(x=0\)或\(x=-1\)，交点为\((0,-1)\)、\((-1,0)\)．}
\fi
\end{ansblock}

\begin{ansblock}[2章-导-课时09-G5]
% ans:2章-导-课时09-G5
\ansitem{21}{(1) 内切；(2) 外切；(3) 相交}
\ifansbookdetail
\ansnote{详解}{(1)圆心距\(\sqrt{9+16}=5=6-1\)，内切；(2)\(d=\sqrt{16+9}=5=2+3\)，外切；(3)化标准形：圆心\((-1,2)\)、\((1,0)\)，\(d=\sqrt{4+4}=2\sqrt{2}\)，半径3与1，\(1<2\sqrt{2}<4\)，相交．}
\fi
\end{ansblock}

\jietitle{2.4 曲线与方程}

\begin{ansblock}[2章-练-课时10-简1]
% ans:2章-练-课时10-简1
\ansitem{1}{P在曲线上；Q不在曲线上}
\ifansbookdetail
\ansnote{详解}{曲线上的点的坐标必为方程的解，反之方程的解对应曲线上的点，故代入即判．\(P(\sqrt{3},0)\)：\(4\times(\sqrt{3})^{2}+3\times 0^{2}=12+0=12\)，是方程的解，\(P\)在曲线上．\(Q(-2,3)\)：\(4\times(-2)^{2}+3\times 3^{2}=16+27=43\neq 12\)，不是方程的解，\(Q\)不在曲线上．}
\fi
\end{ansblock}

\begin{ansblock}[2章-练-课时10-简2]
% ans:2章-练-课时10-简2
\ansitem{2}{r²=5}
\ifansbookdetail
\ansnote{详解}{曲线通过点\((-1,5)\)，即\((-1,5)\)的坐标是方程的解：\((-1+2)^{2}+(5-3)^{2}=r^{2}\)，得\(r^{2}=1+4=5\)．}
\fi
\end{ansblock}

\begin{ansblock}[2章-练-课时10-简3]
% ans:2章-练-课时10-简3
\ansitem{3}{不是。x−y=0 只表示其中一条直线，轨迹应为 x−y=0 或 x+y=0（即|x|=|y|）；以 x−y=0 为轨迹方程不完备（漏了直线 x+y=0 上的点）}
\ifansbookdetail
\ansnote{详解}{设动点\(M(x,y)\)．到\(x\)轴距离为\(|y|\)、到\(y\)轴距离为\(|x|\)，条件\(|x|=|y|\iff x-y=0\)或\(x+y=0\)，轨迹是两条互相垂直的直线．方程\(x-y=0\)的解对应的点都在轨迹上（纯粹性成立），但轨迹上的点（如\((1,-1)\)）不都是该方程的解（完备性不成立）——它只描述第一、三象限角平分线，漏了第二、四象限角平分线．故“轨迹的方程是\(x-y=0\)”不对．}
\fi
\end{ansblock}

\begin{ansblock}[2章-练-课时10-简7]
% ans:2章-练-课时10-简7
\ansitem{4}{A}
\ifansbookdetail
\ansnote{详解}{A：\(\sqrt{x^{2}}=|x|\)，两方程解集（图象）完全相同，表同一条曲线；B：\(y=x\)为直线，\(y=\sqrt{x^{2}}=|x|\)为折线（\(x<0\)时\(y=-x\)），异；C：\(|y|=x\iff x\geqslant 0\)且\(y=\pm x\)，为两条射线，而\(y=|x|\)还含\(x<0\)部分，异；D：\(x^{2}+y^{2}=0\iff\)点\((0,0)\)，\(y=0\)为整条\(x\)轴，异．故选：A．}
\fi
\end{ansblock}

\begin{ansblock}[2章-练-课时10-简8]
% ans:2章-练-课时10-简8
\ansitem{5}{C}
\ifansbookdetail
\ansnote{详解}{C 与原方程仅差移项，同解变形，解集恒等，表同一条曲线；A：\(\sqrt{4-x^{2}}\) 要求 \(y\geqslant 0\)，仅为上半圆，缺 \(y<0\) 部分不完备，排除；B：圆心 \((2,0)\)、半径 2 的圆，与圆心在原点的圆非同一条曲线，排除；D：\(|x|+|y|=2\) 是以 \((2,0)\)、\((0,2)\)、\((-2,0)\)、\((0,-2)\) 为顶点的正方形边界，点 \((1,1)\) 在其上却不在圆上，排除．故选：C．}
\fi
\end{ansblock}

\begin{ansblock}[2章-练-课时10-简10]
% ans:2章-练-课时10-简10
\ansitem{6}{(1) 假　(2) 真　(3) 真}
\ifansbookdetail
\ansnote{详解}{(1) \(x^{2}+y^{2}=0\) 当且仅当 \(x=y=0\)，表示原点一个点，非直线；(2) \(|x|=|y|\) 当且仅当 \(x^{2}=y^{2}\)（两侧均非负，平方同解），即 \(y^{2}=x^{2}\)，命题真；(3) 由曲线方程定义的完备性（以方程的解为坐标的点都在曲线上），命题真．}
\fi
\end{ansblock}

\begin{ansblock}[2章-练-课时10-简4]
% ans:2章-练-课时10-简4
\ansitem{7}{x+2y−7=0}
\ifansbookdetail
\ansnote{详解}{设\(P(x,y)\)为线段\(AB\)垂直平分线上任意一点，则\(|PA|=|PB|\)：\((x+1)^{2}+(y+1)^{2}=(x-3)^{2}+(y-7)^{2}\)，展开：\(x^{2}+2x+1+y^{2}+2y+1=x^{2}-6x+9+y^{2}-14y+49\)，化简得 \(8x+16y-56=0\)，即 \(x+2y-7=0\)．反向检验：满足 \(x+2y-7=0\) 的点均有\(|PA|=|PB|\)，在\(AB\)的垂直平分线上（两性完备）．验算：\(AB\)中点\((1,3)\)：\(1+6-7=0\)成立；\(k_{AB}=(7+1)/(3+1)=2\)，所求直线斜率\(-1/2\)，\(2\times(-1/2)=-1\)成立．}
\fi
\end{ansblock}

\begin{ansblock}[2章-练-课时10-简5]
% ans:2章-练-课时10-简5
\ansitem{8}{(x−7)²+(y−2)²=32}
\ifansbookdetail
\ansnote{详解}{设\(M(x,y)\)，\(|MA|/|MB|=\sqrt{2}\)，即\(|MA|^{2}=2|MB|^{2}\)：\((x+1)^{2}+(y-2)^{2}=2[(x-3)^{2}+(y-2)^{2}]\)．展开：\(x^{2}+2x+1+(y-2)^{2}=2x^{2}-12x+18+2(y-2)^{2}\)，移项整理：\(x^{2}-14x+(y-2)^{2}+17=0\)，配方得 \((x-7)^{2}+(y-2)^{2}=32\)．轨迹为以\((7,2)\)为圆心、\(4\sqrt{2}\)为半径的圆（阿波罗尼斯圆雏形，此名可不向学生出现）．验算：圆与直线\(AB\)（\(y=2\)）交于 \(M_1(7-4\sqrt{2},2)\)、\(M_2(7+4\sqrt{2},2)\)．对\(M_1\)：\(|M_1A|=|7-4\sqrt{2}+1|=8-4\sqrt{2}=4(2-\sqrt{2})\)，\(|M_1B|=|4\sqrt{2}-4|=4(\sqrt{2}-1)\)（取绝对值），比值\(=(2-\sqrt{2})/(\sqrt{2}-1)=\sqrt{2}(\sqrt{2}-1)/(\sqrt{2}-1)=\sqrt{2}\)成立；对\(M_2\)同法比值亦\(\sqrt{2}\)成立．}
\fi
\end{ansblock}

\begin{ansblock}[2章-练-课时10-简6]
% ans:2章-练-课时10-简6
\ansitem{9}{x²/81+y²/36=1（x≠0）}
\ifansbookdetail
\ansnote{详解}{设\(A(x,y)\)，\(x\neq 0\)（\(x=0\) 时 \(A\) 在 \(y\) 轴上与 \(B\)、\(C\) 共线，不能成三角形且斜率不存在）．\(k_{AB}\cdot k_{AC}=[(y-6)/x]\cdot[(y+6)/x]=(y^{2}-36)/x^{2}=-4/9\)，去分母：\(9(y^{2}-36)=-4x^{2}\)，整理得 \(4x^{2}+9y^{2}=324\)，即 \(x^{2}/81+y^{2}/36=1\)（\(x\neq 0\)）．}
\fi
\end{ansblock}

\begin{ansblock}[2章-练-课时10-简9]
% ans:2章-练-课时10-简9
\ansitem{10}{x²=8(y−2)}
\ifansbookdetail
\ansnote{详解}{设 \(M(x,y)\)．由题 \(|y|=\sqrt{x^{2}+(y-4)^{2}}\)．两边平方（两侧非负，等价）：\(y^{2}=x^{2}+y^{2}-8y+16\)，即 \(x^{2}=8(y-2)\)．检验：方程上任一点均满足定义（平方等价可逆），无杂点；\(A(0,4)\) 代入 \(16\neq 0\) 不在轨迹上，无需剔除．故轨迹方程为 \(x^{2}=8(y-2)\)．}
\fi
\end{ansblock}

\begin{ansblock}[2章-练-课时10-中2]
% ans:2章-练-课时10-中2
\ansitem{11}{x²/4+y²=1（x≠±2）}
\ifansbookdetail
\ansnote{详解}{设\(P(x,y)\)．斜率存在要求\(x\neq\pm 2\)；\(y^{2}/(x^{2}-4)=-1/4\Rightarrow 4y^{2}=4-x^{2}\)，即\(x^{2}/4+y^{2}=1\)（\(x\neq\pm 2\)）．纯粹性：轨迹上任一点的坐标推导步步可逆，都满足斜率方程（成立）；完备性：满足题设的点都满足\(x^{2}/4+y^{2}=1\)且\(x\neq\pm 2\)，都在轨迹上（成立）（\(y=0\)时斜率积为\(0\neq -1/4\)，端点自然排除）．故轨迹为椭圆\(x^{2}/4+y^{2}=1\)除去长轴两端点\((\pm 2,0)\)．}
\fi
\end{ansblock}

\begin{ansblock}[2章-练-课时10-中1]
% ans:2章-练-课时10-中1
\ansitem{12}{(1)两个外切于原点的单位圆（「8」字形）；(2)关于x轴、y轴、原点均对称；x∈[−2,2]，y∈[−1,1]；与y轴仅交于(0,0)，与x轴交于(±2,0)与(0,0)；S=2π}
\ifansbookdetail
\ansnote{详解}{(1)按\(|x|\)分段：\(x\geqslant 0\)时\((x-1)^{2}+y^{2}=1\)，\(x<0\)时\((x+1)^{2}+y^{2}=1\)，两圆圆心\((\pm 1,0)\)、半径1，外切于原点．(2)\(|x|\)使\(C\)关于\(y\)轴对称；两圆连心线在\(x\)轴上，故关于\(x\)轴对称；由两者得关于原点对称．范围：\(x\in[-2,2]\)，\(y\in[-1,1]\)．交点：\(x=0\Rightarrow y^{2}=0\)，仅\((0,0)\)；\(y=0\Rightarrow|x|=2\)或\(0\)，得\((\pm 2,0)\)与\((0,0)\)．两圆仅切于一点无重叠，\(S=2\times\pi\times 1^{2}=2\pi\)．}
\fi
\end{ansblock}

\begin{ansblock}[2章-练-课时10-中3]
% ans:2章-练-课时10-中3
\ansitem{13}{m∈[−6,−4]∪[0,2]}
\ifansbookdetail
\ansnote{详解}{圆心\(C_1(0,m)\)、\(C_2(0,-2)\)，半径\(r_1=3\)、\(r_2=1\)，圆心距\(|C_1C_2|=|m+2|\)．两圆有公共点当且仅当\(|r_1-r_2|\leqslant|C_1C_2|\leqslant r_1+r_2\)，即\(2\leqslant|m+2|\leqslant 4\)，解得\(2\leqslant m+2\leqslant 4\)或\(-4\leqslant m+2\leqslant -2\)，即\(m\in[-6,-4]\cup[0,2]\)．}
\fi
\end{ansblock}

\begin{ansblock}[2章-练-课时10-中4]
% ans:2章-练-课时10-中4
\ansitem{14}{k<9：椭圆；k=9：单点(1,2)；k>9：不表示任何曲线}
\ifansbookdetail
\ansnote{详解}{配方：\((x-1)^{2}+2(y-2)^{2}=9-k\)．\(k<9\)时右边为正，方程化为\((x-1)^{2}/(9-k)+(y-2)^{2}/\bigl((9-k)/2\bigr)=1\)，表示椭圆；\(k=9\)时方程为\((x-1)^{2}+2(y-2)^{2}=0\)，表示单点\((1,2)\)（退化椭圆）；\(k>9\)时左边非负而右边为负，无实数解，不表示任何曲线．}
\fi
\end{ansblock}

\begin{ansblock}[2章-练-课时10-难1]
% ans:2章-练-课时10-难1
\ansitem{15}{BCD}
\ifansbookdetail
\ansnote{详解}{A：令\(y=0\)得\(x^{4}=4x^{2}\)，\(x=0\)或\(\pm 2\)，整点\((2,0)\)、\((-2,0)\)、\((0,0)\)三个；若整点满足\(y\neq 0\)，则由B项结论\(x^{2}+y^{2}\leqslant 4\)只能\(y=\pm 1\)、\(x\in\{0,\pm 1,\pm 2\}\)：代入\((x^{2}+1)^{2}=4(x^{2}-1)\)，记\(X=x^{2}\)得\(X^{2}-2X+5=0\)，判别式\(4-20<0\)无解——\(y\neq 0\)无整点．故曲线仅过3个整点，“5个整点”错．B：由\((x^{2}+y^{2})^{2}=4(x^{2}-y^{2})\leqslant 4(x^{2}+y^{2})\)得\(x^{2}+y^{2}\leqslant 4\)，即\(|OP|\leqslant 2\)（成立）．C：方程中\(x\)、\(y\)互换得\((x^{2}+y^{2})^{2}=4(y^{2}-x^{2})\)，即关于\(y=x\)对称的曲线（成立）．D：\(y=kx\)代入：\(x^{4}(1+k^{2})^{2}=4x^{2}(1-k^{2})\)．\(x=0\)恒给交点\((0,0)\)；\(|x|>0\)的解须\(x^{2}=4(1-k^{2})/(1+k^{2})^{2}>0\)，当且仅当\(|k|<1\)．故\(|k|\geqslant 1\)时只有原点一个交点；\(|k|<1\)时另有\(\pm\)两交点（\((0,0)\)之外），不止一个（成立）（\(k=\pm 1\)时\(x^{4}\cdot 2=0\)仅原点，含入界点正确）．故选：BCD．}
\fi
\end{ansblock}

\begin{ansblock}[2章-练-课时10-难2]
% ans:2章-练-课时10-难2
\ansitem{16}{②④}
\ifansbookdetail
\ansnote{详解}{①\(m=1\)时\((x-1)^{2}+(y-2)^{2}=n^{2}/2\)，但\(n=0\)时退化为点\((1,2)\)，不（恒）表示圆——以偏概全，错误．②\(m=0\)、\(n=2\)：\(C\)：\(x^{2}+y^{2}=2\)．记\(D(3,3)\)，\(|DO|=3\sqrt{2}\)，切线长\(|DA|=|DB|=\sqrt{18-2}=4\)，切点\(A\)、\(B\)即\(C\)与圆\((x-3)^{2}+(y-3)^{2}=16\)（圆心\(D\)半径4）的公共点，公共弦方程由两圆方程相减：\(x^{2}+y^{2}-2-(x^{2}+y^{2}-6x-6y+2)=0\)，即\(3x+3y-2=0\)，正确．③\(m=1\)、\(n=\sqrt{2}\)：\(C\)：\((x-1)^{2}+(y-2)^{2}=1\)．过\((2,0)\)的切线：竖直直线\(x=2\)到圆心\((1,2)\)距离为\(1=\)半径，已是切线；另设\(y=k(x-2)\)：\(|k-2-2k|/\sqrt{k^{2}+1}=1\Rightarrow(k+2)^{2}=k^{2}+1\Rightarrow k=-3/4\)，即\(3x+4y-6=0\)也是切线．结论只给后者、漏\(x=2\)，“切线方程为”不完备，错误．④\(n=m\neq 0\)：\((x-m)^{2}+(y-2m)^{2}=m^{2}/2\)，圆心\((m,2m)\)恒在\(y=2x\)上，半径\((\sqrt{2}/2)|m|\)．圆心到\(y=x\)距离\(=|m-2m|/\sqrt{2}=(\sqrt{2}/2)|m|=\)半径；到\(y=7x\)距离\(=|7m-2m|/\sqrt{50}=(\sqrt{2}/2)|m|=\)半径——两直线都是圆系公切线．穷尽性补核（源未给）：设公切线\(y=kx+b\)对一切\(m\)成立，则\([m(k-2)+b]^{2}=\frac{1}{2}m^{2}(k^{2}+1)\)恒成立，比较系数：\(b=0\)且\((k-2)^{2}=(k^{2}+1)/2\Rightarrow k^{2}-8k+7=0\Rightarrow k=1\)或\(7\)；竖直线\(x=c\)代入\(|m-c|=(\sqrt{2}/2)|m|\)不能对一切\(m\)成立，无．故公切线恰为\(y=x\)、\(y=7x\)，正确．⑤\(n=4\)、\(m=0\)：\(C\)：\(x^{2}+y^{2}=8\)，直线\(kx-y+1-2k=0\)即\(k(x-2)-(y-1)=0\)恒过定点\((2,1)\)，而\(4+1=5<8\)，定点在圆内，直线恒与圆相交而非相离，错误．综上，正确的是②④．}
\fi
\end{ansblock}

\begin{ansblock}[2章-导-课时10-G1]
% ans:2章-导-课时10-G1
\ansitem{17}{D}
\ifansbookdetail
\ansnote{详解}{设\(M(x,y)\)，\(P(x_0,y_0)\)，则\(x_0^{2}+y_0^{2}=4\)．由\(PH\perp y\)轴知\(H(0,y_0)\)；\(M\)为\(PH\)中点，得\(x=x_0/2\)、\(y=y_0\)，即\(x_0=2x\)、\(y_0=y\)．代入圆方程：\(4x^{2}+y^{2}=4\)，整理得\(x^{2}+y^{2}/4=1\)．故选：D．}
\fi
\end{ansblock}

\begin{ansblock}[2章-导-课时10-G2]
% ans:2章-导-课时10-G2
\ansitem{18}{A}
\ifansbookdetail
\ansnote{详解}{\(F(-2,0)\)、\(C(2,0)\)，记\(F\)关于折痕\(l\)的对称点为\(A\)，\(l\)与\(AC\)交于点\(P\)，则\(A\)在圆周上，\(l\)为\(AF\)的垂直平分线，故\(|PA|=|PF|\)．于是\(|PF|+|PC|=|PA|+|PC|=|AC|=8>|FC|=4\)，点\(P\)的轨迹是以\(F\)、\(C\)为左、右焦点的椭圆，\(2a=8\)、\(2c=4\)，\(a=4\)、\(c=2\)、\(b=2\sqrt{3}\)，方程为\(x^{2}/16+y^{2}/12=1\)．故选：A．}
\fi
\end{ansblock}

\begin{ansblock}[2章-导-课时10-G3]
% ans:2章-导-课时10-G3
\ansitem{19}{D}
\ifansbookdetail
\ansnote{详解}{曲线定义域：\(x-1\geqslant 0\)且\(x^{2}+y^{2}-1>0\)．乘积为零当且仅当\(\sqrt{x-1}=0\)或\(\ln(x^{2}+y^{2}-1)=0\)．前者给\(x=1\)（此时\(x^{2}+y^{2}-1=y^{2}>0\)须\(y\neq 0\)），即直线\(x=1\)除去\((1,0)\)；后者给\(x^{2}+y^{2}=2\)且\(x\geqslant 1\)，即圆\(x^{2}+y^{2}=2\)的右半弧（含端点\((1,\pm 1)\)，该两点同时落在\(x=1\)上）．直线\(y=kx-1\)恒过定点\((0,-1)\)：与\(x=1\)（\(y\neq 0\)）交于\((1,k-1)\)，\(k\neq 1\)时是曲线上一交点；又\((0,-1)\)到原点距离\(1<\sqrt{2}\)，直线与整圆恒交于两点，其中落在右半弧上的：临界为直线过弧端点\((1,-1)\)（\(k=0\)）与\((1,1)\)（\(k=2\)），端点处弧上点与\(x=1\)上点重合、两交并一，故须\(0<k<2\)（严格）；再排除\(k=1\)（此时\((1,k-1)=(1,0)\)不在曲线上，曲线只剩弧上一个交点，\(k=1\)时直线与圆\(x\geqslant 1\)半弧恰一交点）．综上\(0<k<2\)且\(k\neq 1\)，故选：D．}
\fi
\end{ansblock}

\begin{ansblock}[2章-导-课时10-G4]
% ans:2章-导-课时10-G4
\ansitem{20}{x²+y²−8x−4y+18=0（除去点(3,1)与(5,3)）}
\ifansbookdetail
\ansnote{详解}{设\(C(x,y)\)．由\(|AC|=|AB|\)得\(\sqrt{(x-4)^{2}+(y-2)^{2}}=\sqrt{(4-5)^{2}+(2-3)^{2}}=\sqrt{2}\)，平方化简得\(x^{2}+y^{2}-8x-4y+18=0\)（圆心\((4,2)\)、半径\(\sqrt{2}\)的圆）．剔除：\(C\)与\(B\)重合，去\((5,3)\)；\(A\)、\(B\)、\(C\)三点共线时不构成三角形——\(k_{AB}=(3-2)/(5-4)=1\)，直线\(AB\)：\(y-2=1\cdot(x-4)\)即\(x-y-2=0\)，与圆联立解得交点\((3,1)\)与\((5,3)\)，故须除去\((3,1)\)、\((5,3)\)．轨迹方程为\(x^{2}+y^{2}-8x-4y+18=0\)（\(x-y-2\neq 0\)，即除去\((3,1)\)、\((5,3)\)）．}
\fi
\end{ansblock}

\begin{ansblock}[2章-导-课时10-G5]
% ans:2章-导-课时10-G5
\ansitem{21}{x=4 或 x=−4（即 x=±4，两条平行于y轴的直线）}
\ifansbookdetail
\ansnote{详解}{设动点\(M(x,y)\)．\(M\)到\(y\)轴距离为\(|x|\)，条件\(|x|=4\)即\(x=4\)或\(x=-4\)．检验两性：直线\(x=\pm 4\)上任意一点到\(y\)轴距离均为4（纯粹性成立）；到\(y\)轴距离为4的点横坐标为\(\pm 4\)，都在两直线上（完备性成立）．故轨迹方程为\(x=\pm 4\)．}
\fi
\end{ansblock}

\jietitle{2.5.1 椭圆的标准方程}

\begin{ansblock}[2章-练-课时11-简7]
% ans:2章-练-课时11-简7
\ansitem{1}{(1)选①：椭圆x²/36+y²/20=1；选②：线段F₁F₂，方程y=0（−4≤x≤4）；(2)0＜a＜4无轨迹；a=4线段y=0（−4≤x≤4）；a＞4椭圆x²/a²+y²/(a²−16)=1}
\ifansbookdetail
\ansnote{详解}{(1)选①：\(12>|F_1F_2|=8\)，轨迹是以\(F_1\)、\(F_2\)为焦点的椭圆，\(c=4\)、\(a=6\)、\(b^{2}=36-16=20\)：\(x^{2}/36+y^{2}/20=1\)．选②：\(8=|F_1F_2|\)，距离和等于焦距，\(M\)只能在线段\(F_1F_2\)上（三角形不等式取等），轨迹为线段，方程\(y=0(-4\leqslant x\leqslant 4)\)．(2)按\(2a\)与\(|F_1F_2|=8\)比较三段：\(0<a<4\)时距离和小于焦距，无点满足，无轨迹；\(a=4\)同上为线段；\(a>4\)为椭圆，\(c=4\)、\(b^{2}=a^{2}-16\)：\(x^{2}/a^{2}+y^{2}/(a^{2}-16)=1\)．综上如答案．}
\fi
\end{ansblock}

\begin{ansblock}[2章-练-课时11-简8]
% ans:2章-练-课时11-简8
\ansitem{2}{x²/9＋y²/5＝1}
\ifansbookdetail
\ansnote{详解}{\(|F_1F_2|=4<6\)，由椭圆定义：\(2a=6\)，\(a=3\)；\(c=2\)；\(b^{2}=a^{2}-c^{2}=9-4=5\)．焦点在\(x\)轴，方程\(x^{2}/9+y^{2}/5=1\)．验算：取顶点\((3,0)\)：\(|PF_1|+|PF_2|=5+1=6\)成立；取\((0,\sqrt{5})\)：\(2\times\sqrt{4+5}=6\)成立．辨析：若和\(=4\)为线段\(F_1F_2\)，若和\(<4\)无轨迹（定义辨析位已讲，本题为\(>4\)基础位）．}
\fi
\end{ansblock}

\begin{ansblock}[2章-练-课时11-简2]
% ans:2章-练-课时11-简2
\ansitem{3}{10+2√10}
\ifansbookdetail
\ansnote{详解}{\(a=5\)、\(b=3\)、\(c=4\)，\(A(4,0)\)恰为右焦点，记左焦点\(C(-4,0)\)．由定义\(|MA|=2a-|MC|=10-|MC|\)，故\(|MA|+|MB|=10+(|MB|-|MC|)\leqslant 10+|BC|=10+\sqrt{(2+4)^{2}+(2-0)^{2}}=10+\sqrt{40}=10+2\sqrt{10}\)．等号当\(M\)在射线\(BC\)（自\(B\)过\(C\)延长方向）与椭圆的交点处取得（\(C\)在椭圆内，射线必与椭圆交于一点），可取．最大值\(10+2\sqrt{10}\)．}
\fi
\end{ansblock}

\begin{ansblock}[2章-练-课时11-简3]
% ans:2章-练-课时11-简3
\ansitem{4}{A}
\ifansbookdetail
\ansnote{详解}{\(a=4\)、\(b=2\)．定义：\(|AF_1|+|AF_2|=|BF_1|+|BF_2|=8\)，两式相加并注意\(|AB|=|AF_1|+|BF_1|\)（\(A\)、\(F_1\)、\(B\)共线）：\(\triangle ABF_2\)周长\(|AF_2|+|BF_2|+|AB|=16=4a\)．故\(|AF_2|+|BF_2|=16-|AB|\)，要最大须过焦点的弦\(|AB|\)最小．竖直位置：\(l\)：\(x=-c=-2\sqrt{3}\)，\(y^{2}=4(1-12/16)=1\)，\(|AB|=2\)（通径\(2b^{2}/a=2\)）；水平位置：长轴弦\(|AB|=2a=8\)．通径为最短焦点弦（一般证明需弦长联立，2.5.3后；此处按两位置比较简验）．最大值\(=16-2=14\)，故选：A．}
\fi
\end{ansblock}

\begin{ansblock}[2章-练-课时11-简4]
% ans:2章-练-课时11-简4
\ansitem{5}{(m/4)√(n²−m²)}
\ifansbookdetail
\ansnote{详解}{\(n>m\)：\(A\)到两定点\(B\)、\(C\)距离和为常数\(n\)且大于\(|BC|\)，由定义\(A\)在以\(B\)、\(C\)为焦点的椭圆上（\(2a=n\)、\(2c=m\)，\(A\)不在直线\(BC\)上故去长轴端点）．以\(BC\)为底，底长\(m\)固定，面积最大即高最大；椭圆上点到焦点连线所在直线的最大距离为短半轴长\(b=\sqrt{a^{2}-c^{2}}=\sqrt{(n/2)^{2}-(m/2)^{2}}=\frac{1}{2}\sqrt{n^{2}-m^{2}}\)（短轴端点取得，该点不与\(B\)、\(C\)共线，合法）．\(S_{\max}=\frac{1}{2}\cdot m\cdot\frac{1}{2}\sqrt{n^{2}-m^{2}}=(m/4)\sqrt{n^{2}-m^{2}}\)．}
\fi
\end{ansblock}

\begin{ansblock}[2章-练-课时11-简1]
% ans:2章-练-课时11-简1
\ansitem{6}{x²/12+y²/4=1}
\ifansbookdetail
\ansnote{详解}{定义：\(2a=4\sqrt{3}\)，即\(a=2\sqrt{3}\)、\(a^{2}=12\)．设\(A(c,y_0)\)：\(c^{2}/a^{2}+y_0^{2}/b^{2}=1\)，即\(y_0^{2}=b^{2}(1-c^{2}/a^{2})=b^{4}/a^{2}\)，\(|y_0|=b^{2}/a\)；通径\(|AB|=2b^{2}/a=4\sqrt{3}/3\)，即\(b^{2}=a\cdot(2\sqrt{3})/3=2\sqrt{3}\cdot(2\sqrt{3})/3=4\)．椭圆方程\(x^{2}/12+y^{2}/4=1\)．验算：\(c^{2}=12-4=8\)，\(x=c=2\sqrt{2}\)代入得\(y^{2}=4(1-8/12)=4/3\)，\(|AB|=2\times(2/\sqrt{3})=4\sqrt{3}/3\)成立；顶点\((\sqrt{12},0)\)到两焦点\((\pm 2\sqrt{2},0)\)距离和\(=(2\sqrt{2}+\sqrt{12})+(\sqrt{12}-2\sqrt{2})=2\sqrt{12}=4\sqrt{3}\)成立．方程正确．}
\fi
\end{ansblock}

\begin{ansblock}[2章-练-课时11-简9]
% ans:2章-练-课时11-简9
\ansitem{7}{2＜m＜4}
\ifansbookdetail
\ansnote{详解}{椭圆成立需\(m-2>0\)且\(6-m>0\)，即\(2<m<6\)；焦点在\(y\)轴需\(y^{2}\)分母较大：\(6-m>m-2\)，即\(m<4\)．综上\(2<m<4\)．验算：\(m=3\)：\(x^{2}/1+y^{2}/3=1\)，\(a^{2}=3>b^{2}=1\)，焦点\((0,\pm\sqrt{2})\)在\(y\)轴成立；\(m=5\)：\(x^{2}/3+y^{2}/1=1\)焦点在\(x\)轴（排除端外值实证）成立．}
\fi
\end{ansblock}

\begin{ansblock}[2章-练-课时11-简10]
% ans:2章-练-课时11-简10
\ansitem{8}{x²/(16/3)＋y²/4＝1}
\ifansbookdetail
\ansnote{详解}{\(e=c/a=1/2\)，即\(a=2c\)，\(b^{2}=a^{2}-c^{2}=3c^{2}\)．设\(x^{2}/a^{2}+y^{2}/b^{2}=1\)，过\(P\)：\(4/a^{2}+1/b^{2}=1\)．代\(a^{2}=4c^{2}\)、\(b^{2}=3c^{2}\)：\(4/(4c^{2})+1/(3c^{2})=1\)，即\(\frac{1}{c^{2}}+\frac{1}{3c^{2}}=1\)，得\(\frac{4}{3c^{2}}=1\)，\(c^{2}=4/3\)．故\(a^{2}=16/3\)，\(b^{2}=4\)．方程\(x^{2}/(16/3)+y^{2}/4=1\)．验算：\(P(2,1)\)：\(4/(16/3)+1/4=12/16+4/16=1\)成立；\(e\)：\(c=2/\sqrt{3}\)，\(a=4/\sqrt{3}\)，\(e=1/2\)成立．}
\fi
\end{ansblock}

\begin{ansblock}[2章-练-课时11-简5]
% ans:2章-练-课时11-简5
\ansitem{9}{C}
\ifansbookdetail
\ansnote{详解}{定义\(|MF_1|+|MF_2|=2a\)，均值不等式\(|MF_1|\cdot|MF_2|\leqslant((|MF_1|+|MF_2|)/2)^{2}=a^{2}\)，等号当\(|MF_1|=|MF_2|\)（\(M\)为短轴端点，可取）．最大值\(a^{2}=8\)，即\(a=2\sqrt{2}\)（\(a>0\)）．故选：C．}
\fi
\end{ansblock}

\begin{ansblock}[2章-练-课时11-简6]
% ans:2章-练-课时11-简6
\ansitem{10}{ACD}
\ifansbookdetail
\ansnote{详解}{\(\angle F_1PF_2=90^{\circ}\)当且仅当\(P\)在以\(F_1F_2\)为直径的圆\(x^{2}+y^{2}=c^{2}\)上（圆周角定理，\(P\neq F_1\)、\(F_2\)自动满足）．设\(P(x,y)\)在椭圆上，则\(|OP|^{2}=x^{2}+y^{2}=x^{2}+b^{2}(1-x^{2}/a^{2})=b^{2}+x^{2}(1-b^{2}/a^{2})\)，\(x^{2}\in[0,a^{2}]\)且\(1-b^{2}/a^{2}>0\)，故\(|OP|^{2}\in[b^{2},a^{2}]\)，即椭圆上点到原点距离介于\(b\)与\(a\)之间．圆与椭圆有公共点当且仅当\(b^{2}\leqslant c^{2}\leqslant a^{2}\)，右半恒真，条件即\(b\leqslant c\)，即\(b^{2}\leqslant a^{2}-b^{2}\)，即\(a^{2}\geqslant 2b^{2}\)（存在性与上顶点处张角\(\angle F_1BF_2\geqslant 90^{\circ}\)同述）．逐项：A \(25\geqslant 18\)成立；B \(25\geqslant 32\)不成立；C \(18\geqslant 18\)成立；D \(16\geqslant 16\)成立．故选：ACD．}
\fi
\end{ansblock}

\begin{ansblock}[2章-练-课时11-中1]
% ans:2章-练-课时11-中1
\ansitem{11}{2}
\ifansbookdetail
\ansnote{详解}{向量\(PF_1\cdot\)向量\(PF_2=0\)即\(\angle F_1PF_2=90^{\circ}\)，则\(|PF_1|^{2}+|PF_2|^{2}=|F_1F_2|^{2}=4c^{2}\)．记\(m=|PF_1|\cdot|PF_2|\)：面积\(\frac{1}{2}m=4\)，即\(m=8\)．定义：\((|PF_1|+|PF_2|)^{2}=4a^{2}=|PF_1|^{2}+|PF_2|^{2}+2|PF_1||PF_2|=4c^{2}+16\)，即\(a^{2}-c^{2}=4\)，\(b^{2}=4\)，\(b>0\)得\(b=2\)．验算相容性：\(|PF_1|\)、\(|PF_2|\)是\(t^{2}-2at+8=0\)的正根，判别式\(4a^{2}-32\geqslant 0\)即\(a^{2}\geqslant 8\)（题设“存在\(P\)”即椭圆满足此，必要条件下\(b\)恒为2）．}
\fi
\end{ansblock}

\begin{ansblock}[2章-练-课时11-中2]
% ans:2章-练-课时11-中2
\ansitem{12}{√3/3}
\ifansbookdetail
\ansnote{详解}{记左焦点\(F_1\)．\(O\)既是\(PQ\)中点（\(l\)过中心、椭圆关于\(O\)对称）又是\(FF_1\)中点，故四边形\(PF_1QF\)对角线互相平分，为平行四边形．于是\(S_{\triangle PFQ}=S_{\triangle PF_1F}\)（对角线\(F_1F\)所分另一半等积）．平行四边形邻角互补：顶点\(P\)处内角\(\angle F_1PF=\pi-\angle PFQ=\pi/3\)．在\(\triangle PF_1F\)中：\(a=\sqrt{2}\)、\(b=1\)、\(c=1\)，\(|PF_1|+|PF|=2\sqrt{2}\)，\(|F_1F|=2\)，\(\angle F_1PF=\pi/3\)．余弦定理：\(4=|PF_1|^{2}+|PF|^{2}-2|PF_1||PF|\cos(\pi/3)=(2\sqrt{2})^{2}-3|PF_1||PF|\)，即\(|PF_1||PF|=4/3\)．\(S_{\triangle PF_1F}=\frac{1}{2}\cdot\frac{4}{3}\cdot\sin(\pi/3)=\sqrt{3}/3\)，即\(\triangle PFQ\)的面积．验算：积\(4/3\leqslant((2\sqrt{2})/2)^{2}=4\)成立，两半径可正实存在．}
\fi
\end{ansblock}

\begin{ansblock}[2章-练-课时11-中3]
% ans:2章-练-课时11-中3
\ansitem{13}{C}
\ifansbookdetail
\ansnote{详解}{\(A\)、\(C\)恰为椭圆两焦点（\(c=1\)，\(a=2\)）．\(B\)在椭圆上且\(\triangle ABC\)非退化（\(B\)不为长轴端点，\(\sin B\neq 0\)）．正弦定理：\(|BC|=2R'\sin A\)、\(|AB|=2R'\sin C\)、\(|AC|=2R'\sin B\)（\(R'\)为\(\triangle ABC\)外接圆半径），故\((\sin A+\sin C)/\sin B=(|BC|+|AB|)/|AC|\)．由定义\(|AB|+|BC|=2a=4\)，\(|AC|=2c=2\)，比值\(=4/2=2\)，与\(B\)位置无关．故选：C．}
\fi
\end{ansblock}

\begin{ansblock}[2章-练-课时11-中4]
% ans:2章-练-课时11-中4
\ansitem{14}{(0,±1)}
\ifansbookdetail
\ansnote{详解}{\(a^{2}=3\)、\(b^{2}=1\)、\(c^{2}=2\)：\(F_1(-\sqrt{2},0)\)、\(F_2(\sqrt{2},0)\)．设\(A(x_1,y_1)\)、\(B(x_2,y_2)\)．由向量等式：\((x_1+\sqrt{2},y_1)=5(x_2-\sqrt{2},y_2)\)，即\(x_2=(x_1+6\sqrt{2})/5\)、\(y_2=y_1/5\)．两点均在椭圆上：\(x_1^{2}/3+y_1^{2}=1\)且\((x_1+6\sqrt{2})^{2}/75+y_1^{2}/25=1\)．后者乘75：\((x_1+6\sqrt{2})^{2}+3y_1^{2}=75\)；由前者\(3y_1^{2}=3-x_1^{2}\)，代入：\(x_1^{2}+12\sqrt{2}x_1+72+3-x_1^{2}=75\)，即\(12\sqrt{2}x_1=0\)，\(x_1=0\)，\(y_1^{2}=1\)，\(y_1=\pm 1\)．\(A(0,\pm 1)\)．验算：此时\(B=(6\sqrt{2}/5,\pm 1/5)\)：\((72/25)/3+1/25=24/25+1/25=1\)成立，\(B\)确在椭圆上；\(|F_1A|=\sqrt{2+1}=\sqrt{3}\)，\(6\sqrt{2}/5-\sqrt{2}=\sqrt{2}/5\)，\(|F_2B|=\sqrt{2/25+1/25}=\sqrt{3}/5\)，比值\(\sqrt{3}/(\sqrt{3}/5)=5\)成立．}
\fi
\end{ansblock}

\begin{ansblock}[2章-练-课时11-难1]
% ans:2章-练-课时11-难1
\ansitem{15}{2π}
\ifansbookdetail
\ansnote{详解}{\(a=2\)、\(b=\sqrt{3}\)、\(c=1\)，\(|F_1F_2|=2\)．设\(\angle F_1PF_2=\theta\)．①\(\theta\)范围：\(P\)在短轴端点时\(\theta\)最大，此时\(|PF_1|=|PF_2|=2\)、\(|F_1F_2|=2\)为等边三角形，\(\theta_{\max}=\pi/3\)，故\(0<\theta\leqslant\pi/3\)．②外接圆：\(2R=|F_1F_2|/\sin\theta=2/\sin\theta\)，即\(R=1/\sin\theta\)．③内切圆：\(S_{\triangle F_1PF_2}=\frac{1}{2}|PF_1||PF_2|\sin\theta=b^{2}\tan(\theta/2)=3\tan(\theta/2)\)；又\(r=2S/\)周长\(=2\cdot 3\tan(\theta/2)/(2a+2c)=6\tan(\theta/2)/6=\tan(\theta/2)\)．④\(S_1+2S_2=\pi(R^{2}+2r^{2})=\pi[1/\sin^{2}\theta+2\tan^{2}(\theta/2)]\)．记\(t=\tan(\theta/2)\in(0,\sqrt{3}/3]\)（\(\theta\in(0,\pi/3]\)），\(\sin\theta=2t/(1+t^{2})\)：\(1/\sin^{2}\theta=(1+t^{2})^{2}/(4t^{2})\)，\(S_1+2S_2=\pi[(1+t^{2})^{2}/(4t^{2})+2t^{2}]=\pi[1/(4t^{2})+1/2+9t^{2}/4]\geqslant\pi\left[1/2+2\sqrt{\frac{1}{4t^{2}}\cdot\frac{9t^{2}}{4}}\right]=\pi(1/2+3/2)=2\pi\)，等号当\(1/(4t^{2})=9t^{2}/4\)，即\(t^{2}=1/3\)，\(t=\sqrt{3}/3\)为上限\(\theta=\pi/3\)（\(P\)为短轴端点）可取．最小值\(2\pi\)．}
\fi
\end{ansblock}

\begin{ansblock}[2章-练-课时11-难2]
% ans:2章-练-课时11-难2
\ansitem{16}{A}
\ifansbookdetail
\ansnote{详解}{\(O\)为\(F_1F_2\)中点，重心\(G\)在中线\(PO\)上且\(PG:GO=2:1\)．延长\(PI\)交\(x\)轴于\(Q\)（\(IQ\)在\(x\)轴内）．\(IG\parallel x\)轴，则在\(\triangle PQO\)中\(PI:IQ=PG:GO=2:1\)，即\(PQ:IQ=3\)．两三角形\(\triangle PF_1F_2\)与\(\triangle IF_1F_2\)同底\(F_1F_2\)，面积比\(=PQ:IQ=3\)．用内切圆半径\(r\)表示：\(S_{\triangle PF_1F_2}=\frac{1}{2}r(|PF_1|+|PF_2|+|F_1F_2|)\)（内心分三块）、\(S_{\triangle IF_1F_2}=\frac{1}{2}r|F_1F_2|\)，比\(=(|PF_1|+|PF_2|+|F_1F_2|)/|F_1F_2|=3\)，即\((|PF_1|+|PF_2|)/|F_1F_2|=2\)，即\(2a/2c=2\)，\(c/a=1/2\)．离心率为\(1/2\)，故选：A．}
\fi
\end{ansblock}

\begin{ansblock}[2章-导-课时11-G1]
% ans:2章-导-课时11-G1
\ansitem{17}{C}
\ifansbookdetail
\ansnote{详解}{A：\(|F_1F_2|=8\)，距离和\(=8=|F_1F_2|\)，轨迹是线段\(F_1F_2\)，错；B：和\(6<8=|F_1F_2|\)，满足条件的点不存在，错；C：\(M(5,3)\)到两定点距离和\(=\sqrt{9^{2}+3^{2}}+\sqrt{1^{2}+3^{2}}=3\sqrt{10}+\sqrt{10}=4\sqrt{10}>8\)，和为常数且大于\(|F_1F_2|\)，轨迹是椭圆，对；D：到两点距离相等的点轨迹是线段\(F_1F_2\)的垂直平分线（\(y\)轴），错．故选：C．}
\fi
\end{ansblock}

\begin{ansblock}[2章-导-课时11-G2]
% ans:2章-导-课时11-G2
\ansitem{18}{C}
\ifansbookdetail
\ansnote{详解}{比较定值\(a^{2}+1\)与焦距\(2a\)：\(a^{2}+1-2a=(a-1)^{2}\geqslant 0\)，故\(a^{2}+1\geqslant 2a\)恒成立．当\(a>0\)且\(a\neq 1\)时\((a-1)^{2}>0\)，\(|PF_1|+|PF_2|>|F_1F_2|\)，轨迹为椭圆；当\(a=1\)时等号成立，\(|PF_1|+|PF_2|=|F_1F_2|\)，轨迹为线段\(F_1F_2\)．合起来：椭圆或线段，故选：C．}
\fi
\end{ansblock}

\begin{ansblock}[2章-导-课时11-G3]
% ans:2章-导-课时11-G3
\ansitem{19}{C}
\ifansbookdetail
\ansnote{详解}{由定义\(|AF_1|+|AF_2|=2a=10\)，\(|BF_1|+|BF_2|=10\)；\(A\)、\(F_1\)、\(B\)共线故\(|AB|=|AF_1|+|BF_1|\)．周长\(=|AB|+|AF_2|+|BF_2|=|AF_1|+|BF_1|+|AF_2|+|BF_2|=20\)（\(=4a\)，与直线位置无关）．故选：C．}
\fi
\end{ansblock}

\begin{ansblock}[2章-导-课时11-G4]
% ans:2章-导-课时11-G4
\ansitem{20}{25}
\ifansbookdetail
\ansnote{详解}{\(a=5\)，由定义\(|PF_1|+|PF_2|=2a=10\)．均值不等式：\(m=|PF_1|\cdot|PF_2|\leqslant((|PF_1|+|PF_2|)/2)^{2}=25\)，等号当\(|PF_1|=|PF_2|=5\)．存在性验证：短轴端点\(P(0,\pm 3)\)到焦点\((\pm 4,0)\)的距离\(=\sqrt{16+9}=5\)，两边皆5，等号可取．\(m\)最大值为25．}
\fi
\end{ansblock}

\begin{ansblock}[2章-导-课时11-G5]
% ans:2章-导-课时11-G5
\ansitem{21}{±√3/4}
\ifansbookdetail
\ansnote{详解}{\(c^{2}=12-3=9\)，焦点\((\pm 3,0)\)．设另一定点为\(F_2\)．\(M\)为\(PF_1\)中点、\(O\)为\(F_1F_2\)中点，则\(OM\)为\(\triangle PF_1F_2\)的中位线，\(OM\parallel PF_2\)且\(OM\)在\(y\)轴上，故\(PF_2\perp x\)轴，\(P\)横坐标为\(\pm 3\)．代入：\(9/12+y^{2}/3=1\)，即\(y^{2}=3/4\)，\(y_P=\pm\sqrt{3}/2\)．\(M\)纵坐标\(=y_P/2=\pm\sqrt{3}/4\)．}
\fi
\end{ansblock}

\jietitle{2.5.2 椭圆的几何性质}

\begin{ansblock}[2章-练-课时12-简1]
% ans:2章-练-课时12-简1
\ansitem{1}{x²/12＋y²/4＝1}
\ifansbookdetail
\ansnote{详解}{\(2a=4\sqrt{3}\)，即\(a=2\sqrt{3}\)，\(a^{2}=12\)；\(e=c/a=\sqrt{6}/3\)，即\(c^{2}=(6/9)\times 12=8\)，\(c=2\sqrt{2}\)；\(b^{2}=a^{2}-c^{2}=12-8=4\)．题设已限定焦点在\(x\)轴（方程形如\(x^{2}/a^{2}+y^{2}/b^{2}\)，\(a>b>0\)），故椭圆\(C\)的方程为\(x^{2}/12+y^{2}/4=1\)．}
\fi
\end{ansblock}

\begin{ansblock}[2章-练-课时12-简2]
% ans:2章-练-课时12-简2
\ansitem{2}{x²/3＋y²＝1（答案不唯一）}
\ifansbookdetail
\ansnote{详解}{设\(x^{2}/a^{2}+y^{2}/b^{2}=1\)（\(a>b>0\)），\(e^{2}=1-b^{2}/a^{2}=6/9=2/3\)，即\(a^{2}=3b^{2}\)．取\(b^{2}=1\)、\(a^{2}=3\)即得\(x^{2}/3+y^{2}=1\)；任取\(b^{2}=t\)（\(t>0\)）、\(a^{2}=3t\)均合题．}
\fi
\end{ansblock}

\begin{ansblock}[2章-练-课时12-简3]
% ans:2章-练-课时12-简3
\ansitem{3}{(−√2,√2)}
\ifansbookdetail
\ansnote{详解}{点在椭圆内部即代入左端小于1：\(m^{2}/4+1/2<1\)，即\(m^{2}/4<1/2\)，\(m^{2}<2\)，故\(-\sqrt{2}<m<\sqrt{2}\)．}
\fi
\end{ansblock}

\begin{ansblock}[2章-练-课时12-简4]
% ans:2章-练-课时12-简4
\ansitem{4}{A}
\ifansbookdetail
\ansnote{详解}{\(a^{2}=9\)、\(a=3\)，右焦点\(F_2(c,0)\)．椭圆上点到焦点距离的最小值在长轴（同侧）顶点取得：\(|PF_2|_{\min}=a-c\)．由\(3-c=3-2\sqrt{2}\)得\(c=2\sqrt{2}\)，\(b^{2}=a^{2}-c^{2}=9-8=1\)，\(b=1\)（\(b>0\)）．故选：A．}
\fi
\end{ansblock}

\begin{ansblock}[2章-练-课时12-简5]
% ans:2章-练-课时12-简5
\ansitem{5}{C}
\ifansbookdetail
\ansnote{详解}{由蒙日圆结论（本课时讲解位2.5.2.20）：\(x^{2}/a^{2}+y^{2}/b^{2}=1\)的两条互相垂直切线的交点轨迹为\(x^{2}+y^{2}=a^{2}+b^{2}\)．题设圆即该椭圆的蒙日圆，故\(6+b^{2}=8\)，\(b^{2}=2\)（满足\(0<b<\sqrt{6}\)）．故选：C．}
\fi
\end{ansblock}

\begin{ansblock}[2章-练-课时12-简6]
% ans:2章-练-课时12-简6
\ansitem{6}{(1)无数个，如 x²/3＋y²/8＝1、x²/5＋y²/10＝1（任写两个满足 a²−b²＝5 且焦点在 y 轴者均可）；(2)一个，x²/20＋y²/25＝1}
\ifansbookdetail
\ansnote{详解}{\(C\)化标准式\(x^{2}/4+y^{2}/9=1\)，焦点在\(y\)轴，\(c^{2}=9-4=5\)．(1)同焦点即\(a^{2}-b^{2}=5\)且焦点在\(y\)轴，即\(x^{2}/m+y^{2}/(m+5)=1\)（\(m>0\)）对任意\(m>0\)均可，故有无数个；取\(m=3\)、\(m=5\)得\(x^{2}/3+y^{2}/8=1\)、\(x^{2}/5+y^{2}/10=1\)．(2)代入\((4,\sqrt{5})\)：\(16/m+5/(m+5)=1\)，即\(16(m+5)+5m=m(m+5)\)，即\(m^{2}-16m-80=0\)，\((m-20)(m+4)=0\)，\(m=20\)或\(m=-4\)（舍，须\(m>0\)）．故唯一，方程\(x^{2}/20+y^{2}/25=1\)．}
\fi
\end{ansblock}

\begin{ansblock}[2章-练-课时12-简7]
% ans:2章-练-课时12-简7
\ansitem{7}{(−1,1)}
\ifansbookdetail
\ansnote{详解}{设\(A(x_1,y_1)\)、\(B(x_2,y_2)\)，“不关于长轴（\(x\)轴）对称”即\(y_1\neq -y_2\)（且\(A\)、\(B\)为不同两点）．\(|AM|=|BM|\)即\((x_1-m)^{2}+y_1^{2}=(x_2-m)^{2}+y_2^{2}\)，即\((x_1-x_2)(x_1+x_2-2m)=y_2^{2}-y_1^{2}\)．椭圆上\(y^{2}=12-(3/4)x^{2}\)，故\(y_2^{2}-y_1^{2}=(3/4)(x_1^{2}-x_2^{2})=(3/4)(x_1-x_2)(x_1+x_2)\)．若\(x_1=x_2\)，则由上式\(y_1^{2}=y_2^{2}\)，结合“不关于长轴对称”得\(y_1=y_2\)，\(A\)、\(B\)重合，矛盾；故\(x_1\neq x_2\)，可约去\((x_1-x_2)\)：\(x_1+x_2-2m=(3/4)(x_1+x_2)\)，即\(2m=(1/4)(x_1+x_2)\)，\(m=(x_1+x_2)/8\)．又\(-4\leqslant x_i\leqslant 4\)且\(x_1\neq x_2\)，故\(-8<x_1+x_2<8\)（两端点\(x_1=x_2=\pm 4\)不能同时取，等号不可达），得\(-1<m<1\)．}
\fi
\end{ansblock}

\begin{ansblock}[2章-练-课时12-简8]
% ans:2章-练-课时12-简8
\ansitem{8}{x²/4＋y²＝1}
\ifansbookdetail
\ansnote{详解}{\(P_3\)与\(P_4\)关于\(y\)轴对称，而椭圆关于\(y\)轴对称，故\(P_3\)、\(P_4\)同“在”或同“不在”；“恰有三点”须\(P_3\)、\(P_4\)都在其上，第三点为\(P_1\)、\(P_2\)之一．若\(P_1\)在\(C\)上，则与\(P_4\)同横坐标\(x=1\)而纵坐标不同（\(1\neq\sqrt{3}/2\)），同一\(x\)对应椭圆上两点纵坐标互为相反数，矛盾（亦可直接算：\(1/a^{2}+1/b^{2}=1\)与\(1/a^{2}+3/(4b^{2})=1\)相减得\(b=0\)，舍）；故第三点是\(P_2\)．由\(P_2(0,1)\)：\(1/b^{2}=1\)，即\(b^{2}=1\)；由\(P_4(1,\sqrt{3}/2)\)：\(1/a^{2}+3/4=1\)，即\(a^{2}=4\)（\(a>b\)成立）．所以\(C\)：\(x^{2}/4+y^{2}=1\)．回代检验：\(P_1(1,1)\)使\(1/4+1=5/4\neq 1\)，恰不在其上成立．}
\fi
\end{ansblock}

\begin{ansblock}[2章-练-课时12-简9]
% ans:2章-练-课时12-简9
\ansitem{9}{x＋y−1＝0}
\ifansbookdetail
\ansnote{详解}{\(a^{2}=4\)、\(b^{2}=3\)，即\(c^{2}=1\)，右焦点\(F(1,0)\)．设\(l\)的斜率为\(k\)（\(k\)必存在且\(k\neq 0\)：\(P\)在第二象限、\(F\)在\(x\)轴上），\(l\)：\(y=k(x-1)\)．设\(Q(x_2,y_2)\)，则\(Q'(x_2,-y_2)\)，且\(P\)、\(Q\)、\(F\)共线，故直线\(PQ\)即\(l\)，方向向量\((1,k)\)；向量\(FQ'=(x_2-1,-y_2)=(x_2-1)(1,-k)\)（用\(y_2=k(x_2-1)\)），即\(FQ'\)的方向为\((1,-k)\)（\(x_2\neq 1\)，否则\(Q\)在\(x\)轴上、\(Q\)为顶点使\(P\)不在第二象限，舍）．\(PQ\perp FQ'\)即\((1,k)\cdot(1,-k)=0\)，即\(1-k^{2}=0\)，\(k=\pm 1\)．由\(P\)在第二象限定号：\(k=1\)时\(l\)：\(y=x-1\)与\(x^{2}/4+y^{2}/3=1\)联立得\(7x^{2}-8x-8=0\)，两根之积\(-8/7<0\)，其负根对应\(y=x-1<0\)，点在第三象限，不合；\(k=-1\)时\(l\)：\(y=-(x-1)\)同得\(7x^{2}-8x-8=0\)，负根\(x=(4-6\sqrt{2})/7\)对应\(y=1-x>0\)，点在第二象限成立．故\(k=-1\)，\(l\)：\(y=-(x-1)\)，即\(x+y-1=0\)．}
\fi
\end{ansblock}

\begin{ansblock}[2章-练-课时12-简10]
% ans:2章-练-课时12-简10
\ansitem{10}{C}
\ifansbookdetail
\ansnote{详解}{同一中心天体（太阳）即\(a^{3}/T^{2}\)同值，故\(T_2^{2}=(a_2^{3}/a_1^{3})\times T_1^{2}=(60/1.5)^{3}\times 1^{2}=40^{3}=64000\)，\(T_2=\sqrt{64000}=80\sqrt{10}\approx 80\times 3.1=248\)（年）．故选：C．}
\fi
\end{ansblock}

\begin{ansblock}[2章-练-课时12-中1]
% ans:2章-练-课时12-中1
\ansitem{11}{[√5,3]}
\ifansbookdetail
\ansnote{详解}{\(|PA|+|PB|=6>|AB|=4\)，故\(P\)的轨迹是以\(A\)、\(B\)为焦点、\(2a=6\)的椭圆．以\(M\)为原点、\(AB\)所在直线为\(x\)轴建系，则\(A(-2,0)\)、\(B(2,0)\)，\(c=2\)、\(a=3\)、\(b^{2}=a^{2}-c^{2}=9-4=5\)，轨迹\(x^{2}/9+y^{2}/5=1\)．\(M\)即椭圆中心，\(|PM|\)为中心到椭圆上点的距离：\(|PM|^{2}=x^{2}+y^{2}=x^{2}+5(1-x^{2}/9)=5+(4/9)x^{2}\)，\(x^{2}\in[0,9]\)，故\(|PM|^{2}\in[5,9]\)，\(|PM|\in[\sqrt{5},3]\)（最小在短轴端点\(=b\)，最大在长轴端点\(=a\)）．}
\fi
\end{ansblock}

\begin{ansblock}[2章-练-课时12-中2]
% ans:2章-练-课时12-中2
\ansitem{12}{[4＋2√3, 8)}
\ifansbookdetail
\ansnote{详解}{\(a=2\)、\(b=\sqrt{3}\)、\(c=1\)，右焦点\(F(1,0)\)，记左焦点\(F_1(-1,0)\)．\(l\)过原点，\(A\)、\(B\)关于中心\(O\)对称，故\(O\)是\(AB\)的中点，又\(O\)是\(FF_1\)的中点，所以四边形\(AFBF_1\)是平行四边形，\(|BF|=|AF_1|\)．周长\(L=|AB|+|AF|+|BF|=|AB|+(|AF|+|AF_1|)=|AB|+2a=|AB|+4\)．求\(|AB|\)：\(l\)的斜率不为0（斜率为0时\(A\)、\(B\)为长轴端点，与\(F\)共线，三角形退化），设\(l\)：\(x=my\)，代入得\(y^{2}=12/(3m^{2}+4)\)，\(|AB|=2\sqrt{x^{2}+y^{2}}=2\sqrt{(m^{2}+1)y^{2}}=2\sqrt{12(m^{2}+1)/(3m^{2}+4)}=4\sqrt{1-1/(3m^{2}+4)}\)．\(m^{2}\geqslant 0\)，即\(3m^{2}+4\geqslant 4\)，\(1-1/(3m^{2}+4)\in[3/4,1)\)，故\(|AB|\in[2\sqrt{3},4)\)．所以\(L=|AB|+4\in[4+2\sqrt{3},8)\)（上端\(|AB|\to 2a=4\)对应直线趋于长轴、三角形退化，取开）．}
\fi
\end{ansblock}

\begin{ansblock}[2章-练-课时12-中3]
% ans:2章-练-课时12-中3
\ansitem{13}{最大 12(2＋√2)/5，最小 12(2−√2)/5}
\ifansbookdetail
\ansnote{详解}{设\(P(4\cos\theta,3\sin\theta)\)，则\(d=|3\cdot 4\cos\theta-4\cdot 3\sin\theta-24|/\sqrt{3^{2}+4^{2}}=|12\cos\theta-12\sin\theta-24|/5=|12\sqrt{2}\cos(\theta+\pi/4)-24|/5\)．因\(12\sqrt{2}<24\)，\(\cos(\theta+\pi/4)\in[-1,1]\)时\(12\sqrt{2}\cos(\theta+\pi/4)-24\)恒为负，绝对值直接翻号：\(d=(24-12\sqrt{2}\cos(\theta+\pi/4))/5\)．\(\cos(\theta+\pi/4)=-1\)时\(d_{\max}=(24+12\sqrt{2})/5=12(2+\sqrt{2})/5\)；\(\cos(\theta+\pi/4)=1\)时\(d_{\min}=(24-12\sqrt{2})/5=12(2-\sqrt{2})/5\)．（另法：设与\(3x-4y=24\)平行的切线\(3x-4y=t\)，由判别式为0得\(t=\pm 12\sqrt{2}\)，两切线到已知直线的距离\(|24\mp 12\sqrt{2}|/5\)即最小、最大值，结果一致．）}
\fi
\end{ansblock}

\begin{ansblock}[2章-练-课时12-中4]
% ans:2章-练-课时12-中4
\ansitem{14}{(1)2√3 km²；(2)24/7 km}
\ifansbookdetail
\ansnote{详解}{设平行四边形为\(ADBC\)（\(AB\)、\(CD\)为两条对角线），周长8 km即相邻两边之和\(|CA|+|CB|=4\) km，故另两顶点\(C\)（及\(D\)）在以\(A\)、\(B\)为焦点的椭圆上：\(2a=4\)、\(2c=|AB|=2\)，\(a=2\)、\(c=1\)、\(b^{2}=3\)．以\(AB\)所在直线为\(x\)轴、\(AB\)中点为原点建系，\(A(-1,0)\)、\(B(1,0)\)，椭圆\(x^{2}/4+y^{2}/3=1\)（\(y\neq 0\)，退化时不成四边形）．(1)平行四边形\(ADBC\)的面积\(=2S_{\triangle ABC}=2\times(1/2\times|AB|\times|y_C|)=2|y_C|\)，而\(C\)在椭圆上即\(|y_C|\leqslant b=\sqrt{3}\)，当\(C\)为短轴端点时取等，故最大面积\(=2\sqrt{3}\) km\(^{2}\)．(2)水沟过\(A(-1,0)\)且与\(AB\)（\(x\)轴）成45°角，取斜率1，方程\(y=x+1\)（取斜率\(-1\)时由对称性弦长相同）．“被农艺园围住的部分”即直线被椭圆截得的弦（口径已由亲算⑥注记2钉死，非“最大面积菱形”口径）：代入\(x^{2}/4+(x+1)^{2}/3=1\)，即\(3x^{2}+4(x^{2}+2x+1)=12\)，即\(7x^{2}+8x-8=0\)，\(x_1+x_2=-8/7\)，\(x_1x_2=-8/7\)，\(|x_1-x_2|=\sqrt{(8/7)^{2}+4\times 8/7}=\sqrt{288/49}=(12\sqrt{2})/7\)，弦长\(=\sqrt{1+k^{2}}\,|x_1-x_2|=\sqrt{2}\times(12\sqrt{2})/7=24/7\) km．}
\fi
\end{ansblock}

\begin{ansblock}[2章-练-课时12-难1]
% ans:2章-练-课时12-难1
\ansitem{15}{6}
\ifansbookdetail
\ansnote{详解}{\(A(0,b)\)、\(C(0,-b)\)，设\(B(x_1,y_1)\)、\(D(x_2,y_2)\)．\(\angle BAD=\angle BCD=90^{\circ}\)即\(A\)、\(C\)都在以\(BD\)为直径的圆上，即\(A\)、\(B\)、\(C\)、\(D\)四点共圆，圆心为\(BD\)的中点\(((x_1+x_2)/2,(y_1+y_2)/2)\)．又该圆过\(A\)、\(C\)，圆心在弦\(AC\)的垂直平分线（即\(x\)轴）上，故\((y_1+y_2)/2=0\)，即\(y_2=-y_1\)．面积条件：\(\triangle ABC\)与\(\triangle ADC\)共底\(AC\)（在\(y\)轴上），故\(S_{\triangle ABC}/S_{\triangle ADC}=|x_1|/|x_2|=2\)，即\(|x_1|=2|x_2|\)；因\(B\)、\(D\)分居\(y\)轴两侧，取\(x_1=-2x_2\)．于是圆心为\((x_1/4,0)\)，半径平方\(R^{2}=(x_1-x_1/4)^{2}+y_1^{2}=(9/16)x_1^{2}+y_1^{2}\)．把\(A(0,b)\)代入圆方程：\((x_1/4)^{2}+b^{2}=(9/16)x_1^{2}+y_1^{2}\)，即\(b^{2}=(1/2)x_1^{2}+y_1^{2}\)……①又\(B\)在椭圆上：\(x_1^{2}/18+y_1^{2}/b^{2}=1\)，配合①消\(x_1\)（\(x_1^{2}=2(b^{2}-y_1^{2})\)）：\((b^{2}-y_1^{2})/9+y_1^{2}/b^{2}=1\)，乘\(9b^{2}\)得\(b^{4}-b^{2}y_1^{2}+9y_1^{2}-9b^{2}=0\)，即\((b^{2}-9)(b^{2}-y_1^{2})=0\)．\(B\)异于上、下顶点，故\(y_1^{2}<b^{2}\)，\(b^{2}-y_1^{2}\neq 0\)，只能\(b^{2}=9\)，\(b=3\)，短轴长\(2b=6\)．（存在性核验：\(b^{2}=9\)时①化为\(x_1^{2}=18-2y_1^{2}\)，与椭圆\(x^{2}/18+y^{2}/9=1\)同解，构型可行．取\(y_1=0\)：\(B(3\sqrt{2},0)\)、\(D(-3\sqrt{2}/2,0)\)、\(A(0,3)\)、\(C(0,-3)\)，\(AB\cdot AD=(3\sqrt{2},-3)\cdot(-3\sqrt{2}/2,-3)=-9+9=0\)成立，\(\angle BAD=90^{\circ}\)，同理\(\angle BCD=90^{\circ}\)；\(S_{\triangle ABC}=1/2\times 6\times 3\sqrt{2}=9\sqrt{2}\)、\(S_{\triangle ADC}=1/2\times 6\times(3\sqrt{2}/2)=9\sqrt{2}/2\)，比值2成立．）}
\fi
\end{ansblock}

\begin{ansblock}[2章-练-课时12-难2]
% ans:2章-练-课时12-难2
\ansitem{16}{AD}
\ifansbookdetail
\ansnote{详解}{\(e=\sqrt{2}/2\)，即\(c^{2}/a^{2}=1/2\)，\(a^{2}=2b^{2}\)，\(c^{2}=a^{2}-b^{2}=b^{2}\)，即\(a=\sqrt{2}b\)、\(c=b\)．A：蒙日圆半径平方\(r^{2}=a^{2}+b^{2}=3b^{2}\)，圆心为中心，故\(x^{2}+y^{2}=3b^{2}\)，正确．B：把\(P(b,a)\)代入\(l\)：\(b\cdot b+a\cdot a-a^{2}-b^{2}=0\)，故\(P(b,a)\)在\(l\)上；又\(b^{2}+a^{2}=a^{2}+b^{2}=r^{2}\)，故\(P\)在蒙日圆上，过\(P\)的两条切线互相垂直，取\(A\)、\(B\)为两切点时\(PA\cdot PB=0\)，“任意一点恒有\(PA\cdot PB>0\)”不成立，错误．C：\(|AF_1|+|AF_2|=2a\)，故\(|AH|-|AF_2|=|AH|+|AF_1|-2a\geqslant d(F_1,l)-2a\)．\(d(F_1,l)=|b\cdot(-c)+a\cdot 0-a^{2}-b^{2}|/\sqrt{a^{2}+b^{2}}=(bc+a^{2}+b^{2})/\sqrt{a^{2}+b^{2}}=(b^{2}+2b^{2}+b^{2})/(\sqrt{3}b)=4b/\sqrt{3}=(4\sqrt{3}/3)b\)．故\(|AH|-|AF_2|\)的最小值为\((4\sqrt{3}/3)b-2a=(4\sqrt{3}/3)b-2\sqrt{2}b\)，并非\((4\sqrt{3}/3)b\)，错误（\((4\sqrt{3}/3)b\)只是\(F_1\)到\(l\)的距离，漏减\(2a\)）．D：矩形四边均与\(C\)相切，则相邻两边为一对互相垂直的切线，其四个顶点都在蒙日圆上，即蒙日圆是矩形的外接圆，对角线为直径\(2r=2\sqrt{3}b\)．设两邻边为\(p\)、\(q\)，则\(p^{2}+q^{2}=(2r)^{2}=12b^{2}\)，面积\(pq\leqslant(p^{2}+q^{2})/2=6b^{2}\)，当\(p=q=\sqrt{6}b\)（正方形、对角线落在坐标轴上）时取等，正确．故选：AD．}
\fi
\end{ansblock}

\begin{ansblock}[2章-导-课时12-G1]
% ans:2章-导-课时12-G1
\ansitem{17}{A}
\ifansbookdetail
\ansnote{详解}{\(e=\sqrt{1-b^{2}/a^{2}}\)，四个椭圆\(a^{2}\)同为20，故\(b^{2}\)越小\(e\)越大、椭圆越扁．\(b^{2}/a^{2}\)依次为\(9/20<10/20<11/20<12/20\)，所以\(x^{2}/20+y^{2}/9=1\)的离心率最大，形状最扁．故选：A．}
\fi
\end{ansblock}

\begin{ansblock}[2章-导-课时12-G2]
% ans:2章-导-课时12-G2
\ansitem{18}{B}
\ifansbookdetail
\ansnote{详解}{\(2c=|AB|=4\)，\(c=2\)；\(2a=|PA|+|PB|\)，\(e=c/a\)最大即\(a\)最小，即\(|PA|+|PB|\)最小．求折线和最小：\(A\)关于\(l\)：\(x-y+4=0\)的对称点\(A'(x',y')\)，由公式\(x'=x_0-2\cdot(x_0-y_0+4)/2=-2-2=-4\)，\(y'=y_0+2\cdot(x_0-y_0+4)/2=0+2=2\)，即\(A'(-4,2)\)（与“\(AA'\perp l\)、\(AA'\)中点在\(l\)上”两条件验算一致）．于是\(|PA|+|PB|=|PA'|+|PB|\geqslant|A'B|=\sqrt{(-4-2)^{2}+2^{2}}=\sqrt{40}=2\sqrt{10}\)，当\(P\)为\(A'B\)与\(l\)的交点时取等．故\(2a\geqslant 2\sqrt{10}\)，\(a\geqslant\sqrt{10}\)，\(e=c/a\leqslant 2/\sqrt{10}=\sqrt{10}/5\)．故选：B．}
\fi
\end{ansblock}

\begin{ansblock}[2章-导-课时12-G3]
% ans:2章-导-课时12-G3
\ansitem{19}{B}
\ifansbookdetail
\ansnote{详解}{\(e=\sqrt{6}/3\)，即\(c^{2}/a^{2}=2/3\)，\(b^{2}/a^{2}=1/3\)，\(a=\sqrt{3}b\)，直线\(ax-by=0\)即\(y=(a/b)x=\sqrt{3}x\)．圆\(M\)化为\((x-m/2)^{2}+y^{2}=(m^{2}-1)/4\)，圆心\((m/2,0)\)、半径\(r=\frac{1}{2}\sqrt{m^{2}-1}\)（须\(m^{2}>1\)）．相切即圆心到直线\(\sqrt{3}x-y=0\)的距离等于半径：\(|\sqrt{3}\cdot m/2|/2=\frac{1}{2}\sqrt{m^{2}-1}\)，平方得\(3m^{2}/16=(m^{2}-1)/4\)，即\(3m^{2}=4m^{2}-4\)，\(m^{2}=4\)，\(m=\pm 2\)（满足\(m^{2}>1\)）．故选：B．}
\fi
\end{ansblock}

\begin{ansblock}[2章-导-课时12-G4]
% ans:2章-导-课时12-G4
\ansitem{20}{(±√(5−m),0)（0＜m＜5）；(0,±√(m−5))（m＞5）}
\ifansbookdetail
\ansnote{详解}{须\(m>0\)且\(m\neq 5\)．当\(0<m<5\)时焦点在\(x\)轴，\(a^{2}=5\)、\(b^{2}=m\)，\(c^{2}=5-m\)，焦点为\((\pm\sqrt{5-m},0)\)；当\(m>5\)时焦点在\(y\)轴，\(a^{2}=m\)、\(b^{2}=5\)，\(c^{2}=m-5\)，焦点为\((0,\pm\sqrt{m-5})\)．}
\fi
\end{ansblock}

\begin{ansblock}[2章-导-课时12-G5]
% ans:2章-导-课时12-G5
\ansitem{21}{(1)√5/3；(2)√3/2}
\ifansbookdetail
\ansnote{详解}{(1)\(a:b=3:2\)，取\(a=3k\)、\(b=2k\)，\(c^{2}=a^{2}-b^{2}=5k^{2}\)，\(e=c/a=\sqrt{5}k/3k=\sqrt{5}/3\)．(2)短轴端点\((0,\pm b)\)与焦点\((c,0)\)构成等边三角形：底边\(2b\)，腰\(\sqrt{b^{2}+c^{2}}=a\)，故\(a=2b\)；\(c^{2}=a^{2}-b^{2}=3b^{2}\)，\(e=c/a=\sqrt{3}b/2b=\sqrt{3}/2\)．}
\fi
\end{ansblock}

\jietitle{2.6.1 双曲线的标准方程}

\begin{ansblock}[2章-练-课时13-简1]
% ans:2章-练-课时13-简1
\ansitem{1}{(1)存在，轨迹是线段 AB；(2)存在，轨迹是以 B 为端点、指向 A 一侧反向（即线段 AB 的延长线方向）的射线；(3)存在，轨迹是以 A 为端点、沿线段 BA 延长线的射线；(4)不存在；(5)不存在}
\ifansbookdetail
\ansnote{详解}{记\(|AB|=6\)，对动点\(P\)总有\(|PA|+|PB|\geqslant|AB|\)（三角不等式）与\(\big||PA|-|PB|\big|\leqslant|AB|\)（两边之差小于等于第三边，等号当\(P\)、\(A\)、\(B\)共线）．(1)\(|PA|+|PB|=6=|AB|\)：等号成立当且仅当\(P\)在线段\(AB\)上，故轨迹为线段\(AB\)（椭圆定义中\(2a=2c\)的退化位）．(2)\(|PA|-|PB|=6=|AB|\)：须\(\big||PA|-|PB|\big|=|AB|\)，即\(P\)、\(A\)、\(B\)共线且\(P\)不在\(A\)、\(B\)之间；再由\(|PA|>|PB|\)知\(P\)靠\(B\)一侧，即\(P\)在射线（以\(B\)为端点、方向背离\(A\)）上，此即“线段\(AB\)的延长线”．(3)\(|PB|-|PA|=6\)：同(2)对称，\(P\)在以\(A\)为端点、背离\(B\)的射线上（“线段\(BA\)的延长线”）．(4)\(|PA|+|PB|=3<6=|AB|\)：与三角不等式矛盾，满足条件的\(P\)不存在（轨迹不存在）．(5)\(|PB|-|PA|=7>6=|AB|\)：与\(\big||PB|-|PA|\big|\leqslant|AB|\)矛盾，\(P\)不存在（轨迹不存在）．}
\fi
\end{ansblock}

\begin{ansblock}[2章-练-课时13-简9]
% ans:2章-练-课时13-简9
\ansitem{2}{x²/4−y²/5＝1}
\ifansbookdetail
\ansnote{详解}{\(|F_1F_2|=6>4>0\)，由双曲线定义：\(2a=4\)，\(a=2\)；\(c=3\)；\(b^{2}=c^{2}-a^{2}=9-4=5\)．焦点在\(x\)轴，方程\(x^{2}/4-y^{2}/5=1\)．辨析：若差为\(6=|F_1F_2|\)，轨迹为两条射线\(y=0\)（\(x\leqslant -3\)或\(x\geqslant 3\)）；若差大于\(6\)，则无轨迹．}
\fi
\end{ansblock}

\begin{ansblock}[2章-练-课时13-简3]
% ans:2章-练-课时13-简3
\ansitem{3}{(1)k＜−3 或 1＜k＜3；(2)1＜k＜3；(3)k＜−3}
\ifansbookdetail
\ansnote{详解}{两边乘\(-1\)化为标准形：\(x^{2}/(k-1)+y^{2}/(|k|-3)=1\)（须\(k-1\neq 0\)且\(|k|-3\neq 0\)）．(1)表示双曲线当且仅当两项系数异号，即\((k-1)(|k|-3)<0\)：\(k-1>0\)且\(|k|-3<0\)，即\(k>1\)且\(-3<k<3\)，得\(1<k<3\)；\(k-1<0\)且\(|k|-3>0\)，即\(k<1\)且（\(k>3\)或\(k<-3\)），得\(k<-3\)；故\(k<-3\)或\(1<k<3\)．(2)焦点在\(x\)轴上的双曲线当且仅当\(x^{2}\)项系数为正、\(y^{2}\)项为负：\(k-1>0\)且\(|k|-3<0\)，得\(1<k<3\)．(3)焦点在\(y\)轴上的双曲线当且仅当\(y^{2}\)项系数为正、\(x^{2}\)项为负：\(|k|-3>0\)且\(k-1<0\)，得\(k<-3\)（\(k>3\)与\(k<1\)无交）．}
\fi
\end{ansblock}

\begin{ansblock}[2章-练-课时13-简4]
% ans:2章-练-课时13-简4
\ansitem{4}{(1)x²/6 − y²/8 ＝1；(2)y²/9 − x²/16 ＝1；(3)x²/8 − y²/8 ＝1；(4)y²/16 − x²/9 ＝1}
\ifansbookdetail
\ansnote{详解}{(1)点\((\sqrt{6},0)\)在\(x\)轴上且在曲线上，故焦点在\(x\)轴、该点即右顶点，\(a^{2}=6\)；设\(x^{2}/6-y^{2}/b^{2}=1\)，代\((3,2)\)：\(9/6-4/b^{2}=1\)，即\(4/b^{2}=1/2\)，\(b^{2}=8\)，得\(x^{2}/6-y^{2}/8=1\)（核验\(9/6-4/8=3/2-1/2=1\)成立）．(2)焦点在\(y\)轴、\(c=5\)．设\(P(4\sqrt{3}/3,2\sqrt{3})\)，到两焦点\((0,\pm 5)\)的距离平方：\(|PF_1|^{2}=(4\sqrt{3}/3)^{2}+(2\sqrt{3}-5)^{2}=16/3+37-20\sqrt{3}=(10-3\sqrt{3})^{2}/3\)，故\(|PF_1|=10\sqrt{3}/3-3\)；\(|PF_2|^{2}=(10+3\sqrt{3})^{2}/3\)，故\(|PF_2|=10\sqrt{3}/3+3\)；于是\(2a=\big||PF_2|-|PF_1|\big|=6\)，即\(a=3\)、\(b^{2}=c^{2}-a^{2}=25-9=16\)，得\(y^{2}/9-x^{2}/16=1\)；回代\(P\)：\(12/9-(16/3)/16=4/3-1/3=1\)成立（焦点在\(y\)轴成立）．(3)\(a=b\)为等轴双曲线，两轴取向皆可能，设\(x^{2}-y^{2}=t\)（\(t\neq 0\)），代\((3,-1)\)：\(t=9-1=8>0\)，得\(x^{2}-y^{2}=8\)，即\(x^{2}/8-y^{2}/8=1\)．(4)焦点位置未定，设\(mx^{2}-ny^{2}=1\)（\(mn>0\)），代入两点：\(9m-32n=1\)、\((81/16)m-25n=1\)，由第一式\(m=(1+32n)/9\)，代第二式\((9/16)(1+32n)-25n=1\)，即\(9/16+18n-25n=1\)，\(-7n=7/16\)，\(n=-1/16\)、\(m=-1/9\)，方程为\(-x^{2}/9+y^{2}/16=1\)，即\(y^{2}/16-x^{2}/9=1\)（回代\((9/4,5)\)：\(25/16-(81/16)/9=25/16-9/16=1\)成立）．}
\fi
\end{ansblock}

\begin{ansblock}[2章-练-课时13-简10]
% ans:2章-练-课时13-简10
\ansitem{5}{x²/9−y²/16＝1；焦点 (±5,0)；e=5/3}
\ifansbookdetail
\ansnote{详解}{\(2a=6\)即\(a=3\)；\(2b=8\)即\(b=4\)；\(c^{2}=a^{2}+b^{2}=9+16=25\)，\(c=5\)．故标准方程\(x^{2}/9-y^{2}/16=1\)，焦点\((\pm 5,0)\)，\(e=c/a=5/3\)．}
\fi
\end{ansblock}

\begin{ansblock}[2章-练-课时13-简2]
% ans:2章-练-课时13-简2
\ansitem{6}{C}
\ifansbookdetail
\ansnote{详解}{\(a^{2}=25\)、\(b^{2}=23\)，即\(a=5\)、\(2a=10\)，\(c^{2}=48\)、\(c=4\sqrt{3}\)．由定义\(\big||PF_2|-|PF_1|\big|=2a=10\)，即\(\big||PF_2|-8\big|=10\)，得\(|PF_2|=18\)或\(|PF_2|=-2\)（距离非负，舍），故\(|PF_2|=18\)．（存在性核验：\(c-a=4\sqrt{3}-5\approx 1.93\)、\(c+a\approx 11.93\)；\(|PF_1|=8\in[c-a,c+a)\)只能对应\(P\)在左支，此时\(|PF_2|=|PF_1|+2a=18\geqslant c+a\)成立，与选项A、B中“2”所设的右支情形（须\(|PF_2|\geqslant c-a\)且\(|PF_1|\geqslant c+a\)）不合，故不取双解．）故选：C．}
\fi
\end{ansblock}

\begin{ansblock}[2章-练-课时13-简5]
% ans:2章-练-课时13-简5
\ansitem{7}{C}
\ifansbookdetail
\ansnote{详解}{\(A\)、\(B\)在同一支上，由定义\(|AF_2|-|AF_1|=2a\)、\(|BF_2|-|BF_1|=2a\)（该支上的点到另一焦点较远，故取此号），两式相加：\((|AF_2|+|BF_2|)-(|AF_1|+|BF_1|)=4a\)．又弦\(AB\)过焦点\(F_1\)，\(F_1\)在\(A\)、\(B\)之间（焦点在该支围成的开口内），故\(|AF_1|+|BF_1|=|AB|=m\)，于是\(|AF_2|+|BF_2|=4a+m\)，\(\triangle ABF_2\)的周长\(=|AF_2|+|BF_2|+|AB|=4a+m+m=4a+2m\)．故选：C．}
\fi
\end{ansblock}

\begin{ansblock}[2章-练-课时13-简6]
% ans:2章-练-课时13-简6
\ansitem{8}{C}
\ifansbookdetail
\ansnote{详解}{\(a^{2}=9\)、\(b^{2}=16\)，即\(a=3\)、\(b=4\)，\(c^{2}=a^{2}+b^{2}=25\)、\(c=5\)，故左焦点\(F(-5,0)\)，虚轴长\(2b=8\)，\(|PQ|=2\times 8=16\)．点\(A(5,0)\)正是右焦点\(F'\)，且\(A\)在线段\(PQ\)上，故\(|PF|-|PA|=|PF|-|PF'|=2a=6\)（\(P\)在右支），同理\(|QF|-|QA|=6\)，两式相加：\(|PF|+|QF|=12+(|PA|+|QA|)=12+|PQ|=12+16=28\)，所以\(\triangle PFQ\)的周长\(=|PF|+|QF|+|PQ|=28+16=44\)．故选：C．}
\fi
\end{ansblock}

\begin{ansblock}[2章-练-课时13-简7]
% ans:2章-练-课时13-简7
\ansitem{9}{B}
\ifansbookdetail
\ansnote{详解}{\(a^{2}=3-m\)、\(b^{2}=m\)（\(0<m<3\)保证两项均正、焦点在\(x\)轴），\(c^{2}=a^{2}+b^{2}=3\)，即\(c=\sqrt{3}\)、\(|F_1F_2|=2\sqrt{3}\)．\(|OP|=\frac{1}{2}|F_1F_2|=\sqrt{3}=|OF_1|=|OF_2|\)，即\(P\)到线段\(F_1F_2\)中点\(O\)的距离等于该线段的一半，故\(P\)在以\(F_1F_2\)为直径的圆上，\(\angle F_1PF_2=90^{\circ}\)（直径所对圆周角为直角）．在直角\(\triangle F_1PF_2\)中，\(\angle PF_1F_2=30^{\circ}\)，故\(|PF_2|=|F_1F_2|\sin 30^{\circ}=2\sqrt{3}\times\frac{1}{2}=\sqrt{3}\)，\(|PF_1|=|F_1F_2|\cos 30^{\circ}=2\sqrt{3}\times\frac{\sqrt{3}}{2}=3\)．\(P\)在右支上，由定义\(|PF_1|-|PF_2|=2a\)：\(3-\sqrt{3}=2\sqrt{3-m}\)，即\(\sqrt{3-m}=(3-\sqrt{3})/2\)，\(3-m=(9-6\sqrt{3}+3)/4=3-\frac{3\sqrt{3}}{2}\)，得\(m=\frac{3\sqrt{3}}{2}\)．（核验\(0<m<3\)：\(3\sqrt{3}/2\approx 2.598<3\)成立；此时\(a^{2}=3-m\approx 0.402>0\)成立，\(2a=3-\sqrt{3}=|PF_1|-|PF_2|\)成立．）故选：B．}
\fi
\end{ansblock}

\begin{ansblock}[2章-练-课时13-简8]
% ans:2章-练-课时13-简8
\ansitem{10}{A}
\ifansbookdetail
\ansnote{详解}{\(a^{2}=4\)、\(b^{2}=12\)，即\(a=2\)、\(c^{2}=16\)、\(c=4\)，故左焦点\(F(-4,0)\)、右焦点\(F'(4,0)\)．\(P\)在右支上，由定义\(|PF|-|PF'|=2a=4\)，即\(|PF|=|PF'|+4\)，于是\(|PF|+|PA|=|PF'|+|PA|+4\geqslant|AF'|+4\)（三角不等式），\(|AF'|=\sqrt{(1-4)^{2}+4^{2}}=\sqrt{25}=5\)，故\(|PF|+|PA|\geqslant 5+4=9\)．取等核验：须\(P\)在线段\(AF'\)上，联立直线\(AF'\)（过\(A(1,4)\)、\(F'(4,0)\)，方程\(y=(16-4x)/3\)）与右支\(x^{2}/4-y^{2}/12=1\)（\(x\geqslant 2\)）得\(11x^{2}+128x-364=0\)，解得\(x=26/11\)（另一根\(-14\)舍），\(26/11\approx 2.36\in[2,4]\)，对应\(y=24/11\)，交点\((26/11,24/11)\)确实落在线段\(AF'\)上且在右支上，故等号可取，最小值为9．故选：A．}
\fi
\end{ansblock}

\begin{ansblock}[2章-练-课时13-中1]
% ans:2章-练-课时13-中1
\ansitem{11}{4√3}
\ifansbookdetail
\ansnote{详解}{焦点\((\pm 4,0)\)即\(c=4\)、\(c^{2}=16\)；又\(c^{2}=m+4\)，即\(m=12\)，故\(a^{2}=12\)、\(a=2\sqrt{3}\)、\(2a=4\sqrt{3}\)．设\(|MF_1|=t_1\)、\(|MF_2|=t_2\)，由定义\(|t_1-t_2|=2a=4\sqrt{3}\)，平方得\(t_1^{2}+t_2^{2}-2t_1t_2=48\)……①在\(\triangle F_1MF_2\)中用余弦定理：\(t_1^{2}+t_2^{2}-2t_1t_2\cos 60^{\circ}=|F_1F_2|^{2}=64\)，即\(t_1^{2}+t_2^{2}-t_1t_2=64\)……②②\(-\)①得\(t_1t_2=64-48=16\)，所以\(S=\frac{1}{2}t_1t_2\sin 60^{\circ}=\frac{1}{2}\times 16\times\frac{\sqrt{3}}{2}=4\sqrt{3}\)．}
\fi
\end{ansblock}

\begin{ansblock}[2章-练-课时13-中2]
% ans:2章-练-课时13-中2
\ansitem{12}{16/5}
\ifansbookdetail
\ansnote{详解}{\(a^{2}=9\)、\(b^{2}=16\)，即\(c^{2}=25\)、\(c=5\)、\(2a=6\)、\(|F_1F_2|=10\)．设\(|PF_1|=m\)、\(|PF_2|=n\)，由定义\(\big||PF_1|-|PF_2|\big|=2a=6\)，勾股（\(PF_1\perp PF_2\)）\(m^{2}+n^{2}=|F_1F_2|^{2}=100\)，两式并立：\((m-n)^{2}=m^{2}+n^{2}-2mn\)，即\(36=100-2mn\)，得\(mn=32\)．\(\triangle PF_1F_2\)的面积既等于\(\frac{1}{2}mn\)（直角边互乘），又等于\(\frac{1}{2}|F_1F_2|\cdot h\)（以焦轴为底、\(h\)为\(P\)到\(x\)轴的距离）：\(\frac{1}{2}\times 32=\frac{1}{2}\times 10\times h\)，即\(16=5h\)，得\(h=16/5\)．（存在性核验：由\(m-n=\pm 6\)、\(mn=32\)得\((m+n)^{2}=100+64=164\)，\(m+n=2\sqrt{41}\)，故\(\{m,n\}=\{\sqrt{41}+3,\sqrt{41}-3\}\approx\{9.4,3.4\}\)，两值均正且平方和为\(2(41+9)=100\)成立；与双曲线上焦半径的取值范围（近焦点侧\(\geqslant c-a=2\)、远焦点侧\(\geqslant c+a=8\)）比对：\(9.4\geqslant 8\)配远侧、\(3.4\geqslant 2\)配近侧成立，故满足垂直条件的点\(P\)确实存在．）所以点\(P\)到\(x\)轴的距离为\(16/5\)．}
\fi
\end{ansblock}

\begin{ansblock}[2章-练-课时13-中3]
% ans:2章-练-课时13-中3
\ansitem{13}{D}
\ifansbookdetail
\ansnote{详解}{\(a^{2}=16\)、\(b^{2}=9\)，即\(c^{2}=25\)、\(c=5\)，故\(A(-5,0)\)正是左焦点，记右焦点\(F(5,0)\)．\(P\)在右支上，由定义\(|PA|-|PF|=2a=8\)，即\(|PA|=|PF|+8\)，于是\(|PA|-|PB|=8+(|PF|-|PB|)\leqslant 8+|BF|\)（三角形两边之差小于第三边），\(|BF|=\sqrt{(8-5)^{2}+(4-0)^{2}}=\sqrt{25}=5\)，故\(|PA|-|PB|\leqslant 13\)．取等核验：须\(B\)在线段\(PF\)上，即\(P\)在射线\(F\to B\)的延长方向上；该射线参数式\((5+3s,4s)\)（\(s>0\)时经\(B\)，对应\(s=1\)），代入\(9x^{2}-16y^{2}=144\)得\(175s^{2}-270s-81=0\)，正根\(s=9/5>1\)，对应\(P(52/5,36/5)\)（核验\(9\times(52/5)^{2}-16\times(36/5)^{2}=(24336-20736)/25=144\)成立，且\(x=52/5\geqslant 4\)在右支上成立），故等号可取，最大值为13．故选：D．}
\fi
\end{ansblock}

\begin{ansblock}[2章-练-课时13-中4]
% ans:2章-练-课时13-中4
\ansitem{14}{x²/4 − y²/5 ＝1}
\ifansbookdetail
\ansnote{详解}{记\(F_1(-3,0)\)、\(F_2(3,0)\)，两定圆半径均为2，\(|F_1F_2|=6\)．设动圆\(M\)半径为\(r\)．情形一（与圆\(F_1\)内切、与圆\(F_2\)外切）：\(|MF_1|=r-2\)或\(2-r\)；\(|MF_2|=r+2\)．若\(|MF_1|=2-r\)，则\(|MF_1|+|MF_2|=4<|F_1F_2|=6\)，不可能（三角不等式），故\(|MF_1|=r-2\)（须\(r\geqslant 2\)），此时\(|MF_2|-|MF_1|=4\)；情形二（与圆\(F_1\)外切、与圆\(F_2\)内切）：同理\(|MF_1|=r+2\)、\(|MF_2|=r-2\)，\(|MF_1|-|MF_2|=4\)；（动圆含入定圆内的情形如上所验均导致\(|MF_1|+|MF_2|=4<6\)，故皆不出现）合写为\(\big||MF_1|-|MF_2|\big|=4<|F_1F_2|=6\)，由双曲线定义，\(M\)的轨迹是以\(F_1\)、\(F_2\)为焦点的双曲线（情形一给左支、情形二给右支，两支并存），\(2a=4\)、\(2c=6\)，即\(a=2\)、\(c=3\)、\(b^{2}=c^{2}-a^{2}=9-4=5\)，轨迹方程为\(x^{2}/4-y^{2}/5=1\)．}
\fi
\end{ansblock}

\begin{ansblock}[2章-练-课时13-难1]
% ans:2章-练-课时13-难1
\ansitem{15}{9/2}
\ifansbookdetail
\ansnote{详解}{设公共半焦距为\(c\)，椭圆长半轴\(a_1\)、双曲线实半轴\(a_2\)，记\(m=|PF_1|\)、\(n=|PF_2|\)．由椭圆定义\(m+n=2a_1\)、双曲线定义\(|m-n|=2a_2\)，又\(\angle F_1PF_2=90^{\circ}\)得\(m^{2}+n^{2}=4c^{2}\)；将前两式分别平方：\((m+n)^{2}+(m-n)^{2}=2(m^{2}+n^{2})\)，即\(4a_1^{2}+4a_2^{2}=2\times 4c^{2}\)，故\(a_1^{2}+a_2^{2}=2c^{2}\)，两边除\(c^{2}\)：\(1/e_1^{2}+1/e_2^{2}=2\)……①于是\(4e_1^{2}+e_2^{2}=\frac{1}{2}(1/e_1^{2}+1/e_2^{2})(4e_1^{2}+e_2^{2})=\frac{1}{2}\bigl[4+e_2^{2}/e_1^{2}+4e_1^{2}/e_2^{2}+1\bigr]=\frac{1}{2}\bigl[5+e_2^{2}/e_1^{2}+4e_1^{2}/e_2^{2}\bigr]\geqslant\frac{1}{2}\bigl[5+2\sqrt{(e_2^{2}/e_1^{2})\cdot(4e_1^{2}/e_2^{2})}\bigr]=\frac{1}{2}(5+4)=9/2\)（基本不等式）．取等条件：\(e_2^{2}/e_1^{2}=4e_1^{2}/e_2^{2}\)，即\(e_2^{2}=2e_1^{2}\)，代①：\(1/e_1^{2}+1/(2e_1^{2})=2\)，即\(3/(2e_1^{2})=2\)，得\(e_1^{2}=3/4\)、\(e_2^{2}=3/2\)，即\(e_1=\sqrt{3}/2\in(0,1)\)成立（椭圆）、\(e_2=\sqrt{6}/2\in(1,\sqrt{2})\)成立（双曲线），且\(a_1=c/e_1=2c/\sqrt{3}>c\)、\(a_2=c/e_2=\sqrt{2/3}\,c<c\)均合要求；故\(4e_1^{2}+e_2^{2}\)的最小值为\(9/2\)．}
\fi
\end{ansblock}

\begin{ansblock}[2章-练-课时13-难2]
% ans:2章-练-课时13-难2
\ansitem{16}{(1)(4,0)；(2)没有"被抓"的风险}
\ifansbookdetail
\ansnote{详解}{(1)设机器鼠位置为\(P\)，声音传播时间差为距离差除以\(v_0\)：\(|PA|/v_0-|PB|/v_0=8/v_0\)，即\(|PA|-|PB|=8\)（米），而\(|AB|=10\)，\(8<10\)，由双曲线定义\(P\)的轨迹是以\(A\)、\(B\)为焦点的双曲线中靠\(B\)的一支（右支）：\(2a=8\)、\(2c=10\)，即\(a=4\)、\(c=5\)、\(b^{2}=c^{2}-a^{2}=9\)，轨迹方程\(x^{2}/16-y^{2}/9=1\)（\(x\geqslant 4\)）．时刻\(t_0\)有\(|OP|=4\)：对右支上的点\(x\geqslant 4\)，故\(|OP|=\sqrt{x^{2}+y^{2}}\geqslant x\geqslant 4\)，等号仅当\(x=4\)且\(y=0\)，即顶点\((4,0)\)，所以\(t_0\)时机器鼠所在位置坐标为\((4,0)\)．(2)直线\(l\)：\(y=x\)，设与其平行的直线\(l_1\)：\(y=x+m\)与右支相切．代入\(9x^{2}-16y^{2}=144\)：\(9x^{2}-16(x+m)^{2}=144\)，即\(7x^{2}+32mx+16m^{2}+144=0\)，\(\Delta=(32m)^{2}-4\times 7\times(16m^{2}+144)=1024m^{2}-448m^{2}-4032=576m^{2}-4032=0\)，得\(m^{2}=7\)，切点横坐标\(x=-32m/(2\times 7)=-16m/7\)，要在右支须\(x\geqslant 4\)，故取\(m=-\sqrt{7}\)（此时\(x=16\sqrt{7}/7\approx 6.05\geqslant 4\)成立；\(m=+\sqrt{7}\)的切点在左支，舍），即与右支相切且最“靠向”\(l\)的平行线为\(l_1\)：\(y=x-\sqrt{7}\)，切点\(T(16\sqrt{7}/7,9\sqrt{7}/7)\)即右支上到\(l\)最近的点．两平行线\(y=x\)与\(y=x-\sqrt{7}\)的距离\(d=|\sqrt{7}|/\sqrt{1+1}=\sqrt{14}/2\approx 1.87\)（米）\(>1.5\)，（且当\(|m|<\sqrt{7}\)时\(\Delta<0\)，即形如\(y=x+m\)的直线与双曲线无公共点，右支整体落在\(x-y\geqslant\sqrt{7}\)的一侧，故此距离确为最小值而非极值假象）所以机器鼠到直线\(l\)的最小距离为\(\sqrt{14}/2\approx 1.87\)米，大于1.5米，没有“被抓”的风险．}
\fi
\end{ansblock}

\begin{ansblock}[2章-导-课时13-G1]
% ans:2章-导-课时13-G1
\ansitem{17}{C}
\ifansbookdetail
\ansnote{详解}{\(a^{2}=16\)、\(b^{2}=48\)，即\(a=4\)，\(c^{2}=a^{2}+b^{2}=64\)、\(c=8\)，故\(a+c=12\)、\(c-a=4\)．先定支：若\(P\)在右支上，则\(|PF_1|\geqslant a+c=12\)（右支上的点到左焦点的距离以右顶点处最小，为\(a+c\)），与\(|PF_1|=10\)矛盾；所以\(P\)在左支上，左支上\(|PF_2|>|PF_1|\)，由定义\(|PF_2|-|PF_1|=2a=8\)，得\(|PF_2|=10+8=18\)．（另一解2的来源是误按“\(|10-|PF_2||=8\)”取\(|PF_2|=2\)，但此时\(P\)须在右支且\(|PF_2|\geqslant c-a=4\)，不合，舍．）故选：C．}
\fi
\end{ansblock}

\begin{ansblock}[2章-导-课时13-G2]
% ans:2章-导-课时13-G2
\ansitem{18}{D}
\ifansbookdetail
\ansnote{详解}{设动圆\(P\)的半径为\(r\)，则\(|PM|=r\)；圆\(N\)的圆心\(N(4,0)\)、半径4，且\(|MN|=8\)．两圆相切分三种位置：①外切\(|PN|=r+4\)；②圆\(N\)在动圆内（内切）\(|PN|=r-4\)（须\(r\geqslant 4\)）；③动圆在圆\(N\)内（内切）\(|PN|=4-r\)．情形③不可能：此时\(|PM|+|PN|=r+4-r=4<|MN|=8\)，与三角形两边之和不小于第三边矛盾（也即\(M\)在圆\(N\)外，动圆既然过\(M\)就不能整个含在圆\(N\)内）．情形①：\(|PN|-|PM|=4\)；情形②：\(|PM|-|PN|=4\)．合写即\(\big||PN|-|PM|\big|=4<|MN|=8\)，由双曲线定义，\(P\)的轨迹是以\(M\)、\(N\)为焦点的双曲线，\(2a=4\)、\(2c=8\)，即\(a=2\)、\(c=4\)、\(b^{2}=c^{2}-a^{2}=12\)，且两种切法都能实现（在双曲线上任取一点\(P\)，取\(r=|PM|\)即可作出满足条件的动圆），故两支并存、不加\(x\)的限制，轨迹方程为\(x^{2}/4-y^{2}/12=1\)．故选：D．}
\fi
\end{ansblock}

\begin{ansblock}[2章-导-课时13-G3]
% ans:2章-导-课时13-G3
\ansitem{19}{B}
\ifansbookdetail
\ansnote{详解}{\(N\)是\(F_1M\)的中点（对称关系）、\(O\)是\(F_1F_2\)的中点，故\(ON\)是\(\triangle MF_1F_2\)的中位线，即\(|MF_2|=2|ON|=2\)（\(N\)在单位圆上）．\(P\)在\(F_1M\)的垂直平分线上，即\(|PF_1|=|PM|\)；又\(P\)在直线\(F_2M\)上，故\(P\)、\(M\)、\(F_2\)三点共线，于是\(\big||PF_2|-|PF_1|\big|=\big||PF_2|-|PM|\big|=|MF_2|=2<|F_1F_2|=4\)，由双曲线定义，\(P\)的轨迹是以\(F_1\)、\(F_2\)为焦点的双曲线（\(2a=2\)、\(2c=4\)，即\(x^{2}-y^{2}/3=1\)）．故选：B．}
\fi
\end{ansblock}

\begin{ansblock}[2章-导-课时13-G4]
% ans:2章-导-课时13-G4
\ansitem{20}{√2}
\ifansbookdetail
\ansnote{详解}{把三点代入\(x^{2}/a^{2}-y^{2}\)的值：\((-2,1)\)与\((2,-1)\)都得\(4/a^{2}-1\)，\((-2,3)\)得\(4/a^{2}-9\)．先试\((-2,3)\)在\(C\)上：\(4/a^{2}-9=1\)，即\(a^{2}=2/5\)，此时另两点处的值\(4/a^{2}-1=10-1=9\neq 1\)，即只有1个点在\(C\)上，与“恰有两个”矛盾，故\((-2,3)\)不在\(C\)上．于是恰有的两个点只能是\((-2,1)\)与\((2,-1)\)（两点关于原点对称，双曲线关于原点中心对称，故同属\(C\)成立），由\(4/a^{2}-1=1\)，即\(a^{2}=2\)，得\(a=\sqrt{2}\)（\(a>0\)）．（回代核验：\(a^{2}=2\)时\(4/2-1=1\)成立，两点在\(C\)上；\(4/2-9=-7\neq 1\)成立，第三点不在，恰两个成立．）}
\fi
\end{ansblock}

\begin{ansblock}[2章-导-课时13-G5]
% ans:2章-导-课时13-G5
\ansitem{21}{x²/9 − y²/16 ＝1 (x≠±3)}
\ifansbookdetail
\ansnote{详解}{设\(A(x,y)\)，则两斜率存在须\(x\neq\pm 3\)（\(x=\pm 3\)时\(A\)与\(B\)或\(C\)的横坐标相同，直线\(AB\)或\(AC\)垂直于\(x\)轴，斜率不存在，且\(A\)与\(B\)或\(C\)重合不成三角形），\(k_{AB}=y/(x+3)\)、\(k_{AC}=y/(x-3)\)，由题意\(y^{2}/[(x+3)(x-3)]=16/9\)，即\(y^{2}/(x^{2}-9)=16/9\)，去分母：\(9y^{2}=16(x^{2}-9)\)，即\(16x^{2}-9y^{2}=144\)，得\(x^{2}/9-y^{2}/16=1\)．该双曲线上\(x^{2}\geqslant 9\)，等号只在顶点\((\pm 3,0)\)，而这两点使\(A\)与\(B\)、\(C\)重合（三角形退化），故除去\(x=\pm 3\)两点，即\(x\neq\pm 3\)（此时自动满足\(|x|>3\)，两斜率均存在成立）．所以顶点\(A\)的轨迹方程为\(x^{2}/9-y^{2}/16=1\)（\(x\neq\pm 3\)）．}
\fi
\end{ansblock}

\jietitle{2.6.2 双曲线的几何性质}

\begin{ansblock}[2章-练-课时14-01]
% ans:2章-练-课时14-01
\ansitem{1}{(1) 实轴长8√2，虚轴长4，顶点(−4√2,0)和(4√2,0)；(2) 实轴长6，虚轴长18，顶点(−3,0)和(3,0)；(3) 实轴长4，虚轴长4，顶点(0,−2)和(0,2)；(4) 实轴长10，虚轴长14，顶点(0,−5)和(0,5)。}
\ifansbookdetail
\ansnote{详解}{(1)化\(x^{2}/32-y^{2}/4=1\)，焦点在\(x\)轴，\(a^{2}=32\)、\(b^{2}=4\)，\(a=4\sqrt{2}\)、\(b=2\)：实轴长\(2a=8\sqrt{2}\)，虚轴长\(2b=4\)，顶点\((-4\sqrt{2},0)\)和\((4\sqrt{2},0)\)．\par\noindent (2)化\(x^{2}/9-y^{2}/81=1\)，焦点在\(x\)轴，\(a=3\)、\(b=9\)：实轴长6，虚轴长18，顶点\((-3,0)\)和\((3,0)\)．\par\noindent (3)化\(y^{2}/4-x^{2}/4=1\)，焦点在\(y\)轴，\(a=2\)、\(b=2\)：实轴长4，虚轴长4，顶点\((0,-2)\)和\((0,2)\)．\par\noindent (4)化\(y^{2}/25-x^{2}/49=1\)，焦点在\(y\)轴，\(a=5\)、\(b=7\)：实轴长10，虚轴长14，顶点\((0,-5)\)和\((0,5)\)．}
\fi
\end{ansblock}

\begin{ansblock}[2章-练-课时14-02]
% ans:2章-练-课时14-02
\ansitem{2}{实轴长18；虚轴长6；焦点(0,±3√10)；离心率 e=√10/3；渐近线方程 y=±3x。}
\ifansbookdetail
\ansnote{详解}{化\(-9x^{2}+y^{2}=81\)为\(y^{2}/81-x^{2}/9=1\)，焦点在\(y\)轴，\(a^{2}=81\)、\(b^{2}=9\)，\(a=9\)、\(b=3\)，\(c^{2}=a^{2}+b^{2}=90\)，\(c=\sqrt{90}=3\sqrt{10}\)．\par\noindent 于是：实轴长\(2a=18\)；虚轴长\(2b=6\)；焦点\((0,\pm 3\sqrt{10})\)；离心率\(e=c/a=3\sqrt{10}/9=\sqrt{10}/3\)；渐近线\(y=\pm(a/b)x=\pm 3x\)．\par\noindent 验算：\(e=\sqrt{1+(a/b)^{2}}=\sqrt{10}/3\)合；\(2c=6\sqrt{10}\approx 18.97>2a=18\)合；点\((0,9)\)代入方程\(81/81=1\)合．}
\fi
\end{ansblock}

\begin{ansblock}[2章-练-课时14-07]
% ans:2章-练-课时14-07
\ansitem{3}{D}
\ifansbookdetail
\ansnote{详解}{实轴长\(8\)得\(a=4\)；焦点在\(y\)轴，渐近线\(y=\pm(a/b)x=\pm(4/b)x=(4/3)x\)，得\(b=3\)．故\(C\)：\(y^{2}/16-x^{2}/9=1\)，选D．}
\fi
\end{ansblock}

\begin{ansblock}[2章-练-课时14-08]
% ans:2章-练-课时14-08
\ansitem{4}{(1) 实轴长8/3，虚轴长4，焦点(±2√13/3,0)，顶点(±4/3,0)，渐近线 y=±(3/2)x；(2)（m>0）实轴长4，虚轴长2√m，焦点(±√(m+4),0)，顶点(±2,0)，渐近线 y=±(√m/2)x。}
\ifansbookdetail
\ansnote{详解}{(1)\(a^{2}=16/9\)、\(b^{2}=4\)，\(a=4/3\)、\(b=2\)；\(c^{2}=16/9+4=52/9\)，\(c=\sqrt{52/9}=2\sqrt{13}/3\)：实轴长\(8/3\)，虚轴长4，焦点\((\pm 2\sqrt{13}/3,0)\)，顶点\((\pm 4/3,0)\)，渐近线\(y=\pm(b/a)x=\pm(3/2)x\)．\par\noindent (2)双曲线须\(m>0\)：\(a=2\)、\(b=\sqrt{m}\)、\(c^{2}=4+m\)，\(c=\sqrt{m+4}\)：实轴长4，虚轴长\(2\sqrt{m}\)，焦点\((\pm\sqrt{m+4},0)\)，顶点\((\pm 2,0)\)，渐近线\(y=\pm(\sqrt{m}/2)x\)．（\(m<0\)时方程表椭圆，与题设矛盾，仅\(m>0\)一解．）}
\fi
\end{ansblock}

\begin{ansblock}[2章-拓-课时14-14]
% ans:2章-拓-课时14-14
\ansitem{5}{实轴长2；虚轴长4√6；焦点(±5,0)；渐近线方程 y=±2√6x。}
\ifansbookdetail
\ansnote{详解}{\(a^{2}=1\)、\(b^{2}=24\)，\(a=1\)、\(b=2\sqrt{6}\)，\(c^{2}=1+24=25\)，\(c=5\)：实轴长2，虚轴长\(4\sqrt{6}\)，焦点\((\pm 5,0)\)，渐近线\(y=\pm(b/a)x=\pm 2\sqrt{6}x\)．}
\fi
\end{ansblock}

\begin{ansblock}[2章-拓-课时14-15]
% ans:2章-拓-课时14-15
\ansitem{6}{x²/16−y²/9=1；e=5/4。}
\ifansbookdetail
\ansnote{详解}{焦点\((5,0)\)在\(x\)轴：\(b/a=3/4\)，\(c=5\)，\(a^{2}+b^{2}=25\)；\(b^{2}=(9/16)a^{2}\)代入得\(a^{2}(1+9/16)=25\)，\(a^{2}=16\)，\(b^{2}=9\)．标准方程\(x^{2}/16-y^{2}/9=1\)，\(e=c/a=5/4\)．}
\fi
\end{ansblock}

\begin{ansblock}[2章-拓-课时14-16]
% ans:2章-拓-课时14-16
\ansitem{7}{它们的渐近线相同，均为 y=±(b/a)x。}
\ifansbookdetail
\ansnote{详解}{\(\lambda>0\)时方程即\(x^{2}/(a^{2}\lambda)-y^{2}/(b^{2}\lambda)=1\)，焦点在\(x\)轴，渐近线由\(x^{2}/(a^{2}\lambda)-y^{2}/(b^{2}\lambda)=0\)得\(y=\pm(b/a)x\)；\(\lambda<0\)时方程即\(y^{2}/(-b^{2}\lambda)-x^{2}/(-a^{2}\lambda)=1\)，焦点在\(y\)轴，渐近线由同式得\(y=\pm(b/a)x\)．故所有曲线的渐近线相同，均为\(y=\pm(b/a)x\)（渐近线只取决于\(a:b\)，与\(\lambda\)无关）．}
\fi
\end{ansblock}

\begin{ansblock}[2章-拓-课时14-17]
% ans:2章-拓-课时14-17
\ansitem{8}{x²/9−y²/27=1（焦点在 x 轴；焦点在 y 轴时为 y²/9−x²/27=1）。}
\ifansbookdetail
\ansnote{详解}{顶点距\(6\)，即\(2a=6\)，\(a=3\)；焦点段\([-c,c]\)上的分点\(-a\)，\(0\)，\(a\)将其四等分，等分距为\(a\)，故\(a=c/2\)，\(c=6\)，\(b^{2}=c^{2}-a^{2}=36-9=27\)．焦点在\(x\)轴：\(x^{2}/9-y^{2}/27=1\)；题面未指明焦点位置，焦点在\(y\)轴时\(y^{2}/9-x^{2}/27=1\)同法可得．}
\fi
\end{ansblock}

\begin{ansblock}[2章-练-课时14-03]
% ans:2章-练-课时14-03
\ansitem{9}{B}
\ifansbookdetail
\ansnote{详解}{由\(x^{2}-y^{2}/4=1\)知焦点在\(x\)轴，\(a=1\)、\(b=2\)，渐近线\(y=\pm(b/a)x=\pm 2x\)，故选B．}
\fi
\end{ansblock}

\begin{ansblock}[2章-练-课时14-04]
% ans:2章-练-课时14-04
\ansitem{10}{3}
\ifansbookdetail
\ansnote{详解}{\(c^{2}=16+9=25\)，\(c=5\)，焦点\((\pm 5,0)\)；渐近线方程\(3x\pm 4y=0\)．取\(F(5,0)\)与直线\(3x-4y=0\)：\(d=|3\times 5-4\times 0|/\sqrt{3^{2}+4^{2}}=15/5=3\)．由对称性，任一焦点到任一条渐近线的距离均为3（恰为半虚轴长\(b=3\)）．}
\fi
\end{ansblock}

\begin{ansblock}[2章-练-课时14-11]
% ans:2章-练-课时14-11
\ansitem{11}{B}
\ifansbookdetail
\ansnote{详解}{\(F(-c,0)\)，\(A(a,0)\)，取\(B(0,b)\)，直线\(AB\)：\(bx+ay-ab=0\)．\(F\)到\(AB\)的距离\(d=|-bc-ab|/\sqrt{a^{2}+b^{2}}=b(c+a)/c\)（分子提出\(b\)，分母用\(c^{2}=a^{2}+b^{2}\)）．\par\noindent 由\(d=(5/4)b\)得\(c+a=(5/4)c\)，即\(4a=c\)；\(b^{2}=c^{2}-a^{2}=15a^{2}\)，\(b/a=\sqrt{15}\)，渐近线\(y=\pm\sqrt{15}x\)，故选B．验算：\(c=4a\)时\(d=b\cdot 5a/(4a)=(5/4)b\)合．}
\fi
\end{ansblock}

\begin{ansblock}[2章-拓-课时14-07]
% ans:2章-拓-课时14-07
\ansitem{12}{B}
\ifansbookdetail
\ansnote{详解}{共渐近线设\(x^{2}-2y^{2}=m\)（\(m\neq 0\)）；焦点\((0,-6)\)在\(y\)轴上，故\(m<0\)．标准形\(y^{2}/(-m/2)-x^{2}/(-m)=1\)：\(a^{2}=-m/2\)、\(b^{2}=-m\)，\(c^{2}=a^{2}+b^{2}=-3m/2=36\)，得\(m=-24\)，方程\(y^{2}/12-x^{2}/24=1\)，选B．验算：\(c^{2}=12+24=36\)合，焦点含\((0,-6)\)合．}
\fi
\end{ansblock}

\begin{ansblock}[2章-拓-课时14-19]
% ans:2章-拓-课时14-19
\ansitem{13}{x²/(32/5)−y²/(18/5)=1。}
\ifansbookdetail
\ansnote{详解}{椭圆\(a^{2}=10\)、\(b^{2}=5\)，\(c^{2}=5\)，长轴端点\((\pm\sqrt{10},0)\)，即双曲线焦点在\(x\)轴，\(c=\sqrt{10}\)；渐近线\(y=\pm(3/4)x\)，\(b/a=3/4\)．\(a^{2}+b^{2}=10\)，\(a^{2}(1+9/16)=10\)，得\(a^{2}=32/5\)，\(b^{2}=10-32/5=18/5\)，标准方程\(x^{2}/(32/5)-y^{2}/(18/5)=1\)．}
\fi
\end{ansblock}

\begin{ansblock}[2章-拓-课时14-20]
% ans:2章-拓-课时14-20
\ansitem{14}{x²/(64/13)−y²/(144/13)=1。}
\ifansbookdetail
\ansnote{详解}{焦点在\(x\)轴，\(c=4\)；渐近线\(y=(3/2)x\)，\(b/a=3/2\)，\(a^{2}+b^{2}=16\)；\(b^{2}=(9/4)a^{2}\)代入得\(a^{2}(13/4)=16\)，\(a^{2}=64/13\)，\(b^{2}=16-64/13=144/13\)，标准方程\(x^{2}/(64/13)-y^{2}/(144/13)=1\)．}
\fi
\end{ansblock}

\begin{ansblock}[2章-拓-课时14-21]
% ans:2章-拓-课时14-21
\ansitem{15}{x²/12−y²/8=1。}
\ifansbookdetail
\ansnote{详解}{原双曲线\(c^{2}=16+4=20\)，有公共焦点，故所求双曲线\(c^{2}=20\)，焦点在\(x\)轴；半实轴长\(a=2\sqrt{3}\)，\(a^{2}=12\)，\(b^{2}=c^{2}-a^{2}=20-12=8\)，标准方程\(x^{2}/12-y^{2}/8=1\)．}
\fi
\end{ansblock}

\begin{ansblock}[2章-练-课时14-05]
% ans:2章-练-课时14-05
\ansitem{16}{2}
\ifansbookdetail
\ansnote{详解}{焦点在\(x\)轴，渐近线\(y=\pm(b/a)x\)，故\(b/a=\sqrt{3}\)．\(e=c/a=\sqrt{1+(b/a)^{2}}=\sqrt{1+3}=2\)．}
\fi
\end{ansblock}

\begin{ansblock}[2章-练-课时14-06]
% ans:2章-练-课时14-06
\ansitem{17}{5/4（焦点在x轴）或 5/3（焦点在y轴）}
\ifansbookdetail
\ansnote{详解}{焦点在\(x\)轴时，渐近线\(y=\pm(b/a)x\)，\(b/a=3/4\)，\(e=\sqrt{1+(b/a)^{2}}=\sqrt{1+9/16}=\sqrt{25/16}=5/4\)；焦点在\(y\)轴时，渐近线\(y=\pm(a/b)x\)，\(a/b=3/4\)，\(b/a=4/3\)，\(e=\sqrt{1+16/9}=\sqrt{25/9}=5/3\)．题面未限焦点位置，两解并立．}
\fi
\end{ansblock}

\begin{ansblock}[2章-练-课时14-09]
% ans:2章-练-课时14-09
\ansitem{18}{x²/3−y²=1}
\ifansbookdetail
\ansnote{详解}{由渐近线得\(b/a=\sqrt{3}/3\)，即\(a=\sqrt{3}b\)．直线\(l\)过\(P(0,b)\)、\(Q(a,0)\)，方程\(bx+ay-ab=0\)，原点到\(l\)的距离\(ab/\sqrt{a^{2}+b^{2}}=\sqrt{3}/2\)，两边平方得\(4a^{2}b^{2}=3(a^{2}+b^{2})\)．\par\noindent 联立\(a=\sqrt{3}b\)：\(4\times 3b^{4}=3(3b^{2}+b^{2})\)，即\(12b^{4}=12b^{2}\)，\(b^{2}=1\)（\(b>0\)），\(a^{2}=3\)．故双曲线方程为\(x^{2}/3-y^{2}=1\)．验算：\(a=\sqrt{3}\)、\(b=1\)时\(ab/\sqrt{a^{2}+b^{2}}=\sqrt{3}/2\)合，渐近线\(y=\pm x/\sqrt{3}\)合．}
\fi
\end{ansblock}

\begin{ansblock}[2章-练-课时14-10]
% ans:2章-练-课时14-10
\ansitem{19}{(1) x²/9−y²/25=1 与 y²/25−x²/9=1（答案不唯一，凡渐近线为 5x±3y=0 的标准方程均可）；(2) y²/(56/9)−x²/(56/25)=1（即 9y²−25x²=56）。}
\ifansbookdetail
\ansnote{详解}{(1)两渐近线即\(y=\pm(5/3)x\)．焦点在\(x\)轴：\(b/a=5/3\)，取\(a=3\)、\(b=5\)得\(x^{2}/9-y^{2}/25=1\)；焦点在\(y\)轴：\(a/b=5/3\)，取\(a=5\)、\(b=3\)得\(y^{2}/25-x^{2}/9=1\)．\par\noindent (2)以\(l_1\)、\(l_2\)为渐近线的双曲线可设为\(x^{2}/9-y^{2}/25=\lambda\)（\(\lambda\neq 0\)）．代入\(M(1,3)\)：\(\lambda=1/9-9/25=(25-81)/225=-56/225<0\)，故焦点在\(y\)轴，方程即\(y^{2}/25-x^{2}/9=56/225\)，化标准形\(y^{2}/(56/9)-x^{2}/(56/25)=1\)．验算：\(9y^{2}-25x^{2}=56\)代入\(M(1,3)\)：\(81-25=56\)合．提醒：渐近线同为\(5x\pm 3y=0\)的双曲线焦点可在\(x\)轴或\(y\)轴，由所过点的坐标定号（本题\(1/9-9/25<0\)定\(y\)轴）．}
\fi
\end{ansblock}

\begin{ansblock}[2章-拓-课时14-18]
% ans:2章-拓-课时14-18
\ansitem{20}{曲线是双曲线（两支）：以 AB 所在直线为 x 轴、AB 中点为原点，其方程为 x²/9−y²/16=1。}
\ifansbookdetail
\ansnote{详解}{曲线上点\(P\)满足\(\bigl||PA|-|PB|\bigr|=6<|AB|=10\)，由定义轨迹为双曲线，\(2a=6\)、\(2c=10\)，\(a=3\)、\(c=5\)，\(b^{2}=25-9=16\)．以\(AB\)中点为原点、\(AB\)所在直线为\(x\)轴建系，方程为\(x^{2}/9-y^{2}/16=1\)．作图观察：交点连成两支光滑曲线，与双曲线吻合．}
\fi
\end{ansblock}

\begin{ansblock}[2章-拓-课时14-23]
% ans:2章-拓-课时14-23
\ansitem{21}{x²/4−y²/5=1。}
\ifansbookdetail
\ansnote{详解}{（直接法）设\(M(x,y)\)：\(\sqrt{(x-3)^{2}+y^{2}}=(3/2)\left|x-4/3\right|\)．两边平方：\(x^{2}-6x+9+y^{2}=(9/4)\left(x^{2}-(8/3)x+16/9\right)=(9/4)x^{2}-6x+4\)．\par\noindent \(-6x\)抵消：\(y^{2}=(9/4)x^{2}+4-x^{2}-9=(5/4)x^{2}-5\)，即\((5/4)x^{2}-y^{2}=5\)，标准形\(x^{2}/4-y^{2}/5=1\)．验算：\(a=2\)、\(c=3\)与准线位置\(x=a^{2}/c=4/3\)、比值\(c/a=3/2\)相容合．}
\fi
\end{ansblock}

\begin{ansblock}[2章-拓-课时14-24]
% ans:2章-拓-课时14-24
\ansitem{22}{比值等于 c/a（离心率 e）。}
\ifansbookdetail
\ansnote{详解}{（直接法）\(P(x,y)\)在双曲线上：\(y^{2}=b^{2}(x^{2}/a^{2}-1)\)．\(d_1=|PF|=\sqrt{(x+c)^{2}+y^{2}}\)，代入并用\(b^{2}=c^{2}-a^{2}\)化简：\(x^{2}+2cx+c^{2}+(c^{2}-a^{2})x^{2}/a^{2}-(c^{2}-a^{2})=(c^{2}/a^{2})x^{2}+2cx+a^{2}=\left[(c/a)x+a\right]^{2}\)．\par\noindent \(d_2=\left|x+a^{2}/c\right|\)（\(P\)到\(l\)的距离）．故\(d_1/d_2=(c/a)\left|x+a^{2}/c\right|/\left|x+a^{2}/c\right|=c/a\)（即离心率\(e\)）．}
\fi
\end{ansblock}

\begin{ansblock}[2章-拓-课时14-27]
% ans:2章-拓-课时14-27
\ansitem{23}{x²/a²−y²/b²=1，其中 b²=c²−a²。}
\ifansbookdetail
\ansnote{详解}{（直接法）设\(M(x,y)\)：\(\sqrt{(x+c)^{2}+y^{2}}=(c/a)\left|x+a^{2}/c\right|\)．平方：\(x^{2}+2cx+c^{2}+y^{2}=(c^{2}/a^{2})x^{2}+2cx+a^{2}\)，\(2cx\)抵消：\(y^{2}=(c^{2}/a^{2}-1)x^{2}-(c^{2}-a^{2})=(b^{2}/a^{2})x^{2}-b^{2}\)（\(b^{2}=c^{2}-a^{2}\)），即\(x^{2}/a^{2}-y^{2}/b^{2}=1\)．解释：2.6.2 例 3 所得轨迹方程正是\(x^{2}/a^{2}-y^{2}/b^{2}=1\)（\(b^{2}=c^{2}-a^{2}\)），本题表明该双曲线可由“到定点与到定直线距离之比为常数\(c/a\)”直接刻画．}
\fi
\end{ansblock}

\begin{ansblock}[2章-练-课时14-12]
% ans:2章-练-课时14-12
\ansitem{24}{A}
\ifansbookdetail
\ansnote{详解}{设\(M(x_0,y_0)\)，则\(N(-x_0,-y_0)\)，\(P(x,y)\)（\(x\neq\pm x_0\)）．\(k_1k_2=\dfrac{y-y_0}{x-x_0}\cdot\dfrac{y+y_0}{x+x_0}=\dfrac{y^{2}-y_0^{2}}{x^{2}-x_0^{2}}\)（坐标点差法，不引名目）．\(M\)、\(P\)均在双曲线上：\(y^{2}=b^{2}(x^{2}/a^{2}-1)\)，\(y_0^{2}=b^{2}(x_0^{2}/a^{2}-1)\)，相减得\(y^{2}-y_0^{2}=(b^{2}/a^{2})(x^{2}-x_0^{2})\)，故\(k_1k_2=b^{2}/a^{2}\)（定值）．\par\noindent 于是\(|k_1|+|k_2|\geqslant 2\sqrt{|k_1k_2|}=2b/a\)（当\(|k_1|=|k_2|=b/a\)时取等，\(P\)可取到），由最小值2得\(b/a=1\)，\(e=\sqrt{1+(b/a)^{2}}=\sqrt{2}\)，故选A．}
\fi
\end{ansblock}

\begin{ansblock}[2章-练-课时14-13]
% ans:2章-练-课时14-13
\ansitem{25}{1}
\ifansbookdetail
\ansnote{详解}{\(c^{2}=16+9=25\)，\(c=5\)，右焦点\(F(5,0)\)．设\(P(x,y)\)，\(|PF|^{2}=(x-5)^{2}+y^{2}\)．由\(y^{2}=9(x^{2}/16-1)\)（\(|x|\geqslant 4\)）：\(|PF|^{2}=x^{2}-10x+25+(9/16)x^{2}-9=(25/16)x^{2}-10x+16\)．对称轴\(x=10/(2\times 25/16)=16/5<4\)，故在\(|x|\geqslant 4\)上于\(x=4\)取最小：\(|PF|^{2}=(25/16)\times 16-40+16=1\)，\(|PF|_{\min}=1\)（即右顶点处取得，同支焦半径最小值\(c-a=5-4=1\)合）．}
\fi
\end{ansblock}

\begin{ansblock}[2章-练-课时14-14]
% ans:2章-练-课时14-14
\ansitem{26}{证明见解析。}
\ifansbookdetail
\ansnote{详解}{证明：不妨先设焦点在\(x\)轴：双曲线\(x^{2}/a^{2}-y^{2}/b^{2}=1\)（\(a>0\)，\(b>0\)），焦点\(F(c,0)\)，渐近线\(y=\pm(b/a)x\)，即\(bx\mp ay=0\)．焦点到渐近线的距离\(d=|bc-a\cdot 0|/\sqrt{a^{2}+b^{2}}=bc/c=b\)（半虚轴长）．\par\noindent 焦点在\(y\)轴时（\(y^{2}/a^{2}-x^{2}/b^{2}=1\)，焦点\(F(0,c)\)，渐近线\(y=\pm(a/b)x\)即\(ay\mp bx=0\)）：\(d=|a\cdot 0-bc|/\sqrt{a^{2}+b^{2}}=bc/c=b\)，等式仍成立．证毕．}
\fi
\end{ansblock}

\begin{ansblock}[2章-拓-课时14-22]
% ans:2章-拓-课时14-22
\ansitem{27}{2}
\ifansbookdetail
\ansnote{详解}{设\(P(x,y)\)（\(|x|\geqslant 1\)）：\(|PA|^{2}=(x-3)^{2}+y^{2}=(x-3)^{2}+3(x^{2}-1)=4x^{2}-6x+6\)．对称轴\(x=3/4<1\)，故\(x=1\)（右顶点）时取最小：\(4-6+6=4\)，\(|PA|_{\min}=2\)．}
\fi
\end{ansblock}

\begin{ansblock}[2章-拓-课时14-25]
% ans:2章-拓-课时14-25
\ansitem{28}{16}
\ifansbookdetail
\ansnote{详解}{\(a=3\)，\(b=4\)，\(c=5\)．设两焦半径\(r_1=|MF_1|\)，\(r_2=|MF_2|\)（\(r_1>r_2\)）．\(MF_1\perp MF_2\)：\(r_1^{2}+r_2^{2}=|F_1F_2|^{2}=100\)；定义：\(r_1-r_2=2a=6\)．\par\noindent \((r_1-r_2)^{2}=100-2r_1r_2=36\)，得\(r_1r_2=32\)．\(S=(1/2)r_1r_2\sin 90^{\circ}=16\)．}
\fi
\end{ansblock}

\begin{ansblock}[2章-拓-课时14-26]
% ans:2章-拓-课时14-26
\ansitem{29}{(1) ax²−y²=36a（y≠0）；(2) a>0 时为焦点在 x 轴的双曲线（除去顶点）；a<0 时为椭圆：其中 a=−1 为圆 x²+y²=36（y≠0），−1<a<0 为焦点在 x 轴的椭圆，a<−1 为焦点在 y 轴的椭圆（均除去长轴端点——方程中 y≠0）。}
\ifansbookdetail
\ansnote{详解}{(1)过\((6,0)\)斜率\(k_1\)、过\((-6,0)\)斜率\(k_2\)，\(k_1k_2=a\)．设\(M(x,y)\)：\(k_1=y/(x-6)\)，\(k_2=y/(x+6)\)（\(x\neq\pm 6\)，且\(M\)不与\((\pm 6,0)\)重合故\(y\neq 0\)）．\([y/(x-6)]\cdot[y/(x+6)]=a\)，得\(y^{2}=a(x^{2}-36)\)，即\(ax^{2}-y^{2}=36a\)（\(y\neq 0\)）．\par\noindent (2)\(a>0\)：\(x^{2}/36-y^{2}/(36a)=1\)，焦点在\(x\)轴的双曲线（除去顶点）；\(a<0\)：记\(a=-|a|\)，\(x^{2}/36+y^{2}/(36|a|)=1\)，椭圆；其中\(a=-1\)：\(x^{2}+y^{2}=36\)，圆；\(-1<a<0\)：\(36>36|a|\)，焦点在\(x\)轴的椭圆；\(a<-1\)：\(36<36|a|\)，焦点在\(y\)轴的椭圆（均除去\(x\)轴上的点，即\(y\neq 0\)）．}
\fi
\end{ansblock}

\begin{ansblock}[2章-拓-课时14-08]
% ans:2章-拓-课时14-08
\ansitem{30}{证明见解析。}
\ifansbookdetail
\ansnote{详解}{证明：设\(C\)：\(x^{2}/m^{2}-y^{2}/m^{2}=1\)（\(m>0\)），则\(F_1(-\sqrt{2}m,0)\)，\(F_2(\sqrt{2}m,0)\)．设\(P(x,y)\)，满足\(x^{2}-y^{2}=m^{2}\)．\(|PO|^{2}=x^{2}+y^{2}\)．\par\noindent \(|PF_1|\cdot|PF_2|=\sqrt{(x+\sqrt{2}m)^{2}+y^{2}}\cdot\sqrt{(x-\sqrt{2}m)^{2}+y^{2}}\)，根号内合并为\(\sqrt{(x^{2}+2m^{2}+y^{2})^{2}-8m^{2}x^{2}}\)．用\(y^{2}=x^{2}-m^{2}\)化简：\(x^{2}+2m^{2}+y^{2}=2x^{2}+m^{2}\)，\((2x^{2}+m^{2})^{2}-8m^{2}x^{2}=4x^{4}-4m^{2}x^{2}+m^{4}=(2x^{2}-m^{2})^{2}\)．\par\noindent 故\(|PF_1|\cdot|PF_2|=|2x^{2}-m^{2}|=2x^{2}-m^{2}\)（\(|x|\geqslant m\)时为正）\(=x^{2}+y^{2}=|PO|^{2}\)．证毕．验算：\(m=1\)，\(P(\sqrt{2},1)\)：\(|PO|^{2}=3\)，\(|PF_1|=3\)，\(|PF_2|=1\)，积为3合．}
\fi
\end{ansblock}

\begin{ansblock}[2章-拓-课时14-01]
% ans:2章-拓-课时14-01
\ansitem{31}{B}
\ifansbookdetail
\ansnote{详解}{\(PQ\)为通径，\(P(c,\pm b^{2}/a)\)，\(|PQ|=2b^{2}/a\)．\(\triangle PF_1Q\)中\(\angle PF_1Q=90^{\circ}\)，\(F_2\)为斜边\(PQ\)的中点，由直角三角形斜边中线定理\(|F_1F_2|=|PQ|/2\)，即\(2c=b^{2}/a\)，\(2ac=b^{2}=c^{2}-a^{2}\)．两边除以\(a^{2}\)：\(e^{2}-2e-1=0\)，\(e>1\)，\(e=1+\sqrt{2}\)，故选B．}
\fi
\end{ansblock}

\begin{ansblock}[2章-拓-课时14-02]
% ans:2章-拓-课时14-02
\ansitem{32}{2}
\ifansbookdetail
\ansnote{详解}{\(A\)在线段\(F_1B\)上且\(|F_1A|=|AB|\)，故\(A\)为\(F_1B\)的中点；\(O\)为\(F_1F_2\)的中点，故\(OA\)是\(\triangle F_1F_2B\)的中位线：\(OA\parallel F_2B\)，\(|F_2B|=2|OA|\)．\par\noindent 由\(\overrightarrow{F_1B}\cdot\overrightarrow{F_2B}=0\)得\(F_1B\perp F_2B\)，结合\(OA\parallel F_2B\)得\(OA\perp F_1B\)．于是\(\triangle F_1OB\)中，\(OA\)既是边\(F_1B\)上的中线又是高，故\(\triangle F_1OB\)等腰：\(|OB|=|OF_1|=c\)，且\(\angle AOF_1=\angle AOB\)．\par\noindent 又\(A\)、\(B\)分别在两条渐近线上，两渐近线关于\(x\)轴对称，故\(\angle BOF_2=\angle AOF_1\)；而\(\angle AOF_1+\angle AOB+\angle BOF_2=180^{\circ}\)，三者在上述等量关系下相等，各为\(60^{\circ}\)．故\(b/a=\tan 60^{\circ}=\sqrt{3}\)，\(e=\sqrt{1+(b/a)^{2}}=\sqrt{1+3}=2\)．}
\fi
\end{ansblock}

\begin{ansblock}[2章-拓-课时14-03]
% ans:2章-拓-课时14-03
\ansitem{33}{C}
\ifansbookdetail
\ansnote{详解}{六个顶点内接于以\(F_1F_2\)为直径的圆，由直径所对圆周角\(\angle F_1PF_2=90^{\circ}\)；正六边形各边对圆心角\(60^{\circ}\)，得\(\angle PF_2F_1=30^{\circ}\)．解三角形：\(|PF_2|=2c\cos 30^{\circ}=\sqrt{3}c\)，\(|PF_1|=2c\sin 30^{\circ}=c\)．\(P\)在双曲线上，\(|PF_2|-|PF_1|=2a\)，即\(\sqrt{3}c-c=2a\)，\(e=c/a=2/(\sqrt{3}-1)=\sqrt{3}+1\)，故选C．}
\fi
\end{ansblock}

\begin{ansblock}[2章-拓-课时14-04]
% ans:2章-拓-课时14-04
\ansitem{34}{D}
\ifansbookdetail
\ansnote{详解}{圆\(x^{2}+y^{2}=a^{2}+b^{2}\)的半径为\(c\)，其直径恰为\(F_1F_2\)，故\(P\)在圆上，\(\angle F_1PF_2=90^{\circ}\)（直径所对圆周角）．设\(|PF_1|=n\)，由定义\(|PF_2|=2a+n\)．Rt\(\triangle PF_1F_2\)中：\(n^{2}+(2a+n)^{2}=4c^{2}\)，展开得\(n^{2}=2b^{2}-2an\)．\par\noindent 由\(\overrightarrow{F_1P}=(1/2)\overrightarrow{F_2Q}\)（同向向量）得\(|F_2Q|=2n\)，且\(F_1P\parallel F_2Q\)，故\(\angle QF_2F_1=180^{\circ}-\angle PF_1F_2\)，\(\cos\angle QF_2F_1=-n/(2c)\)．由定义\(|F_1Q|=2a+2n\)，\(\triangle F_1F_2Q\)中余弦定理：\((2a+2n)^{2}=(2n)^{2}+(2c)^{2}+4n^{2}\)，化简得\(2an=b^{2}+n^{2}\)，即\(b^{2}=2an-n^{2}\)．\par\noindent 与\(n^{2}=2b^{2}-2an\)联立：\(3n^{2}=2an\)，\(a=(3/2)n\)；\(b^{2}=2n^{2}\)，\(c^{2}=(17/4)n^{2}\)，\(e=c/a=\sqrt{17}/3\)，故选D．}
\fi
\end{ansblock}

\begin{ansblock}[2章-拓-课时14-05]
% ans:2章-拓-课时14-05
\ansitem{35}{√3}
\ifansbookdetail
\ansnote{详解}{\(O(0,0)\)，\(A(-a,0)\)，\(B(a,0)\)，圆\(P\)的圆心为\(P(0,a)\)，半径\(\sqrt{2}a\)．\(|OA|=|OB|=|OP|=a\)，故\(O\)、\(A\)、\(P\)、\(B\)四点共圆（圆心\(O\)、半径\(a\)），其中\(AB\)为该圆的直径，\(P\)在圆上，由直径所对圆周角得\(\angle APB=90^{\circ}\)．\(M\)在圆上且在第一象限，位于弦\(AB\)所对优弧，故\(\angle AMB=(1/2)\angle APB=45^{\circ}\)．\par\noindent 夹角公式：\(\tan\angle AMB=(n-m)/(1+nm)=1\)，\(n-m=3\)，得\(nm=2\)．设\(M(x,y)\)：\(m=y/(x+a)\)，\(n=y/(x-a)\)，\(nm=y^{2}/(x^{2}-a^{2})\)；由\(y^{2}=b^{2}(x^{2}/a^{2}-1)\)得\(nm=b^{2}/a^{2}=2\)．\(e^{2}=1+b^{2}/a^{2}=3\)，\(e=\sqrt{3}\)．}
\fi
\end{ansblock}

\begin{ansblock}[2章-拓-课时14-06]
% ans:2章-拓-课时14-06
\ansitem{36}{2√3/3}
\ifansbookdetail
\ansnote{详解}{设\(A(a,0)\)，\(B(-a,0)\)，\(P(c,y)\)（\(l\)：\(x=c\)）．\(k_{PA}=y/(c-a)\)，\(k_{PB}=y/(c+a)\)．由夹角公式\(\tan\angle APB=\left|[k_{PA}-k_{PB}]/[1+k_{PA}k_{PB}]\right|=\sqrt{3}\)：\(\left|(2ay)/(c^{2}-a^{2})\right|/\left|1+y^{2}/(c^{2}-a^{2})\right|=\sqrt{3}\)，整理为关于\(y\)的方程\(\sqrt{3}y^{2}-2ay+\sqrt{3}(c^{2}-a^{2})=0\)有实根．判别式\(\Delta=4a^{2}-12(c^{2}-a^{2})=4(4a^{2}-3c^{2})\geqslant 0\)，得\(e^{2}\leqslant 4/3\)，\(e\leqslant 2\sqrt{3}/3\)．离心率最大值\(2\sqrt{3}/3\)．验算：\(e=2\sqrt{3}/3\)时\(\Delta=0\)，\(P\)取对称轴位置\(y=a/\sqrt{3}\)合．}
\fi
\end{ansblock}

\begin{ansblock}[2章-拓-课时14-09]
% ans:2章-拓-课时14-09
\ansitem{37}{18}
\ifansbookdetail
\ansnote{详解}{\(PQ\)与\(F_1F_2\)互相平分（\(P\)、\(Q\)关于\(O\)对称，\(F_1\)、\(F_2\)关于\(O\)对称）且\(|PQ|=|F_1F_2|\)，对角线相等又互相平分，故四边形\(PF_1QF_2\)为矩形，其面积\(=|PF_1|\cdot|PF_2|\)．矩形中\(|PF_1|^{2}+|PF_2|^{2}=|F_1F_2|^{2}=4c^{2}=100\)；又\(\bigl||PF_1|-|PF_2|\bigr|=2a=8\)．故\(2|PF_1||PF_2|=100-64=36\)，\(|PF_1||PF_2|=18\)，面积18．}
\fi
\end{ansblock}

\begin{ansblock}[2章-拓-课时14-10]
% ans:2章-拓-课时14-10
\ansitem{38}{3+2√3}
\ifansbookdetail
\ansnote{详解}{左焦点\((-\sqrt{a^{2}+1},0)=(-2,0)\)，得\(a^{2}=3\)．设\(P(x_0,y_0)\)（右支\(x_0\geqslant\sqrt{3}\)），\(y_0^{2}=x_0^{2}/3-1\)．\(\overrightarrow{OP}\cdot\overrightarrow{FP}=(x_0,y_0)\cdot(x_0+2,y_0)=x_0^{2}+2x_0+y_0^{2}=(4/3)x_0^{2}+2x_0-1\)．在\([\sqrt{3},+\infty)\)上单调递增，\(x_0=\sqrt{3}\)时取最小：\((4/3)\times 3+2\sqrt{3}-1=3+2\sqrt{3}\)．}
\fi
\end{ansblock}

\begin{ansblock}[2章-拓-课时14-11]
% ans:2章-拓-课时14-11
\ansitem{39}{(1) x²/400−y²/500=1（x<0）；(2) 20√2 km。}
\ifansbookdetail
\ansnote{详解}{(1)以\(O\)为原点、\(AB\)所在直线为\(x\)轴建系：\(A(-30,0)\)，\(B(30,0)\)，\(C(0,30)\)．设观察员位置\(P(x,y)\)：\(|PB|-|PA|=(40/V_0)\cdot V_0=40<|AB|=60\)，轨迹为双曲线左支．\(2a=40\)，\(2c=60\)，\(a=20\)，\(c=30\)，\(b^{2}=c^{2}-a^{2}=500\)，轨迹方程\(x^{2}/400-y^{2}/500=1\)（\(x<0\)）．\par\noindent (2)设轨迹上一点\(M(x,y)\)：\(|MC|^{2}=x^{2}+(y-30)^{2}=x^{2}+y^{2}-60y+900\)．由\(x^{2}=(4/5)y^{2}+400\)：\(|MC|^{2}=(9/5)y^{2}-60y+1300=(9/5)(y-50/3)^{2}+800\geqslant 800\)．当\(y=50/3\)时\(|MC|_{\min}=20\sqrt{2}\)，扫描半径\(r\)至少\(20\sqrt{2}\) km．}
\fi
\end{ansblock}

\begin{ansblock}[2章-拓-课时14-12]
% ans:2章-拓-课时14-12
\ansitem{40}{CD}
\ifansbookdetail
\ansnote{详解}{A：按定义实虚轴互换，共轭双曲线应为\(y^{2}/b^{2}-x^{2}/a^{2}=1\)，A错；B：原曲线渐近线\(y=\pm(b/a)x\)，共轭曲线（\(y\)轴型）渐近线\(y=\pm(b/a)x\)（\(a/b\)与\(b/a\)互换后方程相同），渐近线相同，B错；C：\(e_1=c/a\)，\(e_2=c/b\)，\(e_1e_2=c^{2}/(ab)=(a^{2}+b^{2})/(ab)=a/b+b/a\geqslant 2\)（当且仅当\(a=b\)取等），C对；D：焦点\((\pm c,0)\)与\((0,\pm c)\)均在圆\(x^{2}+y^{2}=c^{2}\)上，D对．故选CD．}
\fi
\end{ansblock}

\begin{ansblock}[2章-拓-课时14-32-1]
% ans:2章-拓-课时14-32-1
\ansitem{41}{D}
\ifansbookdetail
\ansnote{详解}{设\(P(x_0,y_0)\)、\(Q(-x_0,-y_0)\)，\(\triangle APQ\)三顶点纵坐标和为\(y_0-y_0+0=0\)，坐标和为\((a,0)\)，其重心为\(G(a/3,0)\)．重心在每条中线上，故\(G\)在中线\(QE\)上；\(QE\)与\(x\)轴的交点唯一，故\(M(1,0)\)即重心：\(a/3=1\)，\(a=3\)．\(b^{2}=1\)，\(c=\sqrt{9+1}=\sqrt{10}\)，\(e=c/a=\sqrt{10}/3\)，故选D．}
\fi
\end{ansblock}

\begin{ansblock}[2章-拓-课时14-32-2]
% ans:2章-拓-课时14-32-2
\ansitem{42}{(0, 3√3)}
\ifansbookdetail
\ansnote{详解}{\(e=2\)，\(c^{2}=4a^{2}\)，\(b^{2}=c^{2}-a^{2}=3a^{2}\)，\(b/a=\sqrt{3}\)．\(k_1k_2=\left[y_0/(x_0+a)\right]\cdot\left[y_0/(x_0-a)\right]=y_0^{2}/(x_0^{2}-a^{2})=(b^{2}/a^{2})(x_0^{2}-a^{2})/(x_0^{2}-a^{2})=b^{2}/a^{2}=3\)（坐标点差法，同例24）．\(P\)在第一象限右支：\(k_3=y_0/x_0\in(0,b/a)=(0,\sqrt{3})\)（\(P\)趋近顶点时\(k_3\to 0\)，沿渐近线方向\(k_3\to\sqrt{3}\)）．故\(m=3k_3\in(0,3\sqrt{3})\)．}
\fi
\end{ansblock}

\begin{ansblock}[2章-拓-课时14-33-1]
% ans:2章-拓-课时14-33-1
\ansitem{43}{[√3, 2√3)}
\ifansbookdetail
\ansnote{详解}{机器人到\(P\)的距离比到\(Q\)大2：距离差\(2<|PQ|=4\)，且距\(P\)更远，轨迹为双曲线右支：\(a=1\)，\(c=2\)，\(b^{2}=3\)，即\(x^{2}-y^{2}/3=1\)（\(x\geqslant 1\)）．直线\(y-3=k(x-\sqrt{3})\)．\(k=\sqrt{3}\)时直线恰为渐近线\(y=\sqrt{3}x\)，符合“接触不到”；\(3-k^{2}\neq 0\)时由\(\Delta<0\)解得\(\sqrt{3}<k<2\sqrt{3}\)（\(k=-\sqrt{3}\)不合）．综上\(k\in[\sqrt{3},2\sqrt{3})\)．验算：\(k=2\sqrt{3}\)时\(\Delta=0\)，直线与右支相切（可触），故右端开合．}
\fi
\end{ansblock}

\begin{ansblock}[2章-拓-课时14-33-2]
% ans:2章-拓-课时14-33-2
\ansitem{44}{A}
\ifansbookdetail
\ansnote{详解}{\(e=2\)，\(b^{2}/a^{2}=3\)，直线\(MN\)：\(y=\sqrt{3}(x+a)\)．设\(P(t,\sqrt{3}t+\sqrt{3}a)\)，\(t\in[-a,0]\)．\(\overrightarrow{PF_1}\cdot\overrightarrow{PF_2}=(-c-t,-\sqrt{3}t-\sqrt{3}a)\cdot(c-t,-\sqrt{3}t-\sqrt{3}a)=t^{2}-c^{2}+3(t+a)^{2}=4t^{2}+6at-a^{2}=4(t+3a/4)^{2}-13a^{2}/4\)．\(t=-3a/4\)取最小（\(y_P=\sqrt{3}a/4\)），\(t=0\)取最大（\(y_P=\sqrt{3}a\)）．面积比即纵坐标绝对值之比：\(S_1/S_2=\sqrt{3}a/(\sqrt{3}a/4)=4\)，故选A．}
\fi
\end{ansblock}

\begin{ansblock}[2章-拓-课时14-33-3]
% ans:2章-拓-课时14-33-3
\ansitem{45}{2}
\ifansbookdetail
\ansnote{详解}{\(D(a,b)\)，\(E(a,-b)\)，\(|DE|=2b\)，\(O\)到直线\(x=a\)距离为\(a\)：\(S=(1/2)\cdot a\cdot 2b=ab\)．焦距\(4\)，\(c=2\)，\(a^{2}+b^{2}=4\geqslant 2ab\)，得\(ab\leqslant 2\)（\(a=b=\sqrt{2}\)取等）．\(\triangle ODE\)面积最大值2．}
\fi
\end{ansblock}

\begin{ansblock}[2章-拓-课时14-33-12]
% ans:2章-拓-课时14-33-12
\ansitem{46}{√5}
\ifansbookdetail
\ansnote{详解}{以最细处直径\(A_1A_2\)为\(x\)轴、中垂线为\(y\)轴建系：\(2a=16\)，\(a=8\)．上瓶口点\(B(10,12)\)在双曲线\(x^{2}/64-y^{2}/b^{2}=1\)上：\(100/64-144/b^{2}=1\)，\(144/b^{2}=36/64\)，\(b^{2}=256\)．\(c^{2}=a^{2}+b^{2}=64+256=320\)，\(e^{2}=320/64=5\)，\(e=\sqrt{5}\)．}
\fi
\end{ansblock}

\begin{ansblock}[2章-拓-课时14-33-13]
% ans:2章-拓-课时14-33-13
\ansitem{47}{在 A 的北偏东 30° 方向，相距 10 海里处。}
\ifansbookdetail
\ansnote{详解}{以\(AB\)为\(x\)轴、\(AB\)中点为原点建系：\(A(3,0)\)，\(B(-3,0)\)，\(C(-5,2\sqrt{3})\)．\(|PB|-|PA|=4<6=|AB|\)，\(P\)在以\(A\)、\(B\)为焦点、实轴长4的双曲线右支上：\(x^{2}/4-y^{2}/5=1\)（\(x\geqslant 2\)）．\par\noindent \(|PB|=|PC|\)（\(B\)、\(C\)同时收到），\(P\)在\(BC\)中垂线上：\(BC\)中点\((-4,\sqrt{3})\)，\(k_{BC}=2\sqrt{3}/(-5+3)=-\sqrt{3}\)，中垂线\(y-\sqrt{3}=(\sqrt{3}/3)(x+4)\)，即\(x-\sqrt{3}y+7=0\)．与双曲线联立（\(y=(x+7)/\sqrt{3}\)代入）：\(11x^{2}-56x-256=0\)，解得\(x=8\)或\(x=-32/11\)（左支交点，不合\(x\geqslant 2\)，舍），\(y=5\sqrt{3}\)．\(P(8,5\sqrt{3})\)：\(k_{AP}=5\sqrt{3}/(8-3)=\sqrt{3}\)，倾斜角\(60^{\circ}\)，\(|PA|=10\)．故\(P\)在\(A\)的北偏东\(30^{\circ}\)方向，相距10海里处．验算：\(|PB|=\sqrt{11^{2}+75}=14\)，\(|PB|-|PA|=4\)合．}
\fi
\end{ansblock}

\begin{ansblock}[2章-拓-课时14-34-1]
% ans:2章-拓-课时14-34-1
\ansitem{48}{2}
\ifansbookdetail
\ansnote{详解}{两者焦点均在\(x\)轴：\(9-k^{2}=k+3\)（\(k>0\)），即\(k^{2}+k-6=0\)，\((k+3)(k-2)=0\)，\(k=2\)（\(k=-3\)舍）．验算：\(k=2\)时椭圆\(c^{2}=9-4=5\)，双曲线\(c^{2}=2+3=5\)合．}
\fi
\end{ansblock}

\begin{ansblock}[2章-拓-课时14-34-2]
% ans:2章-拓-课时14-34-2
\ansitem{49}{C}
\ifansbookdetail
\ansnote{详解}{以\(O\)为原点、\(AD\)所在直线为\(x\)轴建系．八等分结合\(AB=1\)、\(BC=2\)、\(CD=1\)：圆半径\(R=2\)，双曲线过点\((\sqrt{2},\sqrt{2})\)，\(a=1\)．\(2/1-2/b^{2}=1\)，得\(b^{2}=2\)，\(c^{2}=1+2=3\)，\(c=\sqrt{3}\)，焦距\(2\sqrt{3}\)，故选C．}
\fi
\end{ansblock}

\begin{ansblock}[2章-拓-课时14-34-3]
% ans:2章-拓-课时14-34-3
\ansitem{50}{B}
\ifansbookdetail
\ansnote{详解}{\(\overrightarrow{P_0F_1}\cdot\overrightarrow{PF_1}=0\)，即\(P_0F_1\perp PF_1\)；\(P_0\)与\(P\)关于原点对称，\(|P_0F_2|=|PF_1|=a\)且\(P_0F_1\perp P_0F_2\)（\(P_0F_2\parallel PF_1\)）．\(P_0\)在右支：\(|P_0F_1|-|P_0F_2|=2a\)，得\(|P_0F_1|=3a\)．Rt\(\triangle P_0F_1F_2\)（直角在\(P_0\)）：\((3a)^{2}+a^{2}=4c^{2}\)，\(10a^{2}=4(a^{2}+b^{2})\)，\(b^{2}/a^{2}=3/2\)．渐近线\(y=\pm(\sqrt{6}/2)x\)，故选B．}
\fi
\end{ansblock}

\begin{ansblock}[2章-拓-课时14-34-4]
% ans:2章-拓-课时14-34-4
\ansitem{51}{(−∞,−2]}
\ifansbookdetail
\ansnote{详解}{设\(P(x_0,y_0)\)（\(|x_0|\geqslant\sqrt{3}\)），\(Q(-x_0,-y_0)\)．\(\overrightarrow{MP}=(x_0,y_0-1)\)，\(\overrightarrow{MQ}=(-x_0,-y_0-1)\)，\(\overrightarrow{MP}\cdot\overrightarrow{MQ}=-x_0^{2}-y_0^{2}+1\)．由\(y_0^{2}=x_0^{2}/3-1\)：\(\overrightarrow{MP}\cdot\overrightarrow{MQ}=-x_0^{2}-x_0^{2}/3+1+1=2-(4/3)x_0^{2}\leqslant 2-(4/3)\times 3=-2\)．取值范围\((- \infty,-2]\)．}
\fi
\end{ansblock}

\begin{ansblock}[2章-拓-课时14-34-5]
% ans:2章-拓-课时14-34-5
\ansitem{52}{最大值为 5/3。}
\ifansbookdetail
\ansnote{详解}{\(|PF_1|-|PF_2|=3|PF_2|=2a\)，得\(|PF_2|=(2/3)a\)．右支焦半径\(|PF_2|\geqslant c-a\)：\((2/3)a\geqslant c-a\)，\(c\leqslant (5/3)a\)，\(e\leqslant 5/3\)；又\(e>1\)．故\(e\)的最大值\(5/3\)．}
\fi
\end{ansblock}

\begin{ansblock}[2章-拓-课时14-34-6]
% ans:2章-拓-课时14-34-6
\ansitem{53}{A}
\ifansbookdetail
\ansnote{详解}{\(e^{2}=3/2=c^{2}/a^{2}=(a^{2}+b^{2})/a^{2}\)，得\(b^{2}=a^{2}/2\)．过\((1,-2)\)：\(4/a^{2}-1/b^{2}=1\)，\(4/a^{2}-2/a^{2}=1\)，\(a^{2}=2\)，\(b^{2}=1\)．方程\(y^{2}/2-x^{2}=1\)，故选A．验算：\(4/2-1/1=1\)合，\(e=\sqrt{3/2}=\sqrt{6}/2\)合．}
\fi
\end{ansblock}

\begin{ansblock}[2章-练-课时14-15]
% ans:2章-练-课时14-15
\ansitem{54}{(1) e=√5/2；(2) 证明见解析，定点坐标为 (−10/3,0)。}
\ifansbookdetail
\ansnote{详解}{(1)\(a=2\)，\(b=1\)，\(c=\sqrt{4+1}=\sqrt{5}\)，\(e=c/a=\sqrt{5}/2\)．\par\noindent (2)设\(A(x_1,y_1)\)，\(B(x_2,y_2)\)．联立\(y=kx+m\)与\(x^{2}/4-y^{2}=1\)：\((1-4k^{2})x^{2}-8mkx-4(m^{2}+1)=0\)，其中\(1-4k^{2}\neq 0\)（否则\(l\)与渐近线平行，不合两交点），\(\Delta=64m^{2}k^{2}+16(1-4k^{2})(m^{2}+1)>0\)；\(x_1+x_2=8mk/(1-4k^{2})\)，\(x_1x_2=-4(m^{2}+1)/(1-4k^{2})\)；\(y_1y_2=(kx_1+m)(kx_2+m)=k^{2}x_1x_2+mk(x_1+x_2)+m^{2}=(m^{2}-4k^{2})/(1-4k^{2})\)．\par\noindent 以\(AB\)为直径的圆过\(D(-2,0)\Leftrightarrow\overrightarrow{DA}\cdot\overrightarrow{DB}=0\Leftrightarrow(x_1+2)(x_2+2)+y_1y_2=0\)，即\(x_1x_2+2(x_1+x_2)+4+y_1y_2=0\)．通分（分母\(1-4k^{2}\)）：\([m^{2}-4k^{2}]+[-4(m^{2}+1)]+[16mk]+[4(1-4k^{2})]=-3m^{2}+16mk-20k^{2}=0\)，即\(3m^{2}-16mk+20k^{2}=0\)，解得\(m=2k\)或\(m=(10/3)k\)．\par\noindent 当\(m=2k\)时，\(l\)：\(y=k(x+2)\)过\(D(-2,0)\)，此时\(A\)、\(B\)之一即\(D\)，与“\(A\)、\(B\)均异于左、右顶点”矛盾，舍去；当\(m=(10/3)k\)时（\(k\neq 0\)），\(l\)：\(y=k(x+10/3)\)恒过定点\((-10/3,0)\)，且不过\(D\)，符合题意．故直线\(l\)过定点，定点坐标为\((-10/3,0)\)．验算：\(3m^{2}-16mk+20k^{2}=(3m-10k)(m-2k)\)合；定点\((-10/3,0)\)在左顶点\(D(-2,0)\)左侧，直线族绕该点转动时与双曲线恒有两交点（\(k\)充分小时），满足\(\Delta>0\)合．}
\fi
\end{ansblock}

\begin{ansblock}[2章-练-课时14-16]
% ans:2章-练-课时14-16
\ansitem{55}{(1) y=±√3x；(2) 存在，定直线方程 x=1/2。}
\ifansbookdetail
\ansnote{详解}{(1)\(|OF_1|=|OF_2|=c=2\)，\(|OP|=|OF_1|\)说明\(\triangle PF_1F_2\)中边\(F_1F_2\)上的中线等于该边一半，故\(\angle F_1PF_2=90^{\circ}\)，即\(PF_1\perp PF_2\)．\(P\)在右支：\(|PF_1|-|PF_2|=2a\)；面积：\(|PF_1|\cdot|PF_2|=2S=6\)．\(|PF_1|^{2}+|PF_2|^{2}=|F_1F_2|^{2}=(2c)^{2}=16=(|PF_1|-|PF_2|)^{2}+2|PF_1||PF_2|=4a^{2}+12\)，得\(4a^{2}=4\)，\(a=1\)．\(b^{2}=c^{2}-a^{2}=4-1=3\)，渐近线\(y=\pm(b/a)x=\pm\sqrt{3}x\)．\par\noindent (2)双曲线\(C\)：\(x^{2}-y^{2}/3=1\)．（i）\(l\)斜率不存在：\(l\)：\(x=2\)，\(M(2,3)\)，\(N(2,-3)\)．直线\(A_1M\)：\(y=x+1\)；直线\(A_2N\)：\(y=-3x+3\)．联立得\(Q(1/2,3/2)\)，横坐标\(1/2\)．\par\noindent （ii）\(l\)斜率存在：\(l\)过\(F_2(2,0)\)且不与渐近线平行，设\(l\)：\(y=k(x-2)\)（\(k\neq 0\)，\(k\neq\pm\sqrt{3}\)），\(M(x_1,y_1)\)，\(N(x_2,y_2)\)．联立得\((3-k^{2})x^{2}+4k^{2}x-4k^{2}-3=0\)，\(x_1+x_2=4k^{2}/(k^{2}-3)\)，\(x_1x_2=(4k^{2}+3)/(k^{2}-3)\)．直线\(A_1M\)：\(y=[y_1/(x_1+1)](x+1)\)；直线\(A_2N\)：\(y=[y_2/(x_2-1)](x-1)\)．交点\(Q(x,y)\)满足\((x+1)/(x-1)=y_2(x_1+1)/[y_1(x_2-1)]\)（\(x\neq 1\)）．两边平方，以\(y_1^{2}=3(x_1^{2}-1)\)、\(y_2^{2}=3(x_2^{2}-1)\)代入：\([(x+1)/(x-1)]^{2}=[3(x_2^{2}-1)(x_1+1)^{2}]/[3(x_1^{2}-1)(x_2-1)^{2}]=[(x_2+1)(x_1+1)]/[(x_1-1)(x_2-1)]=[x_1x_2+(x_1+x_2)+1]/[x_1x_2-(x_1+x_2)+1]=[(4k^{2}+3)+4k^{2}+(k^{2}-3)]/[(4k^{2}+3)-4k^{2}+(k^{2}-3)]=9k^{2}/k^{2}=9\)．故\((x+1)/(x-1)=\pm 3\)，得\(x=1/2\)或\(x=2\)．\par\noindent 取\(x=2\)需\(y_2(x_1+1)/[y_1(x_2-1)]=+3>0\)；但\(M\)、\(N\)是过右焦点的弦的两端点且均非顶点：\(0<k^{2}<3\)时\(M\)、\(N\)分居两支（\(x_1x_2<0\)），\(x_2-1<0\)而\(y_1y_2>0\)；\(k^{2}>3\)时\(M\)、\(N\)同在右支且分居\(x\)轴两侧，\(y_1y_2<0\)而\((x_1+1)(x_2-1)>0\)——两种情形该式恒为负，只能取\(-3\)．故\(x=2\)为平方增根，舍去，\(x=1/2\)．综上，交点\(Q\)恒在定直线\(x=1/2\)上．验算：（i）的\(Q(1/2,3/2)\)落在\(x=1/2\)上合；（ii）韦达分子分母化简\(4k^{2}+3+4k^{2}+k^{2}-3=9k^{2}\)、\(4k^{2}+3-4k^{2}+k^{2}-3=k^{2}\)，比值9合．}
\fi
\end{ansblock}

\begin{ansblock}[2章-拓-课时14-28]
% ans:2章-拓-课时14-28
\ansitem{56}{C}
\ifansbookdetail
\ansnote{详解}{设\(PQ\)：\(x=my+b\)，联立\(x^{2}-y^{2}/2=1\)：\((m^{2}-1/2)y^{2}+2bmy+b^{2}-1=0\)．设\(P(x_1,y_1)\)，\(Q(x_2,y_2)\)：\(y_1+y_2=-2bm/(m^{2}-1/2)\)，\(y_1y_2=(b^{2}-1)/(m^{2}-1/2)\)．\(AP\perp AQ\)（\(A(-1,0)\)）：\((x_1+1)(x_2+1)+y_1y_2=0\)，以\(x_i=my_i+b\)代入：\((1+m^{2})y_1y_2+(b+1)m(y_1+y_2)+(b+1)^{2}=0\)．\par\noindent 通分（乘\(m^{2}-1/2\)）：\(2(b^{2}-1)(1+m^{2})-4bm^{2}(b+1)+(2m^{2}-1)(b+1)^{2}=0\)．按\(m\)整理——\(m^{2}\)系数：\(2(b^{2}-1)-4b(b+1)+2(b+1)^{2}=2b^{2}-2-4b^{2}-4b+2b^{2}+4b+2=0\)（恒为0）；常数项：\(2(b^{2}-1)-(b+1)^{2}=b^{2}-2b-3=0\)，得\(b=3\)或\(b=-1\)．\(b=-1\)时\(PQ\)过\(A(-1,0)\)（\(P\)、\(Q\)之一与\(A\)重合或直线退化），舍去；\(b=3\)，\(PQ\)恒过\((3,0)\)．故选C．\par\noindent 验算（特例）：\(PQ\perp x\)轴即\(m=0\)、\(x=3\)：与\(x^{2}-y^{2}/2=1\)交于\((3,\pm 4)\)；\(\overrightarrow{AP}=(4,4)\)、\(\overrightarrow{AQ}=(4,-4)\)，\(\overrightarrow{AP}\cdot\overrightarrow{AQ}=16-16=0\)合．}
\fi
\end{ansblock}

\begin{ansblock}[2章-拓-课时14-29]
% ans:2章-拓-课时14-29
\ansitem{57}{(1) x²−y²/3=1；(2) 证明见解析，M(−1,0)。}
\ifansbookdetail
\ansnote{详解}{(1)取渐近线\(y=(b/a)x\)，与\(x=a^{2}/c\)交于\(P(a^{2}/c,ab/c)\)．\(|PF|^{2}=(c-a^{2}/c)^{2}+(ab/c)^{2}\)，其中\(c-a^{2}/c=(c^{2}-a^{2})/c=b^{2}/c\)，故\(|PF|^{2}=b^{4}/c^{2}+a^{2}b^{2}/c^{2}=b^{2}(a^{2}+b^{2})/c^{2}=b^{2}c^{2}/c^{2}=b^{2}\)．由\(|PF|=\sqrt{3}\)得\(b^{2}=3\)．\(e^{2}=c^{2}/a^{2}=4=(a^{2}+b^{2})/a^{2}\)，得\(a^{2}=1\)．\(C\)：\(x^{2}-y^{2}/3=1\)．\par\noindent (2)\(F(2,0)\)．设\(Q(x_0,y_0)\)（\(x_0\geqslant 1\)），\(3x_0^{2}-y_0^{2}=3\)．①\(x_0=2\)时\(y_0=\pm 3\)：\(\angle QFM=90^{\circ}\)需\(\angle QMF=45^{\circ}\)，\(|MF|=|QF|=|y_0|=3\)，得\(M(-1,0)\)，符合．②\(x_0\neq 2\)时设\(M(t,0)\)（\(t<0\)）．有向正切：\(\tan\angle QFM=-y_0/(x_0-2)\)，\(\tan\angle QMF=y_0/(x_0-t)\)．\(\angle QFM=2\angle QMF\)，由倍角公式：\(-y_0/(x_0-2)=2[y_0/(x_0-t)]/(1-[y_0/(x_0-t)]^{2})\)．\(y_0\neq 0\)时约去\(y_0\)、交叉相乘：\(-[(x_0-t)^{2}-y_0^{2}]=2(x_0-2)(x_0-t)\)，展开整理：\(3x_0^{2}-y_0^{2}-(4+4t)x_0+4t+t^{2}=0\)．以\(3x_0^{2}-y_0^{2}=3\)代入：\(-(4+4t)x_0+(t^{2}+4t+3)=0\)，对任意\(x_0\geqslant 1\)（\(x_0\neq 2\)）恒成立，故\(-(4+4t)=0\)且\(t^{2}+4t+3=0\)，得\(t=-1\)（两式公共解）．\(y_0=0\)（\(Q\)为右顶点）时直接验证\(M(-1,0)\)亦满足\(\angle QFM=2\angle QMF\)．综上，存在定点\(M(-1,0)\)．验算：\(t^{2}+4t+3=(t+1)(t+3)\)与\(-(4+4t)\)公共根\(t=-1\)合；取\(Q(2,3)\)：\(\angle QFM=90^{\circ}\)、\(\angle QMF=45^{\circ}\)，\(2\times 45^{\circ}=90^{\circ}\)合．}
\fi
\end{ansblock}

\begin{ansblock}[2章-拓-课时14-30]
% ans:2章-拓-课时14-30
\ansitem{58}{(1) x²−y²/3=1（x>0）；(2) (−∞,−√3)∪(√3,+∞)；(3) 存在，M(−1,0)。}
\ifansbookdetail
\ansnote{详解}{(1)\(|PF_1|-|PF_2|=2<|F_1F_2|=4\)，轨迹为以\(F_1\)、\(F_2\)为焦点、实轴长2的双曲线右支：\(c=2\)，\(a=1\)，\(b^{2}=3\)，\(\Gamma\)：\(x^{2}-y^{2}/3=1\)（\(x>0\)）．\par\noindent (2)\(l\)：\(y=k(x-2)\)，代入\(\Gamma\)：\((3-k^{2})x^{2}+4k^{2}x-4k^{2}-3=0\)．\(A\)、\(B\)为两交点（右支上\(x>0\)）：需\(3-k^{2}\neq 0\)、\(\Delta>0\)、\(x_1+x_2=-4k^{2}/(3-k^{2})>0\)、\(x_1x_2=(-4k^{2}-3)/(3-k^{2})>0\)．由\(x_1+x_2>0\)得\(k^{2}>3\)；此时\(x_1x_2>0\)、\(\Delta=16k^{4}-4(3-k^{2})(-4k^{2}-3)>0\)亦成立．故\(k\in(-\infty,-\sqrt{3})\cup(\sqrt{3},+\infty)\)．\par\noindent (3)存在性：设\(M(m,0)\)．\(l\)与\(\Gamma\)联立同上：\((3-k^{2})x^{2}+4k^{2}x-4k^{2}-3=0\)（\(k^{2}>3\)），\(x_1+x_2=-4k^{2}/(3-k^{2})\)，\(x_1x_2=(-4k^{2}-3)/(3-k^{2})\)．\(MA\perp MB\Leftrightarrow(x_1-m)(x_2-m)+y_1y_2=0\)，\(y_i=k(x_i-2)\)：\((1+k^{2})x_1x_2-(m+2k^{2})(x_1+x_2)+m^{2}+4k^{2}=0\)．乘\((3-k^{2})\)展开：\((1+k^{2})(-4k^{2}-3)+4k^{2}(2k^{2}+m)+(m^{2}+4k^{2})(3-k^{2})=0\)，即\(k^{2}(-m^{2}+4m+5)+3m^{2}-3=0\)，对一切\(k^{2}>3\)恒成立．故\(-m^{2}+4m+5=0\)且\(3m^{2}-3=0\)，得\(m=-1\)（\(m=1\)不满足第一式）．斜率不存在时\(l\)：\(x=2\)，\(A(2,3)\)、\(B(2,-3)\)：\(\overrightarrow{MA}=(3,3)\)、\(\overrightarrow{MB}=(3,-3)\)，\(\overrightarrow{MA}\cdot\overrightarrow{MB}=0\)合．综上，存在\(M(-1,0)\)．验算：\(k^{2}(-m^{2}+4m+5)+3m^{2}-3\)在\(m=-1\)时两项恒为0合；(3)的韦达代入式与(2)一致合．}
\fi
\end{ansblock}

\begin{ansblock}[2章-拓-课时14-33-4]
% ans:2章-拓-课时14-33-4
\ansitem{59}{(1) e=2；(2) 3/2。}
\ifansbookdetail
\ansnote{详解}{(1)右支上\(|PF_1|_{\min}=c+a\)（左顶点处），\(|PF_2|_{\min}=c-a\)（右顶点处）．\(m_1m_2=c^{2}-a^{2}=b^{2}=3a^{2}\)，得\(c^{2}=4a^{2}\)，\(e=2\)．\par\noindent (2)\(a=2\)，\(c=4\)，\(b^{2}=12\)，双曲线\(x^{2}/4-y^{2}/12=1\)，\(F_1(-4,0)\)．设\(AB\)：\(y=k(x+4)\)（\(k\neq\pm\sqrt{3}\)）．联立：\((3-k^{2})x^{2}-8k^{2}x-16k^{2}-12=0\)，\(x_1+x_2=8k^{2}/(3-k^{2})\)，\(x_1x_2=(-16k^{2}-12)/(3-k^{2})\)．弦长\(|AB|=\sqrt{1+k^{2}}\cdot\sqrt{(x_1+x_2)^{2}-4x_1x_2}=12(1+k^{2})/|3-k^{2}|\)．中点\((4k^{2}/(3-k^{2}),12k/(3-k^{2}))\)（\(y_1+y_2=k(x_1+x_2)+8k=24k/(3-k^{2})\)）．垂直平分线：\(y-12k/(3-k^{2})=-(1/k)(x-4k^{2}/(3-k^{2}))\)，令\(x=0\)得\(y=16k/(3-k^{2})\)，\(|OD|=16|k|/|3-k^{2}|\)．\par\noindent \(|AB|/|OD|=[12(1+k^{2})/|3-k^{2}|]\cdot[|3-k^{2}|/(16|k|)]=(3/4)(|k|+1/|k|)\geqslant 3/2\)（\(|k|=1\)取等），最小值\(3/2\)．验算：\((x_1-x_2)^{2}=(8k^{2}/(3-k^{2}))^{2}+4(16k^{2}+12)/(3-k^{2})=[64k^{4}+4(16k^{2}+12)(3-k^{2})]/(3-k^{2})^{2}=144(k^{2}+1)/(3-k^{2})^{2}\)，\(|AB|=\sqrt{1+k^{2}}\cdot 12\sqrt{k^{2}+1}/|3-k^{2}|\)合．}
\fi
\end{ansblock}

\begin{ansblock}[2章-拓-课时14-33-5]
% ans:2章-拓-课时14-33-5
\ansitem{60}{(1) x²/2−y²=1；(2) 过定点 (0,1)。}
\ifansbookdetail
\ansnote{详解}{(1)\(e^{2}=3/2=c^{2}/a^{2}=(a^{2}+b^{2})/a^{2}\)，得\(b^{2}/a^{2}=1/2\)，即\(a^{2}=2b^{2}\)，双曲线方程为\(x^{2}/(2b^{2})-y^{2}/b^{2}=1\)．代入\(P\)：\(3/(2b^{2})-(1/2)/b^{2}=1\)，\((3-1)/(2b^{2})=1\)，\(b^{2}=1\)，\(a^{2}=2\)．\(C\)：\(x^{2}/2-y^{2}=1\)．\par\noindent (2)\(AB\)斜率不存在时：\(A(x_0,y_0)\)、\(B(x_0,-y_0)\)，\(k_1+k_2=(y_0-1)/(x_0-2)+(-y_0-1)/(x_0-2)=-2/(x_0-2)=1\)，得\(x_0=0\)，不在曲线上，舍；故斜率存在．设\(AB\)：\(y=kx+t\)，代入：\((2k^{2}-1)x^{2}+4ktx+2t^{2}+2=0\)（\(2k^{2}-1\neq 0\)），\(x_1+x_2=-4kt/(2k^{2}-1)\)，\(x_1x_2=(2t^{2}+2)/(2k^{2}-1)\)．\par\noindent \(k_1+k_2=1\)：\((kx_1+t-1)/(x_1-2)+(kx_2+t-1)/(x_2-2)=1\)，即\((2k-1)x_1x_2+(t-2k+1)(x_1+x_2)-4t=0\)．代入：\((2k-1)(2t^{2}+2)-4kt(t-2k+1)-4t(2k^{2}-1)=0\)（通分后乘\(2k^{2}-1\)），整理\(t^{2}+(2k-2)t-1+2k=0\)，即\((t-1)(t+2k-1)=0\)．\(t=1\)：\(y=kx+1\)过\((0,1)\)；\(t=1-2k\)：\(y=k(x-2)+1\)过\(Q(2,1)\)，不合（\(A\)、\(B\)异于\(Q\)且\(Q\)为弦端情形）．故直线\(AB\)过定点\((0,1)\)．验算：\(t^{2}+(2k-2)t+2k-1=(t-1)(t+2k-1)\)展开核对合．}
\fi
\end{ansblock}

\begin{ansblock}[2章-拓-课时14-33-6]
% ans:2章-拓-课时14-33-6
\ansitem{61}{(1) x²/4−y²=1；(2) 存在，Q(23/8,0)，QM·QN=273/64。}
\ifansbookdetail
\ansnote{详解}{(1)\(b/a=1/2\)且\(16/a^{2}-3/b^{2}=1\)：\(b^{2}=a^{2}/4\)，\(16/a^{2}-12/a^{2}=1\)，\(a^{2}=4\)，\(b^{2}=1\)．\(C\)：\(x^{2}/4-y^{2}=1\)．\par\noindent (2)设\(l\)：\(x=my+1\)，联立：\((m^{2}-4)y^{2}+2my-3=0\)（\(m^{2}\neq 4\)，\(\Delta>0\)得\(m^{2}>3\)），\(y_1+y_2=-2m/(m^{2}-4)\)，\(y_1y_2=-3/(m^{2}-4)\)；\(x_1+x_2=m(y_1+y_2)+2=-8/(m^{2}-4)\)，\(x_1x_2=(my_1+1)(my_2+1)=-(4m^{2}+4)/(m^{2}-4)\)．设\(Q(t,0)\)：\(\overrightarrow{QM}\cdot\overrightarrow{QN}=(x_1-t)(x_2-t)+y_1y_2=x_1x_2-t(x_1+x_2)+t^{2}+y_1y_2=-4+t^{2}+(8t-23)/(m^{2}-4)\)．与\(m\)无关\(\Leftrightarrow 8t-23=0\)，得\(t=23/8\)，常数\(-4+t^{2}=-4+529/64=273/64\)．验算：\(x_1x_2=(-4m^{2}-4)/(m^{2}-4)=-4-20/(m^{2}-4)\)；三项合并分子\(-20-3+8t=8t-23\)，故数量积\(-4+t^{2}+(8t-23)/(m^{2}-4)\)合．}
\fi
\end{ansblock}

\begin{ansblock}[2章-拓-课时14-33-7]
% ans:2章-拓-课时14-33-7
\ansitem{62}{(1) 15/4；(2) 双曲线中 k_AB·k_OM=b²/a²（定值），证明见解析。}
\ifansbookdetail
\ansnote{详解}{(1)联立得\(16x^{2}+12x-9=0\)，\(x_1+x_2=-3/4\)，\(x_1x_2=-9/16\)．根差\((x_1-x_2)^{2}=(-3/4)^{2}-4\times(-9/16)=45/16\)，故\(|AB|=\sqrt{5}\cdot 3\sqrt{5}/4=15/4\)．\par\noindent (2)性质：双曲线\(x^{2}/a^{2}-y^{2}/b^{2}=1\)上任意两点\(A\)、\(B\)（斜率存在时），\(M\)为\(AB\)中点，则\(k_{AB}\cdot k_{OM}=b^{2}/a^{2}\)．证明（点差法）：\(A(x_3,y_3)\)、\(B(x_4,y_4)\)、\(M(x_0,y_0)\)．两式相减：\((x_3+x_4)(x_3-x_4)/a^{2}=(y_3+y_4)(y_3-y_4)/b^{2}\)，即\(2x_0(x_3-x_4)/a^{2}=2y_0(y_3-y_4)/b^{2}\)，变形得\(k_{AB}\cdot k_{OM}=[(y_3-y_4)/(x_3-x_4)]\cdot(y_0/x_0)=b^{2}/a^{2}\)．定值证毕．验算：\(k_{AB}\cdot k_{OM}=b^{2}/a^{2}\)与例24的\(k_1k_2=b^{2}/a^{2}\)相容（中心对称两点为其特例）合．}
\fi
\end{ansblock}

\begin{ansblock}[2章-拓-课时14-33-8]
% ans:2章-拓-课时14-33-8
\ansitem{63}{(1) x²−y²/8=1；(2) 定值 2√2。}
\ifansbookdetail
\ansnote{详解}{(1)\(c=3\)，\(b/a=2\sqrt{2}\)，\(a^{2}+b^{2}=9\)，得\(a^{2}(1+8)=9\)，\(a^{2}=1\)，\(b^{2}=8\)．\(C\)：\(x^{2}-y^{2}/8=1\)．\par\noindent (2)\(l\)：\(y=kx+m\)（\(k\neq\pm 2\sqrt{2}\)，\(m\neq 0\)，斜率存在且不为0）．\(D(-m/k,0)\)，\(|OD|=|m/k|\)．相切：联立\((8-k^{2})x^{2}-2kmx-m^{2}-8=0\)，\(\Delta=0\)，得\(4k^{2}m^{2}+4(8-k^{2})(m^{2}+8)=0\)，即\(m^{2}=k^{2}-8\)（\(k^{2}>8\)）．\(M\)与\(y=2\sqrt{2}x\)交：\((m/(2\sqrt{2}-k),2\sqrt{2}m/(2\sqrt{2}-k))\)；\(N\)与\(y=-2\sqrt{2}x\)交：\((-m/(2\sqrt{2}+k),2\sqrt{2}m/(2\sqrt{2}+k))\)．\(S_{\triangle MON}=S_{\triangle MOD}+S_{\triangle NOD}=(1/2)|OD|\cdot|y_M-y_N|=(1/2)|m|\cdot|x_M-x_N|=(1/2)|m|\cdot|4\sqrt{2}m/(8-k^{2})|=2\sqrt{2}m^{2}/(k^{2}-8)=2\sqrt{2}\)（定值）．验算：\(m^{2}=k^{2}-8>0\)与\(k^{2}>8\)一致合；\(S\)只依赖\(m^{2}/(k^{2}-8)\)，化简后为常数合．}
\fi
\end{ansblock}

\begin{ansblock}[2章-拓-课时14-33-9]
% ans:2章-拓-课时14-33-9
\ansitem{64}{(1) y=±√3x；(2) 定值 k₁/k₂=−1/3。}
\ifansbookdetail
\ansnote{详解}{(1)\(a=1\)，\(c=2\)，\(b^{2}=3\)，\(\Gamma\)：\(x^{2}-y^{2}/3=1\)，渐近线\(y=\pm\sqrt{3}x\)．\par\noindent (2)\(l\)斜率不存在：\(P(2,3)\)、\(Q(2,-3)\)，\(k_1=3/3=1\)、\(k_2=-3/1=-3\)，比值\(-1/3\)．\(l\)：\(y=k(x-2)\)（\(k\neq 0\)）：\((k^{2}-3)x^{2}-4k^{2}x+4k^{2}+3=0\)，\(x_1+x_2=4k^{2}/(k^{2}-3)\)，\(x_1x_2=(4k^{2}+3)/(k^{2}-3)\)．\(3k_1+k_2=3y_1/(x_1+1)+y_2/(x_2-1)=k[4x_1x_2-5(x_1+x_2)+4]/[(x_1+1)(x_2-1)]\)．\(4x_1x_2-5(x_1+x_2)+4=[4(4k^{2}+3)-20k^{2}+4(k^{2}-3)]/(k^{2}-3)=0\)．故\(3k_1+k_2=0\)，\(k_1/k_2=-1/3\)（定值），综上得证．验算：分子\(16k^{2}+12-20k^{2}+4k^{2}-12=0\)合．}
\fi
\end{ansblock}

\begin{ansblock}[2章-拓-课时14-33-10]
% ans:2章-拓-课时14-33-10
\ansitem{65}{(1) 2x²−y²/2=1；(2) 定直线 12x−y−2=0。}
\ifansbookdetail
\ansnote{详解}{(1)\(2b=2\sqrt{2}\)，得\(b^{2}=2\)；\(e^{2}=5=1+b^{2}/a^{2}\)，得\(a^{2}=1/2\)．\(C\)：\(2x^{2}-y^{2}/2=1\)．\par\noindent (2)设\(M(x,y)\)、\(A(x_1,y_1)\)、\(B(x_2,y_2)\)（\(x_1<x_2<3\)）．\(|AM|/|MB|=|AP|/|PB|\)结合共线定比：\((3-x_1)/(x_2-3)=-(x-x_1)/(x_2-x)\)，即\([6-(x_1+x_2)]x=3(x_1+x_2)-2x_1x_2\)．①设\(l\)：\(y-1=k(x-3)\)，代入\(C\)：\((4-k^{2})x^{2}+(6k^{2}-2k)x-9k^{2}+6k-3=0\)，\(x_1+x_2=(2k-6k^{2})/(4-k^{2})\)，\(x_1x_2=(-9k^{2}+6k-3)/(4-k^{2})\)．代入①整理：\(12x-3=k(x-3)\)，与\(y-1=k(x-3)\)联立消\(k\)：\(12x-3=y-1\)，即\(12x-y-2=0\)．点\(M\)恒落在定直线\(12x-y-2=0\)上．验算：由定比\(x=[3(x_1+x_2)-2x_1x_2]/[6-(x_1+x_2)]\)（分比不变）合；韦达代入化简得\(12x-3=k(x-3)\)，代回\(y-1=k(x-3)\)消\(k\)得\(y=12x-2\)合．}
\fi
\end{ansblock}

\begin{ansblock}[2章-拓-课时14-33-11]
% ans:2章-拓-课时14-33-11
\ansitem{66}{(1) 2√5x−y−4√5=0 或 2√5x+y−4√5=0；(2) 定直线 x=1/2。}
\ifansbookdetail
\ansnote{详解}{设\(l\)：\(x=my+2\)，\(C(x_1,y_1)\)、\(D(x_2,y_2)\)，联立：\((4m^{2}-1)y^{2}+16my+12=0\)，\(y_1+y_2=-16m/(4m^{2}-1)\)，\(y_1y_2=12/(4m^{2}-1)\)．\par\noindent (1)\(|CN|=3|ND|\)（\(N\)为\(F_2\)，\(C\)、\(D\)同向分段），得\(y_1=-3y_2\)，\(y_2=8m/(4m^{2}-1)\)；代入\(y_1y_2=-3y_2^{2}=12/(4m^{2}-1)\)：\(-3\times 64m^{2}/(4m^{2}-1)^{2}=12/(4m^{2}-1)\)，得\(-16m^{2}=4m^{2}-1\)，\(m^{2}=1/20\)，\(m=\pm\sqrt{5}/10\)．直线\(l\)：\(x=\pm(\sqrt{5}/10)y+2\)，即\(2\sqrt{5}x-y-4\sqrt{5}=0\)或\(2\sqrt{5}x+y-4\sqrt{5}=0\)．\par\noindent (2)\(AC\)：\(y=[y_1/(x_1+1)](x+1)\)；\(BD\)：\(y=[y_2/(x_2-1)](x-1)\)．以\(x_1=my_1+2\)、\(x_2=my_2+2\)代入，得\(x_1+1=my_1+3\)、\(x_2-1=my_2+1\)．联立：\((x+1)/(x-1)=(my_1y_2+3y_2)/(my_1y_2+y_1)\)．由韦达\(y_1y_2=12/(4m^{2}-1)\)、\(y_1+y_2=-16m/(4m^{2}-1)\)，得\(y_1y_2=-[3/(4m)](y_1+y_2)\)．于是\((x+1)/(x-1)=[-(3/4)(y_1+y_2)+3y_2]/[-(3/4)(y_1+y_2)+y_1]=(-3y_1+9y_2)/(y_1-3y_2)=-3\)，解得\(x=1/2\)．点\(P\)恒在定直线\(x=1/2\)上．验算：\((-3y_1+9y_2)/(y_1-3y_2)=-3(y_1-3y_2)/(y_1-3y_2)=-3\)合；(1)中\(-16m^{2}=4m^{2}-1\)即\(20m^{2}=1\)合．}
\fi
\end{ansblock}

\begin{ansblock}[2章-导-课时14-G1]
% ans:2章-导-课时14-G1
\ansitem{67}{A}
\ifansbookdetail
\ansnote{详解}{\(e^{2}=1+b^{2}/a^{2}=1+m/9=13/9\)，得\(m/9=4/9\)，\(m=4\)．故\(C\)：\(x^{2}/9-y^{2}/4=1\)，\(a=3\)、\(b=2\)，渐近线\(y=\pm(b/a)x=\pm(2/3)x\)，选A．验算：\(e=\sqrt{1+4/9}=\sqrt{13}/3\)合．}
\fi
\end{ansblock}

\begin{ansblock}[2章-导-课时14-G2]
% ans:2章-导-课时14-G2
\ansitem{68}{A}
\ifansbookdetail
\ansnote{详解}{逐项核两条件：A：\(a^{2}=4\)、\(b^{2}=2\)，渐近线\(y=\pm(b/a)x=\pm(\sqrt{2}/2)x\)，\(c^{2}=4+2=6\)，焦点\((\pm\sqrt{6},0)\)在\(x\)轴，两条件皆合；B：渐近线\(y=\pm(a/b)x=\pm(\sqrt{2}/2)x\)与A相同，但实轴在\(y\)轴上，焦点\((0,\pm\sqrt{6})\)，不合“焦点在\(x\)轴”；C：渐近线\(y=\pm(b/a)x=\pm\sqrt{2}x\)，不合；D：渐近线\(y=\pm(a/b)x=\pm(1/2)x\)且焦点\((0,\pm 2\sqrt{5})\)在\(y\)轴，两条件皆不合．故选A．}
\fi
\end{ansblock}

\begin{ansblock}[2章-导-课时14-G3]
% ans:2章-导-课时14-G3
\ansitem{69}{A}
\ifansbookdetail
\ansnote{详解}{由\(e^{2}=1+(b/a)^{2}\)得\((b/a)^{2}=e^{2}-1\)．\(e\in[\sqrt{3},2]\)，则\(e^{2}\in[3,4]\)，\((b/a)^{2}\in[2,3]\)；又\(b/a>0\)，故\(b/a\in[\sqrt{2},\sqrt{3}]\)，选A．验算：\(e=\sqrt{3}\)时\(b/a=\sqrt{2}\)，\(e=2\)时\(b/a=\sqrt{3}\)，端点可取合．}
\fi
\end{ansblock}

\begin{ansblock}[2章-导-课时14-G4]
% ans:2章-导-课时14-G4
\ansitem{70}{x²/16−y²/4=1}
\ifansbookdetail
\ansnote{详解}{虚轴长\(2b=4\)，得\(b=2\)；焦点在\(x\)轴时渐近线为\(y=\pm(b/a)x\)，由\(b/a=1/2\)得\(a=2b=4\)．故标准方程为\(x^{2}/16-y^{2}/4=1\)．验算：\(a^{2}=16\)、\(b^{2}=4\)，渐近线\(y=\pm(2/4)x=\pm(1/2)x\)合，虚轴\(2b=4\)合，\(c^{2}=20\)，焦点\((\pm 2\sqrt{5},0)\)在\(x\)轴合．}
\fi
\end{ansblock}

\begin{ansblock}[2章-导-课时14-G5]
% ans:2章-导-课时14-G5
\ansitem{71}{m=14；e=2√2}
\ifansbookdetail
\ansnote{详解}{\(C_2\)（\(a^{2}=1\)、\(b^{2}=7\)）的渐近线为\(y=\pm\sqrt{7}x\)．\(C_1\)（\(a^{2}=2\)、\(b^{2}=m\)，焦点同在\(x\)轴）的渐近线为\(y=\pm(\sqrt{m}/\sqrt{2})x\)．渐近线相同即\(\sqrt{m}/\sqrt{2}=\sqrt{7}\)，得\(\sqrt{m}=\sqrt{14}\)，\(m=14\)．此时\(e^{2}=1+b^{2}/a^{2}=1+14/2=8\)，故\(e=2\sqrt{2}\)．验算：\(m=14\)时\(C_1\)渐近线\(y=\pm(\sqrt{14}/\sqrt{2})x=\pm\sqrt{7}x\)合，\(e=\sqrt{8}=2\sqrt{2}\)合．}
\fi
\end{ansblock}

\jietitle{2.7.1 抛物线的标准方程}

\begin{ansblock}[2章-练-课时15-01]
% ans:2章-练-课时15-01
\ansitem{1}{不是。当定点 F 在直线 l 上时，动点轨迹是与 l 垂直的直线（过 F）。}
\ifansbookdetail
\ansnote{详解}{取\(l\)为\(x\)轴、\(F\)为原点．动点\(P(x,y)\)到\(F\)的距离为\(\sqrt{x^{2}+y^{2}}\)，到\(l\)的距离为\(|y|\)．\par\noindent “到\(F\)与到\(l\)距离相等”：\(\sqrt{x^{2}+y^{2}}=|y|\)，平方得\(x^{2}+y^{2}=y^{2}\)，即\(x=0\)．轨迹为整条\(y\)轴：过\(F\)且垂直于\(l\)的直线，不是抛物线（\(p=0\)时抛物线无意义）．验算：点\((0,5)\)到\(F\)距离5、到\(x\)轴距离5，相等合．}
\fi
\end{ansblock}

\begin{ansblock}[2章-练-课时15-03]
% ans:2章-练-课时15-03
\ansitem{2}{x²=−12y}
\ifansbookdetail
\ansnote{详解}{\(F(0,-3)\)不在直线\(y=3\)上，由定义轨迹为抛物线：焦点\(F(0,-3)\)，准线\(y=3\)．开口向下，设\(x^{2}=-2py\)（\(p>0\)）：\(p/2=3\)，\(p=6\)，方程\(x^{2}=-12y\)．验算：顶点\((0,0)\)到\(F\)距3、到准线距3合．}
\fi
\end{ansblock}

\begin{ansblock}[2章-练-课时15-05]
% ans:2章-练-课时15-05
\ansitem{3}{C}
\ifansbookdetail
\ansnote{详解}{左端为\(P(x,y)\)到定点\(F(0,p/2)\)的距离；右端\(|y+p/2|\)为\(P\)到定直线\(l\)：\(y=-p/2\)的距离．\(F(0,p/2)\)不在\(l\)上（\(p>0\)），由抛物线定义，轨迹为抛物线，故选C．}
\fi
\end{ansblock}

\begin{ansblock}[2章-练-课时15-14]
% ans:2章-练-课时15-14
\ansitem{4}{y²=16x。}
\ifansbookdetail
\ansnote{详解}{设\(M(x,y)\)．“\(|MF|\)比到\(l\)的距离小2”即\(|x+6|-|MF|=2\)．当\(x\geqslant -6\)时：\(x+6-\sqrt{(x-4)^{2}+y^{2}}=2\)，即\(\sqrt{(x-4)^{2}+y^{2}}=x+4\)（隐含\(x\geqslant -4\)，自然满足），平方得\((x-4)^{2}+y^{2}=(x+4)^{2}\)，化简为\(y^{2}=16x\)（此时\(x\geqslant 0\geqslant -4\)，与所设一致）；\par\noindent 当\(x<-6\)时：\(-x-6-\sqrt{(x-4)^{2}+y^{2}}=2\)，即\(\sqrt{(x-4)^{2}+y^{2}}=-x-4\)，平方同样得\(y^{2}=16x\)，但该抛物线上\(x\geqslant 0\)，与\(x<-6\)矛盾，此段无轨迹．故轨迹方程\(y^{2}=16x\)（顶点在原点、开口向右的抛物线，焦点\(F(4,0)\)、准线\(x=-4\)）．验算：\(y^{2}=16x\)上任一点\(|MF|=x+4\)，到直线\(x=-6\)的距离为\(x+6\)，差恰为2合．\par\noindent 辨析：本题须防“大2”误读——若误作“\(|MF|\)比到直线\(x=-6\)的距离大2”，则\(|MF|=|x+8|\)，平方化简为\(y^{2}=24(x+2)\)，为平移型抛物线（顶点\((-2,0)\)），与标准形不符；按原题“小2”，平移准线得\(|MF|\)等于到直线\(x=-4\)的距离，才落到标准抛物线\(y^{2}=16x\)．}
\fi
\end{ansblock}

\begin{ansblock}[2章-拓-课时15-15]
% ans:2章-拓-课时15-15
\ansitem{5}{D}
\ifansbookdetail
\ansnote{详解}{由定义，抛物线上点到焦点的距离等于它到准线\(x=-p/2\)的距离，其最小值在顶点处取得，为\(p/2\)．恒大于\(1\Leftrightarrow p/2>1\Leftrightarrow p>2\)，故选D．验算：\(p=2\)时顶点到焦点距离\(=1\)，不满足“恒大于”，开区间正确合．}
\fi
\end{ansblock}

\begin{ansblock}[2章-练-课时15-02]
% ans:2章-练-课时15-02
\ansitem{6}{y²=3x 或 y²=−3x 或 x²=3y 或 x²=−3y。}
\ifansbookdetail
\ansnote{详解}{焦准距即\(p\)：\(p=3/2\)，\(2p=3\)．开口未指定，四种标准方程并立：\(y^{2}=3x\)、\(y^{2}=-3x\)、\(x^{2}=3y\)、\(x^{2}=-3y\)．}
\fi
\end{ansblock}

\begin{ansblock}[2章-练-课时15-07]
% ans:2章-练-课时15-07
\ansitem{7}{(1) y²=8x；(2) y²=6x。}
\ifansbookdetail
\ansnote{详解}{(1)焦点\((2,0)\)在\(x\)正半轴：\(p/2=2\)，\(p=4\)，\(2p=8\)，方程\(y^{2}=8x\)．(2)准线\(x=-3/2\)得焦点在\(x\)正半轴：\(p/2=3/2\)，\(p=3\)，\(2p=6\)，方程\(y^{2}=6x\)．验算：(1)准线\(x=-2\)合；(2)焦点\((3/2,0)\)合．}
\fi
\end{ansblock}

\begin{ansblock}[2章-拓-课时15-07]
% ans:2章-拓-课时15-07
\ansitem{8}{(a,0)（a≠0）。}
\ifansbookdetail
\ansnote{详解}{按标准形\(y^{2}=2px\)对照：\(2p=4a\)，\(p=2a\)，焦点\((p/2,0)=(a,0)\)．\(a>0\)时开口向右，焦点\((a,0)\)；\(a<0\)时开口向左，焦点仍为\((a,0)\)（即\(x\)轴负半轴上）．故焦点\((a,0)\)（\(a\neq 0\)）．验算：\(a=-1\)时\(y^{2}=-4x\)，焦点\((-1,0)\)合．}
\fi
\end{ansblock}

\begin{ansblock}[2章-拓-课时15-08]
% ans:2章-拓-课时15-08
\ansitem{9}{(1) 焦点 (0,1/8)，准线 y=−1/8；(2) 焦点 (0,1/(4a))，准线 y=−1/(4a)（a≠0）。}
\ifansbookdetail
\ansnote{详解}{(1)\(y=2x^{2}\)即\(x^{2}=(1/2)y=2\cdot (1/4)y\)：\(2p=1/2\)，\(p=1/4\)，开口向上，焦点\((0,1/8)\)，准线\(y=-1/8\)．\par\noindent (2)\(y=ax^{2}\)即\(x^{2}=(1/a)y=2\cdot [1/(2a)]y\)：\(2p=1/a\)，\(p=1/(2a)\)．\(a>0\)时开口向上，焦点\((0,1/(4a))\)，准线\(y=-1/(4a)\)；\(a<0\)时开口向下，焦点、准线表达式不变（\(1/(4a)<0\)，准线在\(x\)轴上方）．验算：\(a=2\)时与(1)一致，\(1/(4\times 2)=1/8\)合．}
\fi
\end{ansblock}

\begin{ansblock}[2章-拓-课时15-10]
% ans:2章-拓-课时15-10
\ansitem{10}{(0,−1/3)；y=1/3。}
\ifansbookdetail
\ansnote{详解}{\(3x^{2}+4y=0\)即\(x^{2}=-(4/3)y=-2\cdot (2/3)y\)：\(2p=4/3\)，\(p=2/3\)，开口向下，焦点\((0,-p/2)=(0,-1/3)\)，准线\(y=p/2=1/3\)．验算：顶点到焦点、到准线等距\(1/3\)合．}
\fi
\end{ansblock}

\begin{ansblock}[2章-拓-课时15-11]
% ans:2章-拓-课时15-11
\ansitem{11}{x²=8y}
\ifansbookdetail
\ansnote{详解}{准线\(y=-2\)在\(x\)轴下方，开口向上：设\(x^{2}=2py\)（\(p>0\)），\(-p/2=-2\)得\(p=4\)，\(2p=8\)，方程\(x^{2}=8y\)．验算：焦点\((0,4)\)、准线\(y=-2\)合．}
\fi
\end{ansblock}

\begin{ansblock}[2章-练-课时15-08]
% ans:2章-练-课时15-08
\ansitem{12}{y²=24x}
\ifansbookdetail
\ansnote{详解}{焦点在\(x\)正半轴，开口向右，设\(y^{2}=2px\)（\(p>0\)），准线\(x=-p/2\)．准线与\(y\)轴距离\(|-p/2|=6\)得\(p=12\)，\(2p=24\)，方程\(y^{2}=24x\)．验算：焦点\((6,0)\)、准线\(x=-6\)，与\(y\)轴距离均为6合．}
\fi
\end{ansblock}

\begin{ansblock}[2章-练-课时15-04]
% ans:2章-练-课时15-04
\ansitem{13}{y²=−16x}
\ifansbookdetail
\ansnote{详解}{焦点在\(x\)轴，设\(y^{2}=mx\)（\(m\neq 0\)）．代入\((-4,-8)\)：\(64=m\times (-4)\)，\(m=-16\)，方程\(y^{2}=-16x\)（开口向左，与横坐标为负相符）．}
\fi
\end{ansblock}

\begin{ansblock}[2章-练-课时15-09]
% ans:2章-练-课时15-09
\ansitem{14}{y²=8x}
\ifansbookdetail
\ansnote{详解}{设\(y^{2}=2px\)（\(p>0\)）．代入\((2,-4)\)：\(16=4p\)，\(p=4\)，\(2p=8\)，方程\(y^{2}=8x\)．验算：焦点\((2,0)\)、准线\(x=-2\)；\(M(2,-4)\)到焦点距离4、到准线距离4合．}
\fi
\end{ansblock}

\begin{ansblock}[2章-练-课时15-11]
% ans:2章-练-课时15-11
\ansitem{15}{y²=8x 或 x²=−y}
\ifansbookdetail
\ansnote{详解}{\(A(2,-4)\)在第四象限，开口向右或向下．开口向右：设\(y^{2}=2px\)（\(p>0\)），\((-4)^{2}=2p\times 2\)得\(p=4\)，方程\(y^{2}=8x\)；开口向下：设\(x^{2}=-2py\)（\(p>0\)），\(4=-2p\times (-4)\)得\(p=1/2\)，方程\(x^{2}=-y\)．综上\(y^{2}=8x\)或\(x^{2}=-y\)．验算：\(4^{2}=8\times 2\)合；\(2^{2}=-(-4)\)合．}
\fi
\end{ansblock}

\begin{ansblock}[2章-拓-课时15-12]
% ans:2章-拓-课时15-12
\ansitem{16}{D}
\ifansbookdetail
\ansnote{详解}{椭圆\(x^{2}/2+y^{2}=1\)：\(a^{2}=2\)、\(b^{2}=1\)，\(c^{2}=1\)，焦点\((\pm 1,0)\)，对称中心\(O(0,0)\)．以\(O\)为顶点、焦点为\((1,0)\)或\((-1,0)\)：\(p/2=1\)，\(2p=4\)，\(y^{2}=4x\)或\(y^{2}=-4x\)，故选D．验算：\(y^{2}=\pm 4x\)的焦点为\((\pm 1,0)\)合．}
\fi
\end{ansblock}

\begin{ansblock}[2章-拓-课时15-13]
% ans:2章-拓-课时15-13
\ansitem{17}{D}
\ifansbookdetail
\ansnote{详解}{由\((2\sqrt{2})^{2}=2px_0\)得\(x_0=4/p\)．\(M\)到准线的距离\(=x_0+p/2=4/p+p/2=3\)，得\(p^{2}-6p+8=0\)，\(p=2\)或\(p=4\)，抛物线\(y^{2}=4x\)或\(y^{2}=8x\)，故选D．验算：\(p=2\)时\(x_0=2\)，\(M(2,2\sqrt{2})\)到准线\(x=-1\)距离3合；\(p=4\)时\(x_0=1\)，\(M(1,2\sqrt{2})\)到准线\(x=-2\)距离3合．}
\fi
\end{ansblock}

\begin{ansblock}[2章-练-课时15-06]
% ans:2章-练-课时15-06
\ansitem{18}{3}
\ifansbookdetail
\ansnote{详解}{由抛物线定义，抛物线上点到焦点的距离与到准线的距离相等，故\(M\)到准线的距离为\(|MF|=3\)．}
\fi
\end{ansblock}

\begin{ansblock}[2章-练-课时15-10]
% ans:2章-练-课时15-10
\ansitem{19}{M(6,6√2) 或 M(6,−6√2)。}
\ifansbookdetail
\ansnote{详解}{\(y^{2}=12x\)：\(2p=12\)，\(p=6\)，准线\(x=-3\)．由定义\(|MF|=x_M+3=9\)，得\(x_M=6\)；代入得\(y^{2}=72\)，\(y=\pm 6\sqrt{2}\)．故\(M(6,6\sqrt{2})\)或\(M(6,-6\sqrt{2})\)（两点均在抛物线上，关于\(x\)轴对称）．验算：\(\sqrt{(6-3)^{2}+72}=\sqrt{81}=9\)合．}
\fi
\end{ansblock}

\begin{ansblock}[2章-练-课时15-12]
% ans:2章-练-课时15-12
\ansitem{20}{x=0；1/2。}
\ifansbookdetail
\ansnote{详解}{设\(P(x,y)\)：\(\sqrt{x^{2}+(y-1)^{2}}+|y+1|=3\)．①\(y>2\)或\(y<-4\)（\(|y+1|>3\)）：左边\(>3\)，无解；②\(-1\leqslant y\leqslant 2\)：\(\sqrt{x^{2}+(y-1)^{2}}=2-y\)，平方得\(x^{2}=-2y+3\)（\(-1\leqslant y\leqslant 3/2\)），对称轴\(x=0\)，此时\(|PF|=2-y\geqslant 1/2\)；\par\noindent ③\(-4\leqslant y<-1\)：\(\sqrt{x^{2}+(y-1)^{2}}=y+4\)，平方得\(x^{2}=10y+15\)（\(-3/2\leqslant y<-1\)），对称轴\(x=0\)，此时\(|PF|=y+4\geqslant 5/2\)．综上，对称轴方程\(x=0\)，\(|PF|_{\min}=1/2\)．验算：②段平方展开\(x^{2}+(y-1)^{2}=(2-y)^{2}\)得\(x^{2}=-2y+3\)合；②段与③段区间衔接经复核无缝合．}
\fi
\end{ansblock}

\begin{ansblock}[2章-拓-课时15-14]
% ans:2章-拓-课时15-14
\ansitem{21}{答案见解析：x 的系数为正数且越大时，抛物线的开口向右且开口越大。}
\ifansbookdetail
\ansnote{详解}{三条抛物线同以\(x\)轴为对称轴、顶点在原点、开口向右．作图观察（或取同一\(y\)值比较横坐标\(x=y^{2}/(2p)\)）：系数1、2、4递增时，同一高度处横坐标\(y^{2}/(2p)\)递减，曲线横向收拢更快、开口更大．故：\(x\)的系数为正且越大，开口向右且开口越大．}
\fi
\end{ansblock}

\begin{ansblock}[2章-练-课时15-13]
% ans:2章-练-课时15-13
\ansitem{22}{B}
\ifansbookdetail
\ansnote{详解}{以顶端\(O\)为原点建系，设\(x^{2}=-2py\)（\(p>0\)）．由对称性\(D(15,h)\)（\(h<0\)），\(B(30,h-150)\)．联立：\(225=-2ph\)，\(900=-2p(h-150)\)．两式相除：\(4=(h-150)/h\)，得\(h=-50\)；代入\(225=-2p\times (-50)\)得\(p=2.25\)．顶端\(O\)到\(AB\)的距离\(=|h|+150=200\)(m)，故选B．验算：\(p=2.25\)，\(2p=4.5\)，\(900=4.5\times 200\)合．}
\fi
\end{ansblock}

\begin{ansblock}[2章-练-课时15-16]
% ans:2章-练-课时15-16
\ansitem{23}{不能安全通过隧道。}
\ifansbookdetail
\ansnote{详解}{以抛物线顶点为原点、对称轴为\(y\)轴建系．由图：拱高3m、半宽3m，抛物线过\(A(3,-3)\)；矩形部分高2m、总高（顶点至地面）5m、地面\(y=-5\)．设\(x^{2}=-2py\)（\(p>0\)）：\(9=-2p\times (-3)\)，\(2p=3\)，方程\(x^{2}=-3y\)．\par\noindent 车居中：\(x=\pm 1.5\)处拱高\(y=-(1/3)\times 2.25=-0.75\)，净空\(=5-0.75=4.25\)(m)．\(4.25<4.5\)（车与集装箱总高），即使集装箱处于隧道正中，总高仍高于该处拱高，故不能安全通过．验算：\((\pm 1.5)^{2}=-3y\)得\(y=-0.75\)合；\(2p=3\)时准线\(y=3/4\)合．}
\fi
\end{ansblock}

\begin{ansblock}[2章-拓-课时15-01]
% ans:2章-拓-课时15-01
\ansitem{24}{桥拱所在抛物线 y=−(1/80)x²；溢流孔所在抛物线分别为 y=−(5/36)(x+14)²、y=−(5/36)(x+7)²、y=−(5/36)(x−7)²、y=−(5/36)(x−14)²；交点 A((140/13),−(245/169))、B(10,−5/4)、C((70/13),−(245/676))。}
\ifansbookdetail
\ansnote{详解}{设桥拱\(y=ax^{2}\)．由图注，桥拱过点\((20,-5)\)：\(-5=a\times 20^{2}\)，\(a=-1/80\)，桥拱\(y=-(1/80)x^{2}\)．①\par\noindent 设\(A\)所在溢流孔\(y=a_1(x-14)^{2}\)（顶点\((14,0)\)），亦过\((20,-5)\)：\(-5=a_1\times 36\)，\(a_1=-5/36\)．4孔轮廓相同、顶点\((\pm 7,0)\)、\((\pm 14,0)\)：\(y=-(5/36)(x\mp 7)^{2}\)、\(y=-(5/36)(x\mp 14)^{2}\)．②③\par\noindent 联立①②：\(400(x-14)^{2}=36x^{2}\)（两边乘2880），即\(91x^{2}-2800x+19600=0\)，\(\Delta=705600\)，\(\sqrt{\Delta}=840\)，\(x=20\)（\(y=-5\)）或\(x=140/13\)（\(y=-245/169\)），取交点\(A(140/13,-245/169)\)；联立①③：\(400(x-7)^{2}=36x^{2}\)，即\(91x^{2}-1400x+4900=0\)，\(\Delta=176400\)，\(\sqrt{\Delta}=420\)，\(x=10\)（\(y=-5/4\)）或\(x=70/13\)（\(y=-245/676\)），得\(C(70/13,-245/676)\)、\(B(10,-5/4)\)．验算：\(x=140/13\)：\(y=-(1/80)(140/13)^{2}=-19600/(80\times 169)=-245/169\)合；\(x=10\)：\(-100/80=-5/4\)合；\(x=70/13\)：\(-4900/(80\times 169)=-245/676\)合．}
\fi
\end{ansblock}

\begin{ansblock}[2章-拓-课时15-09]
% ans:2章-拓-课时15-09
\ansitem{25}{(1) y=−2x+6；(2) 6m。}
\ifansbookdetail
\ansnote{详解}{(1)\(|BF|=p/2=1/2\)，灯柱\(CD\)过焦点\(F(1/2,0)\)且\(A\)到灯柱距离1.5，得\(x_A=1/2+3/2=2\)；代入\(y^{2}=2x\)得\(y_A=2\)（取上半支），\(A(2,2)\)．设\(A\)处切线\(y-2=k(x-2)\)：与\(y^{2}=2x\)联立消\(x\)得\(ky^{2}-2y+4-4k=0\)，相切\(\Delta=4-4k(4-4k)=0\)，即\(16k^{2}-16k+4=0\)，\(k=1/2\)．灯罩轴线与切线垂直：斜率\(-2\)，方程\(y-2=-2(x-2)\)，即\(y=-2x+6\)．\par\noindent (2)路宽\(|DH|=10\)，灯罩轴线过路面中线，轴线与路面交点到\(y\)轴距离\(1/2+5=11/2\)．对\(y=-2x+6\)，\(x=11/2\)时\(y=-5\)，交点\((11/2,-5)\)，\(|FD|=5\)；又\(x=1/2\)代入\(y^{2}=2x\)得\(y=\pm 1\)，\(|CF|=1\)．灯柱高\(|CD|=|DF|+|FC|=6\)(m)．验算：\(16k^{2}-16k+4=4(2k-1)^{2}=0\)合；\(|CD|=6\)合．}
\fi
\end{ansblock}

\begin{ansblock}[2章-拓-课时15-03]
% ans:2章-拓-课时15-03
\ansitem{26}{BCD}
\ifansbookdetail
\ansnote{详解}{A：\(D_1D\perp\)底面，\(DN=\sqrt{MN^{2}-MD^{2}}=\sqrt{4-1}=\sqrt{3}\)，\(N\)的轨迹是以\(D\)为圆心、\(\sqrt{3}\)为半径的圆．设\(MN\)中点\(P\)、\(MD\)中点\(Q\)，则\(PQ\parallel ND\)且\(PQ=\sqrt{3}/2\)，\(P\)的轨迹是以\(Q\)为圆心\(\sqrt{3}/2\)为半径的圆，面积\(\pi\times 3/4=3\pi/4\neq \pi\)，A错；\par\noindent B：\(N\)到\(B_1B\)的距离\(=NB\)（\(B_1B\perp\)底面），即\(N\)到定点\(B\)与定直线\(DC\)距离相等（\(B\)不在\(DC\)上），轨迹为抛物线，B对；\par\noindent C：\(AB\parallel D_1C_1\)，\(D_1N\)与\(D_1C_1\)所成角\(\pi/3\)，即\(D_1N\)为以\(D_1C_1\)为轴、轴截面顶角\(2\pi/3\)的圆锥母线，底面\(ABCD\)与轴\(D_1C_1\)平行，交线为双曲线，C对；\par\noindent D：\(D_1N\)为以\(D_1B\)为轴、顶角\(\pi/3\)的圆锥母线，平面\(ABCD\)与轴\(D_1B\)斜交，交线为椭圆，D对．故选BCD．}
\fi
\end{ansblock}

\begin{ansblock}[2章-拓-课时15-04]
% ans:2章-拓-课时15-04
\ansitem{27}{BCD}
\ifansbookdetail
\ansnote{详解}{以\(D\)为原点、\(DA\)、\(DC\)、\(DD_1\)为\(x\)、\(y\)、\(z\)轴建系：\(A(2,0,0)\)，\(C(0,2,0)\)，\(B(2,2,0)\)，\(M(1,0,2)\)．\par\noindent A：\(\overrightarrow{AC}=(-2,2,0)\)，\(\overrightarrow{BM}=(-1,-2,2)\)，\(|\cos\langle\overrightarrow{AC},\overrightarrow{BM}\rangle|=|\overrightarrow{AC}\cdot\overrightarrow{BM}|/(|\overrightarrow{AC}||\overrightarrow{BM}|)=|2-4+0|/(2\sqrt{2}\times 3)=\sqrt{2}/6\neq\sqrt{2}/3\)，A错；\par\noindent B：\(S=(1/2)|AC_1|\cdot h=\sqrt{3}\)，\(|AC_1|=2\sqrt{3}\)，得\(h=1\)：\(P\)在以\(AC_1\)为轴、半径1的圆柱面上，该圆柱面与侧面\(BCC_1B_1\)（与轴斜交）的交线为椭圆的一部分，B对；\par\noindent C：\(P\)在侧面\(BCC_1B_1\)上时，到\(C_1D_1\)的距离\(=|PC_1|\)，条件即\(|PC_1|\)等于\(P\)到直线\(BC\)的距离——到定点与定直线等距且\(C_1\)不在\(BC\)上，轨迹为抛物线的一部分，C对；\par\noindent D：取\(AD\)中点\(H\)，\(MH=2\)且\(MH\perp\)底面；当交线与\(AC\)平行（\(HN\perp AC\)）时二面角最小，截面为过\(M\)、\(N\)及四棱中点的正六边形，边长\(\sqrt{2}\)，面积\(6\times(\sqrt{3}/4)\times 2=3\sqrt{3}\)，D对．故选BCD．}
\fi
\end{ansblock}

\begin{ansblock}[2章-拓-课时15-05]
% ans:2章-拓-课时15-05
\ansitem{28}{C}
\ifansbookdetail
\ansnote{详解}{折叠不变量：\(PB=AB\)，\(PD=AD\)，故\(PB+PD=10>BD=8\)，\(P\)在以\(B\)、\(D\)为焦点、长轴10、焦距8的椭圆上（不在长轴端点）．\par\noindent \(\triangle BPD\cong\triangle BCD\)（三边相等），得\(\angle PBD=\angle CBD\)．过\(P\)作\(PE\perp BD\)于\(E\)，连\(CE\)：\(\triangle BPE\cong\triangle BCE\)，得\(CE\perp BD\)且\(PE=CE\)，故\(BD\perp\)平面\(PCE\)．\(V_{P-BCD}=V_{B-PCE}+V_{D-PCE}=(1/3)S_{\triangle PCE}\cdot BD=(8/3)S_{\triangle PCE}\)．取\(PC\)中点\(F\)，\(EF\perp PC\)（\(PE=CE\)）：\(S_{\triangle PCE}=(1/2)\cdot PC\cdot EF=\sqrt{PE^{2}-1}\)（\(PC=2\)得\(CF=1\)）．\(PE\)的最大值为椭圆短半轴\(\sqrt{5^{2}-4^{2}}=3\)，故\(S_{\triangle PCE}\)最大\(=\sqrt{9-1}=2\sqrt{2}\)．\(V_{\max}=(8/3)\times 2\sqrt{2}=16\sqrt{2}/3\)，故选C．验算：\(EF=\sqrt{PE^{2}-CF^{2}}=\sqrt{PE^{2}-1}\)合；\(PE=3\)时\(EF=2\sqrt{2}\)，\(S=(1/2)\times 2\times 2\sqrt{2}=2\sqrt{2}\)合．}
\fi
\end{ansblock}

\begin{ansblock}[2章-拓-课时15-06]
% ans:2章-拓-课时15-06
\ansitem{29}{BCD}
\ifansbookdetail
\ansnote{详解}{过\(P\)作\(PO\perp\)平面\(ABCD\)于\(O\)，\(OH\perp AD\)于\(H\)．\(AD\perp\)平面\(POH\)，故\(AD\perp PH\)，\(\angle PHO=\alpha\)；又\(\angle PBO=\beta\)，\(\alpha=\beta\)得\(OH=OB\)．\par\noindent \(O\)的轨迹：底面内到定点\(B\)与定直线\(AD\)等距的点集，即以\(B\)为焦点、\(AD\)为准线的抛物线在四边形\(ABCD\)内（含边界）的部分；\(P\)的轨迹为其上平移（以\(B_1\)为焦点、\(A_1D_1\)为准线），是抛物线的一部分而非整条，A错；\par\noindent B：\(|PB|^{2}=(2\sqrt{2})^{2}+|OB|^{2}\)，\(O\)为\(AB\)中点时\(|OB|_{\min}=1\)，故\(|PB|_{\min}=\sqrt{8+1}=3\)，B对；\par\noindent C：\(PA_1\)与\(CD\)所成角即\(OA\)与\(AB\)所成角\(\angle OAB\)．以\(AB\)中点\(M\)为原点建系，抛物线\(y^{2}=4x\)，\(A(-1,0)\)．设切线\(x=ty-1\)：代入得\(y^{2}-4ty+4=0\)，\(\Delta=16t^{2}-16=0\)，\(t=\pm 1\)，取\(t=1\)（\(O\)在四边形内），\(\angle OAB=\pi/4\)，最大值\(\pi/4\)，C对；\par\noindent D：\(V_{P-A_1BC_1}=V_{B-PA_1C_1}=(1/3)S_{\triangle PA_1C_1}\cdot BB_1\)，\(A_1C_1=2\sqrt{2}\)．\(P\)到\(A_1C_1\)距离最大\(\Leftrightarrow O\)到\(AC\)距离最大：\(O\)为\(AB\)中点时最大，值\((1/4)|BD|=\sqrt{2}/2\)．\(V_{\max}=(2\sqrt{2}/3)\times (1/2)\times 2\sqrt{2}\times (\sqrt{2}/2)=2\sqrt{2}/3\)，D对．故选BCD．验算：\(y^{2}=4x\)（焦点\(B(1,0)\)，准线\(x=-1\)即\(AD\)所在直线）：\(|OB|\)与\(|OH|\)分别为到焦点、准线的距离，等距即轨迹合；切线\(y=x+1\)过\(A(-1,0)\)，斜率1，\(\angle OAB=45^{\circ}\)合．}
\fi
\end{ansblock}

\begin{ansblock}[2章-拓-课时15-02]
% ans:2章-拓-课时15-02
\ansitem{30}{A}
\ifansbookdetail
\ansnote{详解}{\(y^{2}=8x\)：\(p=4\)，准线\(x=-2\)，焦点\(F(2,0)\)．\(P\)在准线\(x=-2\)上，又在\(x-y+6=0\)上：\(P(-2,4)\)．\(k_{PF}=(0-4)/(2+2)=-1\)；由特征(3)\(PF\perp AB\)，得\(k_{AB}=1\)．\(AB\)过\(F(2,0)\)：\(y=x-2\)，即\(x-y-2=0\)，故选A．验算：\(P(-2,4)\)代入\(x-y+6=0\)：\(-2-4+6=0\)合．}
\fi
\end{ansblock}

\begin{ansblock}[2章-练-课时15-15]
% ans:2章-练-课时15-15
\ansitem{31}{(1) y²=4x；(2) 存在，此时 Q(9,6) 或 Q(4,−4)。}
\ifansbookdetail
\ansnote{详解}{(1)\(16=8p\)得\(p=2\)，\(C\)：\(y^{2}=4x\)．\par\noindent (2)\(F(1,0)\)．设\(N(t,2t+6)\)、\(Q(x_0,y_0)\)．四边形\(PQFN\)为平行四边形\(\Leftrightarrow\overrightarrow{PQ}=\overrightarrow{NF}\Leftrightarrow Q-P=F-N\)：\((4-t,-2t-2)=(x_0-1,y_0)\)，得\(x_0=5-t\)，\(y_0=-2t-2\)．\(Q\)在\(C\)上：\((-2t-2)^{2}=4(5-t)\)，即\(4t^{2}+8t+4=20-4t\)，\(t^{2}+3t-4=0\)，\(t=1\)或\(t=-4\)．\(t=1\)：\(N(1,8)\)，\(Q(4,-4)\)；\(t=-4\)：\(N(-4,-2)\)，\(Q(9,6)\)．\par\noindent 非退化核验（\(P\)、\(Q\)、\(F\)、\(N\)不共线，\(N\neq F\)）：\(t=1\)时\(\overrightarrow{PQ}=(0,-8)=\overrightarrow{NF}\)合；\(\overrightarrow{QF}=(-3,4)\)，\(\overrightarrow{PN}=\overrightarrow{P}-\overrightarrow{N}=(3,-4)\)，平行且相等合；四点不共线合，\(N(1,8)\neq F(1,0)\)合；\par\noindent \(t=-4\)时\(\overrightarrow{PQ}=(5,2)\)，\(\overrightarrow{NF}=\overrightarrow{F}-\overrightarrow{N}=(5,2)\)合；\(\overrightarrow{QF}=(-8,-6)\)，\(\overrightarrow{PN}=(8,6)\)合；斜率\(k_{PQ}=2/5\neq k_{QF}=3/4\)，四点不共线合，\(N\neq F\)合．两解均成立：存在\(N\)使\(PQFN\)为平行四边形，\(Q(9,6)\)或\(Q(4,-4)\)．验算：\(t^{2}+3t-4=(t-1)(t+4)\)合；\(Q(9,6)\)：\(36=4\times 9\)合；\(Q(4,-4)\)：\(16=16\)合．}
\fi
\end{ansblock}

\begin{ansblock}[2章-导-课时15-G1]
% ans:2章-导-课时15-G1
\ansitem{32}{B}
\ifansbookdetail
\ansnote{详解}{设\(M(x_0,y_0)\)．到对称轴的距离\(|y_0|=2\sqrt{2}\)，故\(y_0^{2}=8\)；到准线\(x=-p/2\)的距离\(x_0+p/2=3\)，得\(x_0=3-p/2\)．代入\(y_0^{2}=2px_0\)：\(8=2p(3-p/2)=6p-p^{2}\)，即\(p^{2}-6p+8=0\)，解得\(p=2\)或\(p=4\)．两值均有\(x_0=3-p/2\geqslant 0\)（\(p=2\)时\(x_0=2\)；\(p=4\)时\(x_0=1\)），皆合题意，故选B．验算：\(p=2\)时准线\(x=-1\)，\(M(2,\pm 2\sqrt{2})\)，\(x_0+1=3\)合；\(p=4\)时准线\(x=-2\)，\(M(1,\pm 2\sqrt{2})\)，\(x_0+2=3\)合．}
\fi
\end{ansblock}

\begin{ansblock}[2章-导-课时15-G2]
% ans:2章-导-课时15-G2
\ansitem{33}{B}
\ifansbookdetail
\ansnote{详解}{\(y=-(1/8)x^{2}\)即\(x^{2}=-8y\)，\(2p=8\)、\(p=4\)，开口向下，准线\(y=p/2=2\)，故选B．验算：焦点\((0,-2)\)、准线\(y=2\)，顶点\((0,0)\)到两者等距2合．}
\fi
\end{ansblock}

\begin{ansblock}[2章-导-课时15-G3]
% ans:2章-导-课时15-G3
\ansitem{34}{D}
\ifansbookdetail
\ansnote{详解}{两边同除以5：\(\sqrt{(x-1)^{2}+(y-2)^{2}}=|3x-4y+12|/5\)，而\(|3x-4y+12|/\sqrt{3^{2}+4^{2}}\)恰为\(M\)到直线\(3x-4y+12=0\)的距离．故动点\(M\)到定点\((1,2)\)的距离等于它到定直线\(3x-4y+12=0\)的距离；又\(3\times 1-4\times 2+12=7\neq 0\)，定点不在定直线上，由抛物线的定义，轨迹是以\((1,2)\)为焦点、\(3x-4y+12=0\)为准线的抛物线，故选D．说明：系数5的作用仅是把右端配成点到直线距离的标准形式（\(5=\sqrt{3^{2}+(-4)^{2}}\)），并非离心率记号．}
\fi
\end{ansblock}

\begin{ansblock}[2章-导-课时15-G4]
% ans:2章-导-课时15-G4
\ansitem{35}{y²=16x 或 x²=−8y}
\ifansbookdetail
\ansnote{详解}{直线\(x-2y-4=0\)与坐标轴的交点：令\(y=0\)得\(A(4,0)\)；令\(x=0\)得\(B(0,-2)\)．①焦点为\(A(4,0)\)：开口向右，设\(y^{2}=2px\)（\(p>0\)），\(p/2=4\)得\(p=8\)，方程\(y^{2}=16x\)；②焦点为\(B(0,-2)\)：开口向下，设\(x^{2}=-2py\)（\(p>0\)），\(p/2=2\)得\(p=4\)，方程\(x^{2}=-8y\)．综上，所求抛物线的方程为\(y^{2}=16x\)或\(x^{2}=-8y\)．验算：\(y^{2}=16x\)的焦点\((4,0)\)满足\(4-0-4=0\)，在直线上合；\(x^{2}=-8y\)的焦点\((0,-2)\)满足\(0+4-4=0\)，在直线上合．}
\fi
\end{ansblock}

\begin{ansblock}[2章-导-课时15-G5]
% ans:2章-导-课时15-G5
\ansitem{36}{y²=−8x；m=±2√6}
\ifansbookdetail
\ansnote{详解}{\(M(-3,m)\)的横坐标为负，抛物线开口向左，设\(y^{2}=-2px\)（\(p>0\)），则焦点\(F(-p/2,0)\)，准线\(x=p/2\)．由抛物线的定义，\(|MF|\)等于\(M\)到准线的距离：\(p/2+3=5\)，得\(p=4\)，故标准方程为\(y^{2}=-8x\)．将\(M(-3,m)\)代入：\(m^{2}=24\)，\(m=\pm 2\sqrt{6}\)．验算：焦点\(F(-2,0)\)：\(|MF|=\sqrt{(-3+2)^{2}+m^{2}}=\sqrt{1+24}=5\)合（两点距与定义两法互证）．}
\fi
\end{ansblock}

\jietitle{2.7.2 抛物线的几何性质}

\begin{ansblock}[2章-练-课时16-04]
% ans:2章-练-课时16-04
\ansitem{1}{1}
\ifansbookdetail
\ansnote{详解}{\(2p=8\)、\(p=4\)，焦点 \(F(2,0)\)．由点到直线的距离公式，\(d=|2-\sqrt{3}\times 0|/\sqrt{1^{2}+(\sqrt{3})^{2}}=2/\sqrt{4}=1\)．}
\fi
\end{ansblock}

\begin{ansblock}[2章-练-课时16-03]
% ans:2章-练-课时16-03
\ansitem{2}{(1) 焦点(3/2,0)，准线 x=−3/2；(2) 焦点(0,−5/2)，准线 y=5/2；(3) 焦点(0,2/3)，准线 y=−2/3；(4) 焦点(−a/4,0)，准线 x=a/4}
\ifansbookdetail
\ansnote{详解}{(1) \(y^{2}=6x\)：\(2p=6\)、\(p=3\)，开口向右，焦点 \((p/2,0)=(3/2,0)\)，准线 \(x=-3/2\)．(2) \(x^{2}=-10y\)：\(2p=10\)、\(p=5\)，开口向下，焦点 \((0,-5/2)\)，准线 \(y=5/2\)．(3) \(x^{2}=(8/3)y\)：\(2p=8/3\)、\(p=4/3\)，开口向上，焦点 \((0,2/3)\)，准线 \(y=-2/3\)．(4) \(y^{2}=-ax\)（\(a\neq 0\)）：一次项系数为 \(-a\)，无论 \(a>0\)（开口向左）或 \(a<0\)（开口向右），焦点均为 \((-a/4,0)\)、准线均为 \(x=a/4\)（\(a>0\) 时焦点在 \(x\) 轴负半轴，\(a<0\) 时在正半轴，公式统一）．}
\fi
\end{ansblock}

\begin{ansblock}[2章-练-课时16-01]
% ans:2章-练-课时16-01
\ansitem{3}{AD}
\ifansbookdetail
\ansnote{详解}{对照标准形 \(y^{2}=-2px\)（\(p>0\)）：\(2p=2\)、\(p=1\)，一次项系数为负，开口向左，A 正确；焦点为 \((-p/2,0)=(-1/2,0)\)，B 所给 \((-1,0)\) 错（把 \(p\) 误作 \(p/2\)）；准线为 \(x=p/2=1/2\)，C 所给 \(x=1\) 错；对称轴为 \(x\) 轴，D 正确．故选：AD．}
\fi
\end{ansblock}

\begin{ansblock}[2章-练-课时16-05]
% ans:2章-练-课时16-05
\ansitem{4}{6}
\ifansbookdetail
\ansnote{详解}{\(y^{2}=8x\) 中 \(2p=8\)、\(p=4\)，准线 \(x=-2\)．点 \(P\) 到 \(y\) 轴距离为 4，而抛物线上点横坐标 \(x\geqslant 0\)，故 \(x_{P}=4\)．由定义，\(P\) 到焦点的距离等于 \(P\) 到准线的距离：\(|PF|=x_{P}+p/2=4+2=6\)．（校验：\(P(4,\pm 4\sqrt{2})\)，\(|PF|=\sqrt{(4-2)^{2}+32}=\sqrt{36}=6\) \ding{51}．）}
\fi
\end{ansblock}

\begin{ansblock}[2章-练-课时16-02]
% ans:2章-练-课时16-02
\ansitem{5}{17/16}
\ifansbookdetail
\ansnote{详解}{化标准形：\(y=4x^{2}\) 即 \(x^{2}=(1/4)y\)，故 \(2p=1/4\)、\(p=1/8\)，焦点 \(F(0,1/16)\)，准线 \(y=-1/16\)．由抛物线定义，点 \(M\) 到焦点的距离等于它到准线的距离，\(|MF|=y_{M}+p/2=1+1/16=17/16\)．（回代校验：\(x_{M}^{2}=1/4\)、\(M(\pm 1/2,1)\)，\(|MF|^{2}=x_{M}^{2}+(y_{M}-1/16)^{2}=1/4+225/256=289/256\)，\(|MF|=17/16\) \ding{51}．）}
\fi
\end{ansblock}

\begin{ansblock}[2章-练-课时16-06]
% ans:2章-练-课时16-06
\ansitem{6}{M(3p/2, ±√3 p)}
\ifansbookdetail
\ansnote{详解}{设 \(M(x_{0},y_{0})\)，\(x_{0}\geqslant 0\)．由定义，\(|MF|\) 等于 \(M\) 到准线 \(x=-p/2\) 的距离，即 \(x_{0}+p/2=2p\)，得 \(x_{0}=3p/2\)．代入方程：\(y_{0}^{2}=2p\cdot(3p/2)=3p^{2}\)，\(y_{0}=\pm\sqrt{3}p\)．故 \(M(3p/2,\sqrt{3}p)\) 或 \(M(3p/2,-\sqrt{3}p)\)．}
\fi
\end{ansblock}

\begin{ansblock}[2章-练-课时16-07]
% ans:2章-练-课时16-07
\ansitem{7}{y²=−4x 或 x²=√2 y}
\ifansbookdetail
\ansnote{详解}{\(P(-2,2\sqrt{2})\) 在第二象限，开口只可能向左或向上．情形一，对称轴为 \(x\) 轴：设 \(y^{2}=-2px\)（\(p>0\)），代入得 \((2\sqrt{2})^{2}=-2p\cdot(-2)=4p\)、\(p=2\)，即 \(y^{2}=-4x\)．情形二，对称轴为 \(y\) 轴：设 \(x^{2}=2py\)（\(p>0\)），代入得 \(4=2p\cdot 2\sqrt{2}\)、\(2p=\sqrt{2}\)，即 \(x^{2}=\sqrt{2}y\)．（回验：\(8=-4\times(-2)\) \ding{51}；\(4=\sqrt{2}\times 2\sqrt{2}\) \ding{51}．）}
\fi
\end{ansblock}

\begin{ansblock}[2章-练-课时16-08]
% ans:2章-练-课时16-08
\ansitem{8}{p=4}
\ifansbookdetail
\ansnote{详解}{椭圆中 \(c^{2}=6-2=4\)，右焦点 \((2,0)\)．抛物线焦点 \((p/2,0)\)，由 \(p/2=2\) 得 \(p=4\)．}
\fi
\end{ansblock}

\begin{ansblock}[2章-练-课时16-10]
% ans:2章-练-课时16-10
\ansitem{9}{C}
\ifansbookdetail
\ansnote{详解}{顶点在原点、对称轴为 \(x\) 轴，方程形如 \(y^{2}=\pm 2px\)（\(p>0\)）．通径（过焦点且垂直于对称轴的弦）长为 \(2p\)，由 \(2p=8\) 得 \(p=4\)，故方程为 \(y^{2}=8x\)（焦点在 \(x\) 轴正半轴）或 \(y^{2}=-8x\)（焦点在负半轴），两向各一，选 C．（排除 D：其对称轴为 \(y\) 轴．）}
\fi
\end{ansblock}

\begin{ansblock}[2章-练-课时16-09]
% ans:2章-练-课时16-09
\ansitem{10}{A（1/3）}
\ifansbookdetail
\ansnote{详解}{\(y^{2}=x\) 中 \(2p=1\)，焦点 \(F(1/4,0)\)，准线 \(x=-1/4\)．直线 \(l\)：\(y=\sqrt{3}(x-1/4)\)．代入 \(y^{2}=x\) 得 \(3(x-1/4)^{2}=x\)，整理为 \(48x^{2}-40x+3=0\)，故 \(x_{1}+x_{2}=40/48=5/6\)，\(x_{1}x_{2}=3/48=1/16\)．由定义，\(|AF|=x_{1}+1/4\)、\(|BF|=x_{2}+1/4\)，\(|AF|\cdot|BF|=x_{1}x_{2}+(1/4)(x_{1}+x_{2})+1/16=1/16+5/24+1/16=16/48=1/3\)．故选 A．}
\fi
\end{ansblock}

\begin{ansblock}[2章-练-课时16-11]
% ans:2章-练-课时16-11
\ansitem{11}{B}
\ifansbookdetail
\ansnote{详解}{\(y^{2}=4x\) 的焦点 \(F(1,0)\)，准线 \(x=-1\)，即直线 \(x+1=0\)．分别过 \(A\)、\(B\)、\(M\) 作准线的垂线，垂足为 \(C\)、\(D\)、\(N\)．由定义，\(|AC|=|AF|\)、\(|BD|=|BF|\)，故 \(|AC|+|BD|=|AF|+|BF|=|AB|=8\)．又 \(AC\parallel MN\parallel BD\) 且 \(M\) 为 \(AB\) 中点，\(MN\) 为直角梯形 \(ABDC\) 的中位线，\(|MN|=(|AC|+|BD|)/2=4\)，即点 \(M\) 到准线 \(x=-1\) 的距离为 4．故选：B．}
\fi
\end{ansblock}

\begin{ansblock}[2章-练-课时16-12]
% ans:2章-练-课时16-12
\ansitem{12}{P(−2,1/2)（最小值 6）}
\ifansbookdetail
\ansnote{详解}{\(x^{2}=8y\) 的准线为 \(l\)：\(y=-2\)．设 \(P\) 在 \(l\) 上的射影为 \(Q\)，过 \(A\) 作 \(AB\perp l\) 于 \(B\)，则 \(|AB|=y_{A}-(-2)=4+2=6\)．由定义 \(|PF|=|PQ|\)，故 \(|PF|+|PA|=|PQ|+|PA|\geqslant|AQ|\geqslant|AB|=6\)：第一个不等号当 \(A\)、\(P\)、\(Q\) 共线（\(P\) 与 \(A\) 同取横坐标 \(-2\)）时取等，第二个不等号当 \(Q\) 与 \(B\) 重合时取等，两条件同时满足．把 \(x=-2\) 代入 \(x^{2}=8y\) 得 \(y=1/2\)，故所求点 \(P(-2,1/2)\)，此时 \((|PF|+|PA|)_{\min}=|AB|=y_{A}+2=6\)．}
\fi
\end{ansblock}

\begin{ansblock}[2章-练-课时16-13]
% ans:2章-练-课时16-13
\ansitem{13}{|MA| 的最小值为 4，此时 M(0,0)}
\ifansbookdetail
\ansnote{详解}{设 \(M(x,y)\)，则 \(x\geqslant 0\)，\(y^{2}=16x\)．\par\noindent \(|MA|^{2}=(x-4)^{2}+y^{2}=(x-4)^{2}+16x=x^{2}+8x+16=(x+4)^{2}\)．\par\noindent 由 \(x\geqslant 0\) 得 \((x+4)^{2}\geqslant 16\)，当且仅当 \(x=0\) 时等号成立，此时 \(y=0\)，\(|MA|_{\min}=4\)，\(M(0,0)\)．（注：\(2p=16\)、\(p/2=4\)，\(A(4,0)\) 恰为该抛物线的焦点，\(|MA|=|MF|=x+4\geqslant 4\) 与配方路径互证．）}
\fi
\end{ansblock}

\begin{ansblock}[2章-练-课时16-14]
% ans:2章-练-课时16-14
\ansitem{14}{4√3 p}
\ifansbookdetail
\ansnote{详解}{设正三角形 \(OAB\) 的顶点 \(A(x_{1},y_{1})\)、\(B(x_{2},y_{2})\) 在抛物线上，则 \(y_{1}^{2}=2px_{1}\)、\(y_{2}^{2}=2px_{2}\)．由 \(|OA|=|OB|\) 得 \(x_{1}^{2}+y_{1}^{2}=x_{2}^{2}+y_{2}^{2}\)，即 \((x_{1}^{2}-x_{2}^{2})+2p(x_{1}-x_{2})=0\)，\((x_{1}-x_{2})(x_{1}+x_{2}+2p)=0\)．因 \(x_{1}\geqslant 0\)、\(x_{2}\geqslant 0\) 且 \(2p>0\)（若 \(x_{1}\)、\(x_{2}\) 之一为 0 则三角形退化），括号第二因子为正，故 \(x_{1}=x_{2}\)，从而 \(y_{2}=-y_{1}\neq 0\)，线段 \(AB\) 关于 \(x\) 轴对称．\(x\) 轴垂直平分 \(AB\) 且平分 \(\angle AOB\)，\(\angle AOx=30^{\circ}\)，故 \(y_{1}/x_{1}=\tan 30^{\circ}=\sqrt{3}/3\)、\(x_{1}=\sqrt{3}y_{1}\)．代入 \(y_{1}^{2}=2p\cdot\sqrt{3}y_{1}\) 得 \(y_{1}=2\sqrt{3}p\)，边长 \(|AB|=2y_{1}=4\sqrt{3}p\)．}
\fi
\end{ansblock}

\begin{ansblock}[2章-练-课时16-15]
% ans:2章-练-课时16-15
\ansitem{15}{(1) |AB|=4p；(2) 定值 cos∠AOB=−3√41/41（与 p 无关）}
\ifansbookdetail
\ansnote{详解}{(1) 过 \(F(p/2,0)\)、倾斜角 \(\pi/4\) 的直线为 \(y=x-p/2\)．代入 \(y^{2}=2px\)：\((x-p/2)^{2}=2px\)，即 \(x^{2}-3px+p^{2}/4=0\)（\(\Delta=9p^{2}-p^{2}=8p^{2}>0\)）．设 \(A(x_{A},y_{A})\)、\(B(x_{B},y_{B})\)，则 \(x_{A}+x_{B}=3p\)，\(x_{A}x_{B}=p^{2}/4\)．由定义（焦半径），\par\noindent \(|AB|=|AF|+|BF|=(x_{A}+p/2)+(x_{B}+p/2)=x_{A}+x_{B}+p=4p\)．\par\noindent (2) 该弦所在直线与 \(x\) 轴交于 \(F\)，两端点 \(A\)、\(B\) 分居 \(x\) 轴两侧，纵坐标一正一负，其积为负：\par\noindent \(y_{A}y_{B}=(x_{A}-p/2)(x_{B}-p/2)=x_{A}x_{B}-(p/2)(x_{A}+x_{B})+p^{2}/4=p^{2}/4-3p^{2}/2+p^{2}/4=-p^{2}\)（与焦点弦结论 \(y_{1}y_{2}=-p^{2}\) 相符，此即符号判据）．\par\noindent 于是 \(x_{A}x_{B}+y_{A}y_{B}=p^{2}/4-p^{2}=-3p^{2}/4\)．又 \(|OA|^{2}=x_{A}^{2}+y_{A}^{2}=x_{A}^{2}+2px_{A}=x_{A}(x_{A}+2p)\)，同理 \(|OB|^{2}=x_{B}(x_{B}+2p)\)．\par\noindent 故 \(|OA|\cdot|OB|=\sqrt{x_{A}x_{B}}\cdot\sqrt{x_{A}x_{B}+2p(x_{A}+x_{B})+4p^{2}}=(p/2)\cdot\sqrt{p^{2}/4+6p^{2}+4p^{2}}=(p/2)\cdot(\sqrt{41}/2)p=\sqrt{41}p^{2}/4\)．\par\noindent 在 \(\triangle AOB\) 中由余弦定理，\(\cos\angle AOB=(|OA|^{2}+|OB|^{2}-|AB|^{2})/(2|OA|\cdot|OB|)\)．\par\noindent 而 \(|OA|^{2}+|OB|^{2}-|AB|^{2}=(x_{A}-x_{B})^{2}+(y_{A}-y_{B})^{2}=2(x_{A}x_{B}+y_{A}y_{B})=-3p^{2}/2\)，\par\noindent 故 \(\cos\angle AOB=(-3p^{2}/4)/(\sqrt{41}p^{2}/4)=-3/\sqrt{41}\)，分母有理化得 \(-3\sqrt{41}/41\)，\(p^{2}\) 全部约去，与 \(p\) 无关．}
\fi
\end{ansblock}

\begin{ansblock}[2章-练-课时16-16]
% ans:2章-练-课时16-16
\ansitem{16}{24}
\ifansbookdetail
\ansnote{详解}{\(x^{2}=4y\) 的焦点 \(F(0,1)\)，准线 \(y=-1\)，抛物线上点之焦半径 \(|AF|=y_{A}+1\)（到准线距离，焦半径公式见例 36）．\(l_{1}\)：\(y=k_{1}x+1\) 代入 \(x^{2}=4y\) 得 \(x^{2}-4k_{1}x-4=0\)（\(\Delta=16k_{1}^{2}+16>0\) 恒成立）．设 \(A(x_{1},y_{1})\)、\(B(x_{2},y_{2})\)，则 \(x_{1}+x_{2}=4k_{1}\)，\(y_{1}+y_{2}=k_{1}(x_{1}+x_{2})+2=4k_{1}^{2}+2\)．于是 \par\noindent \(|AB|=|AF|+|BF|=(y_{1}+1)+(y_{2}+1)=y_{1}+y_{2}+2=4k_{1}^{2}+4\)．\par\noindent 同理 \(|DE|=4k_{2}^{2}+4\)．故 \par\noindent \(|AB|+|DE|=4(k_{1}^{2}+k_{2}^{2})+8\geqslant 8|k_{1}k_{2}|+8=8\times 2+8=24\)，\par\noindent 当且仅当 \(|k_{1}|=|k_{2}|\) 且 \(k_{1}k_{2}=-2\)，即 \((k_{1},k_{2})=(\sqrt{2},-\sqrt{2})\) 或 \((-\sqrt{2},\sqrt{2})\) 时等号成立．所求最小值为 24．}
\fi
\end{ansblock}

\begin{ansblock}[2章-拓-课时16-01]
% ans:2章-拓-课时16-01
\ansitem{17}{B}
\ifansbookdetail
\ansnote{详解}{点在 \(y^{2}=4x\) 上，故 \(x\geqslant 0\) 且 \(y^{2}=4x\)．代入得 \(z=x^{2}+(1/2)\cdot 4x+3=x^{2}+2x+3=(x+1)^{2}+2\)．由 \(x\geqslant 0\)，\(x=0\) 时 \(z\) 最小，最小值为 3．故选：B．}
\fi
\end{ansblock}

\begin{ansblock}[2章-拓-课时16-03]
% ans:2章-拓-课时16-03
\ansitem{18}{C}
\ifansbookdetail
\ansnote{详解}{\(x^{2}=4y\) 的焦点 \(F(0,1)\)，准线 \(y=-1\)．延长 \(QP\) 交准线于 \(M\)，则 \(|PM|=y_{P}+1=|PQ|+1\)；由定义 \(|PM|=|PF|\)，故 \(|PQ|=|PF|-1\)．于是 \par\noindent \(|PA|+|PQ|=|PA|+|PF|-1\geqslant|AF|-1=\sqrt{8^{2}+(7-1)^{2}}-1=10-1=9\)，\par\noindent 当且仅当 \(A\)、\(P\)、\(F\) 三点共线时取等（\(A(8,7)\) 在抛物线外：\(64>4\times 7=28\)，线段 \(AF\) 与抛物线相交，取等可行）．最小值为 9，选 C．}
\fi
\end{ansblock}

\begin{ansblock}[2章-拓-课时16-04]
% ans:2章-拓-课时16-04
\ansitem{19}{D}
\ifansbookdetail
\ansnote{详解}{化标准形 \(x^{2}=-4y\)：焦点 \(F(0,-1)\)，准线 \(l\)：\(y=1\)．第一根式是 \(P\) 到 \(F(0,-1)\) 的距离 \(|PF|\)，第二根式是 \(P\) 到 \(A(4,-5)\) 的距离 \(|PA|\)．由定义，\(|PF|\) 等于 \(P\) 到准线 \(y=1\) 的距离（\(n\leqslant 0\)，即 \(1-n\)），故 \(|PF|+|PA|\) 不小于 \(A\) 到准线的垂线段长 \(1-(-5)=6\)（\(A\) 在抛物线内部：\(4^{2}=16<-4\times(-5)=20\) \ding{51}，\(P\) 取垂线段与抛物线交点即可取等）．最小值 6，选 D．}
\fi
\end{ansblock}

\begin{ansblock}[2章-拓-课时16-02]
% ans:2章-拓-课时16-02
\ansitem{20}{42 或 22}
\ifansbookdetail
\ansnote{详解}{按 \(M\) 与抛物线的位置分类．(1) \(M\) 在抛物线内部或上：\(40^{2}\leqslant 2p\cdot 20\)，即 \(p\geqslant 40\)．设 \(P\) 在准线 \(x=-p/2\) 上的射影为 \(D\)，\(|PF|=|PD|\)，\(M\)、\(P\)、\(D\) 共线时取最小，最小值为 \(M\) 到准线的距离 \(20+p/2=41\)，得 \(p=42\)（\(40\leqslant 42\) \ding{51}）．(2) \(M\) 在抛物线外部：\(p<40\)，\(P\)、\(M\)、\(F\) 共线时取最小，最小值为 \par\noindent \(|MF|=\sqrt{40^{2}+(20-p/2)^{2}}=41\)，\((20-p/2)^{2}=81\)、\(p/2=11\) 或 \(29\)，\(p=22\) 或 \(p=58\)（\(58>40\) 与外部前提矛盾，舍）．\par\noindent 综上 \(p=42\) 或 \(22\)．}
\fi
\end{ansblock}

\begin{ansblock}[2章-拓-课时16-21]
% ans:2章-拓-课时16-21
\ansitem{21}{7/2；P(2,2)}
\ifansbookdetail
\ansnote{详解}{\(A(3,2)\)：\(2^{2}=4<2\times 3=6\)，\(A\) 在抛物线内部．准线 \(l\)：\(x=-1/2\)．设 \(P\) 在 \(l\) 上的射影为 \(Q\)，由定义 \(|PF|=|PQ|\)，\(|PA|+|PF|=|PA|+|PQ|\geqslant A\) 到准线的水平距离 \(=3-(-1/2)=7/2\)，当 \(A\)、\(P\)、\(Q\) 共线（\(y_{P}=y_{A}=2\)）时取等．代入 \(y^{2}=2x\) 得 \(x=2\)，即 \(P(2,2)\)，最小值 \(7/2\)．}
\fi
\end{ansblock}

\begin{ansblock}[2章-拓-课时16-09]
% ans:2章-拓-课时16-09
\ansitem{22}{(0,0)（最小值 8）}
\ifansbookdetail
\ansnote{详解}{设 \(P(x_{0},y_{0})\)，\(y_{0}^{2}=-4x_{0}\)（\(x_{0}\leqslant 0\)）．\par\noindent \(\overrightarrow{AP}\cdot\overrightarrow{BP}=(x_{0}-2)(x_{0}-4)+y_{0}^{2}=x_{0}^{2}-6x_{0}+8-4x_{0}=x_{0}^{2}-10x_{0}+8=(x_{0}-5)^{2}-17\)．\par\noindent 由 \(x_{0}\leqslant 0\)，\((x_{0}-5)^{2}\geqslant 25\) 且随 \(x_{0}\) 增大而减小，故 \(x_{0}=0\) 时 \(\overrightarrow{AP}\cdot\overrightarrow{BP}\) 最小，最小值 8，此时 \(y_{0}=0\)，点 \(P(0,0)\)．}
\fi
\end{ansblock}

\begin{ansblock}[2章-拓-课时16-22]
% ans:2章-拓-课时16-22
\ansitem{23}{D}
\ifansbookdetail
\ansnote{详解}{\(P\) 在抛物线上：\(16=4m\)、\(m=4\)，\(P(4,-4)\)．焦半径（定义）：准线 \(x=-1\)，\(|PF|=x_{P}+1=4+1=5\)．（两点距复核：\(F(1,0)\)，\(|PF|=\sqrt{9+16}=5\) \ding{51}．）选 D．}
\fi
\end{ansblock}

\begin{ansblock}[2章-拓-课时16-06]
% ans:2章-拓-课时16-06
\ansitem{24}{B}
\ifansbookdetail
\ansnote{详解}{由对称性不妨设 \(P(x,y)\)（\(y>0\)）．特殊位：\(x=1\) 时 \(P(1,2)\)，\(PA\perp x\) 轴，\(\tan\angle APB=|AB|/|PA|=2/2=1\)，\(\angle APB=\pi/4\)；\(x=3\) 时 \(P(3,2\sqrt{3})\)，\(PB\perp x\) 轴，\(\tan\angle APB=|AB|/|PB|=2/(2\sqrt{3})=\sqrt{3}/3\)，\(\angle APB=\pi/6\)．一般位 \(x\neq 1\) 且 \(x\neq 3\)：\par\noindent \(\tan\angle APB=(k_{PB}-k_{PA})/(1+k_{PA}\cdot k_{PB})\)，其中 \(k_{PA}=y/(x-1)\)、\(k_{PB}=y/(x-3)\)，代入化简得 \(\tan\angle APB=2y/(x^{2}-4x+3+y^{2})\)；\par\noindent 以 \(y^{2}=4x\) 代 \(x=y^{2}/4\)：\(\tan\angle APB=2y/(y^{4}/16+3)=2/((1/16)y^{3}+3/y)=2/((1/16)y^{3}+1/y+1/y+1/y)\)．\par\noindent 分母四项用均值不等式：\((1/16)y^{3}+1/y+1/y+1/y\geqslant 4\sqrt[4]{(1/16)y^{3}\cdot(1/y)^{3}}=4\cdot(1/16)^{1/4}=2\)，故 \(\tan\angle APB\leqslant 2/2=1\)；等号要求 \((1/16)y^{3}=1/y\) 即 \(y=2\)（\(P(1,2)\)，落在 \(x=1\) 特殊位，不在一般位内），故一般位 \(0<\tan\angle APB<1\)、\(0<\angle APB<\pi/4\)．综上 \(\angle APB\) 最大值 \(\pi/4\)，选 B．}
\fi
\end{ansblock}

\begin{ansblock}[2章-拓-课时16-05]
% ans:2章-拓-课时16-05
\ansitem{25}{D}
\ifansbookdetail
\ansnote{详解}{\(M(x,y)\) 在以 \(B(6,0)\) 为圆心、2 为半径的圆上：\((x-6)^{2}+y^{2}=4\)；\(N(x_{0},y_{0})\) 满足 \(y_{0}^{2}=8x_{0}\)（\(x_{0}\geqslant 0\)）．在圆约束下加权配凑：\par\noindent \((1/4)|MA|^{2}=(1/4)[(x-2)^{2}+y^{2}]=(1/4)[(x-2)^{2}+y^{2}]+(3/4)[(x-6)^{2}+y^{2}-4]=(x-5)^{2}+y^{2}\)（所加项在圆上为 0），即 \((1/2)|MA|=|MC|\)，其中 \(C(5,0)\)．\par\noindent 于是 \(|MN|+(1/2)|MA|=|MN|+|MC|\geqslant|CN|\)．\par\noindent 又 \(|CN|^{2}=(x_{0}-5)^{2}+y_{0}^{2}\)，代入 \(y_{0}^{2}=8x_{0}\) 得 \(|CN|^{2}=(x_{0}-1)^{2}+24\geqslant 24\)，即 \(|CN|\geqslant 2\sqrt{6}\)，当且仅当 \(x_{0}=1\)（\(N(1,\pm 2\sqrt{2})\)）且 \(C\)、\(M\)、\(N\) 三点共线（\(M\) 为线段 \(CN\) 与圆的交点）时取等——取等可行：\(C(5,0)\) 在圆内（\(|CB|=1<2\)）而 \(N\) 在圆外（\(|NB|^{2}=(x_{0}-6)^{2}+8x_{0}=(x_{0}-2)^{2}+32>4\) 恒成立），线段 \(CN\) 必与圆相交．最小值 \(2\sqrt{6}\)，选 D．}
\fi
\end{ansblock}

\begin{ansblock}[2章-拓-课时16-19]
% ans:2章-拓-课时16-19
\ansitem{26}{B}
\ifansbookdetail
\ansnote{详解}{连接 \(PF\)．由抛物线定义，\(P\) 到焦点的距离等于 \(P\) 到准线的距离，即 \(|PF|=|PQ|\)，\(\triangle QPF\) 为等腰三角形，\(P\) 到线段 \(FQ\) 两端点等距，故 \(P\) 在线段 \(FQ\) 的垂直平分线上，选 B．（反例排 A：取 \(y^{2}=4x\)、\(P(1,2)\)，则 \(F(1,0)\)、\(Q(-1,2)\)，\(|OF|=1\) 而 \(|OQ|=\sqrt{5}\)，\(O\) 不在垂直平分线上；C、D 随之不恒成立．）}
\fi
\end{ansblock}

\begin{ansblock}[2章-拓-课时16-11]
% ans:2章-拓-课时16-11
\ansitem{27}{ABC}
\ifansbookdetail
\ansnote{详解}{准线交 \(x\) 轴于 \(P\)，\(|PF|=p\)．分别过 \(A\)、\(B\) 作准线的垂线，垂足 \(E\)、\(M\)．倾斜角 \(60^{\circ}\)：\(AE\parallel x\) 轴 \(\Rightarrow\angle EAF=60^{\circ}\)；由定义 \(|AE|=|AF|\) \(\Rightarrow\triangle AEF\) 为等边三角形 \(\Rightarrow\angle EFP=\angle AEF=60^{\circ}\)、\(\angle FEP=30^{\circ}\)（直角 \(\triangle EPF\) 中），于是 \(|AF|=|EF|=2|PF|=2p=8\)，\(p=4\)，A 对．又 \(|AE|=|EF|=2|PF|\) 且 \(PF\parallel AE\) \(\Rightarrow F\) 为 \(AD\) 中点，\(|DF|=|FA|\)，B 对．\(\angle DAE=60^{\circ}\) 的补算：\(\triangle ADE\) 中 \(\angle AED=90^{\circ}\)、\(\angle ADE=30^{\circ}\) \(\Rightarrow|BD|=2|BM|=2|BF|\)（末步用定义 \(|BM|=|BF|\)），C 对．由 \(|BD|=2|BF|\) 得 \(|BF|=(1/3)|DF|=(1/3)|AF|=8/3\neq 4\)，D 错．故选 ABC．}
\fi
\end{ansblock}

\begin{ansblock}[2章-拓-课时16-12]
% ans:2章-拓-课时16-12
\ansitem{28}{A}
\ifansbookdetail
\ansnote{详解}{圆化标准形：\((x-1)^{2}+y^{2}=1\)，圆心 \(F(1,0)\)、半径 1，恰为抛物线 \(y^{2}=4x\) 的焦点．设 \(A(x_{1},y_{1})\)、\(D(x_{2},y_{2})\)（\(A\)、\(D\) 为直线与抛物线的两交点，\(B\)、\(C\) 为与圆的两交点，自上而下顺次），由焦半径 \(|AF|=x_{1}+1\)、\(|DF|=x_{2}+1\)，得 \(|AB|=|AF|-|BF|=x_{1}+1-1=x_{1}\)，\(|CD|=|DF|-|CF|=x_{2}\)．设 \(l\)：\(x=my+1\) 代入 \(y^{2}=4x\) 得 \(y^{2}-4my-4=0\)，\(y_{1}y_{2}=-4\)，故 \par\noindent \(|AB|\cdot|CD|=x_{1}x_{2}=(y_{1}^{2}/4)(y_{2}^{2}/4)=(y_{1}y_{2})^{2}/16=16/16=1\)，\par\noindent 为定值 1（与 \(m\) 无关，最值之说随之落空）．选 A．}
\fi
\end{ansblock}

\begin{ansblock}[2章-拓-课时16-13]
% ans:2章-拓-课时16-13
\ansitem{29}{B（−4）}
\ifansbookdetail
\ansnote{详解}{联立 \(x^{2}=4(kx+2)\)：\(x^{2}-4kx-8=0\)（\(\Delta=16k^{2}+32>0\) 恒成立），\(x_{1}+x_{2}=4k\)、\(x_{1}x_{2}=-8\)．\par\noindent \(y_{1}y_{2}=(kx_{1}+2)(kx_{2}+2)=k^{2}x_{1}x_{2}+2k(x_{1}+x_{2})+4=-8k^{2}+8k^{2}+4=4\)（即过 \((0,2)\) 的弦两端纵坐标之积为定值 4，\(A\)、\(B\) 分居 \(y\) 轴两侧故 \(x_{1}x_{2}<0\) 与 \(-8\) 相符）．\par\noindent \(\overrightarrow{OA}\cdot\overrightarrow{OB}=x_{1}x_{2}+y_{1}y_{2}=-8+4=-4\)，与 \(k\) 无关．值 \(-4\) 对应选项 B，选 B．}
\fi
\end{ansblock}

\begin{ansblock}[2章-拓-课时16-26]
% ans:2章-拓-课时16-26
\ansitem{30}{A}
\ifansbookdetail
\ansnote{详解}{\(F(1,0)\)．设 \(l\)：\(x=ty+1\)（含斜率不存在），代入 \(y^{2}=4x\)：\(y^{2}-4ty-4=0\)，\(y_{1}+y_{2}=4t\)、\(y_{1}y_{2}=-4\)，\(x_{1}+x_{2}=t(y_{1}+y_{2})+2=4t^{2}+2\)、\(x_{1}x_{2}=(y_{1}y_{2})^{2}/16=1\)．\(\angle AMB\) 为锐角 \(\iff\overrightarrow{MA}\cdot\overrightarrow{MB}>0\)（且 \(A\)、\(M\)、\(B\) 不共线，\(M\) 在 \(x\) 轴负半轴而弦过 \(F\) 恒成立）：\par\noindent \(\overrightarrow{MA}\cdot\overrightarrow{MB}=(x_{1}-m)(x_{2}-m)+y_{1}y_{2}=x_{1}x_{2}-m(x_{1}+x_{2})+m^{2}-4=m^{2}-2m-3-4mt^{2}\)．\par\noindent 按题意，对一切过 \(F\) 的弦（\(\forall t\in\mathrm{R}\)）均存在这样的锐角：\(m<0\) 时 \(-4mt^{2}\geqslant 0\) 且随 \(t^{2}\) 递增，最严处 \(t=0\)，须 \(m^{2}-2m-3>0\) \(\Rightarrow m<-1\) 或 \(m>3\)，配 \(m<0\) 得 \(m\in(-\infty,-1)\)，选 A．}
\fi
\end{ansblock}

\begin{ansblock}[2章-拓-课时16-27]
% ans:2章-拓-课时16-27
\ansitem{31}{B}
\ifansbookdetail
\ansnote{详解}{\(p=4\)，准线 \(x=-2\)．分别过 \(A\)、\(M\)、\(B\) 作准线的垂线，垂足 \(C\)、\(D\)、\(E\)，则 \(MD\) 为梯形（退化情形为中点公式仍合）\(ACEB\) 的中位线，\(|MD|=(|AC|+|BE|)/2=(|AF|+|BF|)/2=|AB|/2\)（定义）．设 \(l\)：\(x=my+2\) 代入 \(y^{2}=8x\)：\(y^{2}-8my-16=0\)，\par\noindent \(|AB|=\sqrt{1+m^{2}}\,|y_{1}-y_{2}|=\sqrt{1+m^{2}}\cdot\sqrt{64m^{2}+64}=8(1+m^{2})\geqslant 8\)（\(m=0\) 即通径取等）．\par\noindent 故 \(|MD|=|AB|/2\geqslant 4\)，最小值 4，选 B．}
\fi
\end{ansblock}

\begin{ansblock}[2章-拓-课时16-28]
% ans:2章-拓-课时16-28
\ansitem{32}{A}
\ifansbookdetail
\ansnote{详解}{由定义 \(d_{1}=|MF|\)；设 \(MN\perp l'\) 于 \(N\)，则 \(d_{2}=|MN|\)，\(d_{1}+d_{2}=|MF|+|MN|\geqslant d(F,l')\)（\(F\) 到 \(l'\) 的垂线段长），当 \(M\) 落在「过 \(F\) 垂直于 \(l'\) 的直线」与抛物线的交点上（该垂线与 \(l'\) 的垂足 \(N\) 在 \(F\) 同侧射影之间，可行）时取最小．\par\noindent \(d(F,l')=|1-0+2|/\sqrt{2}=3\sqrt{2}/2\)．\par\noindent 设 \(l'\) 与 \(x\) 轴交于 \(D(-2,0)\)，\(|FD|=3\)．\par\noindent 在 Rt\(\triangle DNF\)（\(\angle DNF=90^{\circ}\)）中，\(\cos\angle NFD=|FN|/|FD|=(3\sqrt{2}/2)/3=\sqrt{2}/2\)；取最小值时 \(M\) 在线段 \(FN\) 上、射线 \(FO\) 与 \(FD\) 同向，\(\cos\angle MFO=\sqrt{2}/2\)，选 A．（\(M\) 坐标核：\(M(t^{2},2t)\) 且 \(FM\parallel(1,-1)\) \(\Rightarrow t=\sqrt{2}-1\)，\(M(3-2\sqrt{2},2\sqrt{2}-2)\) 在 \(FN\) 段内 \ding{51}．）}
\fi
\end{ansblock}

\begin{ansblock}[2章-拓-课时16-24]
% ans:2章-拓-课时16-24
\ansitem{33}{C}
\ifansbookdetail
\ansnote{详解}{\(|OA|=|OB|\)：\(x_{1}^{2}+y_{1}^{2}=x_{2}^{2}+y_{2}^{2}\)，而 \(y_{i}=x_{i}^{2}/2\)，\((x_{1}^{2}-x_{2}^{2})[1+(x_{1}^{2}+x_{2}^{2})/4]=0\) \(\Rightarrow x_{1}^{2}=x_{2}^{2}\)，\(A\)、\(B\) 异点故 \(x_{2}=-x_{1}\)，两点关于 \(y\) 轴对称．设 \(B(a,a^{2}/2)\)、\(A(-a,a^{2}/2)\)（\(a>0\)），\(S=(1/2)\cdot 2a\cdot(a^{2}/2)=a^{3}/2=12\sqrt{3}\) \(\Rightarrow a^{3}=24\sqrt{3}=(2\sqrt{3})^{3}\) \(\Rightarrow a=2\sqrt{3}\)，\(B(2\sqrt{3},6)\)．记 \(OB\) 与 \(y\) 轴夹角 \(\theta\)：\(\tan\theta=a/(a^{2}/2)=2\sqrt{3}/6=\sqrt{3}/3\) \(\Rightarrow\theta=30^{\circ}\)，\(\angle AOB=2\theta=60^{\circ}\)．选 C．}
\fi
\end{ansblock}

\begin{ansblock}[2章-拓-课时16-15]
% ans:2章-拓-课时16-15
\ansitem{34}{y²=(p/2)x（x≥0），轨迹是顶点在原点、对称轴为 x 轴、开口向右的抛物线（2p̃=p/2，焦点 (p/8,0)）}
\ifansbookdetail
\ansnote{详解}{设抛物线上点 \((x_{0},y_{0})\)，\(y_{0}^{2}=2px_{0}\)（\(x_{0}\geqslant 0\)），垂线段中点 \((x,y)=(x_{0},y_{0}/2)\)．消参：\(x_{0}=x\)、\(y_{0}=2y\) 代入得 \(4y^{2}=2px\)，即 \(y^{2}=(p/2)x\)（\(x\geqslant 0\)）．这是顶点在原点、对称轴为 \(x\) 轴、开口向右的抛物线，其 \(2\tilde{p}=p/2\)（\(\tilde{p}=p/4\)），焦点 \((p/8,0)\)．（原点处垂线段退化为点 \(O\)，中点仍为 \(O\)，轨迹含顶点 \ding{51}．）}
\fi
\end{ansblock}

\begin{ansblock}[2章-拓-课时16-16]
% ans:2章-拓-课时16-16
\ansitem{35}{y²=8x}
\ifansbookdetail
\ansnote{详解}{点 \(A\) 横坐标 \(2>0\) 在 \(y\) 轴右侧，故抛物线开口向右，设 \(y^{2}=2px\)（\(p>0\)），\(F(p/2,0)\)．设 \(A(2,y_{A})\)，\(y_{A}^{2}=4p\)．\(\overrightarrow{FA}=(2-p/2,y_{A})\)、\(\overrightarrow{OA}=(2,y_{A})\)，\(\overrightarrow{FA}\cdot\overrightarrow{OA}=2(2-p/2)+y_{A}^{2}=4-p+4p=4+3p=16\)，解得 \(p=4\)，抛物线为 \(y^{2}=8x\)．（回验：\(A(2,\pm 4)\)，\(\overrightarrow{FA}\cdot\overrightarrow{OA}=2(2-2)+16=16\) \ding{51}．）}
\fi
\end{ansblock}

\begin{ansblock}[2章-拓-课时16-17]
% ans:2章-拓-课时16-17
\ansitem{36}{证明见详解；类似结论：y²=−2px 上 |MF|=p/2−x₀（x₀≤0）；x²=2py 上 |MF|=y₀+p/2；x²=−2py 上 |MF|=p/2−y₀}
\ifansbookdetail
\ansnote{详解}{证明：准线 \(l\)：\(x=-p/2\)．由抛物线定义，\(|MF|\) 等于 \(M\) 到 \(l\) 的距离；\(M\) 横坐标 \(x_{0}\geqslant 0\)，故距离为 \(x_{0}-(-p/2)=x_{0}+p/2\)．类似结论（同法：焦半径＝到准线距离，注意坐标符号范围）：①\(y^{2}=-2px\)（\(x_{0}\leqslant 0\)，准线 \(x=p/2\)）：\(|MF|=p/2-x_{0}\)；②\(x^{2}=2py\)（\(y_{0}\geqslant 0\)，准线 \(y=-p/2\)）：\(|MF|=y_{0}+p/2\)；③\(x^{2}=-2py\)（\(y_{0}\leqslant 0\)，准线 \(y=p/2\)）：\(|MF|=p/2-y_{0}\)．统一记法：焦半径＝\(p/2\)＋（沿开口方向的点到顶点有向距离）．}
\fi
\end{ansblock}

\begin{ansblock}[2章-拓-课时16-07]
% ans:2章-拓-课时16-07
\ansitem{37}{y²=3x 或 y²=−3x}
\ifansbookdetail
\ansnote{详解}{设方程为 \(y^{2}=2px\) 或 \(y^{2}=-2px\)（\(p>0\)），交点 \(A(x_{1},y_{1})\)、\(B(x_{2},y_{2})\)（\(y_{1}>0>y_{2}\)）．两曲线均关于 \(x\) 轴对称，公共弦被 \(x\) 轴垂直平分：\(x_{1}=x_{2}\)、\(y_{2}=-y_{1}\)，弦长 \(y_{1}-y_{2}=2y_{1}=2\sqrt{3}\) \(\Rightarrow y_{1}=\sqrt{3}\)．代入圆：\(x_{1}^{2}=4-3=1\)、\(x_{1}=\pm 1\)．点 \((1,\sqrt{3})\) 在 \(y^{2}=2px\) 上：\(3=2p\)、\(p=3/2\) \(\Rightarrow y^{2}=3x\)；点 \((-1,\sqrt{3})\) 在 \(y^{2}=-2px\) 上：同理 \(p=3/2\) \(\Rightarrow y^{2}=-3x\)．故方程为 \(y^{2}=3x\) 或 \(y^{2}=-3x\)．}
\fi
\end{ansblock}

\begin{ansblock}[2章-拓-课时16-08]
% ans:2章-拓-课时16-08
\ansitem{38}{BCD}
\ifansbookdetail
\ansnote{详解}{圆 \(C\) 圆心 \((0,1)\)、半径 4；抛物线 \(x^{2}=4y\) 的焦点恰为 \(F(0,1)\)、准线 \(y=-1\)，圆心与焦点重合．联立 \(x^{2}=4y\) 与 \(x^{2}+(y-1)^{2}=16\)：\(4y+(y-1)^{2}=16\) \(\Rightarrow y^{2}+2y-15=0\) \(\Rightarrow y=3\)（\(y=-5\) 舍，\(x^{2}=-20\) 无解），交点 \((\pm 2\sqrt{3},3)\)．劣弧 \(AB\) 为弦 \(y=3\) 上方含点 \((0,5)\) 的短弧（圆心到弦距离 \(2<\) 半径 4），其上动点 \(P\) 纵坐标 \(y_{P}\in(3,5)\)，A 项误以横坐标 \(2\sqrt{3}\) 为下界，错．B：\(P\)、\(N\) 同横坐标且 \(P\) 在 \(N\) 上方，\(|PN|=y_{P}-y_{N}\)；由定义 \(|NF|=y_{N}+1\)（\(N\) 到准线距离），故 \(|PN|+|NF|=y_{P}+1\)，恰为 \(P\) 到准线 \(y=-1\) 的距离，对（特位 \(x=0\)：\(P(0,4)\)、\(N(0,0)\)，\(4+1=5=y_{P}+1\) \ding{51}）．C：圆心 \((0,1)\) 到准线 \(y=-1\) 距离为 2，对．D：\(|PF|=\) 半径 4，\(\triangle PFN\) 周长 \(=|PF|+|PN|+|NF|=4+(y_{P}-y_{N})+(y_{N}+1)=y_{P}+5\in(8,10)\)，对．故选 BCD．}
\fi
\end{ansblock}

\begin{ansblock}[2章-拓-课时16-10]
% ans:2章-拓-课时16-10
\ansitem{39}{ABD}
\ifansbookdetail
\ansnote{详解}{A：由定义，\(|PF|\) 等于 \(P\) 到准线 \(l\) 的距离，而 \(PE\perp l\) 于 \(E\)，故 \(|PE|=|PF|\)，对．B：\(PQ\) 平分 \(\angle EPF\) 的外角 \(\Rightarrow\angle FPQ=\angle NPQ\)；又 \(EN\parallel x\) 轴（同垂直于准线）\(\Rightarrow\angle NPQ=\angle PQF\)（内错角），得 \(\angle FPQ=\angle PQF\)，\(\triangle FPQ\) 中 \(|PF|=|QF|\)，对．C：若 \(|PN|=|MF|\)，又 \(|QM|=|QN|\)（\(Q\) 在外角平分线上，到角两边等距）、\(\angle QMF=\angle QNP=90^{\circ}\)，则 \(\triangle FMQ\cong\triangle PNQ\)，得 \(|FQ|=|PQ|\)，结合 B 得 \(\triangle PFQ\) 为正三角形、\(\angle PFQ=\pi/3\)，则 \(P\) 被锁定为定点，与「\(P\) 为异于 \(O\) 的任意一点」矛盾，C 不恒成立．D：四边形 \(KQNE\) 为矩形（\(\angle KEQ=\angle EQN=\angle K=90^{\circ}\)）\(\Rightarrow|EN|=|KQ|\)；\(|PN|=|EN|-|EP|=|KQ|-|PF|=|KF|+|FQ|-|PF|=|KF|\)（用 \(|EP|=|PF|=|FQ|\)），对．故选 ABD．}
\fi
\end{ansblock}

\begin{ansblock}[2章-拓-课时16-14]
% ans:2章-拓-课时16-14
\ansitem{40}{AD}
\ifansbookdetail
\ansnote{详解}{抛物线 \(x^{2}=2y\) 的焦点 \((0,1/2)\)、准线 \(y=-1/2\)．设 \(l\)：\(y=kx+1/2\)，代入 \(x^{2}=2y\) 得 \(x^{2}-2kx-1=0\)，\(x_{1}+x_{2}=2k\)、\(x_{1}x_{2}=-1\)；\(y_{1}+y_{2}=k(x_{1}+x_{2})+1=2k^{2}+1\)、\par\noindent \(y_{1}y_{2}=(kx_{1}+1/2)(kx_{2}+1/2)=k^{2}x_{1}x_{2}+(k/2)(x_{1}+x_{2})+1/4=-k^{2}+k^{2}+1/4=1/4\)．\par\noindent A：过点 \((1,1/2)\)、斜率 \(t\) 的直线 \(y=t(x-1)+1/2\) 代入 \(x^{2}=2y\) 得 \(x^{2}-2tx+2t-1=0\)，即 \((x-1)(x-(2t-1))=0\)，另一交点横标 \(2t-1\)；相切 \(\iff 2t-1=1\iff t=1\)，切线 \(y=x-1/2\) 即 \(2x-2y-1=0\)，A 对．B：\(\overrightarrow{AM}\cdot\overrightarrow{BM}=(-x_{1},-1/2-y_{1})\cdot(-x_{2},-1/2-y_{2})=x_{1}x_{2}+(1/2+y_{1})(1/2+y_{2})=x_{1}x_{2}+1/4+(1/2)(y_{1}+y_{2})+y_{1}y_{2}=-1+1/4+(1/2)(2k^{2}+1)+1/4=k^{2}\)，仅 \(k=0\)（通径位）时为 0，不恒成立，B 错．C：取 \(l\parallel x\) 轴（\(k=0\)，即通径），\(A(1,1/2)\)、\(B(-1,1/2)\)，\(|OA|=|OB|=\sqrt{1+1/4}=\sqrt{5}/2\)，\(|OA|+|OB|=\sqrt{5}\)，严格不等式不成立，C 错．D：\(BN\perp\) 准线 \(\Rightarrow N(x_{2},-1/2)\)；直线 \(OA\)：\(y=(y_{1}/x_{1})x=(x_{1}/2)x\)（用 \(y_{1}=x_{1}^{2}/2\)），代入 \(x=x_{2}\) 得纵标 \((x_{1}x_{2})/2=-1/2\)，即直线 \(OA\) 恰过 \(N(x_{2},-1/2)\)，\(A\)、\(O\)、\(N\) 共线，D 对．故选 AD．}
\fi
\end{ansblock}

\begin{ansblock}[2章-拓-课时16-20]
% ans:2章-拓-课时16-20
\ansitem{41}{2√3}
\ifansbookdetail
\ansnote{详解}{椭圆左顶点 \((-4,0)\)．设 \(P(-y^{2}/4,y)\)，则 \par\noindent \(d^{2}=(-y^{2}/4+4)^{2}+y^{2}=(1/16)y^{4}-y^{2}+16=(1/16)(y^{2}-8)^{2}+12\geqslant 12\)，\par\noindent 当 \(y^{2}=8\)（\(y=\pm 2\sqrt{2}\)，\(x=-2\)，点 \(P(-2,\pm 2\sqrt{2})\) 在抛物线上 \ding{51}）时取等，\(d\) 的最小值为 \(\sqrt{12}=2\sqrt{3}\)．}
\fi
\end{ansblock}

\begin{ansblock}[2章-拓-课时16-23]
% ans:2章-拓-课时16-23
\ansitem{42}{y²=12x}
\ifansbookdetail
\ansnote{详解}{设圆 \(M\) 半径为 \(R\)．\(\odot C\) 圆心 \(C(3,0)\)、半径 1，外切 \(\Rightarrow|MC|=R+1\)；\(\odot M\) 与直线 \(x=-2\) 相切 \(\Rightarrow\) 圆心 \(M\) 到 \(x=-2\) 的距离为 \(R\)，从而 \(M\) 到直线 \(x=-3\) 的距离为 \(R+1\)．于是动点 \(M\) 到定点 \(C(3,0)\) 的距离等于它到定直线 \(x=-3\) 的距离，轨迹是以 \(C\) 为焦点、\(x=-3\) 为准线的抛物线（定义）：\(p/2=3\)、\(p=6\)，标准方程 \(y^{2}=12x\)．（核验 \(x\geqslant 0\)：\(M\) 到 \(x=-3\) 距离 \(=\) 横坐标 \(+3=R+1>0\) 且 \(|MC|=R+1\geqslant 1>0\) 自洽．）}
\fi
\end{ansblock}

\begin{ansblock}[2章-拓-课时16-25]
% ans:2章-拓-课时16-25
\ansitem{43}{7√3/12}
\ifansbookdetail
\ansnote{详解}{等腰梯形两底 \(AB\parallel CD\)，而抛物线上平行于 \(x\) 轴方向的弦（两交点同纵坐标）不存在（\(y^{2}=2px\) 单值），故两底必为垂直于 \(x\) 轴的弦：\(A\)、\(B\) 关于 \(x\) 轴对称，\(C\)、\(D\) 关于 \(x\) 轴对称．\(|AB|=2\Rightarrow y_{A}=1\)（取上半），\(|CD|=4\Rightarrow y_{D}=2\)．设 \(A(m,1)\)；腰与底夹角 \(\angle ADC=60^{\circ}\)，底铅直，故 \(DA\) 与铅直线夹 \(60^{\circ}\)，纵差 \(2-1=1\) \(\Rightarrow\) 横差 \(=\tan 60^{\circ}=\sqrt{3}\)，\(D(m+\sqrt{3},2)\)．两点在抛物线上：\(1=2pm\)、\(4=2p(m+\sqrt{3})\)，两式相减 \(3=2\sqrt{3}p\) \(\Rightarrow p=\sqrt{3}/2\)、\(m=1/(2p)=\sqrt{3}/3\)．由焦半径（例 36 结论）\(|AF|=m+p/2=\sqrt{3}/3+\sqrt{3}/4=7\sqrt{3}/12\)．}
\fi
\end{ansblock}

\begin{ansblock}[2章-拓-课时16-29]
% ans:2章-拓-课时16-29
\ansitem{44}{(1) y²=4x；(2) 直线 AB 恒过定点 (−1,0)}
\ifansbookdetail
\ansnote{详解}{(1) 设圆心 \(M(x,y)\)，圆过 \((2,0)\) \(\Rightarrow\) 半径\({}^{2}=(x-2)^{2}+y^{2}\)；被 \(y\) 轴截弦长 4 \(\Rightarrow\) 半径\({}^{2}=x^{2}+2^{2}\)（半弦 2、弦心距 \(|x|\)）．两式相等：\((x-2)^{2}+y^{2}=x^{2}+4\) \(\Rightarrow y^{2}=4x\)．\par\noindent (2) 顶点在轨迹 \(y^{2}=4x\) 上．设 \(B(x_{1},y_{1})\)、\(C(x_{2},y_{2})\)，\(BC\) 过 \(F(1,0)\) 且不垂直于坐标轴，设 \(BC\)：\(x=ty+1\)：\(y^{2}-4ty-4=0\) \(\Rightarrow y_{1}y_{2}=-4\)．由 \(A\)、\(C\) 关于 \(x\) 轴对称，\(A(x_{2},-y_{2})\)．设 \(AB\)：\(y=kx+m\)（\(k\neq 0\)）代入 \(ky^{2}-4y+4m=0\)，其两根为 \(y_{1}\)、\(-y_{2}\)：\(y_{1}\cdot(-y_{2})=4m/k\)，即 \(4=4m/k\)、\(m=k\)．故 \(AB\)：\(y=k(x+1)\)，对一切实有 \(k\) 恒过定点 \((-1,0)\)．}
\fi
\end{ansblock}

\begin{ansblock}[2章-拓-课时16-30]
% ans:2章-拓-课时16-30
\ansitem{45}{(1) x²=4y；(2) 点 P 恒在定直线 y=−1（即 C 的准线）上}
\ifansbookdetail
\ansnote{详解}{(1) \(A\) 横坐标 \(2\sqrt{2}\)、在圆上：\(y_{A}=(12-8)/2=2\)，\(A(2\sqrt{2},2)\)．代入 \(x^{2}=2py\)：\(8=4p\)、\(p=2\)，抛物线为 \(x^{2}=4y\)．\par\noindent (2) 先给切线式并验证：点 \(M(x_{1},x_{1}^{2}/4)\) 处，直线 \(y=(x_{1}/2)x-x_{1}^{2}/4\) 与 \(x^{2}=4y\) 联立得 \(x^{2}-2x_{1}x+x_{1}^{2}=(x-x_{1})^{2}=0\)，唯一公共点（重根 \(x_{1}\)），即为过 \(M\) 的切线 \(l_{1}\)；同理 \(N(x_{2},x_{2}^{2}/4)\) 处切线 \(l_{2}\)：\(y=(x_{2}/2)x-x_{2}^{2}/4\)．两切线联立：\((x_{1}-x_{2})x/2=(x_{1}^{2}-x_{2}^{2})/4\) \(\Rightarrow x_{0}=(x_{1}+x_{2})/2\)，\(y_{0}=(x_{1}/2)x_{0}-x_{1}^{2}/4=x_{1}x_{2}/4\)．设 \(l\)：\(y=kx+1\)（过 \(F(0,1)\)；\(l\) 为 \(x=0\) 时与抛物线仅一交点，不合「不同两点」），代入 \(x^{2}=4y\)：\(x^{2}-4kx-4=0\) \(\Rightarrow x_{1}x_{2}=-4\)，故 \(y_{0}=-1\) 恒成立，点 \(P\) 恒在定直线 \(y=-1\)（准线）上．}
\fi
\end{ansblock}

\begin{ansblock}[2章-导-课时16-G1]
% ans:2章-导-课时16-G1
\ansitem{46}{A}
\ifansbookdetail
\ansnote{详解}{\(y=2x^{2}\) 即 \(x^{2}=(1/2)y\)，\(2p=1/2\)、\(p=1/4\)，开口向上，焦点 \((0,p/2)=(0,1/8)\)．故选 A．（验算：准线 \(y=-1/8\)，顶点 \((0,0)\) 到焦点、到准线等距 \(1/8\) \ding{51}．）}
\fi
\end{ansblock}

\begin{ansblock}[2章-导-课时16-G2]
% ans:2章-导-课时16-G2
\ansitem{47}{D}
\ifansbookdetail
\ansnote{详解}{\(x^{2}=4y\) 的焦点在 \(y\) 轴上，图象关于 \(y\) 轴对称，对称轴为直线 \(x=0\)．故选 D．（验算：点 \((2,1)\) 在曲线上（\(4=4\times 1\) \ding{51}），其关于 \(y\) 轴的对称点 \((-2,1)\) 也在曲线上 \ding{51}；而关于 \(x\) 轴的对称点 \((2,-1)\) 不在曲线上（\(4\neq -1\)）——对称轴确为 \(y\) 轴 \ding{51}．）}
\fi
\end{ansblock}

\begin{ansblock}[2章-导-课时16-G3]
% ans:2章-导-课时16-G3
\ansitem{48}{D}
\ifansbookdetail
\ansnote{详解}{对圆上一点 \(Q\)，\(|PQ|\geqslant|PM|-|MQ|=|PM|-1\)（\(M\) 为圆心），且当 \(P\)、\(Q\)、\(M\) 共线、\(Q\) 在线段 \(PM\) 上时取等．设 \(P(x_{0},y_{0})\) 在 \(y^{2}=x\) 上，则 \(x_{0}=y_{0}^{2}\)，圆心 \(M(3,0)\)：\par\noindent \(|PM|^{2}=(y_{0}^{2}-3)^{2}+y_{0}^{2}\)．令 \(t=y_{0}^{2}\)（\(t\geqslant 0\)）：\(|PM|^{2}=(t-3)^{2}+t=t^{2}-5t+9=(t-5/2)^{2}+11/4\geqslant 11/4\)，\par\noindent 故 \(|PM|_{\min}=\sqrt{11}/2\)（\(t=5/2\) 时取得，此时 \(y_{0}=\pm\sqrt{10}/2\)，\(P(5/2,\pm\sqrt{10}/2)\) 在抛物线上 \ding{51}），从而 \(|PQ|_{\min}=\sqrt{11}/2-1\)，选 D．（验算：\(t=5/2\)：\(|PM|^{2}=(5/2-3)^{2}+5/2=1/4+5/2=11/4\) \ding{51}；\(|PM|_{\min}=\sqrt{11}/2>1\)，圆内不含抛物线点，\(|PQ|=|PM|-1\) 恒正、等号可取 \ding{51}．）}
\fi
\end{ansblock}

\begin{ansblock}[2章-导-课时16-G4]
% ans:2章-导-课时16-G4
\ansitem{49}{0<a≤1}
\ifansbookdetail
\ansnote{详解}{设 \(P(x,y)\) 为抛物线上任意一点，则 \(y=x^{2}/2\)（\(y\geqslant 0\)）．\par\noindent \(|AP|^{2}=x^{2}+(y-a)^{2}=2y+(y-a)^{2}=y^{2}+2(1-a)y+a^{2}\)．\par\noindent 这是关于 \(y\) 的二次函数，对称轴 \(y=a-1\)、开口向上．「最近点恰为顶点 \((0,0)\)」\(\iff|AP|^{2}\) 在 \(y=0\) 处取最小值 \(\iff\) 对称轴不在开区间 \((0,+\infty)\) 内 \(\iff a-1\leqslant 0\)，即 \(a\leqslant 1\)．又 \(a>0\)，故 \(0<a\leqslant 1\)．（验算：\(a=1\)：\(|AP|^{2}=y^{2}+1\) 在 \(y=0\) 处最小 \ding{51}（顶点最近，边界值取得到）；\(a=1.1\)：对称轴 \(y=0.1>0\)，最近点 \(y=0.1\neq 0\) 非顶点 \ding{55}——故右端取闭 \ding{51}．）}
\fi
\end{ansblock}

\begin{ansblock}[2章-导-课时16-G5]
% ans:2章-导-课时16-G5
\ansitem{50}{3√2/2}
\ifansbookdetail
\ansnote{详解}{设 \(P(t^{2}/4,t)\) 为抛物线上任意一点，则 \par\noindent \(d=|t^{2}/4-t+4|/\sqrt{1^{2}+1^{2}}=|t^{2}/4-t+4|/\sqrt{2}\)．\par\noindent 配方：\(t^{2}/4-t+4=(1/4)(t-2)^{2}+3>0\) 恒成立，故 \(d=(1/4)(t-2)^{2}/\sqrt{2}+3/\sqrt{2}\geqslant 3/\sqrt{2}=3\sqrt{2}/2\)，当 \(t=2\) 即 \(P(1,2)\) 时取得．（验算：\(P(1,2)\)：\(d=|1-2+4|/\sqrt{2}=3/\sqrt{2}=3\sqrt{2}/2\) \ding{51}；联立 \(x=y^{2}/4\) 与直线得 \(y^{2}-4y+16=0\)，\(\Delta=16-64<0\)，直线与抛物线无公共点，最小值为正、在切点型位置取得 \ding{51}．）}
\fi
\end{ansblock}

\jietitle{2.8 直线与圆锥曲线的位置关系（一）}

\begin{ansblock}[2章-练-课时17-01]
% ans:2章-练-课时17-01
\ansitem{1}{有公共点；两公共点 (1,2)、(4,−4)；以之为端点的线段长 3√5}
\ifansbookdetail
\ansnote{详解}{以 \(y\) 为主元消元（直线已给 \(y\)、抛物线给 \(x=y^{2}/4\)）：\(y=-2\times(y^{2}/4)+4=-y^{2}/2+4\)，两边乘 2 整理得 \(y^{2}+2y-8=0\)（二次项系数 \(1\neq 0\)，故确为二次、不会退化）．\(\Delta=2^{2}+4\times 8=36>0\)，故有两个公共点．\(y=(-2\pm 6)/2\)，即 \(y=2\) 或 \(y=-4\)；代回 \(x=y^{2}/4\) 得 \(x=1\) 或 \(x=4\)．\par\noindent 故公共点为 \((1,2)\) 与 \((4,-4)\)，所求线段长 \(=\sqrt{(4-1)^{2}+(-4-2)^{2}}=\sqrt{9+36}=\sqrt{45}=3\sqrt{5}\)．（回验：\(2=-2\times 1+4\) \ding{51}；\(-4=-2\times 4+4\) \ding{51}；\(1=1\)、\(16=4\times 4\) \ding{51}．）}
\fi
\end{ansblock}

\begin{ansblock}[2章-练-课时17-04]
% ans:2章-练-课时17-04
\ansitem{2}{C}
\ifansbookdetail
\ansnote{详解}{联立 \(y=x+1\) 与 \(x^{2}+y^{2}/2=1\)，消去 \(y\)：\(x^{2}+(x+1)^{2}/2=1\)，乘 2 得 \(2x^{2}+x^{2}+2x+1=2\)，即 \(3x^{2}+2x-1=0\)．\(\Delta=2^{2}-4\times 3\times(-1)=4+12=16>0\)，且二次项系数 \(3\neq 0\)，故有两个公共点，直线与椭圆相交，选 C．（易错点：消元后须检查二次项系数是否为 0；本题 \(x^{2}\) 项由椭圆方程给出、恒在，无降次之虞．）}
\fi
\end{ansblock}

\begin{ansblock}[2章-练-课时17-02]
% ans:2章-练-课时17-02
\ansitem{3}{反例：直线 y=1 与抛物线 y²=4x 只有一个公共点 (1/4,1)，但二者不相切（该点处抛物线的切线是 y=2x+1/2）}
\ifansbookdetail
\ansnote{详解}{取抛物线 \(C\)：\(y^{2}=4x\) 与直线 \(y=1\)（平行于 \(C\) 的对称轴 \(x\) 轴）．把 \(y=1\) 代入得 \(1=4x\)，唯一解 \(x=1/4\)，故恰有一个公共点 \((1/4,1)\)．\par\noindent 再验它不相切：设过该点的直线 \(y=k(x-1/4)+1\)，记 \(t=x-1/4\)（则 \(4x=4t+1\)），代入 \(y^{2}=4x\) 得 \((1+kt)^{2}=1+4t\)，展开整理为 \(k^{2}t^{2}+(2k-4)t=0\)．\(k\neq 0\) 时方程有二根 \(t=0\) 与 \(t=(4-2k)/k^{2}\)，相切须二根重合，即 \((4-2k)/k^{2}=0\) \(\iff k=2\)，切线为 \(y=2(x-1/4)+1\)，即 \(y=2x+1/2\)；而 \(k=0\)（即直线 \(y=1\)）时联立式退化为一次方程 \(-4t=0\)，只有一个公共点却不是重根，按「相切＝联立有重根」的判定，它不是切线、直线与抛物线不相切．\par\noindent 故「直线与抛物线只有一个公共点」成立时相切未必成立，即它不是相切的充分条件．（本席事理与变式 2、例 5 之退化特位三处互证．）}
\fi
\end{ansblock}

\begin{ansblock}[2章-练-课时17-05]
% ans:2章-练-课时17-05
\ansitem{4}{A}
\ifansbookdetail
\ansnote{详解}{正向：相切即联立所得二次方程有重根（\(\Delta=0\)），重根只对应一个公共点，故「相切」\(\Rightarrow\)「只有一个公共点」，充分性成立．\par\noindent 反向：抛物线中消元可能降次——直线平行于对称轴时（如 \(y^{2}=4x\) 与 \(y=1\)，见例 2），联立式退化为一次方程，恰有一个公共点却不相切，故「只有一个公共点」未必能推「相切」，必要性不成立．\par\noindent 于是「相切」是「只有一个公共点」的充分不必要条件，选 A．}
\fi
\end{ansblock}

\begin{ansblock}[2章-练-课时17-07]
% ans:2章-练-课时17-07
\ansitem{5}{(1)错误　(2)错误　(3)正确　(4)正确}
\ifansbookdetail
\ansnote{详解}{(1) 错．椭圆外一点可作两条切线：如过 \((4,0)\) 对椭圆 \(x^{2}/4+y^{2}=1\)，设 \(y=k(x-4)\) 代入得 \((1+4k^{2})x^{2}-32k^{2}x+64k^{2}-4=0\)，令 \(\Delta=16-192k^{2}=0\) 得 \(k=\pm\sqrt{3}/6\)，恰两条，故「只能一条」错．\par\noindent (2) 错．原点 \((0,0)\) 在椭圆内部（\(0+0<1\)），过内点的直线必与椭圆交于两点；直接联立 \(b^{2}x^{2}+a^{2}y^{2}=a^{2}b^{2}\) 与 \(y=x\) 得 \((a^{2}+b^{2})x^{2}=a^{2}b^{2}\)，即 \(x^{2}=a^{2}b^{2}/(a^{2}+b^{2})>0\)，确有两交点，故「不一定相交」错．\par\noindent (3) 对．\((3,0)\) 满足 \(9/9+0/16=1\)，是椭圆 \(x^{2}/9+y^{2}/16=1\) 的右顶点，在椭圆上；过椭圆上一点有且仅有一条切线（此处为 \(x=3\)），故对．\par\noindent (4) 对．椭圆是封闭曲线，直线与它联立消元后二次项系数恒不为 0（不存在像抛物线平行对称轴、双曲线平行渐近线那样的降次情形），故「恰一个公共点」\(\iff\)「方程有重根」\(\iff\)「相切」，与变式 2 中抛物线的反例恰成对照．\par\noindent 故四判断依次为「错、错、对、对」．}
\fi
\end{ansblock}

\begin{ansblock}[2章-练-课时17-11]
% ans:2章-练-课时17-11
\ansitem{6}{一个公共点：k∈\{√2,−√2,3/2\} 或 l 的斜率不存在；两个公共点：k∈(−∞,−√2)∪(−√2,√2)∪(√2,3/2)；没有公共点：k∈(3/2,+∞)}
\ifansbookdetail
\ansnote{详解}{① \(l\) 垂直 \(x\) 轴（斜率不存在）：直线 \(x=1\) 代入 \(2x^{2}-y^{2}=2\) 得 \(y=0\)，唯一公共点 \((1,0)\)（该处切线恰为 \(x=1\)），有一个公共点．\par\noindent ② \(l\) 不垂直 \(x\) 轴：设 \(y-2=k(x-1)\)，代入 \(2x^{2}-y^{2}=2\) 并整理得 \((2-k^{2})x^{2}+2(k^{2}-2k)x-k^{2}+4k-6=0\)．\par\noindent (a) \(2-k^{2}=0\)，即 \(k=\pm\sqrt{2}\)（\(l\) 平行渐近线 \(y=\pm\sqrt{2}x\)）：方程降为一次 \(2(k^{2}-2k)x-k^{2}+4k-6=0\)，一次项系数 \(2(2-2k)=4(1-k)\neq 0\)（\(k=\pm\sqrt{2}\) 时），故恰有一个公共点．\par\noindent (b) \(2-k^{2}\neq 0\)：\(\Delta=4(k^{2}-2k)^{2}+4(2-k^{2})(k^{2}-4k+6)=4(12-8k)=48-32k\)．\(\Delta=0\) \(\iff k=3/2\)，一个公共点（相切）；\(\Delta>0\) \(\iff k<3/2\)，两个公共点（本情形下须剔除 \(k=\pm\sqrt{2}\)，已归 (a)）；\(\Delta<0\) \(\iff k>3/2\)，没有公共点．\par\noindent 综上：有一个公共点 \(\iff k\in\{\sqrt{2},-\sqrt{2},3/2\}\) 或斜率不存在；有两个公共点 \(\iff k\in(-\infty,-\sqrt{2})\cup(-\sqrt{2},\sqrt{2})\cup(\sqrt{2},3/2)\)；没有公共点 \(\iff k\in(3/2,+\infty)\)．（注意 \(3/2\) 不在 \(k=\pm\sqrt{2}\) 之内（\(\sqrt{2}\approx 1.414<1.5\)），分段无冲突．）}
\fi
\end{ansblock}

\begin{ansblock}[2章-练-课时17-13]
% ans:2章-练-课时17-13
\ansitem{7}{两个公共点：k<−√6/3 或 k>√6/3；只有一个公共点：k=±√6/3；没有公共点：−√6/3<k<√6/3}
\ifansbookdetail
\ansnote{详解}{把 \(y=kx+2\) 代入 \(x^{2}/3+y^{2}/2=1\)，乘 6 得 \(2x^{2}+3(kx+2)^{2}=6\)，即 \((2+3k^{2})x^{2}+12kx+6=0\)．二次项系数 \(2+3k^{2}\geqslant 2>0\) 恒成立（椭圆无降次情形），故公共点个数完全由 \(\Delta\) 定：\(\Delta=(12k)^{2}-4(2+3k^{2})\times 6=144k^{2}-48-72k^{2}=72k^{2}-48\)．\par\noindent \(\Delta>0\) \(\iff k^{2}>48/72=2/3\) \(\iff|k|>\sqrt{6}/3\)，两个公共点；\(\Delta=0\) \(\iff k=\pm\sqrt{6}/3\)，只有一个公共点（相切）；\(\Delta<0\) \(\iff k^{2}<2/3\) \(\iff|k|<\sqrt{6}/3\)，没有公共点．（边界核对：直线恒过 \((0,2)\)，而 \((0,2)\) 在椭圆外（\(0+4/2=2>1\)），故存在相离与相切位置、与读数自洽．）}
\fi
\end{ansblock}

\begin{ansblock}[2章-练-课时17-12]
% ans:2章-练-课时17-12
\ansitem{8}{k=1 时相切；k<1 时相交；k>1 时相离}
\ifansbookdetail
\ansnote{详解}{联立 \(y=kx+1\) 与 \(y^{2}=4x\) 消 \(y\)：\((kx+1)^{2}=4x\)，即 \(k^{2}x^{2}+(2k-4)x+1=0\)．\par\noindent ① \(k\neq 0\)：二次项系数 \(k^{2}>0\)，为一元二次方程，\(\Delta=(2k-4)^{2}-4k^{2}=4k^{2}-16k+16-4k^{2}=16(1-k)\)．\(\Delta=0\) \(\iff k=1\)，相切；\(\Delta>0\) \(\iff k<1\) 且 \(k\neq 0\)，相交（两交点）；\(\Delta<0\) \(\iff k>1\)，相离．\par\noindent ② \(k=0\)：直线为 \(y=1\)，平行于对称轴 \(x\) 轴，代入得 \(1=4x\) 即 \(x=1/4\)，恰一个公共点 \((1/4,1)\)，属相交（不切）——此即例 2、变式 2 的退化特位，并入「相交」档．\par\noindent 综上：\(k=1\) 相切；\(k<1\)（含 \(k=0\)）相交；\(k>1\) 相离．}
\fi
\end{ansblock}

\begin{ansblock}[2章-练-课时17-14]
% ans:2章-练-课时17-14
\ansitem{9}{|k|<√3/2（即 k∈(−√3/2, √3/2)）}
\ifansbookdetail
\ansnote{详解}{设 \(l\)：\(y=kx\)（过原点的直线中 \(x=0\) 即 \(y\) 轴与双曲线无公共点，故可只议斜率存在者）．代入 \(x^{2}/4-y^{2}/3=1\)，乘 12 得 \(3x^{2}-4k^{2}x^{2}=12\)，即 \((3-4k^{2})x^{2}=12\)．\par\noindent ① \(3-4k^{2}>0\)（\(|k|<\sqrt{3}/2\)）：\(x^{2}=12/(3-4k^{2})>0\)，解得 \(x=\pm\sqrt{12/(3-4k^{2})}\)，对应两个不同交点（关于原点对称），符合「相交于两点」．\par\noindent ② \(3-4k^{2}\leqslant 0\)（\(|k|\geqslant\sqrt{3}/2\)）：当 \(=0\)（\(l\) 与渐近线 \(y=\pm(\sqrt{3}/2)x\) 平行，且过原点即为渐近线本身）时方程为 \(0=12\) 无解；当 \(<0\) 时 \(x^{2}<0\) 无实数解；两种都无公共点．\par\noindent 故所求范围为 \(|k|<\sqrt{3}/2\)．（几何直观：过原点（双曲线中心）的直线只要比两条渐近线陡即与双曲线交于两支上对称两点；渐近线本身与曲线零公共点，正是例 4 中 \(k=\pm\sqrt{2}\) 型降次的对照位．）}
\fi
\end{ansblock}

\begin{ansblock}[2章-练-课时17-10]
% ans:2章-练-课时17-10
\ansitem{10}{1}
\ifansbookdetail
\ansnote{详解}{\(a^{2}=20\)、\(b^{2}=5\)，\(a=2\sqrt{5}\)、\(c^{2}=25\) 即右焦点 \(F(5,0)\)．过 \(F\) 的直线（除 \(x\) 轴外）可设为 \(x=my+5\)，代入 \(x^{2}-4y^{2}=20\) 得 \((my+5)^{2}-4y^{2}=20\)，整理为 \((m^{2}-4)y^{2}+10my+5=0\)．\par\noindent ① \(m^{2}=4\)（\(|m|=2\)，即直线斜率 \(|k|=1/2=b/a\)，平行渐近线）：二次项消失，化为一次方程 \(10my+5=0\)，只有一个公共点，不构成弦．\par\noindent ② \(m^{2}<4\)（斜率 \(|k|>1/2\)，较渐近线陡，两点同在右支）：\(y_{1}+y_{2}=-10m/(m^{2}-4)\)、\(y_{1}y_{2}=5/(m^{2}-4)\)，故 \((y_{1}-y_{2})^{2}=[100m^{2}-20(m^{2}-4)]/(m^{2}-4)^{2}=80(m^{2}+1)/(m^{2}-4)^{2}\)，\(|AB|=\sqrt{1+m^{2}}\,|y_{1}-y_{2}|=4\sqrt{5}(m^{2}+1)/(4-m^{2})\)．记 \(u=m^{2}\in[0,4)\)，\(f(u)=-1+5/(4-u)\) 随 \(u\) 单调增，即 \(|AB|\) 随 \(u\) 单调增，最小在 \(u=0\)，值 \(4\sqrt{5}\times(1/4)=\sqrt{5}\)，且仅此一处；\(u=0\) 即直线 \(x=5\)（垂直 \(x\) 轴的通径），其长 \(2b^{2}/a=2\times 5/(2\sqrt{5})=\sqrt{5}\) \ding{51}．\par\noindent ③ \(m^{2}>4\)（斜率 \(|k|<1/2\)，较渐近线平，两点分居两支）：\(|AB|=4\sqrt{5}(u+1)/(u-4)\)，而 \((u+1)/(u-4)=1+5/(u-4)\) 随 \(u\) 单调减，值域为 \((4\sqrt{5},+\infty)\)（\(u\to 4^{+}\) 时趋于 \(+\infty\)，\(u\to+\infty\) 时趋于 \(4\sqrt{5}\)＝实轴长但取不到，水平线 \(y=0\) 单列如下），恒大于 \(4\sqrt{5}>\sqrt{5}\)，取不到 \(\sqrt{5}\)．另单独核直线 \(y=0\)（\(x\) 轴，未被 \(x=my+5\) 覆盖）：它与双曲线交于两顶点，弦长 \(2a=4\sqrt{5}\neq\sqrt{5}\)．\par\noindent 综上，被截弦长等于 \(\sqrt{5}\) 的直线只有通径 \(x=5\) 这一条，故填 1．}
\fi
\end{ansblock}

\begin{ansblock}[2章-练-课时17-03]
% ans:2章-练-课时17-03
\ansitem{11}{√5}
\ifansbookdetail
\ansnote{详解}{直线给出 \(x=2y-2\)，代入 \(x^{2}+4y^{2}=4\)：\((2y-2)^{2}+4y^{2}=4\to 4y^{2}-8y+4+4y^{2}=4\to 8y^{2}-8y=0\to 8y(y-1)=0\)，得 \(y=0\) 或 \(y=1\)，对应 \(x=-2\) 或 \(x=0\)．故 \(A(-2,0)\)、\(B(0,1)\)，\(|AB|=\sqrt{(0+2)^{2}+(1-0)^{2}}=\sqrt{5}\)．（此法直接解出两交点，无需弦长公式；若用弦长公式，\(x=my-2\) 型联立给 \(|y_{1}-y_{2}|=1\)、\(\sqrt{1+m^{2}}=\sqrt{5}\) 之倒数形式，结果同为 \(\sqrt{5}\)．）}
\fi
\end{ansblock}

\begin{ansblock}[2章-练-课时17-09]
% ans:2章-练-课时17-09
\ansitem{12}{16√3/5}
\ifansbookdetail
\ansnote{详解}{\(a^{2}=3\)、\(b^{2}=6\)，\(c^{2}=a^{2}+b^{2}=9\)，右焦点 \(F_{2}(3,0)\)．倾斜角 \(30^{\circ}\) 得 \(k=\tan 30^{\circ}=\sqrt{3}/3\)，直线 \(l\)：\(y=(\sqrt{3}/3)(x-3)\)．代入 \(x^{2}/3-y^{2}/6=1\)：以 \(y^{2}=(x-3)^{2}/3\) 代，得 \(x^{2}/3-(x-3)^{2}/18=1\)，乘 18 整理为 \(6x^{2}-(x-3)^{2}=18\)，即 \(5x^{2}+6x-27=0\)（二次项系数 \(5\neq 0\)，故 \(l\) 不与渐近线平行；\(\Delta=36+540=576>0\)，确两交点）．韦达：\(x_{1}+x_{2}=-6/5\)、\(x_{1}x_{2}=-27/5\)，故 \((x_{1}-x_{2})^{2}=(x_{1}+x_{2})^{2}-4x_{1}x_{2}=36/25+108/5=576/25\)，\(|x_{1}-x_{2}|=24/5\)．弦长 \(|AB|=\sqrt{1+k^{2}}\,|x_{1}-x_{2}|=\sqrt{1+1/3}\times 24/5=(2/\sqrt{3})\times(24/5)=48/(5\sqrt{3})=16\sqrt{3}/5\)．}
\fi
\end{ansblock}

\begin{ansblock}[2章-练-课时17-06]
% ans:2章-练-课时17-06
\ansitem{13}{7}
\ifansbookdetail
\ansnote{详解}{\(2p=4\)、\(p=2\)，焦点 \(F(1,0)\)、准线 \(x=-1\)．由抛物线定义，焦半径等于点到准线的距离：\(|AF|=x_{1}+1\)、\(|BF|=x_{2}+1\)．因 \(A\)、\(F\)、\(B\) 共线且 \(F\) 在 \(A\)、\(B\) 之间（焦点弦的两端点分居 \(x\) 轴两侧），\(|AB|=|AF|+|BF|=x_{1}+x_{2}+2=5+2=7\)．（即焦点弦公式 \(|AB|=x_{1}+x_{2}+p=5+2=7\)；按定义补出焦半径一步，防公式误记 \(p/2\)．）}
\fi
\end{ansblock}

\begin{ansblock}[2章-练-课时17-16]
% ans:2章-练-课时17-16
\ansitem{14}{1/3}
\ifansbookdetail
\ansnote{详解}{焦点 \(F(0,p/2)\)、准线 \(y=-p/2\)．倾角 \(30^{\circ}\) 即直线斜率 \(\tan 30^{\circ}=\sqrt{3}/3\)，直线 \(l\)：\(y=(\sqrt{3}/3)x+p/2\)，等价地 \(x=\sqrt{3}(y-p/2)\)．代入 \(x^{2}=2py\) 得 \(3(y-p/2)^{2}=2py\)，展开：\(3y^{2}-3py+3p^{2}/4-2py=0\)，即 \(3y^{2}-5py+3p^{2}/4=0\)，除以 3：\(y^{2}-(5p/3)y+p^{2}/4=0\)（\(\Delta=25p^{2}/9-p^{2}=16p^{2}/9\)，两根 \(y=(5p/3\pm 4p/3)/2\)，即 \(p/6\) 与 \(3p/2\)）．\par\noindent 为免 \(y\) 轴左右两侧混判，改以 \(x\) 为主元：把 \(y=x/\sqrt{3}+p/2\) 代 \(x^{2}=2py\) 得 \(x^{2}=2p(x/\sqrt{3}+p/2)=(2p/\sqrt{3})x+p^{2}\)，即 \(x^{2}-(2p/\sqrt{3})x-p^{2}=0\)．解之：\(x=(p/\sqrt{3})\pm(1/2)\sqrt{16p^{2}/3}=(p/\sqrt{3})\pm(2p/\sqrt{3})\)，故 \(x=3p/\sqrt{3}=\sqrt{3}p\) 或 \(x=-p/\sqrt{3}\)．\par\noindent 定位辨析：题设「\(A\) 在 \(y\) 轴左侧」即 \(x_{A}<0\)，故只能取负根 \(x_{A}=-p/\sqrt{3}\)（此时 \(y_{A}=x_{A}^{2}/(2p)=(p^{2}/3)/(2p)=p/6\)），而 \(x_{B}=\sqrt{3}p\)（\(y_{B}=3p^{2}/(2p)=3p/2\)）——此两步不可互换，否则比值倒置为 3．\par\noindent 焦半径（开口向上：\(|MF|=y_{M}+p/2\)，到准线的距离，见课前预习）：\(|AF|=y_{A}+p/2=p/6+p/2=2p/3\)；\(|FB|=y_{B}+p/2=3p/2+p/2=2p\)．\par\noindent 故 \(|AF|/|FB|=(2p/3)/(2p)=1/3\)．（互核一（代回曲线）：\(A(-p/\sqrt{3},p/6)\) 使 \(x^{2}=p^{2}/3=2p\times p/6\) \ding{51}；\(B(\sqrt{3}p,3p/2)\) 使 \(x^{2}=3p^{2}=2p\times(3p/2)\) \ding{51}．互核二（焦点弦坐标积）：本题抛物线开口向上，焦点弦之对偶形式为 \(x_{1}x_{2}=-p^{2}\)，本处 \((-p/\sqrt{3})\times\sqrt{3}p=-p^{2}\) \ding{51}．互核三（弦长）：\(|AF|+|FB|=2p/3+2p=8p/3\)，而 \(|AB|=y_{A}+y_{B}+p=p/6+3p/2+p=8p/3\)，与两点距 \(\sqrt{(4p/\sqrt{3})^{2}+(4p/3)^{2}}=\sqrt{64p^{2}/9}=8p/3\) 合 \ding{51}．）}
\fi
\end{ansblock}

\begin{ansblock}[2章-练-课时17-08]
% ans:2章-练-课时17-08
\ansitem{15}{C}
\ifansbookdetail
\ansnote{详解}{联立消 \(y\)：\(x^{2}/4+(x+1)^{2}/2=1\)，乘 4 得 \(x^{2}+2(x+1)^{2}=4\)，即 \(3x^{2}+4x-2=0\)（\(\Delta=16-4\times 3\times(-2)=40>0\)，故确相交于两点）．设 \(A(x_{1},y_{1})\)、\(B(x_{2},y_{2})\)、中点 \(M(x_{0},y_{0})\)，韦达得 \(x_{1}+x_{2}=-4/3\)，故 \(x_{0}=(x_{1}+x_{2})/2=-2/3\)；中点在直线上，\(y_{0}=x_{0}+1=1/3\)．所以中点 \((-2/3,1/3)\)，选 C．（本席为点差法综合的铺垫位，正文按常规联立韦达书写．）}
\fi
\end{ansblock}

\begin{ansblock}[2章-练-课时17-15]
% ans:2章-练-课时17-15
\ansitem{16}{中点 (−2/3, 1/3)；|AB|=8√2/3}
\ifansbookdetail
\ansnote{详解}{直线写成 \(y=x+1\)，代入 \(x^{2}/6+y^{2}/3=1\)，乘 6 得 \(x^{2}+2(x+1)^{2}=6\)，即 \(3x^{2}+4x-4=0\)（\(\Delta=16+48=64>0\)，确两交点）．设 \(A(x_{1},y_{1})\)、\(B(x_{2},y_{2})\)，则 \(x_{1}+x_{2}=-4/3\)、\(x_{1}x_{2}=-4/3\)．\par\noindent 中点：\(x_{0}=(x_{1}+x_{2})/2=-2/3\)，\(y_{0}=x_{0}+1=1/3\)，即中点 \((-2/3,1/3)\)．\par\noindent 弦长：\((x_{1}-x_{2})^{2}=(x_{1}+x_{2})^{2}-4x_{1}x_{2}=16/9+16/3=64/9\)，\(|x_{1}-x_{2}|=8/3\)；直线斜率 \(k=1\)，故 \(|AB|=\sqrt{1+k^{2}}\,|x_{1}-x_{2}|=\sqrt{2}\times 8/3=8\sqrt{2}/3\)．\par\noindent（回代验：方程 \(3x^{2}+4x-4=0\) 之根 \(x=2/3\)、\(x=-2\)，对应点 \((2/3,5/3)\)、\((-2,-1)\)：中点 \(((2/3-2)/2,(5/3-1)/2)=(-2/3,1/3)\) \ding{51}，\(|AB|=\sqrt{(2/3+2)^{2}+(5/3+1)^{2}}=(8/3)\sqrt{2}\) \ding{51}．）}
\fi
\end{ansblock}

\begin{ansblock}[2章-拓-课时17-03]
% ans:2章-拓-课时17-03
\ansitem{17}{k<−√5/2 或 k>√5/2（|k|>√5/2）}
\ifansbookdetail
\ansnote{详解}{代入 \(x^{2}-(kx-1)^{2}=4\)，展开 \(x^{2}-k^{2}x^{2}+2kx-1=4\)，整理为 \((1-k^{2})x^{2}+2kx-5=0\)．\par\noindent ① \(1-k^{2}=0\)（\(k=\pm 1\)）：降为一次方程 \(2kx-5=0\)，有解 \(x=5/(2k)\)，即有公共点，不合「没有公共点」．\par\noindent ② \(1-k^{2}\neq 0\)：无公共点 \(\iff\Delta<0\)，\(\Delta=(2k)^{2}-4(1-k^{2})(-5)=4k^{2}+20(1-k^{2})=20-16k^{2}\)．令 \(20-16k^{2}<0\) 得 \(k^{2}>5/4\)，即 \(|k|>\sqrt{5}/2\)（此时 \(1-k^{2}\neq 0\) 自动满足，因 \(5/4>1\)）．\par\noindent 故 \(k\) 的取值范围为 \(k<-\sqrt{5}/2\) 或 \(k>\sqrt{5}/2\)．（几何校：双曲线 \(x^{2}/4-y^{2}/4=1\) 的渐近线斜率 \(\pm 1\)，直线恒过 \((0,-1)\)；\(|k|>\sqrt{5}/2\approx 1.118>1\) 属比渐近线陡、且与实轴交点在双曲线顶点内侧的情形，无交点自洽．）}
\fi
\end{ansblock}

\begin{ansblock}[2章-拓-课时17-17]
% ans:2章-拓-课时17-17
\ansitem{18}{D（4条）}
\ifansbookdetail
\ansnote{详解}{① 斜率不存在：\(l\) 为 \(x=0\)（\(y\) 轴），代入得 \(-y^{2}/3=1\) 无实数解，公共点 0 个，不合．\par\noindent ② 斜率存在：设 \(l\)：\(y=kx+3\)，代入 \(x^{2}/4-y^{2}/3=1\)（乘 12）：\(3x^{2}-4(kx+3)^{2}=12\)，展开得 \((3-4k^{2})x^{2}-24kx-48=0\)．\par\noindent (a) \(3-4k^{2}\neq 0\)（\(l\) 不平行渐近线 \(y=\pm(\sqrt{3}/2)x\)）：一个公共点 \(\iff\Delta=0\)，\(\Delta=(-24k)^{2}+4(3-4k^{2})\times 48=576k^{2}+192(3-4k^{2})=576-192k^{2}\)，令其为 0 得 \(k^{2}=3\)，\(k=\pm\sqrt{3}\)——两条切线，符合．\par\noindent (b) \(3-4k^{2}=0\)，即 \(k=\pm\sqrt{3}/2\)：方程降为一次 \(-24kx-48=0\)，一次项系数 \(-24k\neq 0\)（\(k=\pm\sqrt{3}/2\) 时），有唯一解 \(x=-2/k\)，即 \(l\) 平行渐近线时与双曲线恰一个公共点——又两条，符合．\par\noindent 合计 \(2+2=4\) 条，选 D．（校：\(k=\pm\sqrt{3}\) 时 \(3-4k^{2}=3-12=-9\neq 0\) \ding{51} 两类不重；\(k=\pm\sqrt{3}/2\) 与 \(\pm\sqrt{3}\) 互异 \ding{51} 四类无重复．）}
\fi
\end{ansblock}

\begin{ansblock}[2章-拓-课时17-18]
% ans:2章-拓-课时17-18
\ansitem{19}{D}
\ifansbookdetail
\ansnote{详解}{代入 \(x^{2}-(kx-1)^{2}=1\)，展开 \(x^{2}-k^{2}x^{2}+2kx-1=1\)，整理为 \((1-k^{2})x^{2}+2kx-2=0\)．设两交点横坐标 \(x_{1}\)、\(x_{2}\)，「同在右支且为不同两点」须同时满足：\par\noindent (i) \(1-k^{2}\neq 0\)（保持二次，否则一次方程只一交点）；\par\noindent (ii) \(\Delta=(2k)^{2}-4(1-k^{2})(-2)=4k^{2}+8(1-k^{2})=8-4k^{2}>0\) \(\iff k^{2}<2\)；\par\noindent (iii) 两根之积 \(x_{1}x_{2}=-2/(1-k^{2})>0\) \(\iff 1-k^{2}<0\) \(\iff k^{2}>1\)（两根同号）；\par\noindent (iv) 两根之和 \(x_{1}+x_{2}=-2k/(1-k^{2})=2k/(k^{2}-1)>0\) \(\iff k>0\)（结合 (iii) 同为正，即在右支 \(x\geqslant 1\)）．\par\noindent 合 (ii)(iii)(iv)：\(1<k^{2}<2\) 且 \(k>0\)，得 \(1<k<\sqrt{2}\)，即 \((1,\sqrt{2})\)，选 D．（校：此范围内 \(1-k^{2}\neq 0\) 自动满足；端点 \(k=1\) 使二次项消失（一交点）、\(k=\sqrt{2}\) 使 \(\Delta=0\)（切，一交点），故两端均开，排除 B 之左闭．）}
\fi
\end{ansblock}

\begin{ansblock}[2章-拓-课时17-19]
% ans:2章-拓-课时17-19
\ansitem{20}{(−2,2)}
\ifansbookdetail
\ansnote{详解}{曲线侧：\(2x=\sqrt{4+y^{2}}\) 要求 \(x\geqslant 0\)，两边平方得 \(4x^{2}=4+y^{2}\)，即 \(x^{2}-y^{2}/4=1\) 且 \(x\geqslant 0\)（即该双曲线的右支，含右顶点 \((1,0)\)，渐近线 \(y=\pm 2x\)）．直线 \(y=m(x+1)\) 恒过定点 \((-1,0)\)．代入 \(4x^{2}-y^{2}=4\)：\(4x^{2}-m^{2}(x+1)^{2}=4\)，即 \((4-m^{2})x^{2}-2m^{2}x-(m^{2}+4)=0\)．\par\noindent ① \(|m|<2\)（\(4-m^{2}>0\)）：两根之积 \(=-(m^{2}+4)/(4-m^{2})<0\)，故有一正根；正根满足 \(x^{2}-y^{2}/4=1\) 且 \(x>0\)，即 \(x\geqslant 1\)，落在右支上，有公共点．\par\noindent ② \(|m|=2\)：方程降为 \(-2m^{2}x-(m^{2}+4)=0\)，解得 \(x=-(m^{2}+4)/(2m^{2})=-8/8=-1<0\)，不在右支，无公共点．\par\noindent ③ \(|m|>2\)（\(4-m^{2}<0\)）：积 \(>0\) 且和 \(=2m^{2}/(4-m^{2})<0\)，两根同为负，右支无公共点．\par\noindent 故有公共点 \(\iff|m|<2\)，即 \(m\in(-2,2)\)．}
\fi
\end{ansblock}

\begin{ansblock}[2章-拓-课时17-01]
% ans:2章-拓-课时17-01
\ansitem{21}{C（弦长 8/5）}
\ifansbookdetail
\ansnote{详解}{\(a^{2}=4\)、\(b^{2}=1\)，\(c^{2}=a^{2}-b^{2}=3\)、\(c=\sqrt{3}\)，右焦点 \((\sqrt{3},0)\)．直线 \(l\)：\(y=x-\sqrt{3}\)．代入 \(x^{2}/4+y^{2}=1\)：\(x^{2}/4+(x-\sqrt{3})^{2}=1\)，乘 4 得 \(x^{2}+4(x^{2}-2\sqrt{3}x+3)=4\)，即 \(5x^{2}-8\sqrt{3}x+8=0\)．\(\Delta=(8\sqrt{3})^{2}-4\times 5\times 8=192-160=32>0\)，\(x_{1}+x_{2}=8\sqrt{3}/5\)、\(x_{1}x_{2}=8/5\)．\par\noindent 弦长 \(|AB|=\sqrt{1+k^{2}}\,|x_{1}-x_{2}|=\sqrt{2}\times\sqrt{\Delta}/5=\sqrt{2}\times\sqrt{32}/5=\sqrt{2}\times 4\sqrt{2}/5=8/5\)．（互核：\(|x_{1}-x_{2}|^{2}=(x_{1}+x_{2})^{2}-4x_{1}x_{2}=192/25-32/5=32/25\)，\(|AB|^{2}=2\times 32/25=64/25\)，\(|AB|=8/5\) \ding{51}．）故选 C．}
\fi
\end{ansblock}

\begin{ansblock}[2章-拓-课时17-05]
% ans:2章-拓-课时17-05
\ansitem{22}{−3}
\ifansbookdetail
\ansnote{详解}{焦点 \(F(1,0)\)，直线 \(l\)：\(y=x-1\)．代入 \((x-1)^{2}=4x\) 得 \(x^{2}-2x+1=4x\)，即 \(x^{2}-6x+1=0\)（\(\Delta=36-4=32>0\)）．设 \(A(x_{1},y_{1})\)、\(B(x_{2},y_{2})\)，则 \(x_{1}+x_{2}=6\)、\(x_{1}x_{2}=1\)．又 \(y_{1}y_{2}=(x_{1}-1)(x_{2}-1)=x_{1}x_{2}-(x_{1}+x_{2})+1=1-6+1=-4\)．故 \(\overrightarrow{OA}\cdot\overrightarrow{OB}=x_{1}x_{2}+y_{1}y_{2}=1+(-4)=-3\)．（互核：以 \(y\) 为主元，\(x=y+1\) 代 \(y^{2}=4(y+1)\) 得 \(y^{2}-4y-4=0\)，\(y_{1}y_{2}=-4\) \ding{51}；\(x_{1}x_{2}=(y_{1}+1)(y_{2}+1)=-4+4+1=1\) \ding{51}．）}
\fi
\end{ansblock}

\begin{ansblock}[2章-拓-课时17-07]
% ans:2章-拓-课时17-07
\ansitem{23}{10}
\ifansbookdetail
\ansnote{详解}{\(2p=8\)、\(p=4\)，焦点 \(F(2,0)\)、准线 \(x=-2\)．直线 \(l\)：\(y=2(x-2)\)．代入 \(y^{2}=8x\) 得 \(4(x-2)^{2}=8x\)，即 \((x-2)^{2}=2x\)，展开 \(x^{2}-4x+4=2x\)，得 \(x^{2}-6x+4=0\)（\(\Delta=36-16=20>0\)）．设 \(A(x_{1},y_{1})\)、\(B(x_{2},y_{2})\)，则 \(x_{1}+x_{2}=6\)．由抛物线定义（焦半径＝到准线距离，准线 \(x=-2\)）\(|AF|=x_{1}+2\)、\(|BF|=x_{2}+2\)，焦点弦 \(|AB|=|AF|+|BF|=x_{1}+x_{2}+4=6+4=10\)．（互核：以 \(y\) 为主元，\(x=y/2+2\) 代 \(y^{2}=8(y/2+2)\) 得 \(y^{2}-4y-16=0\)，\(y_{1}+y_{2}=4\)、\(y_{1}y_{2}=-16\)，\((y_{1}-y_{2})^{2}=16+64=80\)，\(|AB|=\sqrt{1+(1/2)^{2}}\,|y_{1}-y_{2}|=(\sqrt{5}/2)\times 4\sqrt{5}=10\) \ding{51}．）}
\fi
\end{ansblock}

\begin{ansblock}[2章-拓-课时17-12]
% ans:2章-拓-课时17-12
\ansitem{24}{25/4}
\ifansbookdetail
\ansnote{详解}{\(2p=8\)、\(p=4\)，焦点 \(F(2,0)\)，准线 \(x=-2\)．由 \(A(8,8)\)、\(F(2,0)\) 得直线 \(AB\) 的斜率 \(=(8-0)/(8-2)=4/3\)，方程 \(y=(4/3)(x-2)\)．代入 \(y^{2}=8x\)：\((16/9)(x-2)^{2}=8x\)，即 \(2(x-2)^{2}=9x\)，展开 \(2x^{2}-8x+8=9x\)，得 \(2x^{2}-17x+8=0\)（\(\Delta=289-64=225=15^{2}\)）．解得 \(x=(17\pm 15)/4\)，即 \(x=8\)（点 \(A\)）或 \(x=1/2\)，故 \(B(1/2,-2)\)（\(y=(4/3)(1/2-2)=-2\)）．\par\noindent 方法一（中点坐标）：中点横坐标 \(x_{0}=(8+1/2)/2=17/4\)，到准线 \(x=-2\) 的距离 \(=17/4+2=25/4\)．\par\noindent 方法二（焦半径/定义）：\(|AF|=x_{A}+p/2=8+2=10\)、\(|BF|=x_{B}+p/2=1/2+2=5/2\)，焦点弦上「中点到准线的距离＝\((|AF|+|BF|)/2=|AB|/2\)」（梯形中位线），得 \((10+5/2)/2=25/4\) \ding{51}．\par\noindent（校：\(y_{1}y_{2}=8\times(-2)=-16=-p^{2}\) \ding{51}，与焦点弦纵积定值相符．）}
\fi
\end{ansblock}

\begin{ansblock}[2章-拓-课时17-13]
% ans:2章-拓-课时17-13
\ansitem{25}{x=3}
\ifansbookdetail
\ansnote{详解}{设直线为 \(x=x_{0}\)（\(x_{0}\geqslant 0\)，否则与抛物线无交点）．代入 \(y^{2}=4x_{0}\)，得 \(y=\pm 2\sqrt{x_{0}}\)，两交点 \(A(x_{0},2\sqrt{x_{0}})\)、\(B(x_{0},-2\sqrt{x_{0}})\)（关于 \(x\) 轴对称，故该弦被 \(x\) 轴垂直平分）．\(|AB|=2\times 2\sqrt{x_{0}}=4\sqrt{x_{0}}=4\sqrt{3}\)，故 \(\sqrt{x_{0}}=\sqrt{3}\)、\(x_{0}=3\)．所以直线 \(AB\) 的方程为 \(x=3\)．（回验：\(y^{2}=12\)、\(y=\pm 2\sqrt{3}\)、\(|AB|=4\sqrt{3}\) \ding{51}；此即通径型垂直弦，长 \(4\sqrt{3}\) 大于通径 \(2p=4\)，点 \((3,\pm 2\sqrt{3})\) 在曲线上 \ding{51}．）}
\fi
\end{ansblock}

\begin{ansblock}[2章-拓-课时17-08]
% ans:2章-拓-课时17-08
\ansitem{26}{(1) |AB|=8√5，中点 (−4,−7)；(2) 不成立（OA·OB=−15≠0）}
\ifansbookdetail
\ansnote{详解}{(1) 代入 \(x^{2}=-8(x-3)\)：\(x^{2}=-8x+24\)，即 \(x^{2}+8x-24=0\)（\(\Delta=64+96=160>0\)）．设 \(A(x_{1},y_{1})\)、\(B(x_{2},y_{2})\)，则 \(x_{1}+x_{2}=-8\)、\(x_{1}x_{2}=-24\)．弦长 \(|AB|=\sqrt{1+1^{2}}\,|x_{1}-x_{2}|=\sqrt{2}\times\sqrt{\Delta}=\sqrt{2}\times\sqrt{160}=\sqrt{320}=8\sqrt{5}\)．中点 \(x_{0}=(x_{1}+x_{2})/2=-4\)，\(y_{0}=x_{0}-3=-7\)，即 \((-4,-7)\)．（根之实值 \(x=-4\pm 2\sqrt{10}\)，对应 \(y=-7\pm 2\sqrt{10}\)，代 \(x^{2}=-8y\) 校：\((-4\pm 2\sqrt{10})^{2}=56\mp 16\sqrt{10}\)，\(-8(-7\pm 2\sqrt{10})=56\mp 16\sqrt{10}\) \ding{51}．）\par\noindent (2) \(y_{1}y_{2}=(x_{1}-3)(x_{2}-3)=x_{1}x_{2}-3(x_{1}+x_{2})+9=-24+24+9=9\)．故 \(\overrightarrow{OA}\cdot\overrightarrow{OB}=x_{1}x_{2}+y_{1}y_{2}=-24+9=-15\neq 0\)，所以 \(\overrightarrow{OA}\) 与 \(\overrightarrow{OB}\) 不垂直，「\(OA\perp OB\)」不成立．（若按斜率复核：\(k_{OA}\cdot k_{OB}=(y_{1}y_{2})/(x_{1}x_{2})=9/(-24)=-3/8\neq -1\)，同得否．）}
\fi
\end{ansblock}

\begin{ansblock}[2章-拓-课时17-11]
% ans:2章-拓-课时17-11
\ansitem{27}{y=2x−2}
\ifansbookdetail
\ansnote{详解}{设 \(l\)：\(y=2x+m\)．因抛物线以 \(x=y^{2}/4\) 表出，改以 \(y\) 为主元：由 \(y=2x+m\) 得 \(x=(y-m)/2\)，代入 \(y^{2}=4x=2(y-m)\)，整理为 \(y^{2}-2y+2m=0\)（二次项系数 \(1\neq 0\)）．相交于两点须 \(\Delta=(-2)^{2}-4\times 2m=4-8m>0\)，即 \(m<1/2\)；此时 \(y_{1}+y_{2}=2\)、\(y_{1}y_{2}=2m\)，\((y_{1}-y_{2})^{2}=(y_{1}+y_{2})^{2}-4y_{1}y_{2}=4-8m\)．\par\noindent 弦长用 \(y\) 的差：\(|AB|=\sqrt{1+(1/k)^{2}}\,|y_{1}-y_{2}|=\sqrt{1+1/4}\times\sqrt{4-8m}=(\sqrt{5}/2)\sqrt{4-8m}\)．令其等于 5：\(\sqrt{4-8m}=10/\sqrt{5}=2\sqrt{5}\)，平方得 \(4-8m=20\)，故 \(m=-2\)（\(<1/2\) \ding{51} 合相交条件）．\par\noindent 所以直线 \(l\) 的方程为 \(y=2x-2\)．（回代以 \(x\) 为主元复核：\((2x-2)^{2}=4x\to 4x^{2}-8x+4=4x\to x^{2}-3x+1=0\)，\(|x_{1}-x_{2}|=\sqrt{9-4}=\sqrt{5}\)，\(|AB|=\sqrt{1+4}\times\sqrt{5}=5\) \ding{51}．）}
\fi
\end{ansblock}

\begin{ansblock}[2章-拓-课时17-04]
% ans:2章-拓-课时17-04
\ansitem{28}{(x−1/2)²+(y−1)²=1 或 \hspace{0pt plus 1.5em}(x−1/2)²+(y+1)²=1（两圆）}
\ifansbookdetail
\ansnote{详解}{\(y^{2}=2x\) 中 \(2p=2\)、\(p=1\)，准线 \(l\)：\(x=-1/2\)．设圆心 \(C(a,b)\) 在抛物线上，则 \(b^{2}=2a\)（故 \(a\geqslant 0\)）、半径 \(r\)．与 \(x\) 轴相切 \(\Rightarrow r=|b|\)；与准线 \(x=-1/2\) 相切 \(\Rightarrow r=a+1/2\)（\(a\geqslant 0\) 故 \(a-(-1/2)=a+1/2\)）．两式相等：\(|b|=a+1/2\)，以 \(a=b^{2}/2\) 代得 \(|b|=b^{2}/2+1/2\)，即 \(|b|^{2}-2|b|+1=0\)，\((|b|-1)^{2}=0\)，故 \(|b|=1\)、\(b=\pm 1\)，进而 \(a=1/2\)、\(r=1\)．\par\noindent 所以圆有两解：\((x-1/2)^{2}+(y-1)^{2}=1\) 与 \((x-1/2)^{2}+(y+1)^{2}=1\)．（回验：圆心 \((1/2,1)\) 在抛物线上 \(1=2\times 1/2\) \ding{51}；到 \(x\) 轴距离 \(1=r\) \ding{51}；到准线距离 \(1/2+1/2=1=r\) \ding{51}．）}
\fi
\end{ansblock}

\begin{ansblock}[2章-拓-课时17-06]
% ans:2章-拓-课时17-06
\ansitem{29}{证明见详解}
\ifansbookdetail
\ansnote{详解}{焦点 \(F(p/2,0)\)．设过 \(F\) 的直线为 \(x=my+p/2\)：此写法已含垂直 \(x\) 轴的弦（\(m=0\)，即通径），而未含的只有 \(y=0\)（抛物线的对称轴），它与抛物线仅一个公共点，与题设「相交（有两个交点）」矛盾，故不损一般性．代入 \(y^{2}=2px\)：\(y^{2}=2p(my+p/2)=2pmy+p^{2}\)，整理为关于 \(y\) 的一元二次方程 \(y^{2}-2pmy-p^{2}=0\)．其二次项系数为 \(1\neq 0\)，且 \(\Delta=(2pm)^{2}+4p^{2}=4p^{2}(m^{2}+1)>0\)（\(p>0\)），故确两交点，与题设相符．由韦达定理，\(y_{1}y_{2}=-p^{2}/1=-p^{2}\)．（特位核验：\(m=0\) 通径时 \(y^{2}=p^{2}\)、\(y=\pm p\)、积 \(-p^{2}\) \ding{51}；\(m=1\)、\(p=2\) 时 \(y^{2}-4y-4=0\)、积 \(-4=-p^{2}\) \ding{51}，与例 13 之 \(y_{1}y_{2}=-4\) 同．）}
\fi
\end{ansblock}

\begin{ansblock}[2章-拓-课时17-09]
% ans:2章-拓-课时17-09
\ansitem{30}{证明见详解}
\ifansbookdetail
\ansnote{详解}{双曲线分两种标准形，证法同，先取 \(C\)：\(x^{2}/a^{2}-y^{2}/b^{2}=1\)（\(a>0,b>0\)），其渐近线为 \(y=\pm(b/a)x\)．题设「\(l\) 的斜率与渐近线的斜率相等」即 \(l\) 与某条渐近线平行或重合，不妨设 \(l\)：\(y=(b/a)x+m\)（\(m\) 为任意实数，\(m=0\) 时 \(l\) 就是渐近线本身）．代入 \(C\)：\(x^{2}/a^{2}-((b/a)x+m)^{2}/b^{2}=1\)，而 \(((b/a)x+m)^{2}/b^{2}=(x/a+m/b)^{2}=x^{2}/a^{2}+2mx/(ab)+m^{2}/b^{2}\)，故二次项 \(x^{2}/a^{2}\) 相消，得 \(-2mx/(ab)-m^{2}/b^{2}=1\)，即 \(-(2m/(ab))x=1+m^{2}/b^{2}\)．\par\noindent ① \(m\neq 0\)：一次项系数 \(-2m/(ab)\neq 0\)，方程有唯一解 \(x=-ab(1+m^{2}/b^{2})/(2m)\)，即 \(l\) 与 \(C\) 恰一个公共点．\par\noindent ② \(m=0\)：等式左端为 0、右端为 1，矛盾，无解，即渐近线与 \(C\) 没有公共点．\par\noindent 两种情形公共点均不超过一个，故「最多只有一个公共点」．焦点在 \(y\) 轴上时（\(C\)：\(y^{2}/a^{2}-x^{2}/b^{2}=1\)，渐近线 \(y=\pm(a/b)x\)）同法：设 \(y=(a/b)x+m\) 代入并以 \(x\) 为主元，二次项 \(x^{2}/b^{2}-x^{2}/b^{2}=0\) 相消，化为一次方程或矛盾式，结论同上．（本条与例 4、变式 4 之降次情形同一事理，可作其原理说明．）}
\fi
\end{ansblock}

\begin{ansblock}[2章-拓-课时17-10]
% ans:2章-拓-课时17-10
\ansitem{31}{三条：x=0；y=p；y=(1/2)x+p（即 x−2y+2p=0）}
\ifansbookdetail
\ansnote{详解}{分「斜率不存在」「斜率为 0」「斜率存在且非 0」三类穷举．\par\noindent ① 斜率不存在：直线 \(x=0\)（即 \(y\) 轴），代 \(y^{2}=2p\times 0=0\) 得 \(y=0\)，唯一公共点 \((0,0)\)（顶点），符合．\par\noindent ② 斜率为 0：直线 \(y=p\)（过 \(A(0,p)\) 且平行对称轴 \(x\) 轴），代 \(p^{2}=2px\) 得 \(x=p/2\)，唯一公共点 \((p/2,p)\)，符合——此即例 2／变式 2 的退化型「一公共点而不切」．\par\noindent ③ 斜率存在且非 0：设 \(y=kx+p\)（\(k\neq 0\)），代 \((kx+p)^{2}=2px\)，得 \(k^{2}x^{2}+(2kp-2p)x+p^{2}=0\)（二次项系数 \(k^{2}\neq 0\)）．只有一个公共点 \(\iff\Delta=0\)：\(\Delta=(2kp-2p)^{2}-4k^{2}p^{2}=4p^{2}[(k-1)^{2}-k^{2}]=4p^{2}(1-2k)=0\)，故 \(k=1/2\)，切线为 \(y=(1/2)x+p\)，唯一公共点为 \(x=(2p-2kp)/(2k^{2})=(2p-p)/(1/2)=2p\)、\(y=2p\)，即切点 \((2p,2p)\)（校：\((2p)^{2}=4p^{2}=2p\times 2p\) \ding{51}）．\par\noindent 综上满足条件的直线共三条：\(x=0\)、\(y=p\)、\(y=(1/2)x+p\)．（几何校：\(A(0,p)\) 在抛物线外（\(p^{2}>0=2p\times 0\)），故恰一条切线；另两条分别由「与对称轴平行」「过顶点的竖直割线」两个退化方向提供，条数与读数合．）}
\fi
\end{ansblock}

\begin{ansblock}[2章-拓-课时17-14]
% ans:2章-拓-课时17-14
\ansitem{32}{证明见详解}
\ifansbookdetail
\ansnote{详解}{取抛物线标准形 \(y^{2}=2px\)（\(p>0\)），顶点 \(O(0,0)\)、焦点 \(F(p/2,0)\)、准线 \(x=-p/2\)、对称轴为 \(x\) 轴．设 \(P(x_{1},y_{1})\)、\(Q(x_{2},y_{2})\)，\(y_{1}y_{2}\neq 0\)（焦点弦的端点不为顶点，且由例 17 有 \(y_{1}y_{2}=-p^{2}\neq 0\)）．\par\noindent 直线 \(OP\) 的斜率 \(=y_{1}/x_{1}\)，而 \(y_{1}^{2}=2px_{1}\) 故 \(x_{1}=y_{1}^{2}/(2p)\)，斜率 \(=y_{1}/(y_{1}^{2}/(2p))=2p/y_{1}\)，即 \(OP\)：\(y=(2p/y_{1})x\)．\par\noindent 令 \(x=-p/2\)（准线），得 \(M\) 的纵坐标 \(y_{M}=(2p/y_{1})(-p/2)=-p^{2}/y_{1}\)．\par\noindent 由 \(y_{1}y_{2}=-p^{2}\) 得 \(y_{2}=-p^{2}/y_{1}\)，于是 \(y_{M}=y_{2}\)．\par\noindent \(M(-p/2,y_{M})\) 与 \(Q(x_{2},y_{2})\) 纵坐标相同，而横坐标不同（\(x_{2}\geqslant 0>-p/2\)），故直线 \(MQ\) 为水平线，即 \(MQ\) 平行于抛物线的对称轴 \(x\) 轴．（通径特位：\(PQ\perp x\) 轴时 \(x_{1}=x_{2}=p/2\)、\(y_{2}=-y_{1}\)，由 \(y_{1}^{2}=2p\times(p/2)=p^{2}\) 得 \(y_{M}=-p^{2}/y_{1}=-y_{1}=y_{2}\) \ding{51} 证法不变；焦点在 \(y\) 轴、开口向左等其他标准形同法（换元即可），结论为「\(MQ\) 平行于该抛物线的对称轴」．）}
\fi
\end{ansblock}

\begin{ansblock}[2章-拓-课时17-15]
% ans:2章-拓-课时17-15
\ansitem{33}{证明见详解}
\ifansbookdetail
\ansnote{详解}{取标准形 \(y^{2}=2px\)（\(p>0\)），准线 \(l\)：\(x=-p/2\)．设 \(P_{1}(x_{1},y_{1})\)、\(P_{2}(x_{2},y_{2})\)，弦 \(P_{1}P_{2}\) 的中点为 \(M\)，则圆心为 \(M\)、半径 \(r=|P_{1}P_{2}|/2\)．\par\noindent 由抛物线定义，焦半径等于点到准线的距离：\(|P_{1}F|=x_{1}+p/2\)、\(|P_{2}F|=x_{2}+p/2\)．因 \(F\) 在线段 \(P_{1}P_{2}\) 上（焦点弦过焦点），\(|P_{1}P_{2}|=|P_{1}F|+|P_{2}F|=x_{1}+x_{2}+p\)，故 \(r=(x_{1}+x_{2}+p)/2\)．\par\noindent 又圆心 \(M\) 的横坐标为 \((x_{1}+x_{2})/2\)，\(M\) 到准线 \(x=-p/2\) 的距离 \(d=(x_{1}+x_{2})/2+p/2=(x_{1}+x_{2}+p)/2\)．\par\noindent 于是 \(d=r\)，即圆心到准线的距离恰等于半径，故以 \(P_{1}P_{2}\) 为直径的圆与准线相切．（另一表述：分别过 \(P_{1}\)、\(P_{2}\)、\(M\) 向准线作垂线，垂足为 \(H_{1}\)、\(H_{2}\)、\(H\)，则 \(|MH|\) 是直角梯形 \(P_{1}H_{1}H_{2}P_{2}\) 的中位线，\(|MH|=(|P_{1}F|+|P_{2}F|)/2=|P_{1}P_{2}|/2=r\)，同得相切——此即「中位线＝焦半径和之半」结构，例 14 方法二同族．）}
\fi
\end{ansblock}

\begin{ansblock}[2章-拓-课时17-16]
% ans:2章-拓-课时17-16
\ansitem{34}{(x−3)²+(y−2)²=16 或 (x−11)²+(y+6)²=144（两圆）}
\ifansbookdetail
\ansnote{详解}{焦点 \(F(1,0)\)、准线 \(x=-1\)，直线 \(l\)：\(y=x-1\)．代入 \((x-1)^{2}=4x\) 得 \(x^{2}-6x+1=0\)（\(\Delta=32>0\)），设 \(A(x_{1},y_{1})\)、\(B(x_{2},y_{2})\)，则 \(x_{1}+x_{2}=6\)、\(x_{1}x_{2}=1\)，\(x=3\pm 2\sqrt{2}\)．\par\noindent 先定必备数据：弦的中点 \(M=(3,2)\)；\(|AB|=\sqrt{2}\,|x_{1}-x_{2}|=\sqrt{2}\times\sqrt{36-4}=\sqrt{2}\times\sqrt{32}=8\)，故 \(|MA|=|MB|=4\)．\par\noindent 所求圆过 \(A\)、\(B\)，其圆心在 \(AB\) 的中垂线上：中垂线斜率 \(-1\)、过 \(M\)，即 \(y=-x+5\)，设圆心 \(C(t,5-t)\)．\par\noindent 圆的半径按「与准线 \(x=-1\) 相切」为 \(r=|t+1|\)；按「过 \(A\)」为 \(r^{2}=|CA|^{2}=|CM|^{2}+|MA|^{2}=2(t-3)^{2}+16\)（因 \(|CM|^{2}=(t-3)^{2}+(5-t-2)^{2}=2(t-3)^{2}\)）．\par\noindent 令两式相等：\((t+1)^{2}=2(t-3)^{2}+16\)，展开 \(t^{2}+2t+1=2t^{2}-12t+18+16\)，移项得 \(t^{2}-14t+33=0\)，\((t-3)(t-11)=0\)，故 \(t=3\) 或 \(t=11\)（均使 \(r=|t+1|>0\) \ding{51}，两解皆合）．\par\noindent \(t=3\)：圆心 \((3,2)=M\)、\(r=4\)，即 \((x-3)^{2}+(y-2)^{2}=16\)——此圆恰为以焦点弦 \(AB\) 为直径的圆，与例 19（以焦点弦为直径的圆与准线相切）相符．\par\noindent \(t=11\)：圆心 \((11,-6)\)、\(r=12\)，即 \((x-11)^{2}+(y+6)^{2}=144\)．\par\noindent 故满足条件的圆有两个：\((x-3)^{2}+(y-2)^{2}=16\) 与 \((x-11)^{2}+(y+6)^{2}=144\)．（验第二圆：取 \(A(3+2\sqrt{2},2+2\sqrt{2})\)，\(|CA|^{2}=(8-2\sqrt{2})^{2}+(-8-2\sqrt{2})^{2}=144=r^{2}\) \ding{51}；圆心到准线距离 \(11+1=12=r\) \ding{51}．）}
\fi
\end{ansblock}

\begin{ansblock}[2章-拓-课时17-20]
% ans:2章-拓-课时17-20
\ansitem{35}{D}
\ifansbookdetail
\ansnote{详解}{若 \(l\) 斜率不存在（\(x=1\)）：代入得 \(2-y^{2}=2\)，\(y=0\)，只一交点，不合「交于两点且 \(P\) 为中点」；故设 \(l\)：\(y=k(x-1)+2\)．代入 \(2x^{2}-y^{2}=2\)：\((2-k^{2})x^{2}+2k(k-2)x-(k^{2}-4k+6)=0\)（其中一次项系数由 \(-2k(2-k)\) 化为 \(2k(k-2)\)、常数项为 \(-(2-k)^{2}-2=-(k^{2}-4k+6)\)）．须 \(2-k^{2}\neq 0\)．\par\noindent 中点条件：\(x_{1}+x_{2}=2\) \(\Rightarrow -2k(k-2)/(2-k^{2})=2\) \(\Rightarrow -2k^{2}+4k=4-2k^{2}\) \(\Rightarrow 4k=4\) \(\Rightarrow k=1\)（合 \(2-k^{2}=1\neq 0\) \ding{51}）．\par\noindent 此时方程为 \(x^{2}-2x-3=0\)，\(x_{1}+x_{2}=2\)、\(x_{1}x_{2}=-3\)，两根 \(3\)、\(-1\)，\(\Delta=4+12=16>0\) \ding{51} 两交点 \(M(3,4)\)、\(N(-1,0)\)（校：\(2\times 9-16=2\) \ding{51}；\(2\times 1-0=2\) \ding{51}）．\par\noindent \(|MN|=\sqrt{1+k^{2}}\,|x_{1}-x_{2}|=\sqrt{2}\times\sqrt{(x_{1}+x_{2})^{2}-4x_{1}x_{2}}=\sqrt{2}\times\sqrt{4+12}=\sqrt{2}\times 4=4\sqrt{2}\)，选 D．}
\fi
\end{ansblock}

\begin{ansblock}[2章-拓-课时17-21]
% ans:2章-拓-课时17-21
\ansitem{36}{A（m=−3）}
\ifansbookdetail
\ansnote{详解}{圆 \(C\) 配方法：\((x+1)^{2}+(y+2)^{2}=5-m\)（须 \(m<5\)），圆心 \((-1,-2)\)、半径 \(r=\sqrt{5-m}\)．\(AB\) 是圆 \(C\) 的直径，故圆心即 \(AB\) 的中点：中点 \((-1,-2)\)．\par\noindent 求弦所在直线（按常规联立韦达，不引中点弦公式名目）：设 \(l\) 斜率为 \(k\)（若 \(l\) 为 \(x=-1\)，代入双曲线得 \(1-y^{2}/2=1\)、\(y=0\)，只一交点，不合），以 \(u=x+1\) 平移（则 \(x=u-1\)、\(y=ku-2\)，中点条件即 \(u_{1}+u_{2}=0\)）．代入 \(2x^{2}-y^{2}=2\)：\(2(u-1)^{2}-(ku-2)^{2}=2\to 2u^{2}-4u+2-k^{2}u^{2}+4ku-4=2\)，即 \((2-k^{2})u^{2}+(4k-4)u-4=0\)．由 \(u_{1}+u_{2}=0\) 得 \(-(4k-4)/(2-k^{2})=0\)，故 \(k=1\)（\(2-k^{2}=1\neq 0\) \ding{51}）．\par\noindent 于是 \(l\)：\(y=(x+1)-2=x-1\)，方程化为 \(u^{2}-4=0\)，\(u=\pm 2\to A(1,0)\)、\(B(-3,-4)\)（校：\(1-0=1\) \ding{51}；\(9-16/2=1\) \ding{51}），\(|AB|=\sqrt{4^{2}+4^{2}}=4\sqrt{2}\)．\par\noindent 由直径 \(|AB|^{2}=D^{2}+E^{2}-4F\)（此处 \(D=2\)、\(E=4\)、\(F=m\)）：\(|AB|^{2}=4+16-4m=20-4m\)，令 \((4\sqrt{2})^{2}=32=20-4m\)，解得 \(m=-3\)（\(<5\) \ding{51} 圆确存在）．故选 A．（互核：\(r=|AB|/2=2\sqrt{2}\)、\(r^{2}=8=5-m\Rightarrow m=-3\) \ding{51}；又 \(|CA|^{2}=(1+1)^{2}+(0+2)^{2}=8=r^{2}\) \ding{51}．）}
\fi
\end{ansblock}

\begin{ansblock}[2章-拓-课时17-02]
% ans:2章-拓-课时17-02
\ansitem{37}{±1}
\ifansbookdetail
\ansnote{详解}{联立消 \(y\)：\(x^{2}/2+(kx+1)^{2}=1\)，乘 2 得 \(x^{2}+2(k^{2}x^{2}+2kx+1)=2\)，即 \((1+2k^{2})x^{2}+4kx=0\)．二次项系数 \(1+2k^{2}>0\) 恒成立．设 \(M(x_{1},y_{1})\)、\(N(x_{2},y_{2})\)，则 \(x_{1}+x_{2}=-4k/(1+2k^{2})\)、\(x_{1}x_{2}=0\)（事实上直线过定点 \((0,1)\)，即椭圆上顶点，一交点恒为它）．\par\noindent 由 \(|MN|=4\sqrt{2}/3\) 得 \((x_{1}-x_{2})^{2}+(y_{1}-y_{2})^{2}=32/9\)，即 \((1+k^{2})(x_{1}-x_{2})^{2}=32/9\)；而 \((x_{1}-x_{2})^{2}=(x_{1}+x_{2})^{2}-4x_{1}x_{2}=16k^{2}/(1+2k^{2})^{2}\)，代入得 \((1+k^{2})\times 16k^{2}/(1+2k^{2})^{2}=32/9\)，约 16 并交叉相乘：\(9k^{2}(1+k^{2})=2(1+2k^{2})^{2}\)，展开 \(9k^{2}+9k^{4}=2+8k^{2}+8k^{4}\)，即 \(k^{4}+k^{2}-2=0\)，\((k^{2}+2)(k^{2}-1)=0\)，故 \(k^{2}=1\)、\(k=\pm 1\)．（回验 \(k=1\)：\((1+2)x^{2}+4x=0\to x=0\)、\(-4/3\)，\(|MN|=\sqrt{2}\times 4/3=4\sqrt{2}/3\) \ding{51}．）}
\fi
\end{ansblock}

\begin{ansblock}[2章-导-课时17-G1]
% ans:2章-导-课时17-G1
\ansitem{38}{B}
\ifansbookdetail
\ansnote{详解}{焦点 \(F(3/2,0)\)．所求直线过焦点且与抛物线交于两点，可设 \(x=my+3/2\)（\(m\in R\)，\(m=0\) 即竖直线）．代入 \(y^{2}=6x\)：\(y^{2}-6my-9=0\)，\(\Delta=36m^{2}+36>0\) 恒成立，恒交于两点．\(x_{1}+x_{2}=m(y_{1}+y_{2})+3=m\cdot 6m+3=6m^{2}+3=3\)，得 \(m^{2}=0\)、\(m=0\)．故直线唯一：\(x=3/2\)，条数为 1，选 B．（验算：\(x=3/2\) 代入 \(y^{2}=6x\) 得 \(y^{2}=9\)，两交点 \(A(3/2,3)\)、\(B(3/2,-3)\)，\(x_{1}+x_{2}=3\) \ding{51}（此即通径位置）．）}
\fi
\end{ansblock}

\begin{ansblock}[2章-导-课时17-G2]
% ans:2章-导-课时17-G2
\ansitem{39}{C}
\ifansbookdetail
\ansnote{详解}{\(a^{2}=2\)，\(2a=2\sqrt{2}\)．\(A\) 在左支：\(|AF_{2}|-|AF_{1}|=2a\)；\(B\) 在右支：\(|BF_{1}|-|BF_{2}|=2a\)．设 \(|AF_{2}|=|BF_{2}|=t\)，则 \(|AF_{1}|=t-2a\)、\(|BF_{1}|=t+2a\)．\(F_{1}(-\sqrt{6},0)\) 在左支顶点 \((-\sqrt{2},0)\) 的左方（位于左支凹侧内），过 \(F_{1}\) 的弦自 \(F_{1}\) 出发先遇左支点 \(A\)、再遇右支点 \(B\)，故 \(|AB|=|BF_{1}|-|AF_{1}|=4a=4\sqrt{2}\)．在 \(\mathrm{Rt}\triangle AF_{2}B\)（\(\angle AF_{2}B=90^{\circ}\)）中：\(|AF_{2}|^{2}+|BF_{2}|^{2}=|AB|^{2}\)，即 \(2t^{2}=(4\sqrt{2})^{2}=32\)，\(t^{2}=16\)，\(t=4\)（\(t>0\)）．故 \(|AF_{2}|=4\)，选 C．（验算：\(t=4\)：\(|AF_{1}|=4-2\sqrt{2}>0\) \ding{51}；\(|AB|=|BF_{1}|-|AF_{1}|=(4+2\sqrt{2})-(4-2\sqrt{2})=4\sqrt{2}\) \ding{51}；勾股 \(2\times 4^{2}=32=(4\sqrt{2})^{2}\) \ding{51}．）}
\fi
\end{ansblock}

\begin{ansblock}[2章-导-课时17-G3]
% ans:2章-导-课时17-G3
\ansitem{40}{C}
\ifansbookdetail
\ansnote{详解}{\(a=4\)，\(2a=8\)．\(F_{2}\) 在椭圆内部，弦 \(AB\) 过 \(F_{2}\)，故 \(F_{2}\) 在 \(A\)、\(B\) 之间，\(|AF_{2}|+|BF_{2}|=|AB|\)．\(\triangle ABF_{1}\) 的周长 \(|AB|+|AF_{1}|+|BF_{1}|=(|AF_{1}|+|AF_{2}|)+(|BF_{1}|+|BF_{2}|)=2a+2a=16\)（椭圆定义，与弦的位置无关）．由 \(|AB|+|AF_{1}|=10\) 得 \(|BF_{1}|=16-10=6\)，选 C．（验算：\(|BF_{1}|=6\) 落在焦半径范围 \([a-c,a+c]=[4-\sqrt{7},4+\sqrt{7}]\approx[1.35,6.65]\) 内 \ding{51}；周长恒等 \(4a\) 不依赖过焦弦的方向 \ding{51}．）}
\fi
\end{ansblock}

\begin{ansblock}[2章-导-课时17-G4]
% ans:2章-导-课时17-G4
\ansitem{41}{8}
\ifansbookdetail
\ansnote{详解}{设 \(A(x_{1},y_{1})\)、\(B(x_{2},y_{2})\)（\(y_{1}>0>y_{2}\)），直线 \(x=my+p/2\) 代入 \(y^{2}=2px\)：\(y^{2}-2pmy-p^{2}=0\)，得 \(y_{1}y_{2}=-p^{2}\)．由抛物线定义及 \(y_{1}^{2}=2px_{1}\)：\(|AF|=x_{1}+p/2=(y_{1}^{2}+p^{2})/(2p)\)，同理 \(|FB|=(y_{2}^{2}+p^{2})/(2p)\)．\(|AF|=2|FB|\) 即 \(y_{1}^{2}+p^{2}=2(y_{2}^{2}+p^{2})\)，得 \(y_{1}^{2}=2y_{2}^{2}+p^{2}\)．两边乘 \(y_{2}^{2}\) 并以 \((y_{1}y_{2})^{2}=p^{4}\) 代入：\(p^{4}=2y_{2}^{4}+p^{2}y_{2}^{2}\)．令 \(u=y_{2}^{2}\)：\(2u^{2}+p^{2}u-p^{4}=0\)，因式分解 \((2u-p^{2})(u+p^{2})=0\)，\(u>0\) 得 \(u=p^{2}/2\)，即 \(y_{2}^{2}=p^{2}/2\)（此时 \(y_{2}=-p/\sqrt{2}\)、\(y_{1}=\sqrt{2}p=-2y_{2}\) \ding{51}）．\(C(-p/2,0)\)、\(F(p/2,0)\)，\(|CF|=p\)，\(B\) 到 \(x\) 轴距离 \(|y_{2}|=p/\sqrt{2}\)，故 \(S_{\triangle CBF}=\frac{1}{2}\cdot p\cdot p/\sqrt{2}=p^{2}/(2\sqrt{2})=4\sqrt{2}\)，得 \(p^{2}=16\)、\(p=4\)，通径长 \(2p=8\)．（验算：\(p=4\)：\(C(-2,0)\)、\(F(2,0)\)、\(|y_{2}|=2\sqrt{2}\)，\(S=\frac{1}{2}\times 4\times 2\sqrt{2}=4\sqrt{2}\) \ding{51}；\(|AF|=(32+16)/8=6\)、\(|FB|=(8+16)/8=3\)，\(6=2\times 3\) \ding{51}．）}
\fi
\end{ansblock}

\begin{ansblock}[2章-导-课时17-G5]
% ans:2章-导-课时17-G5
\ansitem{42}{y=x−1}
\ifansbookdetail
\ansnote{详解}{\(P(2,1)\)：\(4/2-1=1\) \ding{51}，\(P\) 在 \(C\) 上．设切线 \(l\)：\(y-1=k(x-2)\)，即 \(y=kx+(1-2k)\)．代入 \(x^{2}/2-y^{2}=1\)（乘 2）：\(x^{2}-2(kx+1-2k)^{2}=2\)，整理得 \((1-2k^{2})x^{2}-4k(1-2k)x-2[(1-2k)^{2}+1]=0\)．若 \(1-2k^{2}=0\)（\(k=\pm\sqrt{2}/2\)），方程降为一次，\(l\) 平行于渐近线 \(y=\pm(\sqrt{2}/2)x\) 而与 \(C\) 仅一个公共点，非切线，排除．故 \(1-2k^{2}\neq 0\)，\(l\) 与 \(C\) 相切即 \(\Delta=0\)：\par\noindent \(\Delta=16k^{2}(1-2k)^{2}+8(1-2k^{2})[(1-2k)^{2}+1]=0\)，除以 8 并展开（\((1-2k)^{2}=1-4k+4k^{2}\)）：\(2k^{2}(1-4k+4k^{2})+(1-2k^{2})(2-4k+4k^{2})=(2k^{2}-8k^{3}+8k^{4})+(2-4k+8k^{3}-8k^{4})=2k^{2}-4k+2=2(k-1)^{2}=0\)，\par\noindent 得 \(k=1\)（\(1-2\times 1=-2\neq 0\)，保持二次 \ding{51}）．故切线方程为 \(y=x-1\)．（验算：\(y=x-1\) 代入：\(x^{2}/2-(x-1)^{2}=1\) 即 \(-x^{2}/2+2x-2=0\)，亦 \(x^{2}-4x+4=0\)、\((x-2)^{2}=0\)，重根 \(x=2\)，切点即 \(P(2,1)\) \ding{51}——恰一公共点且不平行渐近线，为真切线；源解以导数求斜率系后置工具，本解改走联立 \(\Delta=0\)，读数一致．）}
\fi
\end{ansblock}

\jietitle{2.8 直线与圆锥曲线的位置关系（二）}

\begin{ansblock}[2章-练-课时18-01]
% ans:2章-练-课时18-01
\ansitem{1}{B（|MF₂|=5/2）}
\ifansbookdetail
\ansnote{详解}{\(F_1(-2,0)\)、\(F_2(2,0)\)、\(A(-1,0)\)，故 \(F_1A/F_1F_2=1/4\)．\(NA\parallel MF_2\)，得 \(\triangle F_1AN\sim\triangle F_1MF_2\)，故 \(F_1N/F_1M=1/4\)．设 \(M(x_0,y_0)\)（\(x_0>0\)，\(y_0>0\)），则 \(N=F_1+(1/4)(M-F_1)=((x_0-6)/4,y_0/4)\)．\(M\)、\(N\) 均在 \(E\) 上：\(x_0^{2}-y_0^{2}/3=1\)，且 \(((x_0-6)/4)^{2}-(y_0/4)^{2}/3=1\)，后者乘16得 \((x_0-6)^{2}-y_0^{2}/3=16\)；两式相减 \((x_0-6)^{2}-x_0^{2}=15\)，即 \(-12x_0+36=15\)，解得 \(x_0=7/4\)，进而 \(y_0^{2}=3(x_0^{2}-1)=99/16\)，\(y_0=3\sqrt{11}/4\)．于是 \(|MF_2|=\sqrt{(7/4-2)^{2}+99/16}=\sqrt{100/16}=5/2\)．故选：B．}
\fi
\end{ansblock}

\begin{ansblock}[2章-练-课时18-07]
% ans:2章-练-课时18-07
\ansitem{2}{(1) 4/5；(2) 恒过定点 (0,±10)}
\ifansbookdetail
\ansnote{详解}{(1) 由 \(2b=a+c\) 得 \(c=2b-a\)，代入 \(c^{2}=a^{2}-b^{2}\)：\((2b-a)^{2}=a^{2}-b^{2}\)，展开得 \(4b^{2}-4ab=0\)，故 \(b/a=4/5\)．\par\noindent (2) \(2c=12\)，故 \(c=6\)；由 \(b^{2}=16a^{2}/25\) 与 \(a^{2}-b^{2}=36\) 得 \(9a^{2}/25=36\)，故 \(a^{2}=100\)、\(b^{2}=64\)，\(A(0,8)\)．设 \(P(x_0,y_0)\)（\(x_0\neq 0\)），则 \(Q(-x_0,-y_0)\)．直线 \(AP\)：\(y=(y_0-8)x/x_0+8\)，令 \(y=0\) 得 \(M(8x_0/(8-y_0),0)\)；同理 \(N(-8x_0/(8+y_0),0)\)．以 \(MN\) 为直径的圆：\((x-8x_0/(8-y_0))(x+8x_0/(8+y_0))+y^{2}=0\)，其 \(x\) 的交叉项 \(=8x_0[(8-y_0)-(8+y_0)]/(64-y_0^{2})=-16x_0y_0/(64-y_0^{2})\)．由 \(x_0^{2}/100+y_0^{2}/64=1\) 得 \(64-y_0^{2}=64x_0^{2}/100=16x_0^{2}/25\)，故圆方程化为 \(x^{2}+y^{2}-(25y_0/x_0)x-100=0\)．令 \(x=0\) 得 \(y=\pm 10\)，与 \(P\) 无关，故恒过定点 \((0,10)\) 与 \((0,-10)\)（\(P\) 异于 \(A\) 保证 \(x_0\neq 0\)、\(64-y_0^{2}>0\)）．}
\fi
\end{ansblock}

\begin{ansblock}[2章-练-课时18-08]
% ans:2章-练-课时18-08
\ansitem{3}{(1) x²/4+y²/3=1；(2) 存在，R(4,0)}
\ifansbookdetail
\ansnote{详解}{(1) 圆 \(A\) 即 \((x+1)^{2}+y^{2}=16\)，圆心 \(A(-1,0)\)、半径4．由垂直平分线性质得 \(|NM|=|NB|\)，而 \(N\) 在 \(MA\) 上，故 \(|NA|+|NB|=|NA|+|NM|=|MA|=4>|AB|=2\)，轨迹为以 \(A\)、\(B\) 为焦点的椭圆：\(2a=4\)、\(2c=2\)，即 \(a^{2}=4\)、\(b^{2}=3\)，故 \(C\)：\(x^{2}/4+y^{2}/3=1\)．\par\noindent (2) 设 \(R(t,0)\)．联立得 \((4k^{2}+3)x^{2}-8k^{2}x+4k^{2}-12=0\)，故 \(x_1+x_2=8k^{2}/(4k^{2}+3)\)、\(x_1x_2=(4k^{2}-12)/(4k^{2}+3)\)．\(\angle ORP=\angle ORQ\) 等价于 \(k_{RP}+k_{RQ}=0\)，即 \(y_1/(x_1-t)+y_2/(x_2-t)=0\)．以 \(y_i=k(x_i-1)\) 代入，分子 \(=k[2x_1x_2-(1+t)(x_1+x_2)+2t]=0\)；\(k\neq 0\) 时即 \(2(4k^{2}-12)-8k^{2}(1+t)+2t(4k^{2}+3)=0\)，展开得 \(6t-24=0\)，故 \(t=4\)（\(k=0\) 时 \(P\)、\(Q\) 为长轴两端点，等式显然成立）．故存在 \(R(4,0)\)，与 \(k\) 无关．}
\fi
\end{ansblock}

\begin{ansblock}[2章-练-课时18-10]
% ans:2章-练-课时18-10
\ansitem{4}{B}
\ifansbookdetail
\ansnote{详解}{过焦点的弦分两型：跨支型（\(A\)、\(B\) 各在一支）弦长下限为 \(2a\)，且实轴所在直线（长 \(2a\)）恰1条，长 \(L>2a\) 时对称地有2条；单支型（同在右支）弦长下限为通径 \(2b^{2}/a=6/a\)，通径（\(l\perp x\) 轴）恰1条，\(L>6/a\) 时有2条．且 \(6/a\) 与 \(2a\) 不可能同时等于6（\(a=1\) 时 \(6/a=6\)、\(2a=2\)；\(a=3\) 时 \(2a=6\)、\(6/a=2\)）．令 \(L=6\) 分段计数：\(a<1\) 时跨支 \(6>2a\) 有2条、单支 \(6<6/a\) 有0条，合计2，符合；\(a=1\) 时 \(2+1=3\)，不合；\(1<a<3\) 时 \(2+2=4\)，不合；\(a=3\) 时 \(1+2=3\)，不合；\(a>3\) 时跨支 \(6<2a\) 有0条、单支 \(6>6/a\) 有2条，合计2，符合．故 \(a\in(0,1)\cup(3,+\infty)\)，选 B．}
\fi
\end{ansblock}

\begin{ansblock}[2章-练-课时18-11]
% ans:2章-练-课时18-11
\ansitem{5}{(1) x²/4+y²/2=1；(2) 不存在（QA·QB≡−15/16 恒<1）}
\ifansbookdetail
\ansnote{详解}{(1) \(e=\sqrt{2}/2\)，故 \(a^{2}=2c^{2}\)、\(b^{2}=c^{2}\)；三角形面积 \(=(a-c)b/2=(\sqrt{2}-1)c^{2}/2=\sqrt{2}-1\)，得 \(c^{2}=2\)，故 \(a^{2}=4\)、\(b^{2}=2\)，椭圆 \(C\)：\(x^{2}/4+y^{2}/2=1\)．\par\noindent (2) 直线过 \((1,0)\)（椭圆内一点），恒有两交点．联立得 \((1+2k^{2})x^{2}-4k^{2}x+2k^{2}-4=0\)，故 \(x_1+x_2=4k^{2}/(1+2k^{2})\)、\(x_1x_2=(2k^{2}-4)/(1+2k^{2})\)；\(y_1y_2=k^{2}(x_1-1)(x_2-1)=-3k^{2}/(1+2k^{2})\)．于是 \(QA\cdot QB=(x_1-7/4)(x_2-7/4)+y_1y_2=x_1x_2-(7/4)(x_1+x_2)+49/16+y_1y_2\)，代入化简 \(=[(2k^{2}-4)-7k^{2}-3k^{2}]/(1+2k^{2})+49/16=(-8k^{2}-4)/(1+2k^{2})+49/16=-4+49/16=-15/16\)，与 \(k\) 无关且恒小于1，故不存在这样的 \(k\)．}
\fi
\end{ansblock}

\begin{ansblock}[2章-练-课时18-09]
% ans:2章-练-课时18-09
\ansitem{6}{(1) 两条件均得 x²/4+y²/3=1；(2) DE 恒过定点 (8/5,0)}
\ifansbookdetail
\ansnote{详解}{(1) 条件①：\(e=1/2\)，故 \(a=2c\)；焦点到相应准线距离 \(a^{2}/c-c=3\)，即 \(4c-c=3\)，得 \(c=1\)，故 \(a^{2}=4\)、\(b^{2}=3\)．条件②：\(\odot M\) 圆心 \((6,0)\)、半径4，椭圆在圆内一侧外切，切点必为右顶点，故 \(a+4=6\)，\(a=2\)；\(\odot N\) 圆心 \((0,2\sqrt{3})\)、半径 \(\sqrt{3}\)，故 \(b=\sqrt{3}\)、\(b^{2}=3\)．回验唯一性（外切即恰一公共点）：\(\odot M\) 与椭圆联立消 \(y\) 得 \(x^{2}-48x+92=0\)，\(x=2\) 或 \(x=46\)（46越界，舍）；\(\odot N\) 联立得 \(y^{2}+12\sqrt{3}y-39=0\)，\(y=\sqrt{3}\) 或 \(y=-13\sqrt{3}\)（越界，舍）．两条件同得 \(x^{2}/4+y^{2}/3=1\)．\par\noindent (2) \(F(1,0)\)．设 \(A(x_1,y_1)\)（\(y_1>0\)）、\(B(-x_1,-y_1)\)，满足 \(3x_1^{2}+4y_1^{2}=12\)．直线 \(AF\)：\(x=my+1\)（\(m=(x_1-1)/y_1\)），直线 \(BF\)：\(x=ny+1\)（\(n=(x_1+1)/y_1\)），则 \(n-m=2/y_1\)．\(x=my+1\) 代入 \(3x^{2}+4y^{2}=12\)：\((3m^{2}+4)y^{2}+6my-9=0\)，故 \(y_1y_D=-9/(3m^{2}+4)\)，即 \(y_D=-9/[(3m^{2}+4)y_1]\)；同理 \(y_E=9/[(3n^{2}+4)y_1]\)，且 \(x_D=my_D+1\)、\(x_E=ny_E+1\)．\par\noindent \(D\)、\(E\)、\(P(8/5,0)\) 三点共线等价于 \(y_D(x_E-8/5)=y_E(x_D-8/5)\)，即 \(y_Dy_E(n-m)=(3/5)(y_D-y_E)\)（记为①）．代入 \(y_Dy_E=-81/[(3m^{2}+4)(3n^{2}+4)y_1^{2}]\) 与 \(y_D-y_E=-(9/y_1)(3m^{2}+3n^{2}+8)/[(3m^{2}+4)(3n^{2}+4)]\)，则①等价于 \(30/y_1^{2}=3(m^{2}+n^{2})+8\)，即 \(30=3y_1^{2}(m^{2}+n^{2})+8y_1^{2}=3[(x_1-1)^{2}+(x_1+1)^{2}]+8y_1^{2}=6x_1^{2}+6+8y_1^{2}=2(3x_1^{2}+4y_1^{2})+6=24+6=30\)，恒成立．故 \(DE\) 恒过定点 \((8/5,0)\)．}
\fi
\end{ansblock}

\begin{ansblock}[2章-练-课时18-03]
% ans:2章-练-课时18-03
\ansitem{7}{(1) y²=4x、x²/4+y²/3=1；(2) S_max=2√3/3}
\ifansbookdetail
\ansnote{详解}{(1) 焦点到准线的距离 \(p=2\)，故抛物线 \(y^{2}=4x\)，\(F(1,0)\)；椭圆 \(c=1\)、\(e=c/a=1/2\)，故 \(a=2\)、\(b^{2}=a^{2}-c^{2}=3\)，椭圆 \(x^{2}/4+y^{2}/3=1\)．\par\noindent (2) 设 \(l\)：\(x=my+n\)（不过 \(F\)，故 \(n\neq 1\)）．交抛物线：\(y^{2}-4my-4n=0\)，得 \(y_1y_2=-4n\)、\(x_1x_2=(y_1y_2)^{2}/16=n^{2}\)．由 \(OA\cdot OB=n^{2}-4n=-3\) 得 \(n=1\)（此时 \(l\) 过 \(F(1,0)\)，舍）或 \(n=3\)．交椭圆：\(3(my+3)^{2}+4y^{2}=12\)，即 \((3m^{2}+4)y^{2}+18my+15=0\)，\(\Delta=144m^{2}-240\geqslant 0\)，故 \(m^{2}\geqslant 5/3\)．\(l\) 交 \(x\) 轴于 \(T(3,0)\)，\(|TF|=2\)，故 \(S_{\triangle CDF}=|TF|\cdot|y_C-y_D|/2=|y_C-y_D|\)，\(S^{2}=(144m^{2}-240)/(3m^{2}+4)^{2}\)．令 \(u=3m^{2}+4\)（\(u\geqslant 9\)）：\(S^{2}=48/u-432/u^{2}\)；再令 \(v=1/u\in(0,1/9]\)：\(S^{2}=48v-432v^{2}\) 在 \(v=1/18\) 处取最大（\(1/18\in(0,1/9]\)），\(S^{2}\) 最大值为 \(48/18-432/324=4/3\)，故 \(S\) 的最大值为 \(2\sqrt{3}/3\)，此时 \(3m^{2}+4=18\)，\(m^{2}=14/3\)，满足 \(m^{2}\geqslant 5/3\)．}
\fi
\end{ansblock}

\begin{ansblock}[2章-练-课时18-04]
% ans:2章-练-课时18-04
\ansitem{8}{(1) d=|x₁y₂−x₂y₁|/√(x₁²+y₁²)，证毕；①S=√2；②m=−1/2（此时 S≡√2）}
\ifansbookdetail
\ansnote{详解}{(1) \(l_1\)：\(y_1x-x_1y=0\)（\(x_1=0\) 时即 \(x=0\)，公式仍适用），故 \(d=|y_1x_2-x_1y_2|/\sqrt{x_1^{2}+y_1^{2}}\)．\(O\) 同时平分 \(AB\) 与 \(CD\)，故 \(S=4S_{\triangle OAC}=4\cdot|OA|\cdot d/2=2\sqrt{x_1^{2}+y_1^{2}}\cdot|x_1y_2-x_2y_1|/\sqrt{x_1^{2}+y_1^{2}}=2|x_1y_2-x_2y_1|\)，得证．\par\noindent (2) 记 \(y_i=k_ix_i\)，则 \(x_i^{2}=1/(1+2k_i^{2})\)，\(S=2|x_1||x_2||k_2-k_1|=2|k_2-k_1|/\sqrt{(1+2k_1^{2})(1+2k_2^{2})}\)．\par\noindent ① \(k_1k_2=-1/2\)：分母 \(=(1+2k_1^{2})(1+2k_2^{2})=1+2(k_1^{2}+k_2^{2})+4k_1^{2}k_2^{2}=2(1+k_1^{2}+k_2^{2})\)；\((k_2-k_1)^{2}=k_1^{2}+k_2^{2}-2k_1k_2=1+k_1^{2}+k_2^{2}\)，故 \(S^{2}=4(1+k_1^{2}+k_2^{2})/[2(1+k_1^{2}+k_2^{2})]=2\)，即 \(S=\sqrt{2}\)．\par\noindent ② 记 \(T=k_1^{2}+k_2^{2}\)，则 \((k_2-k_1)^{2}=T-2m\)、分母 \(=2T+1+4m^{2}\)，\(S^{2}=4(T-2m)/(2T+1+4m^{2})\)．要与 \(T\) 无关，须存在 \(\lambda\) 使 \(T-2m\equiv\lambda(2T+1+4m^{2})\)：比较系数得 \(2\lambda=1\)、\(\lambda(1+4m^{2})=-2m\)，故 \((1+4m^{2})/2=-2m\)，即 \(4m^{2}+4m+1=0\)，解得 \(m=-1/2\)（与①自洽，此时对一切 \(T\) 有 \(S^{2}=2\)）．}
\fi
\end{ansblock}

\begin{ansblock}[2章-练-课时18-05]
% ans:2章-练-课时18-05
\ansitem{9}{(1) x²/4+y²/3=1；(2) 存在，t=−2/3，定值 1/2}
\ifansbookdetail
\ansnote{详解}{(1) \(y^{2}=4x\) 的焦点为 \((1,0)\)，故 \(c=1\)；长轴 \(2a=4\)，故 \(a=2\)、\(b^{2}=3\)，椭圆 \(E\)：\(x^{2}/4+y^{2}/3=1\)．\par\noindent (2) \(l\)：\(y=k(x-1)\)（\(k\neq 0\)）．交椭圆：\((3+4k^{2})x^{2}-8k^{2}x+4k^{2}-12=0\)，故 \(|AB|=\sqrt{1+k^{2}}\cdot\sqrt{(x_1+x_2)^{2}-4x_1x_2}=12(1+k^{2})/(3+4k^{2})\)．交抛物线：\(k^{2}x^{2}-(2k^{2}+4)x+k^{2}=0\)，故 \(x_3+x_4=(2k^{2}+4)/k^{2}\)，\(|MN|=x_3+x_4+2=4(1+k^{2})/k^{2}\)（焦点弦）．于是 \(2/|AB|+t/|MN|=(3+4k^{2})/(6(1+k^{2}))+tk^{2}/(4(1+k^{2}))=[(8+3t)k^{2}+6]/(12(1+k^{2}))\)，与 \(k\) 无关等价于 \(8+3t=6\)，故 \(t=-2/3\)，此时定值为 \(6/12=1/2\)．}
\fi
\end{ansblock}

\begin{ansblock}[2章-练-课时18-06]
% ans:2章-练-课时18-06
\ansitem{10}{(1) y=1 或 x=0 或 y=−x+1；(2) 仅 k₁+k₂=−4 为定值，其余三个均不是}
\ifansbookdetail
\ansnote{详解}{(1) 分三类：与对称轴（\(x\) 轴）平行的直线 \(y=1\)，恰一公共点；过顶点的切线 \(x=0\)；斜线 \(y=kx+1\)（\(k\neq 0\)）代入得 \(k^{2}x^{2}+(2k+4)x+1=0\)，令 \(\Delta=(2k+4)^{2}-4k^{2}=16k+16=0\)，得 \(k=-1\)，即 \(y=-x+1\)．共三条．\par\noindent (2) 有两交点，故 \(k\) 存在、\(k\neq 0\) 且 \(\Delta=16k+16>0\)，即 \(k>-1\)．联立 \(k^{2}x^{2}+(2k+4)x+1=0\)，得 \(x_1+x_2=-(2k+4)/k^{2}\)、\(x_1x_2=1/k^{2}\)．\(k_i=y_i/x_i=k+1/x_i\)：\(k_1+k_2=2k+(x_1+x_2)/(x_1x_2)=2k-(2k+4)=-4\)，为定值；\(k_1k_2=y_1y_2/(x_1x_2)=-4k\)，随 \(k\) 变；\(|k_1-k_2|=|x_2-x_1|/|x_1x_2|=4\sqrt{k+1}\)，随 \(k\) 变；\(k_1/k_2\) 同理随 \(k\) 变．故仅 \(k_1+k_2=-4\) 为定值．}
\fi
\end{ansblock}

\begin{ansblock}[2章-练-课时18-15]
% ans:2章-练-课时18-15
\ansitem{11}{(1) x²/4+y²=1；(2) 存在，min=2，此时 x−√2y−√3=0 或 x+√2y−√3=0}
\ifansbookdetail
\ansnote{详解}{(1) 周长 \(=|AB|+|AF_2|+|BF_2|=4a=8\)，故 \(a=2\)；\(e=\sqrt{3}/2\)，故 \(c=\sqrt{3}\)、\(b^{2}=1\)，椭圆 \(E\)：\(x^{2}/4+y^{2}=1\)．\par\noindent (2) 由相切得 \(d(O,l_2)=|m|/\sqrt{1+k^{2}}=1\)，即 \(m^{2}=1+k^{2}\)．焦半径公式（\(|x_i|\leqslant 2\) 保证非负）：\(|MF_2|=a-ex_M=2-(\sqrt{3}/2)x_M\)，同理 \(|NF_2|=2-(\sqrt{3}/2)x_N\)，故和 \(=4-(\sqrt{3}/2)(x_M+x_N)\)．联立 \(y=kx+m\) 与 \(x^{2}+4y^{2}=4\)：\((1+4k^{2})x^{2}+8kmx+4m^{2}-4=0\)，故 \(x_M+x_N=-8km/(1+4k^{2})\)；\(km<0\)，故 \(x_M+x_N=8\sqrt{k^{2}(1+k^{2})}/(1+4k^{2})\)．设 \(u=k^{2}>0\)，则 \(x_M+x_N=8\sqrt{u(1+u)}/(1+4u)\)；对 \(u(1+u)/(1+4u)^{2}\) 求导，分子正比于 \((1+2u)(1+4u)-8(u+u^{2})=1-2u\)，故 \(u=1/2\) 时最大，\(x_M+x_N\) 的最大值为 \(8\cdot(\sqrt{3}/2)/3=4\sqrt{3}/3\)，和的最小值为 \(4-(\sqrt{3}/2)\cdot(4\sqrt{3}/3)=2\)．此时 \(k^{2}=1/2\)、\(m^{2}=3/2\) 且 \(km<0\)，故 \((k,m)=(\sqrt{2}/2,-\sqrt{6}/2)\) 或 \((-\sqrt{2}/2,\sqrt{6}/2)\)，即 \(l_2\)：\(x-\sqrt{2}y-\sqrt{3}=0\) 或 \(x+\sqrt{2}y-\sqrt{3}=0\)．}
\fi
\end{ansblock}

\begin{ansblock}[2章-练-课时18-02]
% ans:2章-练-课时18-02
\ansitem{12}{CD}
\ifansbookdetail
\ansnote{详解}{\(F(2,0)\)，准线 \(x=-2\)，故 \(M(-2,0)\)，\(l\)：\(y=k(x+2)\)．联立得 \(k^{2}x^{2}+(4k^{2}-8)x+4k^{2}=0\)；由两交点得 \(k^{2}\neq 0\) 且 \(\Delta=64(1-k^{2})>0\)，即 \(-1<k<1\) 且 \(k\neq 0\)，A 漏 \(k\neq 0\)，错．\par\noindent 改以 \(y\) 为主元：\(x=y^{2}/8\) 代入 \(l\) 得 \(ky^{2}-8y+16k=0\)，故 \(y_1+y_2=8/k\)、\(y_1y_2=16\)；而 \(x_1x_2=4\)，故 \(8x_1x_2=32\neq 16\)，B 错．\par\noindent \(FA\cdot FB=(x_1-2)(x_2-2)+y_1y_2=x_1x_2-2(x_1+x_2)+4+y_1y_2\)，其中 \(x_1+x_2=(8-4k^{2})/k^{2}\)，代入得 \(FA\cdot FB=32-16/k^{2}\)，当 \(k^{2}=1/2\) 时为0，故 \(\angle AFB\) 可为直角，C 对．\par\noindent \(S=|MF|\cdot|y_1-y_2|/2=2\sqrt{(y_1+y_2)^{2}-4y_1y_2}=2\sqrt{64/k^{2}-64}\)；\(k^{2}=1/2\) 时 \(|y_1-y_2|=8\)，故 \(S=16\)，D 对．故选：CD．}
\fi
\end{ansblock}

\begin{ansblock}[2章-练-课时18-12]
% ans:2章-练-课时18-12
\ansitem{13}{(1) y²=8x；(2) 存在，t=±√2/2}
\ifansbookdetail
\ansnote{详解}{(1) 设 \(y^{2}=ax\)，准线 \(x=-a/4=-2\)，得 \(a=8\)，故抛物线 \(y^{2}=8x\)．\par\noindent (2) 联立 \(x=ty+1\)：\(y^{2}-8ty-8=0\)，得 \(y_1+y_2=8t\)、\(y_1y_2=-8\)；\(x_1+x_2=t(y_1+y_2)+2=8t^{2}+2\)、\(x_1x_2=t^{2}y_1y_2+t(y_1+y_2)+1=1\)．设 \(C(x_0,0)\)（\(x_0<0\)）．\(C\) 在以 \(AB\) 为直径的圆上，故 \(\angle ACB=90^{\circ}\)，即 \(CA\cdot CB=0\)：\((x_1-x_0)(x_2-x_0)+y_1y_2=0\)，得 \(x_0^{2}-(8t^{2}+2)x_0+(x_1x_2+y_1y_2)=0\)，即 \(x_0^{2}-(8t^{2}+2)x_0-7=0\)（记为①）．内心在 \(x\) 轴上等价于 \(x\) 轴平分 \(\angle ACB\)，即 \(k_{CA}+k_{CB}=0\)：\(y_1(x_2-x_0)+y_2(x_1-x_0)=0\)，得 \(2ty_1y_2+(y_1+y_2)(1-x_0)=0\)，即 \(-16t+8t(1-x_0)=0\)，故 \(8t(-1-x_0)=0\)．\(t=0\) 时直线 \(x=1\) 垂直于 \(x\) 轴，与题设矛盾，舍；故 \(x_0=-1<0\)，代入①：\(1+(8t^{2}+2)-7=0\)，故 \(t^{2}=1/2\)，\(t=\pm\sqrt{2}/2\)．}
\fi
\end{ansblock}

\begin{ansblock}[2章-练-课时18-13]
% ans:2章-练-课时18-13
\ansitem{14}{(1) x²/2+y²=1；(2) y=(2−√2)/2·(x+2) 或 y=0}
\ifansbookdetail
\ansnote{详解}{(1) 焦距为2，故 \(c=1\)、\(a^{2}=b^{2}+1\)；点 \((1,\sqrt{2}/2)\) 代入：\(1/a^{2}+(1/2)/b^{2}=1\)．令 \(u=b^{2}\)：\(1/(u+1)+1/(2u)=1\)，即 \(2u+(u+1)=2u(u+1)\)，整理得 \(2u^{2}-u-1=0\)，解得 \(u=1\)，故 \(a^{2}=2\)、\(b^{2}=1\)，椭圆 \(C\)：\(x^{2}/2+y^{2}=1\)．\par\noindent (2) 先取 \(y=0\)：\(A\)、\(B\) 为 \((\pm\sqrt{2},0)\)，\(G\) 在 \(y\) 轴上，\(|GA|=|GB|\) 成立，是一解．\(k\neq 0\) 时设 \(l\)：\(y=k(x+2)\)，联立 \(x^{2}+2y^{2}=2\)：\((1+2k^{2})x^{2}+8k^{2}x+8k^{2}-2=0\)，弦中点 \(M\)：\(x_M=-4k^{2}/(1+2k^{2})\)、\(y_M=k(x_M+2)=2k/(1+2k^{2})\)．\(|GA|=|GB|\) 等价于 \(GM\perp l\)，即 \(x_M+k(y_M+1/2)=0\)，两边乘 \(2(1+2k^{2})\)：\(-8k^{2}+4k^{2}+k(1+2k^{2})=0\)，即 \(k(2k^{2}-4k+1)=0\)（\(k\neq 0\)），故 \(k=(2\pm\sqrt{2})/2\)．两交点条件：\(\Delta=8-16k^{2}>0\)，即 \(k^{2}<1/2\)；\(k=(2+\sqrt{2})/2\) 时 \(k^{2}=(3+2\sqrt{2})/2>1/2\)，舍；\(k=(2-\sqrt{2})/2\) 时 \(k^{2}=(3-2\sqrt{2})/2<1/2\)，符合．故 \(l\)：\(y=(2-\sqrt{2})/2\cdot(x+2)\) 或 \(y=0\)．}
\fi
\end{ansblock}

\begin{ansblock}[2章-练-课时18-14]
% ans:2章-练-课时18-14
\ansitem{15}{(1) a=2、b=1；(2) 存在，8x+3y−8=0}
\ifansbookdetail
\ansnote{详解}{(1) \(C_2\) 与 \(x\) 轴交于 \((\pm 1,0)\)，即 \(A(-1,0)\)、\(B(1,0)\)，且 \(x^{2}/b^{2}=1\) 在 \(y=0\) 处成立，故 \(b=1\)．\(e=c/a=\sqrt{3}/2\)，故 \(c^{2}=3a^{2}/4\)；\(a^{2}-c^{2}=b^{2}=1\)，得 \(a^{2}/4=1\)，故 \(a=2\)．\par\noindent (2) \(C_1\)：\(y^{2}/4+x^{2}=1\)（\(y\geqslant 0\)）．设 \(l\)：\(y=k(x-1)\)（\(k\neq 0\)；\(k=0\) 时 \(l\) 与 \(x\) 轴重合，不合）．交 \(C_1\)：\((k^{2}+4)x^{2}-2k^{2}x+k^{2}-4=0\)，一根 \(x=1\)（点 \(B\)），故 \(x_P=(k^{2}-4)/(k^{2}+4)\)、\(y_P=k(x_P-1)=-8k/(k^{2}+4)\)．交 \(C_2\)：\(k(x-1)=-(x-1)(x+1)\)，故 \((x-1)(k+x+1)=0\)，\(x_Q=-k-1\)（\(x_Q\neq 1\)，故 \(k\neq -2\)）、\(y_Q=k(x_Q-1)=-k^{2}-2k\)．\par\noindent \(\overrightarrow{AP}=(2k^{2}/(k^{2}+4),-8k/(k^{2}+4))=(2k/(k^{2}+4))\cdot(k,-4)\)；\(\overrightarrow{AQ}=(-k,-k^{2}-2k)=-k(1,k+2)\)．以 \(PQ\) 为直径的圆过 \(A\) 等价于 \(AP\perp AQ\)，即 \(\overrightarrow{AP}\cdot\overrightarrow{AQ}=0\)：\((2k/(k^{2}+4))\cdot(-k)\cdot[k-4(k+2)]=0\)，\(k\neq 0\)，故 \(3k+8=0\)，\(k=-8/3\)．回验落位：\(y_P=48/25>0\)（\(P\) 在上半椭圆）、\(y_Q=-16/9\leqslant 0\)（\(Q\) 在下半抛物线），符合．故存在 \(l\)：\(y=-(8/3)(x-1)\)，即 \(8x+3y-8=0\)．}
\fi
\end{ansblock}

\begin{ansblock}[2章-练-课时18-16]
% ans:2章-练-课时18-16
\ansitem{16}{(1) e=2；(2) 恒过定点 (4,0) 与 (−2,0)（两点同时在圆上）}
\ifansbookdetail
\ansnote{详解}{(1) \(l\perp x\) 轴时 \(M\)、\(N\) 的坐标为 \((c,\pm b^{2}/a)\)；\(A\) 在 \(x\) 轴上，\(|AM|=|AN|\) 自动成立，故直角顶点只能是 \(A\)（若直角在 \(M\)，\(MA\) 需水平，即 \(y_M=0\)，不可能），故 \(\angle MAN=90^{\circ}\)，\(k_{AM}=\pm 1\)，即 \(b^{2}/a=a+c\)，故 \(c^{2}-a^{2}=a^{2}+ac\)，即 \(e^{2}-e-2=0\)，得 \(e=2\)（负根舍）．\par\noindent (2) \(2a=4\)，故 \(a=2\)、\(c=2a=4\)、\(b^{2}=12\)，\(C\)：\(x^{2}/4-y^{2}/12=1\)，\(A(-2,0)\)，\(x=a/2=1\)．设 \(l\)：\(x=ty+4\)（\(t\) 为实数，\(t=0\) 即 \(l\perp x\) 轴亦含；\(|t|<1/\sqrt{3}\) 时与右支交两点）．代入得 \((3t^{2}-1)y^{2}+24ty+36=0\)，故 \(y_1y_2=36/(3t^{2}-1)\)、\(y_1+y_2=-24t/(3t^{2}-1)\)、\(x_1+x_2=t(y_1+y_2)+8=-8/(3t^{2}-1)\)、\(x_1x_2=t^{2}y_1y_2+4t(y_1+y_2)+16=(-12t^{2}-16)/(3t^{2}-1)\)．\par\noindent 直线 \(AM\) 过 \(A(-2,0)\)、\(M(x_1,y_1)\)，在 \(x=1\) 处 \(y=3y_1/(x_1+2)\)，故 \(P(1,p)\)、\(Q(1,q)\)，其中 \(p=3y_1/(x_1+2)\)、\(q=3y_2/(x_2+2)\)．而 \((x_1+2)(x_2+2)=x_1x_2+2(x_1+x_2)+4=(-12t^{2}-16-16+12t^{2}-4)/(3t^{2}-1)=-36/(3t^{2}-1)\)，故 \(pq=9y_1y_2/[(x_1+2)(x_2+2)]=9\cdot[36/(3t^{2}-1)]\cdot[(3t^{2}-1)/(-36)]=-9\)，与 \(t\) 无关．\par\noindent 以 \(PQ\) 为直径的圆：\((x-1)^{2}+(y-p)(y-q)=0\)，即 \((x-1)^{2}+y^{2}-(p+q)y+pq=0\)．令 \(y=0\)：\((x-1)^{2}=-pq=9\)，得 \(x=4\) 或 \(x=-2\)，即对每一条满足题设的 \(l\)，\((4,0)\) 与 \((-2,0)\) 同时在该圆上．特例回验：\(t=0\) 时 \(M\)、\(N\) 为 \((4,\pm 6)\)，\(P(1,3)\)、\(Q(1,-3)\)，圆 \((x-1)^{2}+y^{2}=9\)，两点均在圆上；\(3t^{2}-1=0\) 时 \(l\) 平行于渐近线仅一交点，不在题设内．}
\fi
\end{ansblock}

\begin{ansblock}[2章-导-课时18-G1]
% ans:2章-导-课时18-G1
\ansitem{17}{C}
\ifansbookdetail
\ansnote{详解}{\(F(p/2,0)\)，准线 \(x=-p/2\)，记准线与 \(x\) 轴交点 \(A(-p/2,0)\)，\(|AF|=p\)．\(l\) 过 \(F\) 与准线交于 \(P\)，弦 \(MN\) 过 \(F\)，沿直线点序为 \(P\)、\(N\)、\(F\)、\(M\)（\(N\) 近 \(F\)、\(M\) 远 \(F\)，\(|PN|<|PF|<|PM|\)）．由 \(|PM|=2|PF|\) 及 \(|PM|=|PF|+|FM|\) 得 \(|FM|=|PF|\)，即 \(F\) 为 \(PM\) 中点．过 \(M\) 作 \(MB\perp\) 准线于 \(B\)，由定义 \(|FM|=|BM|\)；\(A\) 为 \(PB\) 中点（\(P\)、\(B\) 同在准线上，纵坐标和为0）、\(F\) 为 \(PM\) 中点，故 \(AF\) 为 \(\triangle PMB\) 中位线，\(|BM|=2|AF|=2p\)，故 \(|FM|=|PF|=2p\)．\(P\) 与 \(F\) 的水平距离为 \(p\)，设倾斜角 \(\theta\)，则 \(|PF|\cdot|\cos\theta|=p\)，故 \(2p|\cos\theta|=p\)，\(|\cos\theta|=1/2\)，\(\theta\in(0,\pi)\)，得 \(\theta=\pi/3\) 或 \(\theta=2\pi/3\)．故选：C．\par\noindent 特例回验：取 \(p=2\)（\(y^{2}=4x\)）、\(\theta=\pi/3\)，\(l\)：\(y=\sqrt{3}(x-1)\)，联立得 \(3x^{2}-10x+3=0\)，\(M(3,2\sqrt{3})\)、\(N(1/3,-2\sqrt{3}/3)\)、\(P(-1,-2\sqrt{3})\)；\(|PF|=\sqrt{4+12}=4\)、\(|PM|=\sqrt{16+48}=8\)，\(|PM|=2|PF|\) 且 \(|FM|=x_M+1=4=|PF|\)，符合；\(\theta=2\pi/3\) 与 \(\pi/3\) 关于 \(x\) 轴对称，同成立．}
\fi
\end{ansblock}

\begin{ansblock}[2章-导-课时18-G2]
% ans:2章-导-课时18-G2
\ansitem{18}{D}
\ifansbookdetail
\ansnote{详解}{\(a=2\)、\(b=\sqrt{3}\)，故 \(c=\sqrt{4-3}=1\)，\(F_2(1,0)\)．通径：\(x=1\) 代入得 \(y^{2}=3(1-1/4)=9/4\)，\(|y|=3/2\)，故 \(|AB|=3\)（\(=2b^{2}/a=3\)，两法互证）．以 \(|F_1F_2|=2\) 为底：\(S_{\triangle ABF_1}=|AB|\cdot 2/2=3\)．周长：\(|AF_1|+|BF_1|=(2a-|AF_2|)+(2a-|BF_2|)=8-|AB|=5\)，故周长为 \(5+3=8\)．内切圆半径 \(r=2\times 3/8=3/4\)．故选：D．}
\fi
\end{ansblock}

\begin{ansblock}[2章-导-课时18-G3]
% ans:2章-导-课时18-G3
\ansitem{19}{C}
\ifansbookdetail
\ansnote{详解}{设直线 \(y=kx-2\)（过 \((0,-2)\)），\(A(x_1,y_1)\)、\(B(x_2,y_2)\)．代入 \(y^{2}=8x\)：\(k^{2}x^{2}-4(k+2)x+4=0\)，\(\Delta=16(k+2)^{2}-16k^{2}=64(k+1)>0\)，故 \(k>-1\)．中点横坐标 \((x_1+x_2)/2=2(k+2)/k^{2}=2\)，故 \(k^{2}=k+2\)，即 \(k^{2}-k-2=0\)，得 \(k=2\) 或 \(k=-1\)；\(k=-1\) 时 \(\Delta=0\)（相切不成弦），舍，故 \(k=2\)．此时 \(x^{2}-4x+1=0\)，\(x_1+x_2=4\)、\(x_1x_2=1\)，弦长 \(|AB|=\sqrt{1+k^{2}}\cdot\sqrt{(x_1+x_2)^{2}-4x_1x_2}=\sqrt{5}\cdot\sqrt{12}=2\sqrt{15}\)．故选：C．}
\fi
\end{ansblock}

\begin{ansblock}[2章-导-课时18-G4]
% ans:2章-导-课时18-G4
\ansitem{20}{81π/16}
\ifansbookdetail
\ansnote{详解}{\(a=4\)、\(b=3\)，故 \(c=\sqrt{16-9}=\sqrt{7}\)，\(F(\sqrt{7},0)\)．\(x=\sqrt{7}\) 代入：\(y^{2}=9(1-7/16)=81/16\)，\(|y|=9/4\)，故 \(|AB|=9/2\)（与通径 \(2b^{2}/a=18/4=9/2\) 一致）．圆半径 \(r=|AB|/2=9/4\)，面积 \(S=\pi r^{2}=81\pi/16\)．}
\fi
\end{ansblock}

\begin{ansblock}[2章-导-课时18-G5]
% ans:2章-导-课时18-G5
\ansitem{21}{9/4}
\ifansbookdetail
\ansnote{详解}{\(2p=3\)，故 \(p=3/2\)，\(F(3/4,0)\)，\(k=\tan 30^{\circ}=\sqrt{3}/3\)，直线 \(y=(\sqrt{3}/3)(x-3/4)\)．代入 \(y^{2}=3x\)：\((x-3/4)^{2}=9x\)，即 \(x^{2}-(21/2)x+9/16=0\)（\(\Delta=441/4-9/4=108>0\)），\(x_1+x_2=21/2\)．由抛物线定义 \(|AB|=x_1+x_2+p=21/2+3/2=12\)．\(O\) 到直线 \(AB\)（即 \(\sqrt{3}x-3y-3\sqrt{3}/4=0\)）的距离 \(d=(3\sqrt{3}/4)/\sqrt{3+9}=3/8\)，故 \(S_{\triangle ABO}=|AB|\cdot d/2=12\times(3/8)/2=9/4\)．}
\fi
\end{ansblock}

\jietitle{章末总结与复习}

\begin{ansblock}[2章-练-课时19-01]
% ans:2章-练-课时19-01
\ansitem{1}{A}
\ifansbookdetail
\ansnote{详解}{求根公式：\(x=2\pm\sqrt{3}\)，其中 \(2-\sqrt{3}\approx 0.27\in(0,1)\) 可作椭圆离心率，\(2+\sqrt{3}\approx 3.73\in(1,+\infty)\) 可作双曲线离心率．抛物线离心率恒等于 1，两根均不为 1，排除 B、C；两椭圆离心率都须在 \((0,1)\) 内，而 \(2+\sqrt{3}>1\)，排除 D．故选：A．}
\fi
\end{ansblock}

\begin{ansblock}[2章-练-课时19-02]
% ans:2章-练-课时19-02
\ansitem{2}{B}
\ifansbookdetail
\ansnote{详解}{\(k<9\) 得 \(25-k>9-k>0\)，第二曲线为椭圆，其 \(c^{2}=(25-k)-(9-k)=16\)，与 \(k\) 无关；第一曲线 \(c^{2}=25-9=16\)．故两曲线半焦距均为 4、焦距均为 8，选：B．长轴 10 与 \(2\sqrt{25-k}\)、短轴 6 与 \(2\sqrt{9-k}\)、离心率的平方 \(16/25\) 与 \(16/(25-k)\) 均随 \(k\) 变动，排除 A、C、D．}
\fi
\end{ansblock}

\begin{ansblock}[2章-练-课时19-03]
% ans:2章-练-课时19-03
\ansitem{3}{x²/4+y²/3=1（x>0，y≠0），椭圆 x²/4+y²/3=1 的右半部分（不含与 x 轴的交点）}
\ifansbookdetail
\ansnote{详解}{\(|BC|=2\)，故 \(|AB|+|CA|=2|BC|=4>|BC|=2\)，由椭圆定义，\(A\) 的轨迹是以 \(B\)、\(C\) 为焦点、长半轴 \(a=2\)、半焦距 \(c=1\)、短半轴 \(b^{2}=a^{2}-c^{2}=3\) 的椭圆 \(x^{2}/4+y^{2}/3=1\)（去掉长轴端点，因 \(4>|BC|\) 严格成立）．在该椭圆上 \(|AB|=a+ex\)、\(|CA|=a-ex\)（\(e=1/2\)），\(|AB|>|CA|\) 即 \(x>0\)．又 \(A\) 为三角形顶点，不能与 \(B\)、\(C\) 共线，即 \(y\neq 0\)（舍去点 \((2,0)\)）．故轨迹为 \(x^{2}/4+y^{2}/3=1\)（\(x>0\)，\(y\neq 0\)），是椭圆的右半部分（不含与 \(x\) 轴的交点）．}
\fi
\end{ansblock}

\begin{ansblock}[2章-练-课时19-04]
% ans:2章-练-课时19-04
\ansitem{4}{a<−1 或 a>3}
\ifansbookdetail
\ansnote{详解}{双曲线型须二次项系数异号：\((3-a)(a+1)<0\)，即 \(a<-1\) 或 \(a>3\)．此时常数项 \(a^{2}-2a-3=(a-3)(a+1)>0\) 自动相容，方程确为双曲线：\(a>3\) 时移项配方得 \(x^{2}/(a+1)-y^{2}/(a-3)=1\)；\(a<-1\) 时同法得 \(x^{2}/(a+1)-y^{2}/(a-3)=1\)（分子分母同乘 \(-1\) 后仍为平方差形式）．即 \(-1<a<3\) 时为椭圆型或虚轨迹、退化，不取．故 \(a\in(-\infty,-1)\cup(3,+\infty)\)．}
\fi
\end{ansblock}

\begin{ansblock}[2章-练-课时19-05]
% ans:2章-练-课时19-05
\ansitem{5}{k<−√6/2 或 k>√6/2}
\ifansbookdetail
\ansnote{详解}{联立消 \(y\)：\((1+2k^{2})x^{2}+8kx+6=0\)．相交于不同的两点即 \(\Delta>0\)：\(\Delta=64k^{2}-24(1+2k^{2})=16k^{2}-24=8(2k^{2}-3)>0\)，即 \(k^{2}>3/2\)，\(|k|>\sqrt{6}/2\)．故 \(k<-\sqrt{6}/2\) 或 \(k>\sqrt{6}/2\)．}
\fi
\end{ansblock}

\begin{ansblock}[2章-练-课时19-06]
% ans:2章-练-课时19-06
\ansitem{6}{y=4 或 y=x+2}
\ifansbookdetail
\ansnote{详解}{①斜率不存在：\(x=2\) 代入 \(y^{2}=16\) 得两点 \((2,\pm 4)\)，不合，舍．②斜率 \(k=0\)：\(y=4\)，与抛物线恰一公共点 \((2,4)\)，符合．③\(k\neq 0\)：设 \(y-4=k(x-2)\)，以 \(x=y^{2}/8\) 代入并整理：\(ky^{2}-8y+32-16k=0\)，恰一公共点即 \(\Delta=64-4k(32-16k)=64(k-1)^{2}=0\)，得 \(k=1\)，即 \(y=x+2\)．综上共两条：\(y=4\) 与 \(y=x+2\)．易错提醒：与对称轴平行的情形（\(k=0\)）不可漏．}
\fi
\end{ansblock}

\begin{ansblock}[2章-练-课时19-07]
% ans:2章-练-课时19-07
\ansitem{7}{C}
\ifansbookdetail
\ansnote{详解}{\(e=\sqrt{2}\) 得 \(a=b\)，双曲线为等轴双曲线，渐近线 \(y=\pm x\)．准线 \(x=-p/2\)，交渐近线于 \(A(-p/2,p/2)\)、\(B(-p/2,-p/2)\)，\(|AB|=p\)，\(O\) 到 \(AB\) 的距离为 \(p/2\)，故 \(S=\frac{1}{2}\cdot p\cdot\frac{p}{2}=p^{2}/4=4\)，得 \(p=4\)，抛物线的方程为 \(y^{2}=8x\)．故选：C．}
\fi
\end{ansblock}

\begin{ansblock}[2章-练-课时19-08]
% ans:2章-练-课时19-08
\ansitem{8}{(1) 2√2；(2) y=[(2√2+√5)/3]x−(2√10+2)/3；(3) 证毕（O 到 AB 距离恒为 √2）}
\ifansbookdetail
\ansnote{详解}{\(c^{2}=4+4=8\)，故 \(F(2\sqrt{2},0)\)．圆心 \((0,m)\)（\(m>0\)）且过 \(O\)，半径为 \(m\)，圆 \(C\)：\(x^{2}+y^{2}-2my=0\)．\par\noindent (1) \(AF\perp OF\)（\(x\) 轴）得 \(x_A=2\sqrt{2}\)，代入 \(\Gamma\) 得 \(y_A^{2}=8-4=4\)，取 \(A(2\sqrt{2},2)\)，故 \(S=\frac{1}{2}\cdot|OF|\cdot|y_A|=\frac{1}{2}\cdot 2\sqrt{2}\cdot 2=2\sqrt{2}\)．\par\noindent (2) 代 \(A(\sqrt{5},1)\)：\(5+1-2m=0\)，得 \(m=3\)．圆与 \(\Gamma\) 联立（\(x^{2}-y^{2}=4\) 与 \(x^{2}+y^{2}-6y=0\) 相减）得 \(y^{2}-3y+2=0\)，即 \(y=1\)（点 \(A\)）或 \(y=2\)，故 \(B(2\sqrt{2},2)\)．\(k_{AB}=(2-1)/(2\sqrt{2}-\sqrt{5})=(2\sqrt{2}+\sqrt{5})/3\)（分母有理化，\((2\sqrt{2})^{2}-(\sqrt{5})^{2}=3\)）；截距 \(=1-k\sqrt{5}=1-(2\sqrt{10}+5)/3=-(2\sqrt{10}+2)/3\)，故直线 \(AB\)：\(y=\frac{2\sqrt{2}+\sqrt{5}}{3}x-\frac{2\sqrt{10}+2}{3}\)．\par\noindent (3) 一般情形：圆 \(x^{2}+y^{2}-2my=0\) 与 \(\Gamma\) 消 \(x\) 得 \(y^{2}-my+2=0\)，故 \(y_1y_2=2\)、\(y_1^{2}+y_2^{2}=m^{2}-4\)；右支上 \(x_i=\sqrt{y_i^{2}+4}\)，故 \(x_1x_2=\sqrt{(y_1y_2)^{2}+4(y_1^{2}+y_2^{2})+16}=2\sqrt{m^{2}+1}\)、\(x_1^{2}+x_2^{2}=m^{2}+4\)．原点到 \(AB\) 的距离 \(d=|x_1y_2-x_2y_1|/|AB|\)，其中分子平方 \(=x_1^{2}y_2^{2}+x_2^{2}y_1^{2}-2x_1x_2y_1y_2=y_2^{2}(y_1^{2}+4)+y_1^{2}(y_2^{2}+4)-4\sqrt{m^{2}+1}\cdot 2=2(y_1y_2)^{2}+4(y_1^{2}+y_2^{2})-8\sqrt{m^{2}+1}=4m^{2}-8-8\sqrt{m^{2}+1}\)；分母平方 \(=(x_1^{2}+x_2^{2}-2x_1x_2)+(y_1^{2}+y_2^{2}-2y_1y_2)=(m^{2}+4-4\sqrt{m^{2}+1})+(m^{2}-4-4)=2m^{2}-4-4\sqrt{m^{2}+1}\)．分子平方 \(=2\times\)分母平方恒成立（\(m>2\sqrt{2}\) 时分母 \(>0\)），故 \(d^{2}\equiv 2\)，即 \(d=\sqrt{2}\) 恒等于圆 \(x^{2}+y^{2}=2\) 的半径，直线 \(AB\) 与该圆相切，证毕．}
\fi
\end{ansblock}

\begin{ansblock}[2章-练-课时19-09]
% ans:2章-练-课时19-09
\ansitem{9}{A}
\ifansbookdetail
\ansnote{详解}{设 \(|MF_1|>|MF_2|\)，则 \(|MF_1|+|MF_2|=2a_1\)、\(|MF_1|-|MF_2|=2a_2\)，得 \(|MF_1|=a_1+a_2\)、\(|MF_2|=a_1-a_2\)．\(\angle F_1MF_2=90^{\circ}\) 得 \(|MF_1|^{2}+|MF_2|^{2}=4c^{2}\)，即 \((a_1+a_2)^{2}+(a_1-a_2)^{2}=4c^{2}\)，故 \(a_1^{2}+a_2^{2}=2c^{2}\)，同除 \(c^{2}\)：\(1/e_1^{2}+1/e_2^{2}=2\)．\(e_1\in[\sqrt{6}/3,1)\) 得 \(1/e_1^{2}\in(1,3/2]\)，故 \(1/e_2^{2}=2-1/e_1^{2}\in[1/2,1)\)，即 \(e_2^{2}\in(1,2]\)，\(e_2\in(1,\sqrt{2}]\)．故选：A．}
\fi
\end{ansblock}

\begin{ansblock}[2章-练-课时19-10]
% ans:2章-练-课时19-10
\ansitem{10}{2}
\ifansbookdetail
\ansnote{详解}{两曲线关于原点对称，故 \(O\) 为 \(AB\) 中点，四边形 \(AF_1BF_2\) 为平行四边形，其对角线 \(|AB|=2|OF_1|=2c\) 与 \(|F_1F_2|=2c\) 相等，故四边形为矩形，\(\angle F_1AF_2=\angle F_1BF_2=90^{\circ}\)．\(\triangle OF_1B\) 中 \(|OF_1|=|OB|=c\)，\(\angle OF_1B=\pi/6\)，顶角 \(\angle F_1OB=2\pi/3\)；由 \(A\) 与 \(B\) 关于原点对称得 \(\angle F_1OA=\pi-2\pi/3=\pi/3\)，而 \(|OA|=|OF_1|=c\)，故 \(\triangle AOF_1\) 为等边三角形，\(|AF_1|=c\) 且 \(\angle AF_1O=\pi/3\)，\(\angle AF_1B=\pi/3+\pi/6=\pi/2\)，与矩形一致．于是 \(|BF_1|=\sqrt{|AB|^{2}-|AF_1|^{2}}=\sqrt{4c^{2}-c^{2}}=\sqrt{3}c\)，且平行四边形对边 \(|AF_2|=|BF_1|=\sqrt{3}c\)．椭圆定义：\(c+\sqrt{3}c=2a_1\)，故 \(e_1=2/(\sqrt{3}+1)\)；双曲线定义：\(\sqrt{3}c-c=2a_2\)，故 \(e_2=2/(\sqrt{3}-1)\)．\(e_1e_2=4/[(\sqrt{3}+1)(\sqrt{3}-1)]=4/2=2\)．}
\fi
\end{ansblock}

\begin{ansblock}[2章-练-课时19-11]
% ans:2章-练-课时19-11
\ansitem{11}{BCD}
\ifansbookdetail
\ansnote{详解}{A：圆须 \(16+\lambda=9+\lambda\)，矛盾，不存在，不合．B：\(\lambda>-9\) 时 \(16+\lambda>9+\lambda>0\)，方程表示焦点在 \(x\) 轴上的椭圆，符合．C：\(-16<\lambda<-9\) 时 \((16+\lambda)(9+\lambda)<0\)，表示焦点在 \(x\) 轴上的双曲线，符合．D：\(\lambda>-9\) 时为椭圆，\(c^{2}=(16+\lambda)-(9+\lambda)=7\)；\(-16<\lambda<-9\) 时为双曲线，\(c^{2}=(16+\lambda)+[-(9+\lambda)]=7\)，焦距恒为 \(2\sqrt{7}\)，符合．故选：BCD．}
\fi
\end{ansblock}

\begin{ansblock}[2章-练-课时19-12]
% ans:2章-练-课时19-12
\ansitem{12}{BD}
\ifansbookdetail
\ansnote{详解}{\(\overrightarrow{MF_1}\cdot\overrightarrow{MF_2}=0\) 且 \(|MF_1|=|MF_2|=\sqrt{b^{2}+c^{2}}\)，故 \(\triangle MF_1F_2\) 为等腰直角三角形，\(b=c\)，故 \(e_1^{2}=b^{2}/(b^{2}+c^{2})=1/2\)，\(e_1=\sqrt{2}/2\)，且 \(a=\sqrt{2}c\)．在焦点三角形 \(\triangle PF_1F_2\) 中 \(\angle F_1PF_2=\pi/3\)，设 \(|PF_1|=m\)、\(|PF_2|=n\)（\(m>n\)，\(P\) 取双曲线右支），则 \(m+n=2a_1=2\sqrt{2}c\)、\(m-n=2a_2\)．余弦定理：\(4c^{2}=m^{2}+n^{2}-2mn\cos(\pi/3)=(m+n)^{2}-3mn\)，故 \(mn=(8c^{2}-4c^{2})/3=4c^{2}/3\)．于是 \((m-n)^{2}=(m+n)^{2}-4mn=8c^{2}-16c^{2}/3=8c^{2}/3\)，即 \(4a_2^{2}=8c^{2}/3\)，\(a_2^{2}=2c^{2}/3\)，\(e_2^{2}=3/2\)，\(e_2=\sqrt{6}/2\)．逐项：\(e_1e_2=(\sqrt{2}/2)\cdot(\sqrt{6}/2)=\sqrt{3}/2\)，B 符合；\(e_2^{2}-e_1^{2}=3/2-1/2=1\)，D 符合；\(e_2/e_1=\sqrt{3}\neq 2\)，A 不合；\(e_1^{2}+e_2^{2}=2\neq 5/2\)，C 不合．故选：BD．}
\fi
\end{ansblock}

\begin{ansblock}[2章-练-课时19-13]
% ans:2章-练-课时19-13
\ansitem{13}{(1) x²/2−y²/2=1、y=±x；(2) x−2y+3=0}
\ifansbookdetail
\ansnote{详解}{(1) 椭圆 \(c^{2}=18-14=4\)，故双曲线 \(c=2\)、\(a^{2}+b^{2}=4\)；\(A\) 代入：\(9/a^{2}-7/b^{2}=1\)．令 \(u=a^{2}\)：\(9/u-7/(4-u)=1\)，即 \(9(4-u)-7u=u(4-u)\)，整理得 \(u^{2}-20u+36=0\)，解得 \(u=2\) 或 \(u=18\)（\(>c^{2}\) 舍），故 \(a^{2}=b^{2}=2\)，双曲线 \(C\)：\(x^{2}/2-y^{2}/2=1\)，渐近线 \(y=\pm x\)（等轴双曲线）．\par\noindent (2) 点差法：\(x_1^{2}-y_1^{2}=2\)、\(x_2^{2}-y_2^{2}=2\) 相减得 \((x_1+x_2)(x_1-x_2)=(y_1+y_2)(y_1-y_2)\)，故 \(k_{AB}=(y_1-y_2)/(x_1-x_2)=(x_1+x_2)/(y_1+y_2)=2/4=1/2\)，直线为 \(y-2=\frac{1}{2}(x-1)\)，即 \(x-2y+3=0\)．回验：联立消 \(y\) 得 \(3x^{2}-6x-17=0\)，\(\Delta=36+204=240>0\)，为真弦，符合．}
\fi
\end{ansblock}

\begin{ansblock}[2章-练-课时19-14]
% ans:2章-练-课时19-14
\ansitem{14}{A}
\ifansbookdetail
\ansnote{详解}{\(a^{2}=(\sqrt{5}-1)^{2}=6-2\sqrt{5}\)．\(a/c=(\sqrt{5}-1)/2\) 得 \(c=2a/(\sqrt{5}-1)=a(\sqrt{5}+1)/2\)，故 \(c^{2}=a^{2}(3+\sqrt{5})/2\)．\(m=c^{2}-a^{2}=a^{2}[(3+\sqrt{5})/2-1]=a^{2}(1+\sqrt{5})/2=(6-2\sqrt{5})(\sqrt{5}+1)/2=(3-\sqrt{5})(\sqrt{5}+1)=3\sqrt{5}+3-5-\sqrt{5}=2\sqrt{5}-2\)．故选：A．}
\fi
\end{ansblock}

\begin{ansblock}[2章-练-课时19-15]
% ans:2章-练-课时19-15
\ansitem{15}{(1) 三条件均得 x²/9−y²/16=1；(2) |PF₂|=2 或 14}
\ifansbookdetail
\ansnote{详解}{(1) \(c^{2}=49-24=25\)，故 \(c=5\)．选①：左顶点 \((-3,0)\) 得 \(a=3\)．选②：\(18/a^{2}-16/b^{2}=1\) 配 \(a^{2}+b^{2}=25\)，即 \(18(25-a^{2})-16a^{2}=a^{2}(25-a^{2})\)，整理得 \(a^{4}-59a^{2}+450=0\)，解得 \(a^{2}=9\) 或 \(a^{2}=50\)（\(>25\) 舍），故 \(a^{2}=9\)．选③：\(e=c/a=5/3\)，即 \(5/a=5/3\)，得 \(a=3\)．三路均得 \(b^{2}=25-9=16\)，双曲线的方程为 \(x^{2}/9-y^{2}/16=1\)．\par\noindent (2) \(\bigl||PF_1|-|PF_2|\bigr|=2a=6\)，故 \(|PF_2|=2\) 或 \(14\)，均须回验：\(|PF_2|=2\) 时 \(P\) 为右顶点 \((3,0)\)（\(|PF_1|=3+5=8\)，且 \(8-2=6\) 合定义——\(P\)、\(F\) 成线段，退化但合法）；\(|PF_2|=14\) 时 \(P\) 在左支（\(x=-33/5\)，\(|PF_1|=8\)、\(14-8=6\)，且 \(8+14=22>2c=20\) 三角形存在）．故 \(|PF_2|=2\) 或 \(14\)．}
\fi
\end{ansblock}

\begin{ansblock}[2章-练-课时19-16]
% ans:2章-练-课时19-16
\ansitem{16}{(1) p₁=2；(2) x_A=p/(p+1)，斜率范围 (−1,0)∪(0,1)}
\ifansbookdetail
\ansnote{详解}{(1) \(C_1\) 的准线 \(x=-p_1/2=-1\)，得 \(p_1=2\)，即 \(C_1\)：\(y^{2}=4x\)．\par\noindent (2) 设 \(B(-t^{2}/(2p),t)\)（\(C_2\) 上的参数点），\(B\) 为 \(AC\) 中点、\(C(-1,0)\)，故 \(A=2B-C=(1-t^{2}/p,2t)\)．\(A\) 在 \(C_1\) 上：\(4t^{2}=4(1-t^{2}/p)\)，得 \(t^{2}=p/(p+1)\in(0,1)\)（\(p>0\)），故 \(A\) 的横坐标 \(x_A=1-t^{2}/p=1-1/(p+1)=p/(p+1)\)（\(0<x_A<1\)，符合）．斜率：\(k_{AC}=k_{BC}=t/[-t^{2}/(2p)+1]\)；由 \(t^{2}=p/(p+1)\) 得 \(1/p=1/t^{2}-1\)，故分母 \(=-(t^{2}/2)(1/t^{2}-1)+1=(1+t^{2})/2\)，即 \(k=2t/(1+t^{2})\)（\(t\neq 0\)，否则 \(A=C\) 不成弦）．\(k^{2}=4t^{2}/(1+t^{2})^{2}\)，由 \(0<t^{2}<1\) 得 \(0<k^{2}<1\)（等号 \(t^{2}=1\) 不可达），故 \(k\in(-1,0)\cup(0,1)\)．}
\fi
\end{ansblock}

\begin{ansblock}[2章-导-课时19-G1]
% ans:2章-导-课时19-G1
\ansitem{17}{D}
\ifansbookdetail
\ansnote{详解}{\(C\) 即 \(x^{2}/4-y^{2}/4=1\)，实轴在 \(x\) 轴上，渐近线 \(y=\pm x\)．联立 \(l\) 与 \(C\)：\(x^{2}-(kx+2)^{2}=4\)，即 \((1-k^{2})x^{2}-4kx-8=0\)．左支 \(x<0\)、右支 \(x>0\)，故「两支各一交点」等价于方程两实根异号：\(x_1x_2=-8/(1-k^{2})<0\)，即 \(1-k^{2}>0\)，得 \(-1<k<1\)；此时 \(\Delta=16k^{2}+32(1-k^{2})=32-16k^{2}>0\) 自动满足（\(k=\pm 1\) 时方程退化为一元一次、至多一交点，已排除）．几何印证：\(l\) 过定点 \((0,2)\)，\(k=\pm 1\) 时平行于渐近线仅扫一支．故 \(-1<k<1\)，选：D．}
\fi
\end{ansblock}

\begin{ansblock}[2章-导-课时19-G2]
% ans:2章-导-课时19-G2
\ansitem{18}{C}
\ifansbookdetail
\ansnote{详解}{渐近线 \(y=\pm(b/a)x\)．直线 \(y=\sqrt{3}x\) 过原点：当 \(\sqrt{3}<b/a\) 时直线落在含 \(x\) 轴的开口区域内，与右支必有交点；当 \(\sqrt{3}=b/a\) 时直线即渐近线，无交点；当 \(\sqrt{3}>b/a\) 时直线落在含 \(y\) 轴的区域内，双曲线在该区域无点．故「有交点」等价于 \(b/a>\sqrt{3}\)．于是 \(e=\sqrt{1+b^{2}/a^{2}}>\sqrt{1+3}=2\)，即 \(e\in(2,+\infty)\)．故选：C．}
\fi
\end{ansblock}

\begin{ansblock}[2章-导-课时19-G3]
% ans:2章-导-课时19-G3
\ansitem{19}{D}
\ifansbookdetail
\ansnote{详解}{（点差法）设 \(A(x_1,y_1)\)、\(B(x_2,y_2)\)，由中点 \(P(1,4)\)：\(x_1+x_2=2\)、\(y_1+y_2=8\)；由 \(l\) 的斜率为 1：\((y_1-y_2)/(x_1-x_2)=1\)．\(A\)、\(B\) 代入 \(C\) 并相减：\((x_1+x_2)(x_1-x_2)/a^{2}=(y_1+y_2)(y_1-y_2)/b^{2}\)，故 \(b^{2}/a^{2}=(y_1+y_2)(y_1-y_2)/[(x_1+x_2)(x_1-x_2)]=8/2=4\)，即 \(b^{2}=4a^{2}\)，\(e=\sqrt{1+b^{2}/a^{2}}=\sqrt{5}\)．回验：设 \(a^{2}=t\)、\(b^{2}=4t\)，\(y=x+3\) 代入 \(C\)：\(4x^{2}-(x+3)^{2}=4t\)，即 \(3x^{2}-6x-9-4t=0\)，\(\Delta=36+12(9+4t)=144+48t>0\) 恒成立，且 \(x_1+x_2=2\)、\(y_1+y_2=2[(x_1+x_2)/2+3]=8\)，与中点吻合．故选：D．}
\fi
\end{ansblock}

\begin{ansblock}[2章-导-课时19-G4]
% ans:2章-导-课时19-G4
\ansitem{20}{2x−y−15=0}
\ifansbookdetail
\ansnote{详解}{（点差法）设弦端点 \(A(x_1,y_1)\)、\(B(x_2,y_2)\)，则 \(x_1+x_2=16\)、\(y_1+y_2=2\)．\(x_1^{2}-4y_1^{2}=4\) 与 \(x_2^{2}-4y_2^{2}=4\) 相减：\((x_1+x_2)(x_1-x_2)-4(y_1+y_2)(y_1-y_2)=0\)，即 \(16(x_1-x_2)=8(y_1-y_2)\)，故 \(k=(y_1-y_2)/(x_1-x_2)=2\)．直线：\(y-1=2(x-8)\)，即 \(2x-y-15=0\)．回验 \(\Delta>0\)：代入得 \(x^{2}-4(2x-15)^{2}=4\)，即 \(15x^{2}-240x+904=0\)，\(\Delta=57600-54240=3360>0\)，弦存在，符合．}
\fi
\end{ansblock}

\begin{ansblock}[2章-导-课时19-G5]
% ans:2章-导-课时19-G5
\ansitem{21}{4x−y−3=0}
\ifansbookdetail
\ansnote{详解}{（点差法）设弦端点 \(M(x_1,y_1)\)、\(N(x_2,y_2)\)，由中点 \(A(1,1)\)：\(x_1+x_2=2\)、\(y_1+y_2=2\)．若弦垂直于 \(x\) 轴，其中点应在 \(x\) 轴上（纵坐标和为 0），与 \(y_1+y_2=2\) 矛盾，故斜率存在．\(y_1^{2}=8x_1\) 与 \(y_2^{2}=8x_2\) 相减：\((y_1+y_2)(y_1-y_2)=8(x_1-x_2)\)，故 \(k=(y_1-y_2)/(x_1-x_2)=8/(y_1+y_2)=4\)．直线：\(y-1=4(x-1)\)，即 \(4x-y-3=0\)．回验 \(\Delta>0\)：代入得 \((4x-3)^{2}=8x\)，即 \(16x^{2}-32x+9=0\)，\(\Delta=1024-576=448>0\)，弦存在，符合（且 \(x_1+x_2=32/16=2\) 与中点吻合）．}
\fi
\end{ansblock}

\tailfill[30mm]

\end{multicols}
\end{document}
